\chapter{Riemann 积分}\label{ch:6}

在第 1 章中, 我们曾利用以直代曲的方法将由曲线 $y=x^{2}$, $y=0$, $x=1$ 围成的曲边三角形的面积, 化成一个极限来求. 这一章, 我们将进一步通过计算曲边梯形的面积、由速度求位移及由直线的线密度求质量引进 Riemann 积分的概念, 进而给出 Riemann 可积的各种充分必要条件与 Riemann 积分的简单性质, 着重论述微积分学基本定理与基本公式 (Newton-Leibniz 公式), 再应用变量代换、分部积分等方法计算各种典型的 Riemann 积分实例.

\section{Riemann 积分的概念、Riemann 可积的充要条件}\label{sec:6.1}

考虑由直线 $x=a$, $x=b$, $y=0$, $y=f(x)$ 所围曲边梯形 $ABCD$ 的面积. 先将梯形用直线 $x=x_{i}$, $i=0,1,\cdots,n$, 其中 $a=x_{0}<x_{1}<\cdots<x_{n}=b$, 分成 $n$ 个小曲边梯形, 第 $i$ 个小梯形的面积用 $\Delta s_{i}$ 表示, $\Delta s_{i}\approx f(\xi_{i})\Delta x_{i}$, 其中 $\Delta x_{i}=x_{i}-x_{i-1}$, $\xi_{i}$ 取自区间 $[x_{i-1},x_{i}]$. 于是用 $\sum\limits_{i=1}^{n}f(\xi_{i})\Delta x_{i}$ (图~\ref{fig:6.1.1} 中 $n$ 个小矩形面积之和) 表示曲边梯形面积的近似值. 当分点越来越多时, 视曲边梯形 $ABCD$ 的面积的精确值为
\[
S_{ABCD}=\lim_{\|T\|\to 0}\sum_{i=1}^{n}f(\xi_{i})\Delta x_{i}\xlongequal{\text{def}}\int_{a}^{b}f(x)\,\mathrm{d}x,
\]
其中 $f(x)\geqslant 0$, $T:a=x_{0}<x_{1}<\cdots<x_{n}=b$, $\xi_{i}\in[x_{i-1},x_{i}]$, $\|T\|=\max\{\Delta x_{i}\}$, $\Delta x_{i}=x_{i}-x_{i-1}$, $1\leqslant i\leqslant n$.

\begin{figure}[htbp]
\centering
\begin{tikzpicture}[x=1.12cm,y=1.0cm,>=Stealth]
  \draw[->] (-0.3,0) -- (8.6,0) node[below right] {$x$};
  \draw[->] (0,-0.3) -- (0,2.95) node[left] {$y$};
  \node[below left] at (0,0) {$O$};
  \fill[pattern=north east lines] (1.2,0) rectangle (2.1,1.30);
  \draw (1.2,0) rectangle (2.1,1.30);
  \fill[pattern=north east lines] (2.1,0) rectangle (2.85,0.95);
  \draw (2.1,0) rectangle (2.85,0.95);
  \fill[pattern=north east lines] (4.0,0) rectangle (4.7,1.05);
  \draw (4.0,0) rectangle (4.7,1.05);
  \fill[pattern=north east lines] (4.7,0) rectangle (5.4,1.05);
  \draw (4.7,0) rectangle (5.4,1.05);
  \fill[pattern=north east lines] (6.5,0) rectangle (7.35,2.45);
  \draw (6.5,0) rectangle (7.35,2.45);
  \draw[very thick] plot[smooth] coordinates {(1.2,1.30) (1.7,1.08) (2.2,0.88) (2.85,0.74) (3.5,0.62) (4.0,0.64) (4.7,0.86) (5.4,1.15) (6.0,1.52) (6.5,1.90) (7.0,2.22) (7.35,2.45)};
  \node[left] at (1.15,1.30) {$D$};
  \node[right] at (7.4,2.45) {$C$};
  \node[above left] at (1.2,0) {$A$};
  \node[above right] at (7.35,0) {$B$};
  \node[above] at (4.7,1.05) {$f(\xi_{i})$};
  \foreach \xx/\lbl in {1.2/a,2.1/x_{1},2.85/x_{2},4.0/x_{i-1},4.7/\xi_{i},5.4/x_{i},6.5/x_{n-1},7.35/b}
    \node[below] at (\xx,-0.06) {$\lbl$};
\end{tikzpicture}
\caption{}\label{fig:6.1.1}
\end{figure}

对于一般的函数 $f(x)$, $\int_{a}^{b}|f(x)|\,\mathrm{d}x$ 为所围图形 (见图~\ref{fig:6.1.2}) 的总面积, 而 $\int_{a}^{b}f(x)\,\mathrm{d}x$ 为所围图形 (见图~\ref{fig:6.1.3}) 面积的代数和 ($x$ 轴上方的面积为正, 下方的面积为负).

\begin{figure}[htbp]
\centering
\begin{minipage}[c]{0.47\textwidth}
\centering
\begin{tikzpicture}[x=0.80cm,y=0.80cm,>=Stealth]
  \draw[->] (-0.3,0) -- (6.4,0) node[below right] {$x$};
  \draw[->] (0,-1.35) -- (0,2.0) node[left] {$y$};
  \node[below left] at (0,0) {$O$};
  \fill[pattern=north east lines] (1.0,1.45) -- (2.6,0) -- (1.0,0) -- cycle;
  \draw[thick] (1.0,1.45) -- (2.6,0);
  \fill[pattern=north east lines] (2.6,0) .. controls (3.05,0.95) and (4.15,0.95) .. (4.6,0) -- cycle;
  \draw[thick] (2.6,0) .. controls (3.05,0.95) and (4.15,0.95) .. (4.6,0);
  \fill[pattern=north east lines] (4.6,0) -- (5.9,2.25) -- (5.9,0) -- cycle;
  \draw[thick] (4.6,0) -- (5.9,2.25);
  \draw[thick,dotted] (2.6,0) .. controls (3.05,-0.95) and (4.15,-0.95) .. (4.6,0);
  \node[below] at (1.0,-0.06) {$a$};
  \node[below] at (5.9,-0.06) {$b$};
  \node[below] at (3.6,-1.18) {$y=f(x)$};
\end{tikzpicture}
\caption{}\label{fig:6.1.2}
\end{minipage}\hfill
\begin{minipage}[c]{0.47\textwidth}
\centering
\begin{tikzpicture}[x=0.80cm,y=0.80cm,>=Stealth]
  \draw[->] (-0.3,0) -- (6.4,0) node[below right] {$x$};
  \draw[->] (0,-1.35) -- (0,2.0) node[left] {$y$};
  \node[below left] at (0,0) {$O$};
  \draw[very thick] plot[smooth] coordinates {(1.0,1.45) (1.8,0.78) (2.6,0.0) (3.1,-0.70) (3.6,-0.88) (4.1,-0.70) (4.6,0.0) (5.25,1.15) (5.9,2.25)};
  \node[below] at (1.0,-0.06) {$a$};
  \node[below] at (5.9,-0.06) {$b$};
  \node at (1.85,0.36) {$+$};
  \node at (3.6,-0.48) {$-$};
  \node at (5.15,0.55) {$+$};
  \node[below] at (3.6,-1.05) {$y=f(x)$};
\end{tikzpicture}
\caption{}\label{fig:6.1.3}
\end{minipage}
\end{figure}

再考虑质点作直线运动, 时刻 $t$ 时位移为 $x(t)$, 瞬时速度为
\[
v(t)=x'(t)=\lim_{\Delta t\to 0}\frac{x(t+\Delta t)-x(t)}{\Delta t}.
\]
相反, 如果知道瞬时速度, 要求从时刻 $\alpha$ 到 $\beta$ 这段时间内的位移, 可以将 $[\alpha,\beta]$ 分成 $n$ 小段, $\alpha=t_{0}<t_{1}<\cdots<t_{n}=\beta$, 从时刻 $t_{i-1}$ 到 $t_{i}$ 有位移 $\Delta x_{i}$, 它用时刻 $\xi_{i}\in[t_{i-1},t_{i}]$ 的速度与时间间隔 $\Delta t_{i}=t_{i}-t_{i-1}$ 的乘积来近似表达, 即 $\Delta x_{i}\approx v(\xi_{i})\Delta t_{i}$, 故 $x(\beta)-x(\alpha)=\sum\limits_{i=1}^{n}\Delta x_{i}\approx\sum\limits_{i=1}^{n}v(\xi_{i})\Delta t_{i}$. 当分点越来越密时, 视它的极限值是所要求的位移.
\[
x(\beta)-x(\alpha)=\lim_{\|T\|\to 0}\sum_{i=1}^{n}v(\xi_{i})\Delta t_{i}\xlongequal{\text{def}}\int_{\alpha}^{\beta}v(t)\,\mathrm{d}t=\int_{\alpha}^{\beta}x'(t)\,\mathrm{d}t.
\]
此外, 设一根直细棒, 其质量为 $m=m(x)$, 线密度为
\[
\rho(x)=m'(x)=\lim_{\Delta x\to 0}\frac{m(x+\Delta x)-m(x)}{\Delta x}.
\]
同样, 如果知道细棒的线密度为 $\rho(x)$, 那么从 $a$ 到 $b$ 的质量为
\[
m(b)-m(a)=\lim_{\|T\|\to 0}\sum_{i=1}^{n}\rho(\xi_{i})\Delta x_{i}=\int_{a}^{b}\rho(x)\,\mathrm{d}x=\int_{a}^{b}m'(x)\,\mathrm{d}x.
\]
上面一个几何问题与两个物理问题启发我们要引进 Riemann 积分的概念.

\begin{definition}\label{def:6.1.1}
设闭区间 $[a,b]$ 内有 $n+1$ 个点, 依次为 $a=x_{0}<x_{1}<\cdots<x_{n}=b$, 它们将 $[a,b]$ 分成 $n$ 个小区间 $\Delta_{i}=[x_{i-1},x_{i}]$, $i=1,2,\cdots,n$. 这些分点或闭子区间, 构成了 $[a,b]$ 的一个分割, 记为
\[
T: a=x_{0}<x_{1}<\cdots<x_{n}=b
\]
或 $T=\{\Delta_{1},\Delta_{2},\cdots,\Delta_{n}\}$. 小区间 $\Delta_{i}=[x_{i-1},x_{i}]$ 的长度为 $\Delta x_{i}=x_{i}-x_{i-1}$, 并称
\[
\|T\|=\max_{1\leqslant i\leqslant n}\{\Delta x_{i}\}
\]
为分割 $T$ 的模.

设 $f:[a,b]\to\mathbb{R}$ 为一元函数. 对于 $[a,b]$ 的上述分割 $T$, 任取点 $\xi_{i}\in\Delta_{i}=[x_{i-1},x_{i}]$, $i=1,2,\cdots,n$, 称
\[
\sum_{i=1}^{n}f(\xi_{i})\Delta x_{i}
\]
为 $f$ 在 $[a,b]$ 上的一个 Riemann 和或积分和.

如果存在 $J\in\mathbb{R}$, s. t. 对 $\forall\varepsilon>0$, $\exists\delta>0$, 对任何分割 $T$, $\|T\|<\delta$, $\forall\xi_{i}\in\Delta_{i}=[x_{i-1},x_{i}]$, $i=1,2,\cdots,n$, 有
\[
\left|\sum_{i=1}^{n}f(\xi_{i})\Delta x_{i}-J\right|<\varepsilon,
\]
则称 $J$ 为 $\sum\limits_{i=1}^{n}f(\xi_{i})\Delta x_{i}$ 当 $\|T\|\to 0$ 时的极限, 称它为 $f$ 在 $[a,b]$ 上的 Riemann 积分或定积分, 并称 $f$ 在区间 $[a,b]$ 上是 Riemann 可积或是可积的. 记上述 Riemann 积分或定积分为
\[
J=\int_{a}^{b}f(x)\,\mathrm{d}x=\lim_{\|T\|\to 0}\sum_{i=1}^{n}f(\xi_{i})\Delta x_{i},
\]
其中 $f(x)$ 称为被积函数, $x$ 为积分变量, $[a,b]$ 称为积分区间, $a$, $b$ 分别称为积分下限与积分上限.
\end{definition}

显然, $\int_{a}^{b}f(x)\,\mathrm{d}x$ 完全由 $f$ 与区间 $[a,b]$ 所确定, 它与积分变量 $x$ 无关. 因此,
\[
\int_{a}^{b}f(x)\,\mathrm{d}x=\int_{a}^{b}f(u)\,\mathrm{d}u=\int_{a}^{b}f(t)\,\mathrm{d}t.
\]

\begin{theorem}[惟一性]\label{thm:6.1.1}
设 $f$ 在 $[a,b]$ 上 Riemann 可积, 则积分值 $J=\int_{a}^{b}f(x)\,\mathrm{d}x$ 是惟一的.
\end{theorem}

\begin{proof}
假设 $J_{1},J_{2}\in\mathbb{R}$ 都为 $\sum\limits_{i=1}^{n}f(\xi_{i})\Delta x_{i}$ 当 $\|T\|\to 0$ 时的极限, 则对 $\forall\varepsilon>0$, $\exists\delta>0$, 对 $[a,b]$ 的任何分割 $T$, $\|T\|<\delta$, $\forall\xi_{i}\in\Delta_{i}=[x_{i-1},x_{i}]$, $i=1,\cdots,n$, 有
\[
\left|\sum_{i=1}^{n}f(\xi_{i})\Delta x_{i}-J_{j}\right|<\frac{\varepsilon}{2},\quad j=1,2.
\]
于是,
\begin{align*}
0&\leqslant|J_{1}-J_{2}|=\left|J_{1}-\sum_{i=1}^{n}f(\xi_{i})\Delta x_{i}+\sum_{i=1}^{n}f(\xi_{i})\Delta x_{i}-J_{2}\right|\\
&\leqslant\left|\sum_{i=1}^{n}f(\xi_{i})\Delta x_{i}-J_{1}\right|+\left|\sum_{i=1}^{n}f(\xi_{i})\Delta x_{i}-J_{2}\right|\\
&<\frac{\varepsilon}{2}+\frac{\varepsilon}{2}=\varepsilon.
\end{align*}
令 $\varepsilon\to 0^{+}$ 得到
\begin{gather*}
0\leqslant|J_{1}-J_{2}|\leqslant 0,\\
|J_{1}-J_{2}|=0,\quad J_{1}=J_{2}.
\end{gather*}
\end{proof}

\begin{theorem}[可积的必要条件]\label{thm:6.1.2}
设 $f$ 在 $[a,b]$ 上 Riemann 可积, 则 $f$ 在 $[a,b]$ 上有界 (因而无界函数不是 Riemann 可积的), 但反之不真.
\end{theorem}

\begin{proof}
因为 $f$ 在 $[a,b]$ 上 Riemann 可积, 故有 $J\in\mathbb{R}$, 取 $\varepsilon=1$, $\exists\delta>0$, 对 $[a,b]$ 的任意分割 $T:a=x_{0}<x_{1}<\cdots<x_{n}=b$, $\|T\|<\delta$, $\forall\xi_{i}\in\Delta_{i}=[x_{i-1},x_{i}]$, $i=1,2,\cdots,n$, 有
\[
\left|\sum_{i=1}^{n}f(\xi_{i})\Delta x_{i}-J\right|<1. \tag*{(*)}
\]
(反证) 假设 $f$ 在 $[a,b]$ 上无界, 则至少有一个区间 $[x_{i_{0}-1},x_{i_{0}}]$, $f$ 在其上无界, 于是 $\exists\xi_{i_{0}}\in[x_{i_{0}-1},x_{i_{0}}]$, s. t.
\[
\left|f(\xi_{i_{0}})\Delta x_{i_{0}}\right|>1+\left|\sum_{i\neq i_{0}}f(\xi_{i})\Delta x_{i}\right|+|J|. \tag*{(**)}
\]
但由 $(*)$ 式知,
\begin{align*}
\left|f(\xi_{i_{0}})\Delta x_{i_{0}}\right|&=\left|\left(\sum_{i=1}^{n}f(\xi_{i})\Delta x_{i}-J\right)-\sum_{i\neq i_{0}}f(\xi_{i})\Delta x_{i}+J\right|\\
&\leqslant\left|\sum_{i=1}^{n}f(\xi_{i})\Delta x_{i}-J\right|+\left|\sum_{i\neq i_{0}}f(\xi_{i})\Delta x_{i}\right|+|J|\\
&<1+\left|\sum_{i\neq i_{0}}f(\xi_{i})\Delta x_{i}\right|+|J|,
\end{align*}
这与 $(**)$ 式相矛盾.
\end{proof}

反例: Dirichlet 函数 $D(x)$ 在 $[0,1]$ 上有界但非 Riemann 可积.

显然, $|D(x)|\leqslant 1$, 故 $D(x)$ 有界. 现证 $D(x)$ 在 $[0,1]$ 上非 Riemann 可积.

(反证) 假设 $D(x)$ 在 $[0,1]$ 上 Riemann 可积, 则 $\exists J\in\mathbb{R}$, 对 $\varepsilon_{0}=\frac{1}{2}$, $\exists\delta>0$, 当 $[a,b]$ 的分割 $T:a=x_{0}<x_{1}<\cdots<x_{n}=b$, $\|T\|<\delta$ 时, 有
\[
\left|\sum_{i=1}^{n}D(\xi_{i})\Delta x_{i}-J\right|<\varepsilon_{0}=\frac{1}{2},\quad\forall\xi_{i}\in\Delta_{i}=[x_{i-1},x_{i}],\ i=1,2,\cdots,n.
\]
特别取 $\eta_{i}\in[x_{i-1},x_{i}]\cap\mathbb{Q}$, $\zeta_{i}\in[x_{i-1},x_{i}]\cap(\mathbb{R}\setminus\mathbb{Q})$, 有
\begin{align*}
1&=|0-1|=\left|\sum_{i=1}^{n}D(\eta_{i})\Delta x_{i}-\sum_{i=1}^{n}D(\zeta_{i})\Delta x_{i}\right|\\
&=\left|\left(\sum_{i=1}^{n}D(\eta_{i})\Delta x_{i}-J\right)-\left(\sum_{i=1}^{n}D(\zeta_{i})\Delta x_{i}-J\right)\right|\\
&\leqslant\left|\sum_{i=1}^{n}D(\eta_{i})\Delta x_{i}-J\right|+\left|\sum_{i=1}^{n}D(\zeta_{i})\Delta x_{i}-J\right|\\
&<\varepsilon_{0}+\varepsilon_{0}=\frac{1}{2}+\frac{1}{2}=1,
\end{align*}
矛盾.

定理~\ref{thm:6.1.2} 只是 Riemann 可积的必要条件, 它只能否定无界函数是 Riemann 可积的, 但对有界函数是否为 Riemann 可积函数就束手无策了. 要判断一个函数是否 Riemann 可积, 原则上可以根据定义直接考察 Riemann 和是否无限接近某个常数 $J$, 但由于那个常数不易预知及 Riemann 和的复杂性, 这种判断是非常困难的. 下面给出一些 Riemann 可积的充要条件, 它只与被积函数 $f$ 本身有关, 而不涉及 Riemann 积分的值.

以下总假定 $f(x)$ 在 $[a,b]$ 上有界, 记 $M=\sup\limits_{x\in[a,b]}f(x)$, $m=\inf\limits_{x\in[a,b]}f(x)$. $T:a=x_{0}<x_{1}<\cdots<x_{n}=b$ 为 $[a,b]$ 的任一分割, $\Delta_{i}=[x_{i-1},x_{i}]$, $i=1,2,\cdots,n$.
\[
M_{i}^{f}=\sup_{x\in\Delta_{i}}f(x),\quad m_{i}^{f}=\inf_{x\in\Delta_{i}}f(x).
\]
在不致混淆时, 记 $M_{i}^{f}$ 为 $M_{i}$, $m_{i}^{f}$ 为 $m_{i}$.

做和 $S(T,f)=\sum\limits_{i=1}^{n}M_{i}\Delta x_{i}$ 及 $s(T,f)=\sum\limits_{i=1}^{n}m_{i}\Delta x_{i}$, 分别称它们为 $f$ 关于 $T$ 的 Darboux 上和与 Darboux 下和, 统称为 Darboux 和. 对 $\forall\xi_{i}\in\Delta_{i}$, $i=1,2,\cdots,n$, 显然, 有
\[
m(b-a)\leqslant\sum_{i=1}^{n}m_{i}\Delta x_{i}\leqslant\sum_{i=1}^{n}f(\xi_{i})\Delta x_{i}\leqslant\sum_{i=1}^{n}M_{i}\Delta x_{i}\leqslant M(b-a),
\]
即
\[
m(b-a)\leqslant\underbrace{s(T,f)}_{\text{Darboux 下和}}\leqslant\underbrace{S(T,f,\xi)}_{\text{Riemann 和}}\leqslant\underbrace{S(T,f)}_{\text{Darboux 上和}}\leqslant M(b-a),
\]
其中 $\xi=(\xi_{1},\xi_{2},\cdots,\xi_{n})$.

设 $\omega_{i}^{f}=M_{i}^{f}-m_{i}^{f}$, 称它为 $f$ 在 $\Delta_{i}$ 上的振幅, 简记为 $\omega_{i}=M_{i}-m_{i}$. 而
\begin{align*}
S(T)-s(T)&=\sum_{i=1}^{n}(M_{i}-m_{i})\Delta x_{i}\\
&=\sum_{i=1}^{n}\omega_{i}\Delta x_{i}\ \left(\text{或记为 }\sum_{T}\omega_{i}\Delta x_{i}\right)
\end{align*}
称为 $f$ 在 $[a,b]$ 上关于分割 $T$ 的振幅和.

\begin{lemma}\label{lem:6.1.1}
对 $[a,b]$ 的同一分割 $T$, 相对于任何 $\xi=(\xi_{1},\cdots,\xi_{n})$, $\xi_{i}\in\Delta_{i}$, 有
\[
S(T,f)=\sup_{\xi}\sum_{i=1}^{n}f(\xi_{i})\Delta x_{i},\quad s(T,f)=\inf_{\xi}\sum_{i=1}^{n}f(\xi_{i})\Delta x_{i}.
\]
\end{lemma}

\begin{proof}
一方面, 有
\[
s(T,f)\leqslant\inf_{\xi}\sum_{i=1}^{n}f(\xi_{i})\Delta x_{i}\leqslant\sup_{\xi}\sum_{i=1}^{n}f(\xi_{i})\Delta x_{i}\leqslant S(T,f).
\]
另一方面, $\forall\varepsilon>0$, 由上确界定义, 可选 $\xi_{i}\in\Delta_{i}$, s. t.
\[
f(\xi_{i})>M_{i}-\frac{\varepsilon}{b-a}.
\]
于是, 有
\begin{align*}
\sum_{i=1}^{n}f(\xi_{i})\Delta x_{i}&>\sum_{i=1}^{n}\left(M_{i}-\frac{\varepsilon}{b-a}\right)\Delta x_{i}\\
&=\sum_{i=1}^{n}M_{i}\Delta x_{i}-\frac{\varepsilon}{b-a}\sum_{i=1}^{n}\Delta x_{i}\\
&=S(T,f)-\varepsilon,
\end{align*}
\[
\sup_{\xi}\sum_{i=1}^{n}f(\xi_{i})\Delta x_{i}\geqslant S(T,f)-\varepsilon.
\]
令 $\varepsilon\to 0^{+}$ 得到
\[
\sup_{\xi}\sum_{i=1}^{n}f(\xi_{i})\Delta x_{i}\geqslant S(T,f).
\]
所以,
\[
S(T,f)=\sup_{\xi}\sum_{i=1}^{n}f(\xi_{i})\Delta x_{i}.
\]
同理, 或用 $-f$ 代替 $f$ 得到
\[
s(T,f)=\inf_{\xi}\sum_{i=1}^{n}f(\xi_{i})\Delta x_{i}.
\]
\end{proof}

\begin{lemma}\label{lem:6.1.2}
设 $T'$ 为分割 $T$ 添加 $p$ 个新分点所得的分割, 则有
\[
S(T)\geqslant S(T')\geqslant S(T)-(M-m)p\|T\|,
\]
\[
s(T)\leqslant s(T')\leqslant s(T)+(M-m)p\|T\|,
\]
其中 $s(T)=s(T,f)$, $S(T)=S(T,f)$. 这表明增加新分点后, 上和不变, 下和不减.
\end{lemma}

\begin{proof}
当 $p=1$ 时, 在 $T$ 上添加一个新分点, 它必落在 $T$ 的某个小区间 $\Delta_{k}$ 内, 而且将 $\Delta_{k}$ 分为两个小区间, 记为 $\Delta_{k}'$ 与 $\Delta_{k}''$. 但 $T$ 的其他小区间 $\Delta_{i}$ ($i\neq k$) 仍旧是新分割 $T_{1}$ 所属的小区间. 因此, 比较 $S(T)$ 与 $S(T_{1})$ 的各个被加项, 它们之间的差别仅仅是 $S(T)$ 中的 $M_{k}\Delta x_{k}$ 一项换成 $S(T_{1})$ 中的 $M_{k}'\Delta x_{k}'+M_{k}''\Delta x_{k}''$ 两项 (这里 $M_{k}'$ 与 $M_{k}''$ 分别是 $f$ 在 $\Delta_{k}'$ 与 $\Delta_{k}''$ 上的上确界). 显然, $m\leqslant M_{k}'$ (或 $M_{k}''$) $\leqslant M_{k}\leqslant M$, 所以,
\begin{align*}
0&\leqslant S(T)-S(T_{1})=(M_{k}-M_{k}')\Delta x_{k}'+(M_{k}-M_{k}'')\Delta x_{k}''\\
&\leqslant(M-m)\Delta x_{k}'+(M-m)\Delta x_{k}''=(M-m)\Delta x_{k}\\
&\leqslant(M-m)\|T\|.
\end{align*}
即
\[
S(T)\geqslant S(T_{1})\geqslant S(T)-(M-m)\|T\|.
\]
一般地, 对 $T_{i}$ 增加一个新分点得到 $T_{i+1}$, 就有
\[
0\leqslant S(T_{i})-S(T_{i+1})\leqslant(M-m)\|T_{i}\|,\quad i=0,1,\cdots,p-1,
\]
其中 $T_{0}=T$, $T_{p}=T'$. 将这些不等式对 $i$ 依次相加得到
\begin{gather*}
0\leqslant S(T)-S(T')\leqslant(M-m)\sum_{i=0}^{p-1}\|T_{i}\|\\
\leqslant(M-m)p\|T\|,
\end{gather*}
即
\[
S(T)\geqslant S(T')\geqslant S(T)-(M-m)p\|T\|.
\]
同理有
\[
s(T)\leqslant s(T')\leqslant s(T)+(M-m)p\|T\|.
\]
\end{proof}

\begin{lemma}\label{lem:6.1.3}
设 $T'$ 与 $T''$ 为 $[a,b]$ 的任意两个分割, $T'+T''$ 表示将 $T'$ 与 $T''$ 的所有分点合并得到的分割 (重复的分点只取一次), 则
\begin{gather*}
S(T'+T'')\leqslant S(T'),\quad s(T'+T'')\geqslant s(T');\\
S(T'+T'')\leqslant S(T''),\quad s(T'+T'')\geqslant s(T'').
\end{gather*}
\end{lemma}

\begin{proof}
因为 $T'+T''$ 既可看作 $T'$ 添加新分点后得到的分割, 也可看作 $T''$ 添加新分点后得到的分割, 所以根据引理~\ref{lem:6.1.2} 立即推得上述不等式成立.
\end{proof}

\begin{lemma}\label{lem:6.1.4}
对任意两个分割 $T'$ 与 $T''$, 有
\[
s(T')\leqslant S(T'').
\]
\end{lemma}

\begin{proof}
由引理~\ref{lem:6.1.1} 与引理~\ref{lem:6.1.3} 便有
\[
s(T')\leqslant s(T'+T'')\leqslant S(T'+T'')\leqslant S(T'').
\]
\end{proof}

由此立即有下面引理.

\begin{lemma}\label{lem:6.1.5}
\[
m(b-a)\leqslant\sup_{T'}s(T')\leqslant\inf_{T''}S(T'')\leqslant M(b-a).
\]
\end{lemma}

\begin{proof}
由
\[
m(b-a)\leqslant s(T')\leqslant S(T'')\leqslant M(b-a)
\]
知
\[
m(b-a)\leqslant\sup_{T'}s(T')\leqslant\inf_{T''}S(T'')\leqslant M(b-a),
\]
故 $m(b-a)\leqslant s\leqslant S\leqslant M(b-a)$, 其中 $S=\inf\limits_{T}S(T)$ 与 $s=\sup\limits_{T}s(T)$ 分别称为 $f(x)$ 在 $[a,b]$ 上的上积分与下积分.
\end{proof}

\begin{theorem}[Darboux 定理]\label{thm:6.1.3}
\[
\lim_{\|T\|\to 0}S(T)=S,\qquad \lim_{\|T\|\to 0}s(T)=s.
\]
\end{theorem}

\begin{proof}
$\forall\varepsilon>0$, 由 $S$ 的定义知, 必有某个分割 $T'$, s. t.
\[
S(T')<S+\frac{\varepsilon}{2}.
\]
设 $T'$ 由 $p$ 个分点构成, 对于任意另一个分割 $T$, $T+T'$ 至多比 $T$ 多 $p$ 个分点, 由引理~\ref{lem:6.1.2} 与引理~\ref{lem:6.1.3} 得到
\[
S(T)-(M-m)p\|T\|\leqslant S(T+T')\leqslant S(T').
\]
当 $\|T\|<\delta=\frac{\varepsilon}{2(M-m+1)p}$ 时, 有
\begin{align*}
S&\leqslant S(T)\leqslant S(T')+(M-m)p\|T\|\\
&<S(T')+(M-m)p\cdot\frac{\varepsilon}{2(M-m+1)p}\\
&<S(T')+\frac{\varepsilon}{2}<\left(S+\frac{\varepsilon}{2}\right)+\frac{\varepsilon}{2}=S+\varepsilon.
\end{align*}
这就证明了
\[
\lim_{\|T\|\to 0}S(T)=S.
\]
同理, 或用 $-f$ 代替 $f$ 得到
\[
\lim_{\|T\|\to 0}s(T)=s.
\]
\end{proof}

\begin{definition}\label{def:6.1.2}
设 $A\subseteq\mathbb{R}$, 如果对 $\forall\varepsilon>0$, 存在至多可数个 (有限个或可数个) 开区间 $\{I_{i}\mid i=1,2,\cdots\}$, s. t.
\[
A\subset\bigcup_{i}I_{i},\quad\text{且}\ \sum_{i}|I_{i}|<\varepsilon\quad(|I_{i}|\ \text{为区间}\ I_{i}\ \text{的长度}),
\]
则称 $A$ 为 $\mathbb{R}$ 中的 Lebesgue 零测集, 简称为零测集, 也称 $A$ 的 (Lebesgue) 测度为零.
\end{definition}

\begin{lemma}\label{lem:6.1.6}
(1) 有限集为零测集;

(2) 至多可数集为零测集;

(3) 零测集 $A$ 的任何子集 $B$ 为零测集;

(4) 至多可数个零测集 $A_{i}$ ($i=1,2,\cdots$) 的并集仍为零测集.
\end{lemma}

\begin{proof}
(1) 设有限集 $A=\{a_{1},a_{2},\cdots,a_{k}\}$. 对 $\forall\varepsilon>0$, 令 $I_{i}=\left(a_{i}-\frac{\varepsilon}{4k},a_{i}+\frac{\varepsilon}{4k}\right)$, $i=1,2,\cdots,k$. 则
\[
A=\bigcup_{i=1}^{k}\{a_{i}\}\subset\bigcup_{i=1}^{k}\left(a_{i}-\frac{\varepsilon}{4k},a_{i}+\frac{\varepsilon}{4k}\right),
\]
且
\[
\sum_{i=1}^{k}|I_{i}|=\sum_{i=1}^{k}\left[\left(a_{i}+\frac{\varepsilon}{4k}\right)-\left(a_{i}-\frac{\varepsilon}{4k}\right)\right]=\sum_{i=1}^{k}\frac{\varepsilon}{2k}=\frac{\varepsilon}{2}<\varepsilon,
\]
因此, $A$ 为零测集.

(2) 设可数集 $A=\{a_{1},a_{2},\cdots,a_{k},\cdots\}$. 对 $\forall\varepsilon>0$, 令 $I_{i}=\left(a_{i}-\frac{\varepsilon}{2^{i+2}},a_{i}+\frac{\varepsilon}{2^{i+2}}\right)$, $i=1,2,\cdots,k,\cdots$, 则
\[
A=\bigcup_{i=1}^{\infty}\{a_{i}\}\subset\bigcup_{i=1}^{\infty}\left(a_{i}-\frac{\varepsilon}{2^{i+2}},a_{i}+\frac{\varepsilon}{2^{i+2}}\right),
\]
且
\begin{align*}
\sum_{i=1}^{\infty}|I_{i}|&=\sum_{i=1}^{\infty}\left[\left(a_{i}+\frac{\varepsilon}{2^{i+2}}\right)-\left(a_{i}-\frac{\varepsilon}{2^{i+2}}\right)\right]=\sum_{i=1}^{\infty}\frac{\varepsilon}{2^{i+1}}\\
&=\frac{\varepsilon}{2}<\varepsilon.
\end{align*}
因此, $A$ 为零测集.

(3) 因为 $A$ 为零测集, 故对任何 $\varepsilon>0$, 存在开区间 $I_{i}$ ($i=1,2,\cdots$), 使得
\[
A\subset\bigcup_{i}I_{i},\quad\sum_{i}|I_{i}|<\varepsilon.
\]
因此,
\[
B\subset A\subset\bigcup_{i}I_{i},\quad\sum_{i}|I_{i}|<\varepsilon.
\]
由此知, $B$ 为零测集.

(4) 因为每个 $A_{i}$ 为零测集, 则存在开区间 $I_{j}^{i}$ ($j=1,2,\cdots$), 使得
\[
A_{i}\subset\bigcup_{j}I_{j}^{i},\quad\text{且}\ \sum_{j}|I_{j}^{i}|<\frac{\varepsilon}{2^{i}},\quad i=1,2,\cdots.
\]
于是, $\{I_{j}^{i}\mid i=1,2,\cdots;j=1,2,\cdots\}$ 为可数个开区间\linebreak[4](依 $i+j$ 从小到大排列为 $I_{1}^{1},I_{2}^{1},\allowbreak I_{1}^{2},I_{3}^{1},\allowbreak I_{1}^{3},I_{2}^{2},\cdots$), 使得
\[
\bigcup_{i}A_{i}\subset\bigcup_{i}\left(\bigcup_{j}I_{j}^{i}\right)=\bigcup_{i,j}I_{j}^{i},
\]
且
\[
\sum_{i,j}|I_{j}^{i}|=\sum_{i}\left(\sum_{j}|I_{j}^{i}|\right)<\sum_{i}\frac{\varepsilon}{2^{i}}\leqslant\varepsilon,
\]
所以, $\bigcup_{i}A_{i}$ 为零测集.
\end{proof}

\begin{theorem}[Riemann 可积的充要条件]\label{thm:6.1.4}
设函数 $f$ 在 $[a,b]$ 上有界, 即 $|f(x)|\leqslant M$, $\forall x\in[a,b]$, 则下面结论等价.

(1) $f$ 在 $[a,b]$ 上 Riemann 可积.

(2) $f$ 在 $[a,b]$ 上的上积分与下积分相等, 即 $S=s$.

(3) 对任何 $\varepsilon>0$, 存在 $[a,b]$ 的某个分割 $T$, 使得
\[
S(T)-s(T)=\sum_{i=1}^{n}\omega_{i}\Delta x_{i}<\varepsilon.
\]

(4) 存在 $[a,b]$ 的分割串 $T_{m}$, $m=1,2,\cdots$, 使得
\[
\lim_{m\to+\infty}\left[S(T_{m})-s(T_{m})\right]=0.
\]

(5) $\lim\limits_{\|T\|\to 0}\left[S(T)-s(T)\right]=0$.

(6) 对任何 $\varepsilon>0$, 存在 $[a,b]$ 的某个分割 $T$, 使对任何 $\boldsymbol{\xi}=(\xi_{1},\xi_{2},\cdots,\xi_{n})$, $\xi_{i}\in[x_{i-1},x_{i}]=\Delta_{i}$, $i=1,2,\cdots,n$, 有
\[
\left|S(T,f,\boldsymbol{\xi})-J\right|<\varepsilon.
\]

(7) 对任何 $\varepsilon>0$, $\eta>0$, 存在 $[a,b]$ 的某个分割 $T$, 使得
\[
\sum_{\omega_{i}\geqslant\varepsilon}\Delta x_{i}<\eta.
\]

(8) (Lebesgue) $f$ 在 $[a,b]$ 上几乎处处 (a. e. 即 almost everywhere) 连续, 即 $f$ 的不连续点集 $D_{\text{不}}$ 为零测集, 记作
\[
\operatorname{meas}D_{\text{不}}=0.
\]
\end{theorem}

\begin{proof}
(1) $\Rightarrow$ (2) 由 $f$ 在 $[a,b]$ 上 Riemann 可积, 其积分值为 $J$, 故 $\forall\varepsilon>0$, $\exists\delta>0$, 当 $\|T\|<\delta$ 时, 有
\[
J-\frac{\varepsilon}{4}<\sum_{i=1}^{n}f(\xi_{i})\Delta x_{i}<J+\frac{\varepsilon}{4},
\]
于是,
\begin{align*}
J-\frac{\varepsilon}{4}&\leqslant s(T)=\sum_{i=1}^{n}\inf_{x\in[x_{i-1},x_{i}]}f(x)\Delta x_{i}\\
&\leqslant\sum_{i=1}^{n}\sup_{x\in[x_{i-1},x_{i}]}f(x)\Delta x_{i}=S(T)\leqslant J+\frac{\varepsilon}{4},
\end{align*}
\[
J-\frac{\varepsilon}{4}\leqslant s=\sup_{T}s(T)\leqslant\inf_{T}S(T)=S\leqslant J+\frac{\varepsilon}{4}.
\]
令 $\varepsilon\to 0^{+}$ 得到
\[
J\leqslant s\leqslant S\leqslant J,
\]
\[
S=s=J.
\]

(2) $\Rightarrow$ (1) 设 $J'=S=s$. 由 Darboux 定理得
\[
\lim_{\|T\|\to 0}S(T)=\lim_{\|T\|\to 0}s(T)=J'.
\]
则 $\forall\varepsilon>0$, $\exists\delta>0$, 当 $\|T\|<\delta$ 时, 满足
\[
J'-\varepsilon<s(T)\leqslant\sum_{i=1}^{n}f(\xi_{i})\Delta x_{i}\leqslant S(T)<J'+\varepsilon.
\]
从而
\[
\lim_{\|T\|\to 0}\sum_{i=1}^{n}f(\xi_{i})\Delta x_{i}=J',
\]
即 $f$ 在 $[a,b]$ 上 Riemann 可积, 且 $\int_{a}^{b}f(x)\,\mathrm{d}x=J'$.

(2) $\Rightarrow$ (3) 设 $J=S=s$. $\forall\varepsilon>0$, 由 Darboux 定理知, $\exists\delta>0$, 当 $\|T\|<\delta$ 时,
\[
J-\frac{\varepsilon}{4}\leqslant s(T)\leqslant S(T)\leqslant J+\frac{\varepsilon}{4},
\]
则
\[
\sum_{i=1}^{n}\omega_{i}\Delta x_{i}=S(T)-s(T)\leqslant\left(J+\frac{\varepsilon}{4}\right)-\left(J-\frac{\varepsilon}{4}\right)=\frac{\varepsilon}{2}<\varepsilon.
\]

(3) $\Rightarrow$ (2) $\forall\varepsilon>0$, 由 (3), 存在 $[a,b]$ 的某个分割 $T$, 使得
\[
S(T)-s(T)<\varepsilon.
\]
由
\[
s(T)\leqslant s\leqslant S\leqslant S(T)
\]
得到
\[
0\leqslant S-s\leqslant S(T)-s(T)<\varepsilon.
\]
令 $\varepsilon\to 0^{+}$ 有
\[
0\leqslant S-s\leqslant 0,
\]
即
\[
S=s.
\]

(3) $\Rightarrow$ (4) 对任何 $m\in\mathbb{N}$, 由 (3), 必存在 $[a,b]$ 的某个分割 $T_{m}$, 使得
\[
0\leqslant S(T_{m})-s(T_{m})<\frac{1}{m}.
\]
由数列极限的定义或夹逼定理知,
\[
\lim_{m\to+\infty}\left[S(T_{m})-s(T_{m})\right]=0.
\]

(4) $\Rightarrow$ (3) 设 $[a,b]$ 的分割 $T_{m}$, $m=1,2,\cdots$, 使得
\[
\lim_{m\to+\infty}\left[S(T_{m})-s(T_{m})\right]=0,
\]
则 $\forall\varepsilon>0$, $\exists N\in\mathbb{N}$, s. t. 当 $m>N$ 时, 有
\[
S(T_{m})-s(T_{m})<\varepsilon.
\]
令 $T=T_{N+1}$, 则 $T$ 即为所求.

(1) $\Rightarrow$ (5) 由 (1), $f$ 在 $[a,b]$ 上 Riemann 可积, 故 $\forall\varepsilon>0$, $\exists\delta>0$, 当 $\|T\|<\delta$ 时有
\[
J-\frac{\varepsilon}{4}<\sum_{i=1}^{n}f(\xi_{i})\Delta x_{i}=S(T,f,\xi)<J+\frac{\varepsilon}{4}.
\]
由此得到
\[
J-\frac{\varepsilon}{4}\leqslant s(T,f)\leqslant S(T,f)\leqslant J+\frac{\varepsilon}{4},
\]
\[
S(T,f)-s(T,f)\leqslant\left(J+\frac{\varepsilon}{4}\right)-\left(J-\frac{\varepsilon}{4}\right)=\frac{\varepsilon}{2}<\varepsilon,
\]
从而
\[
\lim_{\|T\|\to 0}\left[S(T,f)-s(T,f)\right]=0.
\]

(5) $\Rightarrow$ (3) 由 (5),
\[
\lim_{\|T\|\to 0}\left[S(T,f)-s(T,f)\right]=0,
\]
故 $\forall\varepsilon>0$, $\exists\delta>0$, 当 $\|T\|<\delta$ 时, $S(T)-s(T)<\varepsilon$.

(1) $\Rightarrow$ (6) 显然.

(6) $\Rightarrow$ (3) 由 (6), 对任何 $\varepsilon>0$, 存在 $[a,b]$ 的某个分割 $T$, 对任何 $\xi_{i}\in[x_{i-1},x_{i}]=\Delta_{i}$, $i=1,2,\cdots,n$, 有
\[
|S(T,f,\xi)-J|<\frac{\varepsilon}{4},
\]
故
\begin{align*}
|S(T,f)-J|&<\frac{\varepsilon}{4},\qquad |s(T,f)-J|<\frac{\varepsilon}{4},\\
|S(T,f)-s(T,f)|&\leqslant|S(T,f)-J|+|s(T,f)-J|\\
&\leqslant\frac{\varepsilon}{4}+\frac{\varepsilon}{4}=\frac{\varepsilon}{2}<\varepsilon.
\end{align*}

(3) $\Rightarrow$ (7) $\forall\varepsilon>0$, $\forall\eta>0$, 由 (3), 存在 $[a,b]$ 的某个分割 $T$, 使得
\[
\sum_{i=1}^{n}\omega_{i}\Delta x_{i}<\sigma=\varepsilon\eta.
\]
于是, 便有
\[
\varepsilon\sum_{\omega_{i}\geqslant\varepsilon}\Delta x_{i}\leqslant\sum_{\omega_{i}\geqslant\varepsilon}\omega_{i}\Delta x_{i}\leqslant\sum_{i}\omega_{i}\Delta x_{i}<\varepsilon\eta.
\]
两边约去 $\varepsilon$ 得到
\[
\sum_{\omega_{i}\geqslant\varepsilon}\Delta x_{i}<\eta.
\]

(7) $\Rightarrow$ (3) $\forall\varepsilon>0$, 取 $\varepsilon'=\frac{\varepsilon}{2(b-a)}>0$, $\eta'=\frac{\varepsilon}{2(M-m+1)}>0$. 由 (7), 存在 $[a,b]$ 的某个分割 $T$, 使得
\[
\sum_{\omega_{i}\geqslant\varepsilon'}\Delta x_{i}<\eta'.
\]
于是,
\begin{align*}
S(T)-s(T)&=\sum_{i}\omega_{i}\Delta x_{i}=\sum_{\omega_{i}\geqslant\varepsilon'}\omega_{i}\Delta x_{i}+\sum_{\omega_{i}<\varepsilon'}\omega_{i}\Delta x_{i}\\
&\leqslant(M-m)\sum_{\omega_{i}\geqslant\varepsilon'}\Delta x_{i}+\varepsilon'\sum_{\omega_{i}<\varepsilon'}\Delta x_{i}\\
&\leqslant(M-m)\eta'+\varepsilon'(b-a)\\
&\leqslant(M-m)\frac{\varepsilon}{2(M-m+1)}+\frac{\varepsilon}{2(b-a)}(b-a)\\
&<\frac{\varepsilon}{2}+\frac{\varepsilon}{2}=\varepsilon.
\end{align*}

(3) $\Rightarrow$ (8) 由 (3), 对任何 $\varepsilon>0$, 存在 $[a,b]$ 的某个分割 $T_{\varepsilon}:a=x_{0}<x_{1}<\cdots<x_{n}=b$, 使得
\[
S(T_{\varepsilon},f)-s(T_{\varepsilon},f)<\frac{\delta}{2}\varepsilon,
\]
其中 $\delta>0$ 为任何正数. 令
\[
D_{\delta}=\left\{x\in[a,b]\ \middle|\ \text{振幅}\ \omega_{f}(x)=\lim_{r\to 0^{+}}\left[\sup f(x-r,x+r)-\inf f(x-r,x+r)\right]\geqslant\delta\right\},
\]
则
\[
\sum_{D_{\delta}\cap(x_{i-1},x_{i})\neq\varnothing}\Delta x_{i}<\frac{\varepsilon}{2}.
\]
若不然,
\begin{align*}
\frac{\delta}{2}\varepsilon&>S(T_{\varepsilon},f)-s(T_{\varepsilon},f)\\
&\geqslant\sum_{D_{\delta}\cap(x_{i-1},x_{i})\neq\varnothing}\left[\sup f(x_{i-1},x_{i})-\inf f(x_{i-1},x_{i})\right]\Delta x_{i}\\
&\geqslant\sum_{D_{\delta}\cap(x_{i-1},x_{i})\neq\varnothing}\delta\Delta x_{i}\geqslant\delta\cdot\frac{\varepsilon}{2}=\frac{\delta}{2}\varepsilon,\ \text{矛盾}.
\end{align*}
显然,
\[
D_{\delta}\subset\bigcup_{D_{\delta}\cap(x_{i-1},x_{i})\neq\varnothing}(x_{i-1},x_{i})\cup\bigcup_{i=0}^{n}\left(x_{i}-\frac{\varepsilon}{4(n+1)},x_{i}+\frac{\varepsilon}{4(n+1)}\right),
\]
且
\[
\sum_{D_{\delta}\cap(x_{i-1},x_{i})\neq\varnothing}\Delta x_{i}+\frac{\varepsilon}{4(n+1)}\cdot 2(n+1)<\frac{\varepsilon}{2}+\frac{\varepsilon}{2}=\varepsilon,
\]
所以, $D_{\delta}$ 为零测集. 由此推得 $D_{\text{不}}=\bigcup_{n}D_{\frac{1}{n}}$ 也为零测集, 即 $f$ 在 $[a,b]$ 上几乎处处连续.

(8) $\Rightarrow$ (3) 由 (8), $D_{\text{不}}$ 为零测集, 故对任何 $\varepsilon>0$, 存在开区间集 $\left\{(\alpha_{i},\beta_{i})\mid i=1,2,\cdots\right\}$, 使得
\[
D_{\text{不}}\subset\bigcup_{i}(\alpha_{i},\beta_{i}),\quad\text{且}\ \sum_{i}(\beta_{i}-\alpha_{i})<\frac{\varepsilon}{4M+1}.
\]
由此, $\forall x\in[a,b]-\bigcup_{i}(\alpha_{i},\beta_{i})$, $f$ 在 $x$ 连续. 于是, 存在含 $x$ 的开区间 $I_{x}$, 当 $u\in I_{x}\cap[a,b]$ 时,
\[
|f(u)-f(x)|<\frac{\varepsilon}{2(b-a)}.
\]
显然, $\mathcal{I}=\left\{(\alpha_{i},\beta_{i}),I_{x}\mid i=1,2,\cdots;x\in[a,b]-\bigcup_{i}(\alpha_{i},\beta_{i})\right\}$\linebreak[4]为紧致集 $[a,b]$ 的一个开覆盖, 根据 Heine-Borel 有限覆盖定理, 存在 $\mathcal{I}$ 的有限子集 $\mathcal{I}'=\left\{(\alpha_{i_{k}},\beta_{i_{k}}),\allowbreak I_{x_{l}}\mid k=1,2,\cdots,m;l=1,2,\cdots,n\right\}$ 覆盖了 $[a,b]$. 再根据 Lebesgue 数定理, 存在 Lebesgue 数 $\lambda=\lambda(\mathcal{I}')>0$. 当 $[a,b]$ 的分割 $T:y_{0}<y_{1}<\cdots<y_{l}=b$ 满足 $\|T\|<\lambda$ 时, 必有
\[
[y_{i-1},y_{i}]\subset(\alpha_{i_{k}},\beta_{i_{k}})\ \text{或}\ I_{x_{l}}.
\]
于是,
\begin{align*}
0&\leqslant S(T,f)-s(T,f)\\
&=\sum_{i=1}^{l}\left[\sup f([y_{i-1},y_{i}])-\inf f([y_{i-1},y_{i}])\right]\Delta y_{i}\\
&\leqslant\sum_{1}\left[\sup f([y_{i-1},y_{i}])-\inf f([y_{i-1},y_{i}])\right]\Delta y_{i}\\
&\quad+\sum_{2}\left[\sup f([y_{i-1},y_{i}])-\inf f([y_{i-1},y_{i}])\right]\Delta y_{i}\\
&\leqslant 2M\sum_{1}\Delta y_{i}+\frac{\varepsilon}{2(b-a)}\sum_{2}\Delta y_{i}\\
&\leqslant 2M\sum_{i}(\beta_{i}-\alpha_{i})+\frac{\varepsilon}{2(b-a)}(b-a)\\
&<2M\cdot\frac{\varepsilon}{4M+1}+\frac{\varepsilon}{2}<\varepsilon.
\end{align*}
\end{proof}

\begin{remark}\label{remark:6.1.1}
定理~\ref{thm:6.1.4} 中的 (1)$\Leftrightarrow$(8) 可参阅参考文献 [2] 286 页定理 3.5.5.
\end{remark}

\begin{example}\label{ex:6.1.1}
Dirichlet 函数
\[
D(x)=\begin{cases}1,&x\in\mathbb{Q},\\0,&x\in\mathbb{R}-\mathbb{Q}\end{cases}
\]
在 $[0,1]$ 上非 Riemann 可积.
\end{example}

\begin{solution}[证法 1]
参阅定理~\ref{thm:6.1.2}.
\end{solution}

\begin{solution}[证法 2]
因为
\[
S(T,f)=\sum_{i}\sup D([x_{i-1},x_{i}])\Delta x_{i}=\sum_{i}1\cdot\Delta x_{i}=1,
\]
\[
s(T,f)=\sum_{i}\inf D([x_{i-1},x_{i}])\Delta x_{i}=\sum_{i}0\cdot\Delta x_{i}=0,
\]
\[
S=1\neq 0=s,
\]
所以, 根据定理~\ref{thm:6.1.4}(2) 知, $D(x)$ 在 $[0,1]$ 上不是 Riemann 可积的.
\end{solution}

\begin{solution}[证法 3]
由 $S(T,f)-s(T,f)=1-0=1\not<\frac{1}{2}=\varepsilon_{0}$ 及定理~\ref{thm:6.1.4}(3) 知, $D(x)$ 在 $[0,1]$ 上不是 Riemann 可积的.
\end{solution}

\begin{solution}[证法 4]
因为 $D_{\text{不}}=[0,1]$ 不为零测集, 根据定理~\ref{thm:6.1.4}(8) 知, $D(x)$ 在 $[0,1]$ 上不是 Riemann 可积的.
\end{solution}

\begin{example}\label{ex:6.1.2}
Riemann 函数
\[
R(x)=\begin{cases}\dfrac{1}{q},&x=\dfrac{p}{q},\ p\ \text{与}\ q\ \text{互质 (互素)},\\[2mm]1,&x=0,1,\\[1mm]0,&x\ \text{为}\ (0,1)\ \text{中的无理数}\end{cases}
\]
在 $[0,1]$ 上 Riemann 可积, 且
\[
\int_{0}^{1}R(x)\,\mathrm{d}x=0.
\]
\end{example}

\begin{solution}[证法 1]
由例 2.4.6 知, $R(x)$ 的不连续点全体为 $D_{\text{不}}=\mathbb{Q}\cap[0,1]$, 它为可数集, 因而为零测集. 根据定理~\ref{thm:6.1.4}(8) 知, $R(x)$ 在 $[0,1]$ 上 Riemann 可积, 且
\begin{align*}
\int_{0}^{1}R(x)\,\mathrm{d}x&=s=\sup_{T}s(T,f)=\sup_{T}\sum_{i=1}^{n}\inf_{\xi}R(\xi_{i})\Delta x_{i}\\
&=\sup_{T}\sum_{i=1}^{n}0\cdot\Delta x_{i}=0.
\end{align*}
\end{solution}

\begin{solution}[证法 2]
对任何 $\varepsilon>0$, 在 $[0,1]$ 内使 $\frac{1}{q}\geqslant\frac{\varepsilon}{2}$ 的有理点 $\frac{p}{q}$ 只有有限个, 设它们为 $r_{1},r_{2},\cdots,r_{k}$. 现对 $[0,1]$ 作分割 $T=\{\Delta_{1},\Delta_{2},\cdots,\Delta_{n}\}$, 使得 $\|T\|<\frac{\varepsilon}{4k+1}$, 并将 $T$ 中所有小区间分为 $\{\Delta_{i}'\mid i=1,2,\cdots,m\}$ 与 $\{\Delta_{i}''\mid i=1,2,\cdots,n-m\}$ 两类, 其中 $\{\Delta_{i}'\}$ 为含 $\{r_{i}\mid i=1,2,\cdots,k\}$ 中点的所有小区间, 这类小区间的个数 $m\leqslant 2k$ (当所有 $r_{i}$ 恰好都是 $T$ 的分点时才有 $m=2k$); 而 $\{\Delta_{i}''\}$ 为 $T$ 中所有其余不含 $\{r_{i}\}$ 中点的小区间. 于是,
\begin{align*}
0&\leqslant\sum_{i=1}^{n}R(\xi_{i})\Delta x_{i}\leqslant\sum_{i=1}^{n}M_{i}\Delta x_{i}\leqslant\sum_{i=1}^{m}1\cdot\Delta x_{i}'+\sum_{i=1}^{n-m}\frac{\varepsilon}{2}\Delta x_{i}''\\
&\leqslant 2k\|T\|+\frac{\varepsilon}{2}\cdot 1<2k\frac{\varepsilon}{4k+1}+\frac{\varepsilon}{2}<\varepsilon,
\end{align*}
所以, $R(x)$ 在 $[0,1]$ 上 Riemann 可积, 且
\[
\int_{0}^{1}R(x)\,\mathrm{d}x=0.
\]
\end{solution}

\begin{solution}[证法 3]
由证法 2, 有
\begin{align*}
\sum_{i=1}^{n}\omega_{i}\Delta x_{i}&=\sum_{i=1}^{m}\omega_{i}'\Delta x_{i}'+\sum_{i=1}^{n-m}\omega_{i}''\Delta x_{i}''\\
&\leqslant\sum_{i=1}^{m}1\cdot\Delta x_{i}'+\sum_{i=1}^{n-m}\frac{\varepsilon}{2}\Delta x_{i}''\\
&\leqslant 2k\cdot\|T\|+\frac{\varepsilon}{2}\cdot 1\\
&\leqslant 2k\cdot\frac{\varepsilon}{4k+1}+\frac{\varepsilon}{2}<\varepsilon,
\end{align*}
根据定理~\ref{thm:6.1.4}(3) 知, $R(x)$ 在 $[0,1]$ 上 Riemann 可积. 再由证法 1 有, $\int_{0}^{1}R(x)\,\mathrm{d}x=0$.
\end{solution}

\begin{solution}[证法 4]
$\forall\varepsilon>0$, $\forall\eta>0$, 由于在 $[0,1]$ 中满足 $\frac{1}{q}\geqslant\varepsilon$ (即 $q\leqslant\frac{1}{\varepsilon}$) 的有理点 $\frac{p}{q}$ 只有有限个, 设为 $k$ 个, 因此, 含这类点的小区间至多 $2k$ 个, 在其中 $\omega_{i}\geqslant\varepsilon$. 当 $\|T\|<\frac{\eta}{2k+1}$ 时, 就能保证这些小区间的总长度为
\[
\sum_{\omega_{i}\geqslant\varepsilon}\Delta x_{i}\leqslant 2k\cdot\|T\|\leqslant 2k\cdot\frac{\eta}{2k+1}<\eta.
\]
根据定理~\ref{thm:6.1.4}(7), $R(x)$ 在 $[0,1]$ 上 Riemann 可积. 再由证法 1 有, $\int_{0}^{1}R(x)\,\mathrm{d}x=0$.
\end{solution}

我们用 4 种方法证明 Riemann 函数 $R(x)$ 在 $[0,1]$ 上 Riemann 可积, 这可使读者思路开阔, 此外, 还可看出证法 2 的方法最好.

\begin{example}\label{ex:6.1.3}
设 $f$ 为 $[a,b]$ 上的连续函数, 则 $f$ 在 $[a,b]$ 上 Riemann 可积.
\end{example}

\begin{solution}[证法 1]
显然 $D_{\text{不}}=\varnothing$ 为零测集, 根据定理~\ref{thm:6.1.4}(8), $f$ 在 $[a,b]$ 上 Riemann 可积.
\end{solution}

\begin{solution}[证法 2]
由于 $f$ 在闭区间 $[a,b]$ 上连续, 则 $f$ 在 $[a,b]$ 上一致连续. 即对 $\forall\varepsilon>0$, $\exists\delta>0$, 当 $x',x''\in[a,b]$ 且 $|x'-x''|<\delta$ 时, 有
\[
|f(x')-f(x'')|<\frac{\varepsilon}{2(b-a)}.
\]
所以, 当 $[a,b]$ 的分割 $T:a=x_{0}<x_{1}<\cdots<x_{n}=b$ 满足 $\|T\|<\delta$ 时,
\begin{align*}
\omega_{i}&=M_{i}-m_{i}=\sup f([x_{i-1},x_{i}])-\inf f([x_{i-1},x_{i}])\\
&=\sup_{x',x''\in\Delta_{i}}|f(x')-f(x'')|\leqslant\frac{\varepsilon}{2(b-a)}.
\end{align*}
从而导致 (见图~\ref{fig:6.1.4})
\[
\sum_{i=1}^{n}\omega_{i}\Delta x_{i}\leqslant\frac{\varepsilon}{2(b-a)}\sum_{i=1}^{n}\Delta x_{i}=\frac{\varepsilon}{2(b-a)}(b-a)=\frac{\varepsilon}{2}<\varepsilon.
\]
根据定理~\ref{thm:6.1.4}(3), $f$ 在 $[a,b]$ 上 Riemann 可积.
\end{solution}

\begin{figure}[htbp]
\centering
\begin{tikzpicture}[x=1.12cm,y=1.0cm,>=Stealth]
  \draw[->] (-0.3,0) -- (6.9,0) node[below right] {$x$};
  \draw[->] (0,-0.3) -- (0,2.8) node[left] {$y$};
  \node[below left] at (0,0) {$O$};
  \draw (1.5,0) rectangle (2.3,1.95);
  \draw (2.3,0) rectangle (2.9,1.10);
  \draw (2.9,0) rectangle (3.45,0.97);
  \draw (3.45,0) rectangle (4.0,1.05);
  \draw (4.0,0) rectangle (4.7,1.62);
  \draw (4.7,0) rectangle (5.4,2.08);
  \draw (5.4,0) rectangle (6.1,2.38);
  \fill[pattern=north east lines] (1.5,1.95) -- (2.3,1.95) -- (2.3,1.10) -- (1.5,1.90) -- cycle;
  \fill[pattern=north east lines] (2.3,1.10) -- (2.9,1.10) -- (2.9,0.95) -- (2.3,1.10) -- cycle;
  \fill[pattern=north east lines] (2.9,0.97) -- (3.45,0.97) -- (3.45,0.92) -- (2.9,0.95) -- cycle;
  \fill[pattern=north east lines] (3.45,1.05) -- (4.0,1.05) -- (4.0,0.98) -- (3.45,0.92) -- cycle;
  \fill[pattern=north east lines] (4.0,1.62) -- (4.7,1.62) -- (4.7,1.30) -- (4.0,0.98) -- cycle;
  \fill[pattern=north east lines] (4.7,2.08) -- (5.4,2.08) -- (5.4,1.85) -- (4.7,1.30) -- cycle;
  \fill[pattern=north east lines] (5.4,2.38) -- (6.1,2.38) -- (6.1,2.35) -- (5.4,1.85) -- cycle;
  \draw[very thick] plot[smooth] coordinates {(1.5,1.90) (2.3,1.10) (2.9,0.95) (3.45,0.92) (4.0,0.98) (4.7,1.30) (5.4,1.85) (6.1,2.35)};
  \node[right] at (6.1,2.35) {$y=f(x)$};
\end{tikzpicture}
\caption{}\label{fig:6.1.4}
\end{figure}

\begin{example}\label{ex:6.1.4}
设有界函数 $f$ 在 $[a,b]$ 上只有至多可数个不连续点, 则 $f$ 在 $[a,b]$ 上 Riemann 可积.
\end{example}

\begin{solution}
因为 $D_{\text{不}}$ 为至多可数集, 它为零测集. 根据定理~\ref{thm:6.1.4}(8) 知, $f$ 在 $[a,b]$ 上 Riemann 可积.
\end{solution}

\begin{example}\label{ex:6.1.5}
设 $f$ 为 $[a,b]$ 上只有有限个不连续点的有界函数, 则 $f$ 在 $[a,b]$ 上 Riemann 可积.
\end{example}

\begin{solution}[证法 1]
因 $D_{\text{不}}$ 为有限集, 故它为零测集. 根据定理~\ref{thm:6.1.4}(8), $f$ 在 $[a,b]$ 上 Riemann 可积.
\end{solution}

\begin{solution}[证法 2]
设 $|f(x)|<M$, $\forall x\in[a,b]$, $f$ 在 $[a,b]$ 上的不连续点全体含在下面的分点内: $a=c_{0}<c_{1}<\cdots<c_{n}=b$. 对 $\forall\varepsilon>0$, 取 $[a,b]$ 的一个分割 (图~\ref{fig:6.1.5})
\[
T_{0}: a=c_{0}<a_{1}<b_{1}<c_{1}<a_{2}<b_{2}<c_{2}<\cdots<c_{n-1}<a_{n}<b_{n}<c_{n}=b,\ \text{s. t.}
\]
\[
a_{j}-c_{j-1}<\frac{\varepsilon}{6nM},\quad c_{j}-b_{j}<\frac{\varepsilon}{6nM},\quad j=1,2,\cdots,n.
\]
又 $f$ 在 $[a_{j},b_{j}]$ 内连续, $j=1,2,\cdots,n$, 故由例~\ref{ex:6.1.3} 及定理~\ref{thm:6.1.4}(3) 知, 存在 $[a_{j},b_{j}]$ 的分割
\[
T_{j}: a_{j}=x_{j_{0}}<x_{j_{1}}<\cdots<x_{j_{l}}=b_{j},\ \text{s. t.}\ \sum_{i=1}^{l}\omega_{j_{i}}\Delta x_{j_{i}}<\frac{\varepsilon}{3n}.
\]
于是, 对 $[a,b]$ 的分割 $T=T_{0}+T_{1}+\cdots+T_{n}$, 有
\begin{align*}
S(T,f)-s(T,f)&=\sum_{j=1}^{n}\left[\sup f([c_{j-1},a_{j}])-\inf f([c_{j-1},a_{j}])\right](a_{j}-c_{j-1})\\
&\quad+\sum_{j=1}^{n}\left[\sup f([b_{j},c_{j}])-\inf f([b_{j},c_{j}])\right](c_{j}-b_{j})\\
&\quad+\sum_{j=1}^{n}\left[S(T_{j},f)-s(T_{j},f)\right]\\
&<2M\cdot\frac{\varepsilon}{6nM}\cdot n+2M\cdot\frac{\varepsilon}{6nM}\cdot n+\frac{\varepsilon}{3n}\cdot n=\varepsilon.
\end{align*}
再由定理~\ref{thm:6.1.4}(3) 知, $f$ 在 $[a,b]$ 上 Riemann 可积.
\end{solution}

\begin{figure}[htbp]
\centering
\begin{tikzpicture}[x=1.02cm,y=1.0cm,>=Stealth]
  \draw[->] (0,0) -- (10.8,0) node[below right] {$x$};
  \foreach \x in {0.6,2.4,3.0,4.0,5.4,6.0,7.0,8.8,9.4,10.2}
    \draw (\x,0) -- (\x,0.14);
  \node[above] at (0.6,0.14) {$a=c_{0}$};
  \node[above] at (4.0,0.14) {$c_{1}$};
  \node[above] at (7.0,0.14) {$c_{2}$};
  \node[above] at (8.8,0.14) {$c_{n-1}$};
  \node[above] at (10.3,0.14) {$c_{n}=b$};
  \node[below] at (2.4,-0.06) {$a_{1}$};
  \node[below] at (3.0,-0.06) {$b_{1}$};
  \node[below] at (5.4,-0.06) {$a_{2}$};
  \node[below] at (6.0,-0.06) {$b_{2}$};
  \node[below] at (9.4,-0.06) {$a_{n}$};
  \node[below] at (9.8,-0.06) {$b_{n}$};
\end{tikzpicture}
\caption{}\label{fig:6.1.5}
\end{figure}

\begin{example}\label{ex:6.1.6}
设 $\{c_{i}\mid i=1,2,\cdots\}$ 为 $f$ 在 $[a,b]$ 上的全体不连续点, 且 $f$ 在 $[a,b]$ 上有界. 如果 $\lim\limits_{n\to\infty}c_{n}=c_{0}\in[a,b]$, 则 $f$ 在 $[a,b]$ 上 Riemann 可积.
\end{example}

\begin{solution}[证法 1]
参阅例~\ref{ex:6.1.4} 的证明.
\end{solution}

\begin{solution}[证法 2]
因 $\lim\limits_{n\to+\infty}c_{n}=c_{0}$, 故 $\forall\varepsilon>0$, $\exists N\in\mathbb{N}$, 当 $n>N$ 时, $c_{n}\in\left(c_{0}-\frac{\varepsilon}{4M+1},c_{0}+\frac{\varepsilon}{4M+1}\right)$, 其中 $|f(x)|\leqslant M$, $\forall x\in[a,b]$. 由于在 $[a,b]-\left(c_{0}-\frac{\varepsilon}{8M+1},c_{0}+\frac{\varepsilon}{8M+1}\right)$ 上, $f$ 只有有限个不连续点, 根据例~\ref{ex:6.1.5} 证法 2 知, $f$ 在 $[a,b]-\left(c_{0}-\frac{\varepsilon}{8M+1},c_{0}+\frac{\varepsilon}{8M+1}\right)$ 上 Riemann 可积. 根据定理~\ref{thm:6.1.4}(3), 必有 $[a,b]-\left(c_{0}-\frac{\varepsilon}{8M+1},c_{0}+\frac{\varepsilon}{8M+1}\right)$ 的分割 $T_{0}$, 使得 $\sum\limits_{T_{0}}\omega_{i}\Delta x_{i}<\frac{\varepsilon}{2}$. 由 $T_{0}$ 的分点构成 $[a,b]$ 的分割 $T$, 则有
\begin{align*}
\sum_{T}\omega_{i}\Delta x_{i}&=\sum_{T_{0}}\omega_{i}\Delta x_{i}+\left[\sup f\left(\left[c_{0}-\frac{\varepsilon}{8M+1},c_{0}+\frac{\varepsilon}{8M+1}\right]\right)\right.\\
&\quad\left.-\inf f\left(\left[c_{0}-\frac{\varepsilon}{8M+1},c_{0}+\frac{\varepsilon}{8M+1}\right]\right)\right]\cdot\frac{2\varepsilon}{8M+1}\\
&<\frac{\varepsilon}{2}+2M\cdot\frac{2\varepsilon}{8M+1}<\frac{\varepsilon}{2}+\frac{\varepsilon}{2}=\varepsilon.
\end{align*}
再根据定理~\ref{thm:6.1.4}(3) 可得, $f$ 在 $[a,b]$ 上 Riemann 可积 (见图~\ref{fig:6.1.6}).
\end{solution}

\begin{figure}[htbp]
\centering
\begin{tikzpicture}[x=1.12cm,y=1.0cm,>=Stealth]
  \draw[->] (0,0) -- (10.4,0) node[below right] {$x$};
  \draw (0.8,0) -- (0.8,0.14);
  \node[below] at (0.8,-0.06) {$a$};
  \draw[thick] (3.0,-0.30) -- (3.0,0.30);
  \draw[thick] (4.9,-0.30) -- (4.9,0.30);
  \draw[thick] (3.0,0) -- (4.9,0);
  \node[below] at (3.95,-0.34) {$c_{0}$};
  \node[above] at (3.0,0.32) {$c_{0}-\frac{\varepsilon}{8M+1}$};
  \node[above] at (5.0,0.32) {$c_{0}+\frac{\varepsilon}{8M+1}$};
  \fill (6.7,0) circle (1.3pt);
  \draw (6.7,0) -- (6.7,0.14);
  \node[below] at (6.7,-0.06) {$c_{n_{1}}$};
  \fill (7.7,0) circle (1.3pt);
  \draw (7.7,0) -- (7.7,0.14);
  \node[below] at (7.7,-0.06) {$c_{n_{k}}$};
  \draw (9.0,0) -- (9.0,0.14);
  \node[below] at (9.0,-0.06) {$b$};
\end{tikzpicture}
\caption{}\label{fig:6.1.6}
\end{figure}

从表面上看, 例~\ref{ex:6.1.5} 与例~\ref{ex:6.1.6} 中证法 1 比较简单, 而证法 2 比较复杂, 但实际上, 如果证法 1 加进定理~\ref{thm:6.1.4}(8) 的证明就不比证法 2 简单了. 此外, 证法 2 是为了培养读者有更强的分析能力.

\begin{example}\label{ex:6.1.7}
设 $f$ 为 $[a,b]$ 上的单调函数, 则 $f$ 在 $[a,b]$ 上 Riemann 可积.
\end{example}

\begin{solution}[证法 1]
因 $f$ 为 $[a,b]$ 上的单调函数, 故由定理 2.4.7 知, $D_{\text{不}}$ 为至多可数集, 从而 $D_{\text{不}}$ 为零测集. 由定理~\ref{thm:6.1.4}(8) 可知, $f$ 在 $[a,b]$ 上 Riemann 可积.
\end{solution}

\begin{solution}[证法 2]
不妨设 $f$ 为增函数. $\forall\varepsilon>0$, 取 $\delta=\frac{\varepsilon}{f(b)-f(a)+1}$, 对 $[a,b]$ 的分割 $T:a=x_{0}<x_{1}<\cdots<x_{n}=b$, $\|T\|<\delta$, 有
\begin{align*}
\sum_{T}\omega_{i}\Delta x_{i}&\leqslant\sum_{i=1}^{n}\left[f(x_{i})-f(x_{i-1})\right]\|T\|\\
&=\left[f(b)-f(a)\right]\cdot\frac{\varepsilon}{f(b)-f(a)+1}<\varepsilon,
\end{align*}
根据定理~\ref{thm:6.1.4}(3) 推得 $f$ 在 $[a,b]$ 上 Riemann 可积.
\end{solution}

\begin{example}\label{ex:6.1.8}
设 $[\alpha,\beta]\subset[a,b]$, 如果 $f$ 在 $[a,b]$ 上 Riemann 可积, 则 $f$ 在子区间 $[\alpha,\beta]$ 上也 Riemann 可积.
\end{example}

\begin{solution}[证法 1]
根据定理~\ref{thm:6.1.4}(8) 知, $f$ 在 $[a,b]$ 上 Riemann 可积等价于 $f$ 在 $[a,b]$ 上有界且几乎处处连续. 因此, $f$ 在 $[\alpha,\beta]\subseteq[a,b]$ 上有界且几乎处处连续, 再根据定理~\ref{thm:6.1.4}(8), $f$ 在 $[\alpha,\beta]$ 上也 Riemann 可积.
\end{solution}

\begin{solution}[证法 2]
因 $f$ 在 $[a,b]$ 上 Riemann 可积, 故对 $\forall\varepsilon>0$, $\exists\delta>0$, 当 $\|T\|<\delta$ 时, 有 $\sum\limits_{T}\omega_{i}\Delta x_{i}<\varepsilon$. 因此, 对 $[\alpha,\beta]$ 上任一分割 $T_{0}'$, 只要 $\|T_{0}'\|<\delta$, 必可构造 $[a,b]$ 上的分割 $T_{0}$, 它由 $T_{0}'$ 增加属于 $[a,b]-[\alpha,\beta]$ 的分点得到, 且 $\|T_{0}\|<\delta$ (图~\ref{fig:6.1.7}). 于是,
\[
\sum_{T_{0}'}\omega_{i}\Delta x_{i}\leqslant\sum_{T_{0}}\omega_{i}\Delta x_{i}<\varepsilon.
\]
由此及定理~\ref{thm:6.1.4}(3), 推得 $f$ 在 $[\alpha,\beta]$ 上 Riemann 可积.
\end{solution}

\begin{figure}[htbp]
\centering
\begin{tikzpicture}[x=1.02cm,y=1.0cm,>=Stealth]
  \draw[->] (0,0) -- (10.4,0) node[below right] {$x$};
  \foreach \x in {1.2,1.5,1.8,2.1,2.4,2.7,3.0,3.3} \fill (\x,0) circle (1.2pt);
  \draw[thick] (3.9,-0.26) -- (3.9,0.26);
  \draw[thick] (6.1,-0.26) -- (6.1,0.26);
  \foreach \x in {6.7,7.0,7.3,7.6,7.9,8.2,8.5} \fill (\x,0) circle (1.2pt);
  \node[below] at (1.2,-0.06) {$a$};
  \node[below] at (3.9,-0.30) {$\alpha$};
  \node[below] at (6.1,-0.30) {$\beta$};
  \node[below] at (8.5,-0.06) {$b$};
\end{tikzpicture}
\caption{}\label{fig:6.1.7}
\end{figure}

\begin{example}\label{ex:6.1.9}
设 $f$, $g$ 在 $[a,b]$ 上 Riemann 可积, 则 $|f|$, $fg$ 在 $[a,b]$ 上也 Riemann 可积.
\end{example}

\begin{solution}[证法 1]
因 $f$, $g$ 在 $[a,b]$ 上 Riemann 可积, 故由定理~\ref{thm:6.1.4}(8) 知, $D_{\text{不}}^{f}$ 与 $D_{\text{不}}^{g}$ 均为零测集. 再由引理~\ref{lem:6.1.6}(4) 知, $D_{\text{不}}^{f}\cup D_{\text{不}}^{g}$ 为零测集. 显然,
\[
D_{\text{不}}^{|f|}\subset D_{\text{不}}^{f},\quad D_{\text{不}}^{fg}\subset D_{\text{不}}^{f}\bigcup D_{\text{不}}^{g}.
\]
根据引理~\ref{lem:6.1.6}(3) 推得 $D_{\text{不}}^{|f|}$ 与 $D_{\text{不}}^{fg}$ 都为零测集. 再由定理~\ref{thm:6.1.4}(8), $|f|$ 与 $fg$ 在 $[a,b]$ 上都是 Riemann 可积的.
\end{solution}

\begin{solution}[证法 2]
由于 $f$, $g$ 都在 $[a,b]$ 上 Riemann 可积, 从而都有界. 设
\[
A=\sup_{x\in[a,b]}|f(x)|,\quad B=\sup_{x\in[a,b]}|g(x)|.
\]
$\forall\varepsilon>0$, 由于 $f$, $g$ 在 $[a,b]$ 上均 Riemann 可积, 故必分别存在 $[a,b]$ 的分割 $T_{1}$, $T_{2}$, 使得
\[
\sum_{T_{1}}\omega_{i}^{f}\Delta x_{i}<\frac{\varepsilon}{2B+1},\quad\sum_{T_{2}}\omega_{i}^{g}\Delta x_{i}<\frac{\varepsilon}{2A+1}.
\]
令 $T=T_{1}+T_{2}=\{\Delta_{i}\}$, 它是 $[a,b]$ 的一个新分割. 于是, 对每个 $\Delta_{i}\in T$, 有
\begin{align*}
\omega_{i}^{fg}&=\sup_{x',x''\in\Delta_{i}}|f(x')g(x')-f(x'')g(x'')|\\
&\leqslant\sup_{x',x''\in\Delta_{i}}\left[|g(x')|\cdot|f(x')-f(x'')|+|f(x'')|\cdot|g(x')-g(x'')|\right]\\
&\leqslant B\omega_{i}^{f}+A\omega_{i}^{g}.
\end{align*}
\[
\sum_{T}\omega_{i}^{fg}\Delta x_{i}\leqslant\sum_{T}\left(B\omega_{i}^{f}+A\omega_{i}^{g}\right)\Delta x_{i}=B\sum_{T}\omega_{i}^{f}\Delta x_{i}+A\sum_{T}\omega_{i}^{g}\Delta x_{i}
\]
\[
\leqslant B\sum_{T_{1}}\omega_{i}^{f}\Delta x_{i}+A\sum_{T_{2}}\omega_{i}^{g}\Delta x_{i}\leqslant B\cdot\frac{\varepsilon}{2B+1}+A\cdot\frac{\varepsilon}{2A+1}<\varepsilon.
\]
此外, 由于
\[
\omega_{i}^{|f|}=\sup_{x',x''\in\Delta_{i}}\left||f(x')|-|f(x'')|\right|\leqslant\sup_{x',x''\in\Delta_{i}}|f(x')-f(x'')|=\omega_{i}^{f},
\]
故
\[
\sum_{T_{1}}\omega_{i}^{|f|}\Delta x_{i}\leqslant\sum_{T_{1}}\omega_{i}^{f}\Delta x_{i}<\frac{\varepsilon}{2B+1}<\varepsilon.
\]
根据定理~\ref{thm:6.1.4}(3) 知, $|f|$ 与 $fg$ 在 $[a,b]$ 上 Riemann 可积.
\end{solution}

\begin{example}\label{ex:6.1.10}
设 $f$ 在 $[a,b]$ 上连续, $\varphi$ 在 $[\alpha,\beta]$ 上 Riemann 可积, $a\leqslant\varphi(t)\leqslant b$, $t\in[\alpha,\beta]$. 则 $f\circ\varphi$ 在 $[\alpha,\beta]$ 上仍 Riemann 可积.

如果 $f$, $\varphi$ 都为 Riemann 可积函数, 则 $f\circ\varphi$ 不一定 Riemann 可积.

如果 $f$ Riemann 可积, $\varphi$ 连续, 则 $f\circ\varphi$ 也不一定 Riemann 可积.
\end{example}

\begin{solution}[证法 1]
因为 $\varphi$ 在 $[\alpha,\beta]$ 上 Riemann 可积, 故 $D_{\text{不}}^{\varphi}$ 为零测集. 又因为 $f$ 在 $[a,b]$ 上连续, 且
\[
D_{\text{不}}^{f\circ\varphi}\subset D_{\text{不}}^{\varphi},
\]
根据引理~\ref{lem:6.1.6}(3) 知, $D_{\text{不}}^{f\circ\varphi}$ 也为零测集. 再由定理~\ref{thm:6.1.4}(8) 推得 $f\circ\varphi$ 在 $[\alpha,\beta]$ 上 Riemann 可积.
\end{solution}

\begin{solution}[证法 2]
因为 $f$ 在 $[a,b]$ 上连续, 所以它在 $[a,b]$ 上一致连续. 于是, $\forall\varepsilon>0$, $\exists\delta>0$, 当 $x',x''\in[a,b]$, $|x'-x''|<\delta$ 时, 有 $|f(x')-f(x'')|<\frac{\varepsilon}{2(\beta-\alpha)}$. 对 $[\alpha,\beta]$ 上 Riemann 可积的函数 $\varphi$ 应用定理~\ref{thm:6.1.4}(7), 存在 $[\alpha,\beta]$ 的某个分割 $T:\alpha=t_{0}<t_{1}<\cdots<t_{m}=\beta$, 有
\[
\sum_{\omega_{i}^{\varphi}\geqslant\delta}\Delta t_{i}<\frac{\varepsilon}{4M+1},
\]
其中 $M=\max\limits_{x\in[a,b]}|f(x)|$. 于是
\begin{align*}
\sum_{i=1}^{m}\omega_{i}^{f\circ\varphi}\Delta t_{i}&=\sum_{\omega_{i}^{\varphi}\geqslant\delta}\omega_{i}^{f\circ\varphi}\Delta t_{i}+\sum_{\omega_{i}^{\varphi}<\delta}\omega_{i}^{f\circ\varphi}\Delta t_{i}\\
&\leqslant 2M\sum_{\omega_{i}^{\varphi}\geqslant\delta}\Delta t_{i}+\frac{\varepsilon}{2(\beta-\alpha)}\sum_{\omega_{i}^{\varphi}<\delta}\Delta t_{i}\\
&\leqslant 2M\cdot\frac{\varepsilon}{4M+1}+\frac{\varepsilon}{2(\beta-\alpha)}(\beta-\alpha)<\varepsilon.
\end{align*}
根据定理~\ref{thm:6.1.4}(3), $f\circ\varphi$ 在 $[\alpha,\beta]$ 上 Riemann 可积.
\end{solution}

\begin{solution}[证法 3]
$\forall\varepsilon>0$, $\forall\eta>0$, 由于 $f$ 在 $[a,b]$ 上一致连续, 因此对上述的 $\eta$, $\exists\delta>0$, 当 $x',x''\in[a,b]$ 且 $|x'-x''|<\delta$ 时, $|f(x')-f(x'')|<\eta$.

再由 $\varphi$ 在 $[\alpha,\beta]$ 上 Riemann 可积, 根据定理~\ref{thm:6.1.4}(7), 存在某个分割 $T$, 使得
\[
\sum_{\omega_{i}^{\varphi}\geqslant\delta}\Delta t_{i}<\varepsilon.
\]
而在其余 $\omega_{i}^{\varphi}<\delta$ 的小区间上必有 $\omega_{i}^{f\circ\varphi}<\eta$. 因此,
\[
\sum_{\omega_{i}^{f\circ\varphi}\geqslant\eta}\Delta t_{i}\leqslant\sum_{\omega_{i}^{\varphi}\geqslant\delta}\Delta t_{i}<\varepsilon.
\]
再根据定理~\ref{thm:6.1.4}(7) 得到 $f\circ\varphi$ 在 $[\alpha,\beta]$ 上 Riemann 可积.
\end{solution}

\begin{solution}[反例 1]
设 $f(x)=\operatorname{sgn}x=\begin{cases}1,&x>0,\\0,&x=0,\\-1,&x<0,\end{cases}\quad \varphi(t)=R(t)$ (Riemann 函数), 则 $f$, $\varphi$ 都在 $[0,1]$ 上 Riemann 可积, 但
\[
f\circ\varphi(t)=\operatorname{sgn}(R(t))=\begin{cases}1,&t\ \text{为有理数},\\0,&t\ \text{为无理数}\end{cases}=D(t)\quad(\text{Dirichlet 函数})
\]
在 $[0,1]$ 上不是 Riemann 可积的.
\end{solution}

\begin{solution}[反例 2]
设 $f(x)=\begin{cases}0,&x=0,\\1,&0<x\leqslant 1,\end{cases}$ 它是一个有界函数, 且只有一个不连续点 $0$, 根据例~\ref{ex:6.1.5} 知, $f$ 为 $[0,1]$ 上的 Riemann 可积函数.

为构造 $\varphi$, 先定义类 Cantor 集 $C$, 在 $[0,1]$ 的中点 $\frac{1}{2}$ 处挖去一个以 $\frac{1}{2}$ 为中心, 长度为 $\frac{1}{4}$ 的小区间; 然后, 在余下的两个区间的中点各挖去一个以该中点为中心, 长度为 $\frac{1}{2\times 8}$ 的小区间; 再在余下的 $8$ 个区间的中点各挖去一个以该中点为中心, 长度为 $\frac{1}{2^{2}\times 16}$ 的小区间; 依此类推. 易见, 挖去小区间的总长度为
\[
\frac{1}{4}+2\times\frac{1}{2\times 8}+2^{2}\times\frac{1}{2^{2}\times 16}+\cdots=\frac{1}{4}+\frac{1}{8}+\frac{1}{16}+\cdots=\frac{\frac{1}{4}}{1-\frac{1}{2}}=\frac{1}{2}.
\]
记 $[0,1]$ 在挖去上述可数个小区间后余下的集合为 $C$, 称为类 Cantor 集. 应用反证法可知, $C$ 不为零测集 (在实变函数论中知道, $C$ 的 Lebesgue 测度为 $1-\frac{1}{2}=\frac{1}{2}\neq 0$). 设 $[0,1]-\bigcup\limits_{n=1}^{\infty}(a_{n},b_{n})=C$, 其中 $(a_{n},b_{n})$ 彼此不相交, 它们为 $[0,1]\setminus C$ 的构成区间. 令
\[
\varphi(t)=\begin{cases}0,&t\in C,\\\frac{1}{2}(b_{n}-a_{n})-\left|t-\frac{1}{2}(a_{n}+b_{n})\right|,&t\in(a_{n},b_{n}),\ n\in\mathbb{N},\end{cases}
\]
易见, $\varphi$ 为 $[0,1]$ 上的连续函数, 且
\[
f\circ\varphi(t)=\begin{cases}0,&t\in C,\\1,&t\in[0,1]\setminus C.\end{cases}
\]
它的不连续点恰为类 Cantor 集 $C$, 即 $D_{\text{不}}^{f\circ\varphi}$, 它不为零测集. 根据定理~\ref{thm:6.1.4}(8) 推得 $f\circ\varphi$ 在 $[0,1]$ 上不是 Riemann 可积的.
\end{solution}

\begin{example}\label{ex:6.1.11}
证明: $\int_{a}^{b}c\,\mathrm{d}x=c(b-a)$, 其中 $c$ 为常数.
\end{example}

\begin{solution}
对 $[a,b]$ 的任何分割 $T:a=x_{0}<x_{1}<\cdots<x_{n}=b$, $\forall\xi_{i}\in[x_{i-1},x_{i}]$, $\forall\varepsilon>0$, 有
\[
\sum_{i=1}^{n}f(\xi_{i})\Delta x_{i}=\sum_{i=1}^{n}c\Delta x_{i}=c(b-a),
\]
\[
\left|\sum_{i=1}^{n}f(\xi_{i})\Delta x_{i}-c(b-a)\right|=|c(b-a)-c(b-a)|=0<\varepsilon.
\]
根据定义~\ref{def:6.1.1}, 可知 $\int_{a}^{b}c\,\mathrm{d}x=c(b-a)$.
\end{solution}

\begin{example}\label{ex:6.1.12}
证明: $\int_{a}^{b}x\,\mathrm{d}x=\frac{b^{2}-a^{2}}{2}=\frac{b+a}{2}(b-a)$.
\end{example}

这恰为梯形 $AabB$ 的面积 (图~\ref{fig:6.1.8}).

\begin{solution}[证法 1]
易见,
\begin{align*}
\sum_{i=1}^{n}\eta_{i}\Delta x_{i}&=\sum_{i=1}^{n}\frac{x_{i-1}+x_{i}}{2}(x_{i}-x_{i-1})\\
&=\frac{1}{2}\sum_{i=1}^{n}(x_{i}^{2}-x_{i-1}^{2})\\
&=\frac{1}{2}(b^{2}-a^{2}).
\end{align*}
$\forall\varepsilon>0$, 取 $\delta=\frac{\varepsilon}{b-a}$, 当 $[a,b]$ 的分割 $T:a=x_{0}<x_{1}<\cdots<x_{n}=b$, $\|T\|<\delta$, $\forall\xi_{i}\in[x_{i-1},x_{i}]$, $i=1,2,\cdots,n$, 有
\begin{align*}
\left|\sum_{i=1}^{n}\xi_{i}\Delta x_{i}-\frac{1}{2}(b^{2}-a^{2})\right|&=\left|\sum_{i=1}^{n}(\xi_{i}-\eta_{i})\Delta x_{i}\right|\\
&\leqslant\sum_{i=1}^{n}|\xi_{i}-\eta_{i}|\Delta x_{i}\leqslant\|T\|\sum_{i=1}^{n}\Delta x_{i}\\
&<\frac{\varepsilon}{b-a}(b-a)=\varepsilon,
\end{align*}
根据定义~\ref{def:6.1.1}, 可知 $\int_{a}^{b}x\,\mathrm{d}x=\frac{b^{2}-a^{2}}{2}$.
\end{solution}

\begin{solution}[证法 2]
因为 $f(x)=x$ 连续, 所以它在 $[a,b]$ 上 Riemann 可积, 且
\begin{align*}
\int_{a}^{b}x\,\mathrm{d}x&=\lim_{n\to+\infty}\sum_{i=1}^{n}\left[a+\frac{i}{n}(b-a)\right]\frac{b-a}{n}\\
&=\lim_{n\to+\infty}\left[a(b-a)+\frac{(b-a)^{2}}{n^{2}}\frac{n(n+1)}{2}\right]\\
&=a(b-a)+\frac{(b-a)^{2}}{2}=\frac{b^{2}-a^{2}}{2}.
\end{align*}
\end{solution}

\begin{figure}[htbp]
\centering
\begin{tikzpicture}[x=1.05cm,y=0.80cm,>=Stealth]
  \draw[->] (-0.3,0) -- (6.6,0) node[below right] {$x$};
  \draw[->] (0,-0.3) -- (0,5.9) node[left] {$y$};
  \node[below left] at (0,0) {$O$};
  \draw[thick] (0,0) -- (5.8,5.8);
  \node[right] at (4.3,5.0) {$y=x$};
  \draw[thick] (1.5,0) -- (1.5,1.5);
  \draw[thick] (2.9,0) -- (2.9,2.9);
  \draw[thick] (4.6,0) -- (4.6,4.6);
  \draw[thick] (5.8,0) -- (5.8,5.8);
  \node[right] at (5.82,5.8) {$B$};
  \node[left] at (1.48,1.5) {$A$};
  \draw[dashed] (0,3.6) -- (3.6,3.6);
  \node[left] at (0,3.6) {$\eta_{i}=\dfrac{x_{i-1}+x_{i}}{2}$};
  \fill (3.6,0) circle (1.3pt);
  \fill (4.1,0) circle (1.3pt);
  \node[below] at (1.5,-0.10) {$a$};
  \node[below] at (2.9,-0.10) {$x_{i-1}$};
  \node[below] at (3.6,-0.12) {$\eta_{i}$};
  \node[below] at (4.1,-0.36) {$\xi_{i}$};
  \node[below] at (4.6,-0.12) {$x_{i}$};
  \node[below] at (5.8,-0.10) {$b$};
\end{tikzpicture}
\caption{}\label{fig:6.1.8}
\end{figure}

\begin{example}\label{ex:6.1.13}
求在区间 $[a,b]$ 上, 以抛物线 $y=x^{2}$ 为曲边, $y=0$, $x=a$ 与 $x=b$ 为直边的曲边梯形的面积 $S$.
\end{example}

\begin{solution}
因为 $y=x^{2}$ 在 $[a,b]$ 上连续, 所以在 $[a,b]$ 上 Riemann 可积. 于是, 将 $[a,b]$ $n$ 等分, 所得分割为 $T_{n}:a=x_{0}<x_{1}<\cdots<x_{n}=b$, $x_{i}=a+\frac{i}{n}(b-a)$, $\|T_{n}\|=\frac{b-a}{n}$, $\xi_{i}=x_{i}$, $i=0,1,\cdots,n-1$. 根据定理~\ref{thm:6.1.4}(5) (图~\ref{fig:6.1.9}), 有
\begin{align*}
S&=\int_{a}^{b}x^{2}\,\mathrm{d}x\\
&=\lim_{n\to+\infty}\sum_{i=0}^{n-1}\left[a+\frac{i}{n}(b-a)\right]^{2}\frac{b-a}{n}\\
&=\lim_{n\to+\infty}\left[a^{2}(b-a)+2a(b-a)^{2}\frac{n(n-1)}{2n^{2}}+(b-a)^{3}\frac{(n-1)n(2n-1)}{6n^{3}}\right]\\
&=a^{2}(b-a)+a(b-a)^{2}+\frac{1}{3}(b-a)^{3}=\frac{b^{3}-a^{3}}{3}.
\end{align*}
\end{solution}

\begin{figure}[htbp]
\centering
\begin{tikzpicture}[x=1.0cm,y=0.80cm,>=Stealth]
  \draw[->] (-0.3,0) -- (6.7,0) node[below right] {$x$};
  \draw[->] (0,-0.3) -- (0,5.0) node[left] {$y$};
  \node[below left] at (0,0) {$O$};
  \fill[pattern=north east lines] (1.4,0) rectangle (2.1,0.485);
  \draw (1.4,0) rectangle (2.1,0.485);
  \fill[pattern=north east lines] (2.1,0) rectangle (2.8,0.862);
  \draw (2.1,0) rectangle (2.8,0.862);
  \fill[pattern=north east lines] (2.8,0) rectangle (3.5,1.348);
  \draw (2.8,0) rectangle (3.5,1.348);
  \fill[pattern=north east lines] (3.5,0) rectangle (4.2,1.940);
  \draw (3.5,0) rectangle (4.2,1.940);
  \fill[pattern=north east lines] (4.2,0) rectangle (4.9,2.641);
  \draw (4.2,0) rectangle (4.9,2.641);
  \fill[pattern=north east lines] (4.9,0) rectangle (5.6,3.450);
  \draw (4.9,0) rectangle (5.6,3.450);
  \fill[pattern=north east lines] (5.6,0) rectangle (6.3,4.366);
  \draw (5.6,0) rectangle (6.3,4.366);
  \draw[very thick,domain=0:6.3,smooth,samples=60] plot (\x,{0.11*\x*\x});
  \node[right] at (3.9,2.90) {$y=x^{2}$};
  \node[below] at (1.4,-0.10) {$a$};
  \node[below] at (3.4,-0.10) {$x_{i-1}$};
  \node[below] at (4.1,-0.10) {$x_{i}$};
  \node[below] at (6.3,-0.10) {$b$};
\end{tikzpicture}
\caption{}\label{fig:6.1.9}
\end{figure}

\begin{example}\label{ex:6.1.14}
求 $\int_{a}^{b}\sin x\,\mathrm{d}x$.
\end{example}

\begin{solution}
因为 $\sin x$ 连续, 所以 $\sin x$ 在 $[a,b]$ 上 Riemann 可积. 于是, 将 $[a,b]$ $n$ 等分, 所得分割为 $T_{n}:a=x_{0}<x_{1}<\cdots<x_{n}=b$, $x_{i}=a+\frac{i}{n}(b-a)$, $\|T_{n}\|=\frac{b-a}{n}$, $\xi_{i}=x_{i}$, $i=1,2,\cdots,n$. 根据定理~\ref{thm:6.1.4}(5) 就有
\begin{align*}
\int_{a}^{b}\sin x\,\mathrm{d}x&=\lim_{n\to+\infty}\sum_{i=1}^{n}\sin\left(a+i\frac{b-a}{n}\right)\cdot\frac{b-a}{n}\\
&=\lim_{n\to+\infty}\frac{b-a}{2n\sin\frac{b-a}{2n}}\sum_{i=1}^{n}2\sin\left(a+i\frac{b-a}{n}\right)\sin\frac{b-a}{2n}\\
&=\lim_{n\to+\infty}\frac{\frac{b-a}{2n}}{\sin\frac{b-a}{2n}}\sum_{i=1}^{n}\left[\cos\left(a+(2i-1)\frac{b-a}{2n}\right)\right.\\
&\quad\left.-\cos\left(a+(2i+1)\frac{b-a}{2n}\right)\right]\\
&=\lim_{n\to+\infty}\frac{\frac{b-a}{2n}}{\sin\frac{b-a}{2n}}\left[\cos\left(a+\frac{b-a}{2n}\right)-\cos\left(a+\frac{2n+1}{2n}(b-a)\right)\right]\\
&=1\cdot\left[\cos a-\cos(a+(b-a))\right]=\cos a-\cos b.
\end{align*}
\end{solution}

\begin{remark}\label{remark:6.1.2}
对于 $[a,b]$ 上 Riemann 可积的函数 $f$, 可选取特殊的分割和点, 例如: $n$ 等分; $\xi_{i}$ 可取 $x_{i-1}$, 或 $x_{i}$, 或 $\frac{x_{i-1}+x_{i}}{2}$, 或 $\sqrt{x_{i-1}x_{i}}$.
\end{remark}

从上述例~\ref{ex:6.1.12}, 例~\ref{ex:6.1.13} 与例~\ref{ex:6.1.14} 可看出, 用定义求 Riemann 积分的值很辛苦, 甚至有时无法求出. 6.3 节将给出求 Riemann 积分值非常简单且十分有效的 Newton-Leibniz 公式 $\int_{a}^{b}f(x)\,\mathrm{d}x=F(b)-F(a)$, 其中 $F'(x)=f(x)$.

\begin{exblock}{练习题 6.1}
  \item 通过对积分区间作等分分割, 并取适当的点集 $\{\xi_{i}\}$, 将 Riemann 积分视作 $\sum\limits_{i=1}^{n}f(\xi_{i})\Delta x_{i}$ 的极限. 由此计算下列 Riemann 积分:
  \begin{enumerate}[label=(\arabic*)]
    \item $\int_{0}^{1}x^{3}\,\mathrm{d}x$; (提示: $\sum\limits_{i=1}^{n}i^{3}=\left[\frac{n(n+1)}{2}\right]^{2}$);
    \item $\int_{a}^{b}\frac{\mathrm{d}x}{x^{2}}$, $0<a<b$; (提示: $\xi_{i}=\sqrt{x_{i-1}x_{i}}$);
    \item $\int_{a}^{b}\mathrm{e}^{x}\,\mathrm{d}x$, $a<b$; (提示: $\mathrm{e}^{x_{i}}-\mathrm{e}^{x_{i-1}}=\mathrm{e}^{\eta_{i}}(x_{i}-x_{i-1})$).
  \end{enumerate}

  \item 利用 Riemann 积分的几何意义, 求下列积分:
  \begin{enumerate}[label=(\arabic*)]
    \item $\int_{a}^{b}\left|x-\frac{a+b}{2}\right|\mathrm{d}x$.
    \item $\int_{a}^{b}\sqrt{(x-a)(b-x)}\,\mathrm{d}x$.
  \end{enumerate}

  \item 设 $p:[a,b]\to\mathbb{R}$. 如果有分割
  \[
  a=x_{0}<x_{1}<\cdots<x_{n}=b
\]
  使得在每一个子区间 $(x_{i-1},x_{i})$, $i=1,2,\cdots,n$ 上, $p$ 为常值函数, 则称 $p$ 为 $[a,b]$ 上的阶梯函数. 又设 $f$ 在 $[a,b]$ 上 Riemann 可积, 证明: 对任何 $\varepsilon>0$, 必存在 $[a,b]$ 上的阶梯函数 $p$ 与 $q$, 使得在 $[a,b]$ 上有 $p\leqslant f\leqslant q$, 并且
  \[
  \int_{a}^{b}[q(x)-p(x)]\,\mathrm{d}x<\varepsilon.
\]

  \item 设 $f$ 在 $[a,b]$ 上 Riemann 可积. 证明: 对任何 $\varepsilon>0$, 必存在 $[a,b]$ 上的连续函数 $p$ 与 $q$, 使得在 $[a,b]$ 上有 $p\leqslant f\leqslant q$, 并且
  \[
  \int_{a}^{b}[q(x)-p(x)]\,\mathrm{d}x<\varepsilon.
\]

  \item 设函数 $f$ 在 $[a,b]$ 上 Riemann 可积, 且对 $[a,b]$ 内的任何子区间 $\Delta$ 上有 $\sup\limits_{\Delta}f\geqslant\sigma$, 这里 $\sigma$ 为一常数. 证明:
  \[
  \int_{a}^{b}f(x)\,\mathrm{d}x\geqslant\sigma(b-a).
\]

  \item 设 $f$ 与 $g$ 在 $[a,b]$ 上都 Riemann 可积. 证明: $\max\{f(x),g(x)\}$ 与 $\min\{f(x),g(x)\}$ 在 $[a,b]$ 上也都 Riemann 可积.

  \item 设 $f$ 在 $[a,b]$ 上 Riemann 可积, 且在 $[a,b]$ 上满足 $|f(x)|\geqslant m>0$. 证明: $\frac{1}{f}$ 在 $[a,b]$ 上也 Riemann 可积.

  \item 设
  \[
  f(x)=\begin{cases}x,&x\ \text{为有理数},\\0,&x\ \text{为无理数}.\end{cases}
\]
  试求: $f$ 在 $[0,1]$ 上的上积分与下积分. 并由此及其他方法判断 $f$ 在 $[0,1]$ 上是否 Riemann 可积.

  \item 设 $f$ 在 $[a,b]$ 上 Riemann 可积, 且 $f(x)\geqslant 0$, $x\in[a,b]$. 试问: $\sqrt{f}$ 在 $[a,b]$ 上是否 Riemann 可积? 并说明理由.

  \item (1) 证明引理~\ref{lem:6.1.2} 中第二式: $s(T)\leqslant s(T')\leqslant s(T)+(M-m)p\|T\|$;

  (2) 证明引理~\ref{lem:6.1.3} 中第二式: $\lim\limits_{\|T\|\to 0}s(T)=s$.
\end{exblock}

\begin{exblock}{思考题 6.1}
  \item 设 $m\in\mathbb{N}$, 用积分的定义证明: $\int_{a}^{b}x^{m}\,\mathrm{d}x=\frac{1}{m+1}(b^{m+1}-a^{m+1})$. (提示: (1) 令 $\eta_{i}=\left(\frac{x_{i-1}^{m}+x_{i-1}^{m-1}x_{i}+\cdots+x_{i}^{m}}{m+1}\right)^{\frac{1}{m}}$, $\sum\limits_{i=1}^{n}\eta_{i}^{m}(x_{i}-x_{i-1})=\frac{1}{m+1}\sum\limits_{i=1}^{n}(x_{i}^{m+1}-x_{i-1}^{m+1})$; (2) 选择 $[a,b]$ 诸分点为 $a,aq,aq^{2},\cdots,aq^{n}=b$, $q=\left(\frac{b}{a}\right)^{\frac{1}{n}}$.)

  \item 应用上题提示 (2) 的方法证明: $\int_{a}^{b}\frac{\mathrm{d}x}{x}=\ln\frac{b}{a}$, $0<a<b$.

  \item 设正数列 $\{a_{n}\}$ 满足 $\lim\limits_{n\to+\infty}\int_{0}^{a_{n}}x^{n}\,\mathrm{d}x=2$. 证明 $\lim\limits_{n\to+\infty}a_{n}=1$.

  \item (1) 设 $f$ 在 $[a,b]$ 上 Riemann 可积. 证明: 开区间 $(a,b)$ 内至少有 $f$ 的一个连续点.

  (2) 设 $f\geqslant 0$ 且在 $[a,b]$ 上 Riemann 可积. 证明:
  \[
  \int_{a}^{b}f(x)\,\mathrm{d}x=0\Longleftrightarrow f\ \text{在连续点处都必取零值}\Longleftrightarrow f\ \text{在}\ [a,b]\ \text{上几乎处处为零}.
\]

  (3) 设 $f\geqslant 0$, 且有界, 如果 $f$ 在 $[a,b]$ 上几乎处处为零, 则是否有 $\int_{a}^{b}f(x)\,\mathrm{d}x=0$?

  \item 设 $f$, $g$ 均为 $[a,b]$ 上的 Riemann 可积函数, $f\geqslant g$, $x\in[a,b]$. 证明: $\int_{a}^{b}f(x)\,\mathrm{d}x=\int_{a}^{b}g(x)\,\mathrm{d}x\Longleftrightarrow f$ 与 $g$ 在 $[a,b]$ 上几乎处处相等, 即 $f\xlongequal{\text{a. e.}}g$, $x\in[a,b]$.

  \item 设 $f$, $g$ 均为 $[a,b]$ 上的 Riemann 可积函数. 证明:
  \[
  \int_{a}^{b}[f(x)-g(x)]^{2}\,\mathrm{d}x=0\Longleftrightarrow f\xrightarrow{\text{a. e.}}g,\ x\in[a,b].
\]

  \item 设 $f$, $g$ 均为 $[a,b]$ 上的有界函数, 且仅在 $[a,b]$ 中有限个点处 $f(x)\neq g(x)$. 证明: 当 $f$ 在 $[a,b]$ 上 Riemann 可积时, $g$ 在 $[a,b]$ 上也 Riemann 可积, 且 $\int_{a}^{b}f(x)\,\mathrm{d}x=\int_{a}^{b}g(x)\,\mathrm{d}x$.

  \item 设 $f$ 与 $g$ 在 $[a,b]$ 上均 Riemann 可积. 证明:
  \[
  \lim_{\|T\|\to 0}\sum_{i=1}^{n}f(\xi_{i})g(\eta_{i})\Delta x_{i}=\int_{a}^{b}f(x)g(x)\,\mathrm{d}x,
\]
  其中 $T:a=x_{0}<x_{1}<\cdots<x_{n}=b$ 为 $[a,b]$ 的分割, $\xi_{i},\eta_{i}\in[x_{i-1},x_{i}]$.

  \item 设 $f$ 在 $[a,b]$ 上 Riemann 可积, 直接用闭区间套原理按下面顺序逐一证明: $f$ 在 $[a,b]$ 内必定有无限多个处处稠密的连续点.

  (1) 若 $T$ 为 $[a,b]$ 的一个分割, 使得 $S(T)-s(T)<b-a$, 则在 $T$ 中存在某个小区间
  \[
  \Delta_{i}=[x_{i-1},x_{i}],\ \text{使}\ \omega_{i}^{f}<1.
\]

  (2) 存在闭区间 $[a_{1},b_{1}]\subset(a,b)$, 使得
  \[
  \omega^{f}([a_{1},b_{1}])=\sup_{x\in[a_{1},b_{1}]}f(x)-\inf_{x\in[a_{1},b_{1}]}f(x)<1,\quad b_{1}-a_{1}<\frac{b-a}{2}.
\]

  (3) 存在闭区间 $[a_{2},b_{2}]\subset(a_{1},b_{1})$, 使得
  \[
  \omega^{f}([a_{2},b_{2}])=\sup_{x\in[a_{2},b_{2}]}f(x)-\inf_{x\in[a_{2},b_{2}]}f(x)<\frac{1}{2},\quad b_{2}-a_{2}<\frac{b-a}{2^{2}}.
\]

  (4) 继续以上方法, 求出一闭区间序列 $[a_{n},b_{n}]\subset(a_{n-1},b_{n-1})$, 使得
  \[
  \omega^{f}([a_{n},b_{n}])=\sup_{x\in[a_{n},b_{n}]}f(x)-\inf_{x\in[a_{n},b_{n}]}f(x)<\frac{1}{n},\quad b_{n}-a_{n}<\frac{b-a}{2^{n}}.
\]
  根据闭区间套原理, $\exists!x_{0}\in\bigcap\limits_{n=1}^{\infty}[a_{n},b_{n}]$, 则 $x_{0}$ 为 $f$ 的连续点.

  (5) 上面求得的 $f$ 的连续点在 $[a,b]$ 内处处稠密.

  \item 设 $f$ 在 $[a,b]$ 上 Riemann 可积, 且处处有 $f(x)>0$. 证明: $\int_{a}^{b}f(x)\,\mathrm{d}x>0$.
\end{exblock}

% ====================================================================
%  徐森林、薛春华《数学分析》(第一册)
%  \S6.2  Riemann 积分的性质、积分第一与第二中值定理
%  源: 原书印刷页 353--371(PDF 物理页 359--377)
%  底本: PaddleOCR-VL-1.6, 见 chapters/_ch06_2_raw.md
%  约定: 全西文标点; 行间公式用一对方括号; 多行公式用 align 或 aligned;
%        跨册引用一律用 extref(写作「定理 x.y.z」原文); 数学模式内的中文必须包text
% ====================================================================
\section{Riemann 积分的性质、积分第一与第二中值定理}\label{sec:6.2}

要深入研究 Riemann 积分, 仅仅停留在定义上是不行的, 必须进一步讨论其各种性质, 才能加以展开. 由 Riemann 积分定义知, $\forall a,b\in\mathbb{R}$, $a<b$, 有
\[
\int_{b}^{a}c\,\mathrm{d}x=\lim_{\|T\|\to 0}\sum_{i=1}^{n}f(\xi_{i})\Delta x_{i}=\lim_{\|T\|\to 0}\sum_{i=1}^{n}c\Delta x_{i}=\lim_{\|T\|\to 0}c(b-a)=c(b-a).
\]
为了方便, 我们约定
\[
\int_{b}^{a}f(x)\,\mathrm{d}x=-\int_{a}^{b}f(x)\,\mathrm{d}x,\quad a<b;
\]
\[
\int_{a}^{a}f(x)\,\mathrm{d}x=0.
\]
由此知, 上边等式对任何 $a,b\in\mathbb{R}$ 成立.

\begin{theorem}\label{thm:6.2.1}
设 $f$, $g$ 在 $[a,b]$ 上 Riemann 可积, $\lambda,\mu\in\mathbb{R}$, 则 $\lambda f$, $f\pm g$ 在 $[a,b]$ 上也 Riemann 可积, 且

(1)
\[
\int_{a}^{b}\lambda f(x)\,\mathrm{d}x=\lambda\int_{a}^{b}f(x)\,\mathrm{d}x;
\]

(2)
\[
\int_{a}^{b}[f(x)\pm g(x)]\,\mathrm{d}x=\int_{a}^{b}f(x)\,\mathrm{d}x\pm\int_{a}^{b}g(x)\,\mathrm{d}x;
\]
两式综合表达为
\[
\int_{a}^{b}[\lambda f(x)+\mu g(x)]\,\mathrm{d}x=\lambda\int_{a}^{b}f(x)\,\mathrm{d}x+\mu\int_{a}^{b}g(x)\,\mathrm{d}x.
\]
\end{theorem}

\begin{proof}
对 $\forall\varepsilon>0$, 由 $f$, $g$ 在 $[a,b]$ 上 Riemann 可积, $\exists\delta>0$, 当 $[a,b]$ 的分割 $T$, $\|T\|<\delta$ 时, 有
\[
\left|\sum_{i=1}^{n}f(\xi_{i})\Delta x_{i}-\int_{a}^{b}f(x)\,\mathrm{d}x\right|<\frac{\varepsilon}{|\lambda|+|\mu|+1},
\]
\[
\left|\sum_{i=1}^{n}g(\xi_{i})\Delta x_{i}-\int_{a}^{b}g(x)\,\mathrm{d}x\right|<\frac{\varepsilon}{|\lambda|+|\mu|+1}.
\]
于是, 有
\begin{align*}
&\left|\sum_{i=1}^{n}[\lambda f(\xi_{i})+\mu g(\xi_{i})]\Delta x_{i}-\left(\lambda\int_{a}^{b}f(x)\,\mathrm{d}x+\mu\int_{a}^{b}g(x)\,\mathrm{d}x\right)\right|\\
&\leqslant|\lambda|\left|\sum_{i=1}^{n}f(\xi_{i})\Delta x_{i}-\int_{a}^{b}f(x)\,\mathrm{d}x\right|\\
&\quad+|\mu|\left|\sum_{i=1}^{n}g(\xi_{i})\Delta x_{i}-\int_{a}^{b}g(x)\,\mathrm{d}x\right|\\
&\leqslant\frac{|\lambda|\varepsilon}{|\lambda|+|\mu|+1}+\frac{|\mu|\varepsilon}{|\lambda|+|\mu|+1}<\varepsilon,
\end{align*}
即
\begin{align*}
\int_{a}^{b}[\lambda f(x)+\mu g(x)]\,\mathrm{d}x&=\lim_{\|T\|\to 0}\sum_{i=1}^{n}[\lambda f(\xi_{i})+\mu g(\xi_{i})]\Delta x_{i}\\
&=\lambda\int_{a}^{b}f(x)\,\mathrm{d}x+\mu\int_{a}^{b}g(x)\,\mathrm{d}x.
\end{align*}
\end{proof}

\begin{theorem}\label{thm:6.2.2}
(1) 设 $f$ 为 $[a,b]$ 上的 Riemann 可积函数, 且 $f(x)\geqslant 0$, $\forall x\in[a,b]$, 则
\[
\int_{a}^{b}f(x)\,\mathrm{d}x\geqslant 0;
\]

(2) 设 $f$, $g$ 均为 $[a,b]$ 上的 Riemann 可积函数, 且 $g(x)\leqslant f(x)$, $\forall x\in[a,b]$, 则有
\[
\int_{a}^{b}g(x)\,\mathrm{d}x\leqslant\int_{a}^{b}f(x)\,\mathrm{d}x.
\]

(3) 在 (2) 中, 如果 $\exists x_{0}\in[a,b]$, s. t. $g(x_{0})<f(x_{0})$, 且 $f$, $g$ 均在 $x_{0}$ 连续, 则
\[
\int_{a}^{b}g(x)\,\mathrm{d}x<\int_{a}^{b}f(x)\,\mathrm{d}x.
\]
特别当 $g=0$ 时, $\int_{a}^{b}f(x)\,\mathrm{d}x>0$.
\end{theorem}

\begin{proof}
(1) 因为 $f(x)\geqslant 0$, 故 $\sum\limits_{i=1}^{n}f(\xi_{i})\Delta x_{i}\geqslant 0$, 从而
\[
\int_{a}^{b}f(x)\,\mathrm{d}x=\lim_{\|T\|\to 0}\sum_{i=1}^{n}f(\xi_{i})\Delta x_{i}\geqslant 0.
\]

(2) 由定理~\ref{thm:6.2.1} 及 $f$, $g$ 在 $[a,b]$ 上 Riemann 可积知, $g(x)-f(x)$ 在 $[a,b]$ 上也 Riemann 可积. 再由 $g(x)\leqslant f(x)$ 推得 $f(x)-g(x)\geqslant 0$. 从 (1) 立即得到
\begin{align*}
\int_{a}^{b}f(x)\,\mathrm{d}x-\int_{a}^{b}g(x)\,\mathrm{d}x&=\int_{a}^{b}[f(x)-g(x)]\,\mathrm{d}x\geqslant 0,\\
\int_{a}^{b}g(x)\,\mathrm{d}x&\leqslant\int_{a}^{b}f(x)\,\mathrm{d}x.
\end{align*}

(3) 因为 $f-g$ 在 $x_{0}$ 连续, 且 $f(x_{0})-g(x_{0})>0$, 所以存在区间 $[\alpha,\beta]$ 含 $x_{0}$, s. t. $\forall x\in[\alpha,\beta]\cap[a,b]=[\overline{\alpha},\overline{\beta}]$, 有
\[
f(x)-g(x)>\frac{1}{2}[f(x_{0})-g(x_{0})]>0.
\]
因此,
\begin{align*}
\int_{a}^{b}[f(x)-g(x)]\,\mathrm{d}x&\xlongequal{\text{定理~\ref{thm:6.2.4}}}\int_{a}^{\overline{\alpha}}[f(x)-g(x)]\,\mathrm{d}x\\
&\quad+\int_{\overline{\alpha}}^{\overline{\beta}}[f(x)-g(x)]\,\mathrm{d}x+\int_{\overline{\beta}}^{b}[f(x)-g(x)]\,\mathrm{d}x\\
&\geqslant 0+\int_{\overline{\alpha}}^{\overline{\beta}}\frac{1}{2}[g(x_{0})-g(x_{0})]\,\mathrm{d}x\\
&=\frac{1}{2}[f(x_{0})-g(x_{0})](\overline{\beta}-\overline{\alpha})>0,
\end{align*}
于是,
\[
\int_{a}^{b}g(x)\,\mathrm{d}x<\int_{a}^{b}f(x)\,\mathrm{d}x.
\]
\end{proof}

\begin{theorem}\label{thm:6.2.3}
设 $f$ 在 $[a,b]$ 上 Riemann 可积, 则 $|f|$ 在 $[a,b]$ 上也 Riemann 可积, 且
\[
\int_{a}^{b}f(x)\,\mathrm{d}x\leqslant\int_{a}^{b}|f(x)|\,\mathrm{d}x.
\]
\end{theorem}

\begin{proof}
因 $f$ 在 $[a,b]$ 上 Riemann 可积, 根据例~\ref{ex:6.1.9}, $|f|$ 在 $[a,b]$ 上也 Riemann 可积, 再由不等式
\[
-|f(x)|\leqslant f(x)\leqslant|f(x)|
\]
及定理~\ref{thm:6.2.2}(2) 得到
\begin{align*}
-\int_{a}^{b}|f(x)|\,\mathrm{d}x&\leqslant\int_{a}^{b}f(x)\,\mathrm{d}x\leqslant\int_{a}^{b}|f(x)|\,\mathrm{d}x,\\
\left|\int_{a}^{b}f(x)\,\mathrm{d}x\right|&\leqslant\int_{a}^{b}|f(x)|\,\mathrm{d}x.
\end{align*}
\end{proof}

\begin{example}\label{ex:6.2.1}
设
\[
f(x)=\begin{cases}1,&x\in\mathbb{Q},\\-1,&x\in\mathbb{R}\setminus\mathbb{Q},\end{cases}
\]
则类似 Dirichlet 函数可证 $f(x)$ 在 $[0,1]$ 上不是 Riemann 可积的, 但 $|f(x)|=1$ 在 $[0,1]$ 上 Riemann 可积.
\end{example}

\begin{theorem}\label{thm:6.2.4}
设 $c\in[a,b]$, $f$ 为 $[a,b]$ 上的函数, 则: $f$ 在 $[a,b]$ 上 Riemann 可积 $\Leftrightarrow f$ 在 $[a,c]$ 上与 $[c,b]$ 上均 Riemann 可积. 并且,
\[
\int_{a}^{b}f(x)\,\mathrm{d}x=\int_{a}^{c}f(x)\,\mathrm{d}x+\int_{c}^{b}f(x)\,\mathrm{d}x.
\]
\end{theorem}

\begin{proof}[证法 1]
$f$ 在 $[a,b]$ 上 Riemann 可积
\[
\xLongleftrightarrow{\text{定理~\ref{thm:6.1.4}(8)}}f\ \text{在}\ [a,b]\ \text{上几乎处处连续且有界}
\]
\[
\Leftrightarrow f\ \text{在}\ [a,c]\ \text{及}\ [c,b]\ \text{上均几乎处处连续且有界}
\]
\[
\xLongleftrightarrow{\text{定理~\ref{thm:6.1.4}(8)}}f\ \text{在}\ [a,c]\ \text{与}\ [c,b]\ \text{上均 Riemann 可积}.
\]
现在考虑 $[a,c]$ 的分割 $T'$, $[c,b]$ 的分割 $T''$, 则 $T=T'+T''$ 为 $[a,b]$ 的分割, $c$ 为 $T'$, $T''$, $T$ 的分点. 于是
\begin{align*}
\int_{a}^{b}f(x)\,\mathrm{d}x&=\lim_{\|T\|\to 0}\sum_{T}f(\xi_{i})\Delta x_{i}\\
&=\lim_{\|T'\|\to 0}\sum_{T'}f(\xi_{i})\Delta x_{i}+\lim_{\|T''\|\to 0}\sum_{T''}f(\xi_{i})\Delta x_{i}\\
&=\int_{a}^{c}f(x)\,\mathrm{d}x+\int_{c}^{b}f(x)\,\mathrm{d}x.
\end{align*}
\end{proof}

\begin{proof}[证法 2]
($\Rightarrow$) 因为 $f$ 在 $[a,b]$ 上 Riemann 可积, 根据例~\ref{ex:6.1.8}, $f$ 在 $[a,c]$ 与 $[c,b]$ 上均 Riemann 可积.

($\Leftarrow$) 设 $f$ 在 $[a,c]$ 与 $[c,b]$ 上均 Riemann 可积, 则对 $\forall\varepsilon>0$, $\exists\delta_{1},\delta_{2}>0$, s. t. $\|T'\|<\delta_{1}$, $\|T''\|<\delta_{2}$, 有
\[
\left|\sum_{T'}f(\xi_{i})\Delta x_{i}-\int_{a}^{c}f(x)\,\mathrm{d}x\right|<\frac{\varepsilon}{3},
\]
\[
\left|\sum_{T''}f(\xi_{i})\Delta x_{i}-\int_{c}^{b}f(x)\,\mathrm{d}x\right|<\frac{\varepsilon}{3},
\]
其中 $T'$ 与 $T''$ 分别为 $[a,c]$ 与 $[c,b]$ 的分割. 令 $T$ 为 $[a,b]$ 的分割,\linebreak[4]$\|T\|<\delta=\min\left\{\delta_{1},\delta_{2},\dfrac{\varepsilon}{9M+1}\right\}$, $M=\sup\limits_{x\in[a,b]}|f(x)|$. 则
\begin{align*}
&\left|\sum_{T}f(\xi_{i})\Delta x_{i}-\left[\int_{a}^{c}f(x)\,\mathrm{d}x+\int_{c}^{b}f(x)\,\mathrm{d}x\right]\right|\\
&\leqslant\left|\sum_{T'}f(\xi_{i})\Delta x_{i}-\int_{a}^{c}f(x)\,\mathrm{d}x\right|+\left|\sum_{T''}f(\xi_{i})\Delta x_{i}-\int_{c}^{b}f(x)\,\mathrm{d}x\right|+3M\|T\|\\
&<\frac{\varepsilon}{3}+\frac{\varepsilon}{3}+3M\cdot\frac{\varepsilon}{9M+1}<\varepsilon.
\end{align*}
因而, $f$ 在 $[a,b]$ 上 Riemann 可积, 且
\[
\int_{a}^{b}f(x)\,\mathrm{d}x=\lim_{\|T\|\to 0}\sum_{T}f(\xi_{i})\Delta x_{i}=\int_{a}^{c}f(x)\,\mathrm{d}x+\int_{c}^{b}f(x)\,\mathrm{d}x.
\]
\end{proof}

\begin{remark}\label{remark:6.2.1}
对于任何 $a,b,c\in\mathbb{R}$, 都有
\[
\int_{a}^{b}f(x)\,\mathrm{d}x=\int_{a}^{c}f(x)\,\mathrm{d}x+\int_{c}^{b}f(x)\,\mathrm{d}x.
\]
例如: 当 $a<b<c$ 时, $f$ 在 $[a,c]$ 上 Riemann 可积 $\Leftrightarrow f$ 在 $[a,b]$ 与 $[b,c]$ 上均 Riemann 可积. 并由定理~\ref{thm:6.2.4}, 有
\[
\int_{a}^{c}f(x)\,\mathrm{d}x=\int_{a}^{b}f(x)\,\mathrm{d}x+\int_{b}^{c}f(x)\,\mathrm{d}x,
\]
\[
\int_{a}^{b}f(x)\,\mathrm{d}x=\int_{a}^{c}f(x)\,\mathrm{d}x-\int_{b}^{c}f(x)\,\mathrm{d}x=\int_{a}^{c}f(x)\,\mathrm{d}x+\int_{c}^{b}f(x)\,\mathrm{d}x.
\]
\end{remark}

\begin{theorem}[积分第一中值定理]\label{thm:6.2.5}
设 $f$ 在 $[a,b]$ 上连续, 则必存在 $\xi\in[a,b]$, 使得 (图~\ref{fig:6.2.1})
\[
\int_{a}^{b}f(x)\,\mathrm{d}x=f(\xi)(b-a).
\]
\end{theorem}

\begin{figure}[htbp]
\centering
\begin{tikzpicture}[x=1.0cm,y=0.92cm,>=Stealth]
  \draw[->] (-0.3,0) -- (7.9,0) node[below right] {$x$};
  \draw[->] (0,-0.3) -- (0,4.85) node[left] {$y$};
  \draw (0,2.60) -- (6.6,2.60);
  \fill[pattern=north east lines] (1.6,2.60) -- (3.6,2.60)
    .. controls (3.0,2.28) and (2.2,1.90) .. (1.6,1.60) -- cycle;
  \fill[pattern=crosshatch] (3.6,2.60)
    .. controls (4.1,2.90) and (4.6,3.20) .. (5.0,3.40)
    .. controls (5.8,3.80) and (6.2,4.10) .. (6.6,4.35)
    -- (6.6,2.60) -- cycle;
  \draw[very thick] (1.6,1.60) .. controls (2.2,1.90) and (3.0,2.28) .. (3.6,2.60)
    .. controls (4.1,2.90) and (4.6,3.20) .. (5.0,3.40)
    .. controls (5.8,3.80) and (6.2,4.10) .. (6.6,4.35);
  \draw (1.6,0) -- (1.6,1.60);
  \draw (6.6,0) -- (6.6,4.35);
  \draw[dashed] (3.6,0) -- (3.6,2.60);
  \node[below left] at (0,0) {$O$};
  \node[left] at (-0.05,2.60) {$f(\xi)$};
  \node[above] at (5.1,4.42) {$y=f(x)$};
  \node[below] at (1.6,-0.06) {$a$};
  \node[below] at (3.6,-0.06) {$\xi$};
  \node[below] at (6.6,-0.06) {$b$};
\end{tikzpicture}
\caption{}\label{fig:6.2.1}
\end{figure}

\begin{proof}
由于 $f$ 在 $[a,b]$ 上连续, 根据最值定理, $f$ 达到最大值 $M$ 与最小值 $m$, 则有 $m\leqslant f(x)\leqslant M$, $x\in[a,b]$.

由定理~\ref{thm:6.2.2}(2) 得到
\begin{align*}
m(b-a)&=\int_{a}^{b}m\,\mathrm{d}x\\
&\leqslant\int_{a}^{b}f(x)\,\mathrm{d}x\\
&\leqslant\int_{a}^{b}M\,\mathrm{d}x=M(b-a),
\end{align*}
即
\[
m\leqslant\frac{1}{b-a}\int_{a}^{b}f(x)\,\mathrm{d}x\leqslant M.
\]
由连续函数的介值定理可知, $\exists\xi\in[a,b]$, s. t.
\[
f(\xi)=\frac{1}{b-a}\int_{a}^{b}f(x)\,\mathrm{d}x,
\]
也就是 $\int_{a}^{b}f(x)\,\mathrm{d}x=f(\xi)(b-a)$, $\xi\in[a,b]$.
\end{proof}

$f(\xi)=\frac{1}{b-a}\int_{a}^{b}f(x)\,\mathrm{d}x$ 称为 $f$ 在区间 $[a,b]$ 上的平均值, 记作 $M(f)$.

\begin{theorem}[推广的积分第一中值定理]\label{thm:6.2.6}
设 $f$ 与 $g$ 都在 $[a,b]$ 上连续, 且 $g(x)$ 在 $[a,b]$ 上不变号, 则 $\exists\xi\in[a,b]$, s. t.
\[
\int_{a}^{b}f(x)g(x)\,\mathrm{d}x=f(\xi)\int_{a}^{b}g(x)\,\mathrm{d}x.
\]
特别地, 当 $g(x)\equiv 1$ 时, 它就是定理~\ref{thm:6.2.5}.
\end{theorem}

\begin{proof}
不妨设 $g(x)\geqslant 0$, $x\in[a,b]$, 此时有
\[
mg(x)\leqslant f(x)g(x)\leqslant Mg(x),\quad x\in[a,b],
\]
其中 $M$ 与 $m$ 分别为 $f$ 在 $[a,b]$ 上的最大值与最小值. 根据定理~\ref{thm:6.2.2}(2),
\[
m\int_{a}^{b}g(x)\,\mathrm{d}x\leqslant\int_{a}^{b}f(x)g(x)\,\mathrm{d}x\leqslant M\int_{a}^{b}g(x)\,\mathrm{d}x.
\]

(1) 若 $\int_{a}^{b}g(x)\,\mathrm{d}x=0$, 则由上式知,
\[
0\leqslant\int_{a}^{b}f(x)g(x)\,\mathrm{d}x\leqslant 0,
\]
\[
\int_{a}^{b}f(x)g(x)\,\mathrm{d}x=0=f(\xi)\cdot 0=f(\xi)\int_{a}^{b}g(x)\,\mathrm{d}x,\quad\forall\xi\in(a,b).
\]

(2) 若 $\int_{a}^{b}g(x)\,\mathrm{d}x>0$, 则
\[
m\leqslant\frac{\displaystyle\int_{a}^{b}f(x)g(x)\,\mathrm{d}x}{\displaystyle\int_{a}^{b}g(x)\,\mathrm{d}x}\leqslant M.
\]
由连续函数的介值定理, $\exists\xi\in[a,b]$, s. t.
\[
f(\xi)=\frac{\displaystyle\int_{a}^{b}f(x)g(x)\,\mathrm{d}x}{\displaystyle\int_{a}^{b}g(x)\,\mathrm{d}x},
\]
即
\[
\int_{a}^{b}f(x)g(x)\,\mathrm{d}x=f(\xi)\int_{a}^{b}g(x)\,\mathrm{d}x.
\]
\end{proof}

\begin{theorem}[积分第二中值定理]\label{thm:6.2.7}
设 $f$ 在 $[a,b]$ 上 Riemann 可积.

(1) 如果 $g$ 在 $[a,b]$ 上单调减, 且 $g(x)\geqslant 0$, $x\in[a,b]$, 则 $\exists\xi\in[a,b]$, s. t.
\[
\int_{a}^{b}f(x)g(x)\,\mathrm{d}x=g(a)\int_{a}^{\xi}f(x)\,\mathrm{d}x;
\]

(2) 如果 $g$ 在 $[a,b]$ 上单调增, 且 $g(x)\geqslant 0$, $x\in[a,b]$, 则 $\exists\eta\in[a,b]$, s. t.
\[
\int_{a}^{b}f(x)g(x)\,\mathrm{d}x=g(b)\int_{\eta}^{b}f(x)\,\mathrm{d}x.
\]
\end{theorem}

\begin{proof}
(1) 设
\[
F(x)=\int_{a}^{x}f(t)\,\mathrm{d}t,\quad x\in[a,b].
\]
由于 $f$ 在 $[a,b]$ 上 Riemann 可积, 故从定理~\ref{thm:6.3.1} 知, $F$ 在 $[a,b]$ 上连续, 因此, 它达到最大值 $M$ 和最小值 $m$.

设 $g(a)=0$. 因 $g(x)\geqslant 0$ 且 $g$ 单调减, 必有 $g(x)\equiv 0$, $x\in[a,b]$, 则
\begin{align*}
\int_{a}^{b}f(x)g(x)\,\mathrm{d}x&=\int_{a}^{b}f(x)\cdot 0\,\mathrm{d}x=0=0\cdot\int_{a}^{\xi}f(x)\,\mathrm{d}x\\
&=g(a)\int_{a}^{\xi}f(x)\,\mathrm{d}x.
\end{align*}

设 $g(a)>0$. 因 $f$ 在 $[a,b]$ 上 Riemann 可积, 根据定理~\ref{thm:6.1.2}, $f$ 在 $[a,b]$ 上有界, 设 $|f(x)|<L$, $x\in[a,b]$. 又因 $g$ 单调减, 故 $g$ 在 $[a,b]$ 上 Riemann 可积, 从而 $\forall\varepsilon>0$, 必有 $[a,b]$ 的分割 $T:a=x_{0}<x_{1}<\cdots<x_{n}=b$, s. t.
\[
\sum_{T}\omega_{i}^{g}\Delta x_{i}<\frac{\varepsilon}{L}.
\]
于是,
\begin{align*}
\int_{a}^{b}f(x)g(x)\,\mathrm{d}x&=\sum_{i=1}^{n}\int_{x_{i-1}}^{x_{i}}f(x)g(x)\,\mathrm{d}x\\
&=\sum_{i=1}^{n}\int_{x_{i-1}}^{x_{i}}[g(x)-g(x_{i-1})]f(x)\,\mathrm{d}x+\sum_{i=1}^{n}g(x_{i-1})\int_{x_{i-1}}^{x_{i}}f(x)\,\mathrm{d}x\\
&\leqslant\sum_{i=1}^{n}\int_{x_{i-1}}^{x_{i}}|g(x)-g(x_{i-1})|\cdot|f(x)|\,\mathrm{d}x\\
&\quad+\sum_{i=1}^{n}g(x_{i-1})[F(x_{i})-F(x_{i-1})]\\
&\leqslant L\sum_{i=1}^{n}\omega_{i}^{g}\Delta x_{i}+\sum_{i=1}^{n-1}F(x_{i})[g(x_{i-1})-g(x_{i})]+F(b)g(x_{n-1})\\
&<L\cdot\frac{\varepsilon}{L}+M\sum_{i=1}^{n-1}[g(x_{i-1})-g(x_{i})]+Mg(x_{n-1})=\varepsilon+Mg(a).
\end{align*}
类似地, 有
\begin{align*}
\int_{a}^{b}f(x)g(x)\,\mathrm{d}x&\geqslant-\sum_{i=1}^{n}\int_{x_{i-1}}^{x_{i}}|g(x)-g(x_{i-1})|\cdot|f(x)|\,\mathrm{d}x\\
&\quad+\sum_{i=1}^{n}g(x_{i-1})[F(x_{i})-F(x_{i-1})]\\
&\geqslant-L\cdot\sum_{i=1}^{n}\omega_{i}^{g}\Delta x_{i}+\sum_{i=1}^{n-1}F(x_{i})[g(x_{i-1})-g(x_{i})]+F(b)g(x_{n-1})\\
&>-L\cdot\frac{\varepsilon}{L}+m\sum_{i=1}^{n-1}[g(x_{i-1})-g(x_{i})]+mg(x_{n-1})=-\varepsilon+mg(a).
\end{align*}
综合上述得到
\[
mg(a)-\varepsilon<\int_{a}^{b}f(x)g(x)\,\mathrm{d}x<Mg(a)+\varepsilon.
\]
令 $\varepsilon\to 0^{+}$, 有
\[
mg(a)\leqslant\int_{a}^{b}f(x)g(x)\,\mathrm{d}x\leqslant Mg(a),
\]
\[
m\leqslant\frac{\displaystyle\int_{a}^{b}f(x)g(x)\,\mathrm{d}x}{g(a)}\leqslant M.
\]
根据连续函数的介值定理, $\exists\xi\in[a,b]$, s. t.
\[
\int_{a}^{\xi}f(x)\,\mathrm{d}x=F(\xi)=\frac{1}{g(a)}\int_{a}^{b}f(x)g(x)\,\mathrm{d}x,
\]
即
\[
\int_{a}^{b}f(x)g(x)\,\mathrm{d}x=g(a)\int_{a}^{\xi}f(x)\,\mathrm{d}x.
\]

(2) 设
\[
\widetilde{F}(x)=\int_{x}^{b}f(t)\,\mathrm{d}t,\quad x\in[a,b].
\]
由于 $f$ 在 $[a,b]$ 上 Riemann 可积, 故从定理~\ref{thm:6.3.1} 知, $\widetilde{F}$ 在 $[a,b]$ 上连续, 因此, 它达到最大值 $M$ 与最小值 $m$.

设 $g(b)=0$. 因 $g(x)\geqslant 0$ 且 $g$ 单调增, 必有 $g(x)\equiv 0$, $x\in[a,b]$, 则
\begin{align*}
\int_{a}^{b}f(x)g(x)\,\mathrm{d}x&=\int_{a}^{b}f(x)\cdot 0\,\mathrm{d}x=0=0\cdot\int_{\eta}^{b}f(x)\,\mathrm{d}x\\
&=g(b)\int_{\eta}^{b}f(x)\,\mathrm{d}x.
\end{align*}

设 $g(b)>0$. 因 $f$ 在 $[a,b]$ 上 Riemann 可积, 根据定理~\ref{thm:6.1.2}, $f$ 在 $[a,b]$ 上有界, 设 $|f(x)|<L$, $x\in[a,b]$, 又因为 $g$ 单调增, 故 $g$ 在 $[a,b]$ 上 Riemann 可积, 从而对 $\forall\varepsilon>0$, 必有 $[a,b]$ 的分割 $T:a=x_{0}<x_{1}<\cdots<x_{n}=b$, s. t.
\[
\sum_{T}\omega_{i}^{g}\Delta x_{i}<\frac{\varepsilon}{L}.
\]
于是,
\begin{align*}
\int_{a}^{b}f(x)g(x)\,\mathrm{d}x&=\sum_{i=1}^{n}\int_{x_{i-1}}^{x_{i}}f(x)g(x)\,\mathrm{d}x\\
&=\sum_{i=1}^{n}\int_{x_{i-1}}^{x_{i}}[g(x)-g(x_{i})]f(x)\,\mathrm{d}x+\sum_{i=1}^{n}g(x_{i})\int_{x_{i-1}}^{x_{i}}f(x)\,\mathrm{d}x\\
&\leqslant\sum_{i=1}^{n}\int_{x_{i-1}}^{x_{i}}|g(x)-g(x_{i-1})|\cdot|f(x)|\,\mathrm{d}x\\
&\quad+\sum_{i=1}^{n}g(x_{i})[\widetilde{F}(x_{i-1})-\widetilde{F}(x_{i})]\\
&\leqslant L\sum_{i=1}^{n}\omega_{i}^{g}\Delta x_{i}+g(x_{1})\widetilde{F}(x_{0})+\sum_{i=1}^{n-1}\widetilde{F}(x_{i})[g(x_{i+1})-g(x_{i})]\\
&<L\cdot\frac{\varepsilon}{L}+Mg(x_{1})+M\sum_{i=1}^{n-1}[g(x_{i+1})-g(x_{i})]<\varepsilon+Mg(x_{n})=\varepsilon+Mg(b).
\end{align*}
类似地, 有
\begin{align*}
\int_{a}^{b}f(x)g(x)\,\mathrm{d}x&\geqslant-\sum_{i=1}^{n}\int_{x_{i-1}}^{x_{i}}|g(x)-g(x_{i-1})|\cdot|f(x)|\,\mathrm{d}x\\
&\quad+\sum_{i=1}^{n}g(x_{i})[\widetilde{F}(x_{i-1})-\widetilde{F}(x_{i})]\\
&\geqslant-L\sum_{i=1}^{n}\omega_{i}^{g}\Delta x_{i}+g(x_{1})\widetilde{F}(x_{0})+\sum_{i=1}^{n}\widetilde{F}(x_{i})[g(x_{i+1})-g(x_{i})]\\
&>-L\cdot\frac{\varepsilon}{L}+mg(x_{1})+m\sum_{i=1}^{n-1}[g(x_{i+1})-g(x_{i})]>-\varepsilon+mg(b).
\end{align*}
综合上述得到
\[
mg(b)-\varepsilon<\int_{a}^{b}f(x)g(x)\,\mathrm{d}x<Mg(b)+\varepsilon.
\]
令 $\varepsilon\to 0^{+}$, 有
\[
mg(b)\leqslant\int_{a}^{b}f(x)g(x)\,\mathrm{d}x\leqslant Mg(b),
\]
\[
m\leqslant\frac{\displaystyle\int_{a}^{b}f(x)g(x)\,\mathrm{d}x}{g(b)}\leqslant M.
\]
根据连续函数的介值定理可知, $\exists\eta\in[a,b]$, s. t.
\[
\int_{\eta}^{b}f(x)\,\mathrm{d}x=\widetilde{F}(\eta)=\frac{1}{g(b)}\int_{a}^{b}f(x)g(x)\,\mathrm{d}x,
\]
即
\[
\int_{a}^{b}f(x)g(x)\,\mathrm{d}x=g(b)\int_{\eta}^{b}f(x)\,\mathrm{d}x.
\]
\end{proof}

\begin{theorem}[一般的积分第二中值定理]\label{thm:6.2.8}
设 $f$ 在 $[a,b]$ 上 Riemann 可积, $g$ 为单调函数, 则 $\exists\xi\in[a,b]$, s. t.
\[
\int_{a}^{b}f(x)g(x)\,\mathrm{d}x=g(a)\int_{a}^{\xi}f(x)\,\mathrm{d}x+g(b)\int_{\xi}^{b}f(x)\,\mathrm{d}x.
\]
\end{theorem}

\begin{proof}
(1) 如果 $g$ 单调减, 令 $\widetilde{g}(x)=g(x)-g(b)$, 则 $\widetilde{g}\geqslant 0$, 且 $\widetilde{g}$ 仍单调减. 根据积分第二中值定理 (1), $\exists\xi\in[a,b]$, s. t.
\begin{align*}
\int_{a}^{b}f(x)g(x)\,\mathrm{d}x-g(b)\int_{a}^{b}f(x)\,\mathrm{d}x&=\int_{a}^{b}f(x)[g(x)-g(b)]\,\mathrm{d}x\\
&=\int_{a}^{b}f(x)\widetilde{g}(x)\,\mathrm{d}x=\widetilde{g}(a)\int_{a}^{\xi}f(x)\,\mathrm{d}x\\
&=[g(a)-g(b)]\int_{a}^{\xi}f(x)\,\mathrm{d}x.
\end{align*}
因此,
\[
\int_{a}^{b}f(x)g(x)\,\mathrm{d}x=g(a)\int_{a}^{\xi}f(x)\,\mathrm{d}x+g(b)\int_{\xi}^{b}f(x)\,\mathrm{d}x.
\]

(2) 如果 $g$ 单调增, 令 $\widetilde{g}(x)=g(x)-g(a)$, 则 $\widetilde{g}\geqslant 0$, 且 $\widetilde{g}(x)$ 仍单调增. 根据积分第二中值定理 (2), $\exists\xi\in[a,b]$, s. t.
\begin{align*}
\int_{a}^{b}f(x)g(x)\,\mathrm{d}x-g(a)\int_{a}^{b}f(x)\,\mathrm{d}x&=\int_{a}^{b}f(x)[g(x)-g(a)]\,\mathrm{d}x\\
&=\int_{a}^{b}f(x)\widetilde{g}(x)\,\mathrm{d}x=\widetilde{g}(b)\int_{\xi}^{b}f(x)\,\mathrm{d}x\\
&=[g(b)-g(a)]\int_{\xi}^{b}f(x)\,\mathrm{d}x.
\end{align*}
因此,
\[
\int_{a}^{b}f(x)g(x)\,\mathrm{d}x=g(a)\int_{a}^{\xi}f(x)\,\mathrm{d}x+g(b)\int_{\xi}^{b}f(x)\,\mathrm{d}x.
\]
或者, 令 $\widetilde{g}(x)=g(b)-g(x)$, 由 $g$ 单调增, 可知 $\widetilde{g}(x)$ 单调减, $\widetilde{g}\geqslant 0$. 根据积分第二中值定理 (1), $\exists\xi\in[a,b]$, s. t.
\begin{align*}
\int_{a}^{b}f(x)[g(b)-g(x)]\,\mathrm{d}x&=[g(b)-g(a)]\int_{a}^{\xi}f(x)\,\mathrm{d}x,\\
\int_{a}^{b}f(x)g(x)\,\mathrm{d}x&=g(a)\int_{a}^{\xi}f(x)\,\mathrm{d}x+g(b)\int_{\xi}^{b}f(x)\,\mathrm{d}x.
\end{align*}
或者, 因为 $g$ 单调增, 故 $-g(x)$ 单调减, 由 (1) 知, $\exists\xi\in[a,b]$, s. t.
\begin{align*}
\int_{a}^{b}f(x)[-g(x)]\,\mathrm{d}x&=-g(a)\int_{a}^{\xi}f(x)\,\mathrm{d}x+[-g(b)]\int_{\xi}^{b}f(x)\,\mathrm{d}x,\\
\int_{a}^{b}f(x)g(x)\,\mathrm{d}x&=g(a)\int_{a}^{\xi}f(x)\,\mathrm{d}x+g(b)\int_{\xi}^{b}f(x)\,\mathrm{d}x.
\end{align*}
\end{proof}

\begin{example}\label{ex:6.2.2}
设
\[
f(x)=\begin{cases}\sin x-1,&-1\leqslant x<0,\\x^{2},&0\leqslant x\leqslant 1,\end{cases}
\]
求 $\int_{-1}^{1}f(x)\,\mathrm{d}x$.
\end{example}

\begin{solution}
如图~\ref{fig:6.2.2} 所示.
\begin{align*}
\int_{-1}^{1}f(x)\,\mathrm{d}x&\xlongequal{\text{定理~\ref{thm:6.2.4}}}\int_{-1}^{0}f(x)\,\mathrm{d}x+\int_{0}^{1}f(x)\,\mathrm{d}x\\
&=\int_{-1}^{0}(\sin x-1)\,\mathrm{d}x+\int_{0}^{1}x^{2}\,\mathrm{d}x\\
&=\cos(-1)-\cos 0-[0-(-1)]+\frac{1^{3}-0^{3}}{3}\\
&=\cos 1-1-1+\frac{1}{3}=\cos 1-\frac{5}{3}.
\end{align*}
\end{solution}

\begin{figure}[htbp]
\centering
\begin{tikzpicture}[x=1.15cm,y=0.85cm,>=Stealth]
  \draw[->] (-2.15,0) -- (2.35,0) node[below right] {$x$};
  \draw[->] (0,-2.75) -- (0,1.75) node[left] {$y$};
  \fill[pattern=north east lines] (-1,0) -- (-1,-1.8415)
    .. controls (-0.6,-1.72) and (-0.25,-1.33) .. (0,-1) -- cycle;
  \draw[thick] (-1,-1.8415) .. controls (-0.6,-1.72) and (-0.25,-1.33) .. (0,-1);
  \draw (-1,0) -- (-1,-1.8415);
  \fill[pattern=dots] (0,0) .. controls (0.45,0.20) and (0.80,0.64) .. (1,1) -- (1,0) -- cycle;
  \draw[thick] (0,0) .. controls (0.45,0.20) and (0.80,0.64) .. (1,1);
  \draw (1,0) -- (1,1);
  \fill (0,0) circle (1.5pt);
  \fill (-1,-1.8415) circle (1.5pt);
  \fill (1,1) circle (1.5pt);
  \draw[fill=white] (0,-1) circle (1.5pt);
  \node[below left] at (0,0) {$O$};
  \node[left] at (-1.04,0.12) {$-1$};
  \node[below] at (1,-0.10) {$1$};
  \node[right] at (0.06,-1) {$-1$};
  \node[right] at (0.08,0.47) {$y=x^{2}$};
  \node[left] at (-0.10,-2.02) {$y=\sin x-1$};
\end{tikzpicture}
\caption{}\label{fig:6.2.2}
\end{figure}

\begin{example}\label{ex:6.2.3}
(1) 设 $f$ 在 $[a,b]$ 上连续, 且 $f(x)\geqslant 0$, $\int_{a}^{b}f(x)\,\mathrm{d}x=0$, 则 $f(x)\equiv 0$, $x\in[a,b]$.

(2) 设 $f$ 在 $[a,b]$ 上 Riemann 可积, 且 $f(x)\geqslant 0$, $\int_{a}^{b}f(x)\,\mathrm{d}x=0$, 则 $f(x)$ 在 $[a,b]$ 上几乎处处为 $0$, 即 $\{x\mid x\in[a,b],f(x)\neq 0\}$ 为零测集, 记为 $f(x)\xrightarrow{\text{a. e.}}0$, $x\in[a,b]$.
\end{example}

\begin{solution}
(1) (反证) 假设 $f(x)\not\equiv 0$, $x\in[a,b]$, 则 $\exists x_{0}\in[a,b]$, s. t. $f(x_{0})\neq 0$, 则 $f(x_{0})>0$. 根据定理~\ref{thm:6.2.2}(3), $\int_{a}^{b}f(x)\,\mathrm{d}x>0$, 这与题设 $\int_{a}^{b}f(x)\,\mathrm{d}x=0$ 相矛盾.

(2) 因 $f$ 在 $[a,b]$ 上 Riemann 可积, 根据定理~\ref{thm:6.1.4}(8), $f$ 在 $[a,b]$ 上几乎处处连续, 由定理~\ref{thm:6.2.2}(3), $f$ 在每个连续点 $x_{0}$ 处, $f(x_{0})=0$. 所以, $f(x)$ 在 $[a,b]$ 上几乎处处为 $0$.
\end{solution}

\begin{example}\label{ex:6.2.4}
求 $f(x)=\sin x$ 在 $[0,\pi]$ 上的平均值 $M(f)$.
\end{example}

\begin{solution}
由例~\ref{ex:6.1.14}, 有
\begin{align*}
M(f)&=\frac{1}{\pi}\int_{0}^{\pi}\sin x\,\mathrm{d}x\\
&=\frac{1}{\pi}(\cos 0-\cos\pi)=\frac{1}{\pi}[1-(-1)]=\frac{2}{\pi}.
\end{align*}
\end{solution}

\begin{example}\label{ex:6.2.5}
不计算积分值比较下列积分的大小:

(1) $\int_{0}^{1}x^{3}\,\mathrm{d}x$ 与 $\int_{0}^{1}x^{2}\,\mathrm{d}x$;

(2) $\int_{1}^{2}x^{3}\,\mathrm{d}x$ 与 $\int_{1}^{2}x^{2}\,\mathrm{d}x$;

(3) $\int_{0}^{1}\mathrm{e}^{x^{2}}(1-x)\,\mathrm{d}x$ 与 $\int_{0}^{1}0\,\mathrm{d}x$.
\end{example}

\begin{solution}
(1) 因为 $x^{3}\leqslant x^{2}$, $x\in[0,1]$, 且 $\left(\frac{1}{2}\right)^{3}<\left(\frac{1}{2}\right)^{2}$, $x^{3}$ 与 $x^{2}$ 均连续, 根据定理~\ref{thm:6.2.2}(3), 有
\[
\int_{0}^{1}x^{3}\,\mathrm{d}x<\int_{0}^{1}x^{2}\,\mathrm{d}x.
\]

(2) 因 $x^{3}\geqslant x^{2}$, $x\in[1,2]$, 且 $2^{3}>2^{2}$, $x^{3}$ 与 $x^{2}$ 均连续, 根据定理~\ref{thm:6.2.2}(3), 有
\[
\int_{1}^{2}x^{3}\,\mathrm{d}x>\int_{1}^{2}x^{2}\,\mathrm{d}x.
\]

(3) 因 $\mathrm{e}^{x^{2}}(1-x)\geqslant 0$, $x\in[0,1]$, 且 $\mathrm{e}\left(\frac{1}{2}\right)^{2}\left(1-\frac{1}{2}\right)>0$, $\mathrm{e}^{x^{2}}(1-x)$ 与 $0$ 均连续, 根据定理~\ref{thm:6.2.2}(3), 有
\[
\int_{0}^{1}\mathrm{e}^{x^{2}}(1-x)\,\mathrm{d}x>\int_{0}^{1}0\,\mathrm{d}x.
\]
\end{solution}

\begin{example}\label{ex:6.2.6}
证明: $\frac{2}{\sqrt[4]{\mathrm{e}}}<\int_{0}^{2}\mathrm{e}^{x^{2}-x}\,\mathrm{d}x<2\mathrm{e}^{2}$.
\end{example}

\begin{solution}
令 $f(x)=\mathrm{e}^{x^{2}-x}$, 则由 $-\frac{1}{4}\leqslant x^{2}-x\leqslant 2$, $\forall x\in[0,2]$ 得到
\[
\frac{1}{\sqrt[4]{\mathrm{e}}}=\mathrm{e}^{-\frac{1}{4}}\leqslant\mathrm{e}^{x^{2}-x}\leqslant\mathrm{e}^{2},\quad\forall x\in[0,2].
\]
又因为
\[
\frac{1}{\sqrt[4]{\mathrm{e}}}<f(1)=1<\mathrm{e}^{2},
\]
根据定理~\ref{thm:6.2.2}(3), 有
\[
\frac{2}{\sqrt[4]{\mathrm{e}}}=\int_{0}^{2}\frac{\mathrm{d}x}{\sqrt[4]{\mathrm{e}}}<\int_{0}^{2}\mathrm{e}^{x^{2}-x}\,\mathrm{d}x<\int_{0}^{2}\mathrm{e}^{2}\,\mathrm{d}x=2\mathrm{e}^{2}.
\]
\end{solution}

\begin{example}\label{ex:6.2.7}
设 $f$, $g$ 在 $[a,b]$ 上连续. 证明 Cauchy-Schwarz 不等式:
\[
\left[\int_{a}^{b}f(x)g(x)\,\mathrm{d}x\right]^{2}\leqslant\int_{a}^{b}f^{2}(x)\,\mathrm{d}x\int_{a}^{b}g^{2}(x)\,\mathrm{d}x,
\]
其中等号当且仅当 $f=\lambda g$ 或 $g=\lambda f$ 时成立 $(\lambda\in\mathbb{R})$.
\end{example}

\begin{solution}[证法 1]
$\forall t\in\mathbb{R}$, 有
\begin{align*}
0&\leqslant\int_{a}^{b}[f(x)-\tan(x)]^{2}\,\mathrm{d}x\\
&=\int_{a}^{b}f^{2}(x)\,\mathrm{d}x-2t\int_{a}^{b}f(x)g(x)\,\mathrm{d}x+t^{2}\int_{a}^{b}g^{2}(x)\,\mathrm{d}x.
\end{align*}

(1) 如果 $\int_{a}^{b}g^{2}(x)\,\mathrm{d}x=0$, 由例~\ref{ex:6.2.2} 知, $g(x)=0$, $x\in[a,b]$, 故
\begin{align*}
\left[\int_{a}^{b}f(x)g(x)\,\mathrm{d}x\right]^{2}&=\left[\int_{a}^{b}f(x)\cdot 0\,\mathrm{d}x\right]^{2}=0\\
&=\int_{a}^{b}f^{2}(x)\,\mathrm{d}x\cdot 0\\
&=\int_{a}^{b}f^{2}(x)\,\mathrm{d}x\cdot\int_{a}^{b}g^{2}(x)\,\mathrm{d}x.
\end{align*}

(2) 如果 $\int_{a}^{b}g^{2}(x)\,\mathrm{d}x>0$, 则由上述关于 $t$ 的二次三项式非负, 故其判别式
\[
\Delta=4\left[\int_{a}^{b}f(x)g(x)\,\mathrm{d}x\right]^{2}-4\int_{a}^{b}f^{2}(x)\,\mathrm{d}x\int_{a}^{b}g^{2}(x)\,\mathrm{d}x\leqslant 0,
\]
即
\[
\left[\int_{a}^{b}f(x)g(x)\,\mathrm{d}x\right]^{2}\leqslant\int_{a}^{b}f^{2}(x)\,\mathrm{d}x\cdot\int_{a}^{b}g^{2}(x)\,\mathrm{d}x.
\]
当 $f=\lambda g$ 时, 显然有
\begin{align*}
\left[\int_{a}^{b}f(x)g(x)\,\mathrm{d}x\right]^{2}&=\left[\int_{a}^{b}\lambda g(x)g(x)\,\mathrm{d}x\right]^{2}\\
&=\lambda^{2}\left[\int_{a}^{b}g^{2}(x)\,\mathrm{d}x\right]^{2}\\
&=\int_{a}^{b}[\lambda g(x)]^{2}\,\mathrm{d}x\cdot\int_{a}^{b}g^{2}(x)\,\mathrm{d}x.
\end{align*}
同理, 当 $g=\lambda f$ 时, 等式也成立.

反之, 如果 $\left[\int_{a}^{b}f(x)g(x)\,\mathrm{d}x\right]^{2}=\int_{a}^{b}f^{2}(x)\,\mathrm{d}x\cdot\int_{a}^{b}g^{2}(x)\,\mathrm{d}x$, 则

① $\int_{a}^{b}g^{2}(x)\,\mathrm{d}x=0$, 由例~\ref{ex:6.2.2} 可知, $g=0=0\cdot f$;

② $\int_{a}^{b}g^{2}(x)\,\mathrm{d}x>0$, 由
\[
\Delta=4\left[\int_{a}^{b}f(x)g(x)\,\mathrm{d}x\right]^{2}-4\int_{a}^{b}f^{2}(x)\,\mathrm{d}x\int_{a}^{b}g^{2}(x)\,\mathrm{d}x=0
\]
立刻知道, 二次三项式
\[
t^{2}\int_{a}^{b}g^{2}(x)\,\mathrm{d}x-2t\int_{a}^{b}f(x)g(x)\,\mathrm{d}x+\int_{a}^{b}f^{2}(x)\,\mathrm{d}x
\]
必有实根 $\lambda$, 即
\begin{align*}
0&=\lambda^{2}\int_{a}^{b}g^{2}(x)\,\mathrm{d}x-2\lambda\int_{a}^{b}f(x)g(x)\,\mathrm{d}x+\int_{a}^{b}f^{2}(x)\,\mathrm{d}x\\
&=\int_{a}^{b}[f(x)-\lambda g(x)]^{2}\,\mathrm{d}x,
\end{align*}
再由例~\ref{ex:6.2.2} 知, $f(x)-\lambda g(x)=0$, 即 $f=\lambda g$.
\end{solution}

\begin{solution}[证法 2]
因 $f$, $g$ 在 $[a,b]$ 上连续, 故 $fg$, $f^{2}$, $g^{2}$ 在 $[a,b]$ 上也连续. 从而它们在 $[a,b]$ 上都 Riemann 可积. 对 $[a,b]$ 的任何分割 $T:a=x_{0}<x_{1}<\cdots<x_{n}=b$, 有 Cauchy-Schwarz 不等式:
\begin{align*}
\left(\sum_{i=1}^{n}f(\xi_{i})g(\xi_{i})\Delta x_{i}\right)^{2}&=\left(\sum_{i=1}^{n}f(\xi_{i})\sqrt{\Delta x_{i}}\cdot g(\xi_{i})\sqrt{\Delta x_{i}}\right)^{2}\\
&\leqslant\sum_{i=1}^{n}f^{2}(\xi_{i})\Delta x_{i}\cdot\sum_{i=1}^{n}g^{2}(\xi_{i})\Delta x_{i}.
\end{align*}
令 $\|T\|\to 0$ 得到
\[
\left[\int_{a}^{b}f(x)g(x)\,\mathrm{d}x\right]^{2}\leqslant\int_{a}^{b}f^{2}(x)\,\mathrm{d}x\cdot\int_{a}^{b}g^{2}(x)\,\mathrm{d}x.
\]
\end{solution}

\begin{remark}\label{remark:6.2.2}
根据例~\ref{ex:6.2.3}(2), 请读者证明:

设 $f$, $g$ 在 $[a,b]$ 上 Riemann 可积, 则有 Cauchy-Schwarz 不等式
\[
\left[\int_{a}^{b}f(x)g(x)\,\mathrm{d}x\right]^{2}\leqslant\int_{a}^{b}f^{2}(x)\,\mathrm{d}x\cdot\int_{a}^{b}g^{2}(x)\,\mathrm{d}x,
\]
其中等号当且仅当 $f\xlongequal{\text{a. e.}}\lambda g$ 或 $g\xlongequal{\text{a. e.}}\lambda f$, $x\in[a,b]$ 时成立 $(\lambda\in\mathbb{R})$.
\end{remark}

\begin{example}\label{ex:6.2.8}
设 $f$ 在 $[a,b]$ 上连续, 且 $\int_{a}^{b}f(x)\,\mathrm{d}x=\int_{a}^{b}xf(x)\,\mathrm{d}x=0$, 则在 $(a,b)$ 内至少存在两点 $x_{1},x_{2}$, 使得 $f(x_{1})=f(x_{2})=0$.
\end{example}

\begin{solution}[证法 1]
如果 $f$ 在 $(a,b)$ 内无零点, 由 $f$ 连续与零值定理, $f$ 在 $(a,b)$ 内恒大于 $0$ 或恒小于 $0$. 不妨设 $f$ 在 $(a,b)$ 内恒大于 $0$. 再由 $f$ 在 $[a,b]$ 上连续知, $f(x)\geqslant 0$, $\forall x\in[a,b]$. 因此,
\[
\int_{a}^{b}f(x)\,\mathrm{d}x>0,
\]
这与题设 $\int_{a}^{b}f(x)\,\mathrm{d}x=0$ 相矛盾.

如果 $f$ 在 $(a,b)$ 内恰有一个零点 $x_{0}$, 则 $f(x)$ 在 $(a,x_{0})$ 与 $(x_{0},b)$ 内严格异号, 从而 $(x-x_{0})f(x)$ 在 $(a,x_{0})$ 与 $(x_{0},b)$ 内必严格同号, 因此,
\[
\int_{a}^{b}(x-x_{0})f(x)\,\mathrm{d}x>0\ (\text{或}<0),
\]
这与题设 $\int_{a}^{b}(x-x_{0})f(x)\,\mathrm{d}x=\int_{a}^{b}xf(x)\,\mathrm{d}x-x_{0}\int_{a}^{b}f(x)\,\mathrm{d}x=0-x_{0}\cdot 0=0$ 相矛盾.

综合上述得到 $f$ 在 $(a,b)$ 内至少有两个零点 $x_{1},x_{2}$, 即 $x_{1}\neq x_{2}$, 且 $f(x_{1})=0=f(x_{2})$.
\end{solution}

\begin{solution}[证法 2]
由 $\int_{a}^{b}f(x)\,\mathrm{d}x=0$ 推得 $f$ 在 $(a,b)$ 内必存在零点. (反证) 假设 $f$ 在 $(a,b)$ 内只有一个零点 $x_{0}$, 即 $f(x_{0})=0$. 则由
\[
0=\int_{a}^{b}f(x)\,\mathrm{d}x=\int_{a}^{x_{0}}f(x)\,\mathrm{d}x+\int_{x_{0}}^{b}f(x)\,\mathrm{d}x
\]
可得
\[
\int_{a}^{x_{0}}f(x)\,\mathrm{d}x=-\int_{x_{0}}^{b}f(x)\,\mathrm{d}x\neq 0
\]
(最后一个不等号是因为 $f$ 在 $(x_{0},b)$ 上严格同号). 又因为 $f$ 在 $[a,x_{0}]$ 与 $[x_{0},b]$ 内都不变号, 故由推广的积分第一中值定理可知, $\exists\xi_{1},\xi_{2}$, s. t. $a<\xi_{1}<x_{0}<\xi_{2}<b$, 且
\begin{align*}
0&=\int_{a}^{b}xf(x)\,\mathrm{d}x=\int_{a}^{x_{0}}xf(x)\,\mathrm{d}x+\int_{x_{0}}^{b}xf(x)\,\mathrm{d}x\\
&=\xi_{1}\int_{a}^{x_{0}}f(x)\,\mathrm{d}x+\xi_{2}\int_{x_{0}}^{b}f(x)\,\mathrm{d}x=(\xi_{2}-\xi_{1})\int_{x_{0}}^{b}f(x)\,\mathrm{d}x\neq 0,
\end{align*}
矛盾.
\end{solution}

\begin{exblock}{练习题 6.2}
  \item 证明下列不等式:
  \begin{enumerate}[label=(\arabic*)]
    \item $\int_{0}^{2\pi}|a\cos x+b\sin x|\,\mathrm{d}x\leqslant 2\pi\sqrt{a^{2}+b^{2}}$;
    \item $\int_{0}^{1}x^{m}(1-x^{n})\,\mathrm{d}x\leqslant\frac{m^{m}n^{n}}{(m+n)^{m+n}}$, \quad $m,n\in\mathbb{N}$;
    \item $\int_{0}^{10}\frac{x}{x^{3}+16}\,\mathrm{d}x\leqslant\frac{5}{6}$;
    \item $1<\int_{0}^{\frac{\pi}{2}}\frac{\sin x}{x}\,\mathrm{d}x<\frac{\pi}{2}$;
    \item $1<\int_{0}^{1}\mathrm{e}^{x^{2}}\,\mathrm{d}x<\mathrm{e}$;
    \item $3\sqrt{\mathrm{e}}<\int_{\mathrm{e}}^{4\mathrm{e}}\frac{\ln x}{\sqrt{x}}\,\mathrm{d}x<6$;
    \item $\frac{\pi}{2}<\int_{0}^{\frac{\pi}{2}}\frac{\mathrm{d}x}{\sqrt{1-\frac{1}{2}\sin^{2}x}}<\frac{\pi}{\sqrt{2}}$.
  \end{enumerate}

  \item 求下列极限:
  \begin{enumerate}[label=(\arabic*)]
    \item $\lim\limits_{n\to+\infty}\int_{a}^{b}\mathrm{e}^{-nx^{2}}\,\mathrm{d}x$, \quad $0<a<b$;
    \item $\lim\limits_{n\to+\infty}\int_{0}^{1}\frac{x^{n}}{1+x}\,\mathrm{d}x$.
  \end{enumerate}

  \item 确定下列积分的正负:
  \begin{enumerate}[label=(\arabic*)]
    \item $\int_{\frac{1}{2}}^{1}\mathrm{e}^{x}\ln^{3}x\,\mathrm{d}x$;
    \item $\int_{0}^{\pi}f(x)\,\mathrm{d}x$, \quad $f(x)=\begin{cases}\dfrac{\sin x}{x},&x\neq 0,\\1,&x=0.\end{cases}$
  \end{enumerate}

  \item 比较下列积分的大小:
  \begin{enumerate}[label=(\arabic*)]
    \item $\int_{0}^{1}\mathrm{e}^{-x}\,\mathrm{d}x$ 与 $\int_{0}^{1}\mathrm{e}^{-x^{2}}\,\mathrm{d}x$;
    \item $\int_{0}^{1}\frac{\sin x}{1+x}\,\mathrm{d}x$ 与 $\int_{0}^{1}\frac{\sin x}{1+x^{2}}\,\mathrm{d}x$;
    \item $\int_{0}^{\frac{\pi}{2}}\frac{\sin x}{x}\,\mathrm{d}x$ 与 $\int_{0}^{\frac{\pi}{2}}\frac{\sin^{2}x}{x^{2}}\,\mathrm{d}x$;
    \item $\int_{-1}^{0}\mathrm{e}^{-x^{2}}\,\mathrm{d}x$ 与 $\int_{0}^{1}\mathrm{e}^{-x^{2}}\,\mathrm{d}x$;
    \item $\int_{0}^{1}x\,\mathrm{d}x$ 与 $\int_{0}^{1}x^{2}\,\mathrm{d}x$;
    \item $\int_{0}^{\frac{\pi}{2}}x\,\mathrm{d}x$ 与 $\int_{0}^{\frac{\pi}{2}}\sin x\,\mathrm{d}x$.
  \end{enumerate}

  \item 设函数 $f$ 与 $g$ 在 $[a,b]$ 上连续, 且 $f(x)\leqslant g(x)$, $\forall x\in[a,b]$, 又 $\int_{a}^{b}f(x)\,\mathrm{d}x=\int_{a}^{b}g(x)\,\mathrm{d}x$, 证明: $f=g$.

  \item 设函数 $f$ 在 $[a,b]$ 上连续, 且 $\int_{a}^{b}f(x)g(x)\,\mathrm{d}x=0$ 对一切连续函数 $g$ 成立. 证明: $f=0$.

  \item 证明 Dirichlet 积分
  \[
  \int_{0}^{\pi}\frac{\sin\left(n+\frac{1}{2}\right)x}{\sin\frac{x}{2}}\,\mathrm{d}x=\pi,\quad n=0,1,2,\cdots.
\]
  (提示: 将被积函数表示为余弦函数之和, 或用数学归纳法.)

  \item 证明: Fejèr 积分
  \[
  \int_{0}^{\pi}\left(\frac{\sin\frac{nx}{2}}{\sin\frac{x}{2}}\right)^{2}\,\mathrm{d}x=n\pi,\quad n=0,1,2,\cdots.
\]
  (提示: 将被积函数表示为题 7 中被积函数之和.)

  \item 设 $f$ 为 $[0,1]$ 上的连续函数, 且 $f>0$. 证明:
  \[
  \int_{0}^{1}f(x)\,\mathrm{d}x\cdot\int_{0}^{1}\frac{1}{f(x)}\,\mathrm{d}x\geqslant 1.
\]

  \item 应用 Riemann 积分 $\int_{1}^{n+1}\frac{\mathrm{d}x}{x}<\sum\limits_{i=1}^{n}\frac{1}{i}$ 与 $\sum\limits_{i=1}^{n-1}\frac{1}{i+1}<\int_{1}^{n}\frac{\mathrm{d}x}{x}$ 证明:
  \[
  \ln(1+n)<1+\frac{1}{2}+\cdots+\frac{1}{n}<1+\ln n.
\]
\end{exblock}

\begin{exblock}{思考题 6.2}
  \item 证明:
  \[
  \int_{0}^{\frac{\pi}{2}}\mathrm{e}^{-R\sin x}\,\mathrm{d}x
  \begin{cases}<\dfrac{\pi}{2R}(1-\mathrm{e}^{-R}),&R>0,\\[3mm]>\dfrac{\pi}{2R}(1-\mathrm{e}^{-R}),&R<0,\\[3mm]=\dfrac{\pi}{2},&R=0.\end{cases}
\]

  \item 设 $f$ 为 $[0,\pi]$ 上的连续函数, 它满足
  \[
  \int_{0}^{\pi}f(\theta)\cos\theta\,\mathrm{d}\theta=\int_{0}^{\pi}f(\theta)\sin\theta\,\mathrm{d}\theta.
\]
  证明: $f$ 在 $(0,\pi)$ 内至少有两个零点.

  \item (1) 设 $f(x)$ 在 $[a,b]$ 上连续, 且 $\int_{a}^{b}f(x)\,\mathrm{d}x=0$, $\int_{a}^{b}xf(x)\,\mathrm{d}x=0$. 证明: 至少存在两点 $x_{1},x_{2}\in(a,b)$ 使得 $f(x_{1})=f(x_{2})=0$.

  (2) 设 $f(x)$ 为定义在 $[a,b]$ 上的连续函数. 如果它的前 $n$ 个矩全为 $0$, 即
  \[
  \int_{a}^{b}f(x)\,\mathrm{d}x=\int_{a}^{b}xf(x)\,\mathrm{d}x=\cdots=\int_{a}^{b}x^{n-1}f(x)\,\mathrm{d}x=0.
\]
  证明: 函数 $f(x)$ 或者恒为 $0$ 或者在 $(a,b)$ 上至少改变 $n$ 次符号, 即 $f$ 至少有 $n$ 个零点.

  \item 设 $f(x)$ 在 $[a,b]$ 上 Riemann 可积, 且它的一切矩全为 $0$, 即
  \[
  \int_{a}^{b}f(x)x^{n}\,\mathrm{d}x=0,\quad n=0,1,2,\cdots.
\]
  证明: $f$ 在每一连续点处为 $0$. 特别当 $f$ 连续时, $f=0$.

  \item 设函数 $f$ 在 $[0,+\infty)$ 上连续, 且 $\lim\limits_{x\to+\infty}f(x)=a$. 证明:
  \[
  \lim_{x\to+\infty}\frac{1}{x}\int_{0}^{x}f(t)\,\mathrm{d}t=a.
\]

  \item 设函数 $f$ 在 $[0,\pi]$ 上连续, $n\in\mathbb{N}$. 证明:
  \[
  \lim_{n\to+\infty}\int_{0}^{\pi}f(x)|\sin nx|\,\mathrm{d}x=\frac{2}{\pi}\int_{0}^{\pi}f(x)\,\mathrm{d}x.
\]

  \item 设 $f$ 是 $(-\infty,+\infty)$ 上的周期为 $T$ 的连续函数. 证明:
  \[
  \lim_{x\to+\infty}\frac{1}{x}\int_{0}^{x}f(t)\,\mathrm{d}t=\frac{1}{T}\int_{0}^{T}f(t)\,\mathrm{d}t.
\]

  \item (1) 设 $f$ 与 $g$ 分别为区间 $[a,b]$ 上的非负与正值连续函数. 证明:
  \[
  \lim_{n\to+\infty}\sqrt[n]{\int_{a}^{b}g(x)[f(x)]^{n}\,\mathrm{d}x}=\max_{a\leqslant x\leqslant b}f(x).
\]

  (2) 设 $f$ 与 $g$ 皆为 $[a,b]$ 上的正值连续函数. 证明:
  \[
  \lim_{n\to+\infty}\frac{\displaystyle\int_{a}^{b}g(x)f^{n+1}(x)\,\mathrm{d}x}{\displaystyle\int_{a}^{b}g(x)f^{n}(x)\,\mathrm{d}x}=\max_{a\leqslant x\leqslant b}f(x),
\]
  其中 $f^{n}(x)=[f(x)]^{n}$.

  \item 进一步证明: 积分第一中值定理~\ref{thm:6.2.5} 与推广积分第一中值定理~\ref{thm:6.2.6} 中, 可选 $\xi$ 使得 $\xi\in(a,b)$.

  \item 设 $f$, $g$ 均在 $[a,b]$ 上 Riemann 可积, $g(x)$ 在 $[a,b]$ 上不变号, $M=\sup\limits_{a\leqslant x\leqslant b}f(x)$, $m=\inf\limits_{a\leqslant x\leqslant b}f(x)$. 证明: $\exists\mu\in[m,M]$, s. t.
  \[
  \int_{a}^{b}f(x)g(x)\,\mathrm{d}x=\mu\int_{a}^{b}g(x)\,\mathrm{d}x.
\]

  \item 设 $f$ 与 $g$ 在 $[a,b]$ 上 Riemann 可积, 证明 Cauchy-Schwarz 不等式
  \[
  \left[\int_{a}^{b}f(x)g(x)\,\mathrm{d}x\right]^{2}\leqslant\int_{a}^{b}f^{2}(x)\,\mathrm{d}x\cdot\int_{a}^{b}g^{2}(x)\,\mathrm{d}x,
\]
  其中等号当且仅当 $f\xlongequal{\text{a. e.}}\lambda g$ 或 $g\xlongequal{\text{a. e.}}\lambda f$, $x\in[a,b]$ 时成立.

  \item 应用 Cauchy-Schwarz 不等式证明:

  (1) 设 $f$ 在 $[a,b]$ 上 Riemann 可积, 则
  \[
  \left[\int_{a}^{b}f(x)\,\mathrm{d}x\right]^{2}\leqslant(b-a)\int_{a}^{b}f^{2}(x)\,\mathrm{d}x;
\]

  (2) 设 $f$ 在 $[a,b]$ 上 Riemann 可积, 且 $f(x)\geqslant m>0$, 则
  \[
  \int_{a}^{b}f(x)\,\mathrm{d}x\cdot\int_{a}^{b}\frac{1}{f(x)}\,\mathrm{d}x\geqslant(b-a)^{2};
\]

  (3) 设 $f$ 与 $g$ 都在 $[a,b]$ 上 Riemann 可积, 则有 Minkowski 不等式
  \[
  \left[\int_{a}^{b}(f(x)+g(x))^{2}\,\mathrm{d}x\right]^{\frac{1}{2}}\leqslant\left[\int_{a}^{b}f^{2}(x)\,\mathrm{d}x\right]^{\frac{1}{2}}+\left[\int_{a}^{b}g^{2}(x)\,\mathrm{d}x\right]^{\frac{1}{2}}.
\]

  \item 设 $f$ 为 $(0,+\infty)$ 上的连续减函数, $f(x)>0$; 又设
  \[
  a_{n}=\sum_{k=1}^{n}f(k)-\int_{1}^{n}f(x)\,\mathrm{d}x.
\]
  证明: 数列 $\{a_{n}\}$ 收敛.

  \item 设 $f$ 在 $[a,b]$ 上 Riemann 可积, $\varphi$ 在 $[\alpha,\beta]$ 上单调且连续可导, $\varphi(\alpha)=a$, $\varphi(\beta)=b$.
  证明: $\int_{a}^{b}f(x)\,\mathrm{d}x=\int_{\alpha}^{\beta}f(\varphi(t))\varphi'(t)\,\mathrm{d}t$.

  \item 设 $f$ 为 $[a,b]$ 上的连续函数, $g$ 为连续可导的单调函数, 试用一个比较简单的论述, 证明积分第二中值定理:
  \[
  \int_{a}^{b}f(x)g(x)\,\mathrm{d}x=g(a)\int_{a}^{\xi}f(x)\,\mathrm{d}x+g(b)\int_{\xi}^{b}f(x)\,\mathrm{d}x.
\]
  (提示: 应用分部积分法 $\int_{a}^{b}f(x)g(x)\,\mathrm{d}x=\int_{a}^{b}g(x)\,\mathrm{d}\int_{a}^{x}f(t)\,\mathrm{d}t$, 以及积分第一中值定理.)

  \item 设 $f(x)$ 在 $[a,b]$ 上连续可导, 而
  \[
  \Delta_{n}=\int_{a}^{b}f(x)\,\mathrm{d}x-\frac{b-a}{n}\sum_{k=1}^{n}f\left(a+k\cdot\frac{b-a}{n}\right).
\]
  证明: (1) $-\frac{1}{2}\left(\frac{b-a}{n}\right)^{2}\sum\limits_{k=1}^{n}M_{k}\leqslant\Delta_{n}\leqslant-\frac{1}{2}\left(\frac{b-a}{n}\right)^{2}\sum\limits_{k=1}^{n}m_{k}$, 其中
  \[
  m_{k}=\min\left\{f'(x)\mid a+(k-1)\frac{b-a}{n}\leqslant x\leqslant a+k\frac{b-a}{n}\right\},
\]
  \[
  M_{k}=\max\left\{f'(x)\mid a+(k-1)\frac{b-a}{n}\leqslant x\leqslant a+k\frac{b-a}{n}\right\}.
\]
  (2) $\lim\limits_{n\to+\infty}n\Delta_{n}=-\frac{b-a}{2}[f(b)-f(a)]$.

  \item 设 $f$ 为 $[a,b]$ 上的连续函数, 且对任一满足 $\int_{a}^{b}g(x)\,\mathrm{d}x=0$ 的连续函数 $g$ 有\linebreak[4]$\int_{a}^{b}f(x)g(x)\,\mathrm{d}x=0$. 证明: $f$ 为常值函数.

  \item 求出满足下列条件的所有函数 $f(x)$:

  (1) $f$ 在 $[0,1]$ 上为连续函数;

  (2) $\int_{0}^{1}f(x)\,\mathrm{d}x=1$;

  (3) $\exists\alpha\in\mathbb{R}$, s. t.
  \[
  \int_{0}^{1}xf(x)\,\mathrm{d}x=\alpha,\quad\int_{0}^{1}x^{2}f(x)\,\mathrm{d}x=\alpha^{2}.
\]
\end{exblock}

% ====================================================================
%  徐森林、薛春华《数学分析》(第一册)
%  \S6.3  微积分基本定理、微积分基本公式
%  源: 原书印刷页 371--386(PDF 物理页 377--392)
%  底本: PaddleOCR-VL-1.6, 见 chapters/_ch06_3_raw.md
%  约定: 全西文标点; 行间公式用一对方括号; 多行公式用 align 或 aligned;
%        跨册引用一律用 extref(写作「定理 x.y.z」原文); 数学模式内的中文必须包text
% ====================================================================
\section{微积分基本定理、微积分基本公式}\label{sec:6.3}

设 $f$ 在 $[a,b]$ 上 Riemann 可积, 因此 $f$ 在 $[a,b]$ 的任一闭子区间上也 Riemann 可积. 于是, 由
\[
\Phi(x)=\int_{a}^{x}f(t)\,\mathrm{d}t,\quad x\in[a,b]
\]
定义了一个以积分上限 $x$ 为自变量的函数, 称为\textbf{变上限积分}, 几何上看是变动曲边梯形 $AaxB$ 的面积 (图~\ref{fig:6.3.1}). 在不致混淆时, 记 $\int_{a}^{x}f(t)\,\mathrm{d}t$ 为 $\int_{a}^{x}f(x)\,\mathrm{d}x$, 但应注意的是积分上限的 $x$ 是真正的变量, 而 $f(x)\,\mathrm{d}x$ 中的 $x$ 是积分变量, 它用 $t,\theta,u,\cdots$ 表示都可以.

类似地, 也可定义\textbf{变下限积分}
\[
\Psi(x)=\int_{x}^{b}f(t)\,\mathrm{d}t,\quad x\in[a,b].
\]
$\Phi$ 与 $\Psi$ 统称为\textbf{变限积分}.

\begin{figure}[htbp]
\centering
\begin{tikzpicture}[x=1.05cm,y=1.00cm,>=Stealth]
  \draw[->] (-0.6,0) -- (5.5,0) node[below right] {$x$};
  \draw[->] (0,-0.5) -- (0,4.05) node[left] {$y$};
  \fill[pattern=north east lines] (1,0) -- (1,2.2)
    .. controls (1.45,2.80) and (2.50,2.65) .. (3.2,2.28)
    -- (3.2,0) -- cycle;
  \draw[thick] (1,2.2) .. controls (1.45,2.80) and (2.50,2.65) .. (3.2,2.28)
    .. controls (3.75,2.05) and (4.15,2.30) .. (4.6,3.35);
  \draw (1,0) -- (1,2.2);
  \draw (3.2,0) -- (3.2,2.28);
  \draw (4.6,0) -- (4.6,3.35);
  \node[below left] at (0,0) {$O$};
  \node[left] at (0.94,2.24) {$A$};
  \node[right] at (3.26,2.24) {$B$};
  \node at (2.55,3.05) {$y=f(x)$};
  \node[below] at (1,-0.10) {$a$};
  \node[below] at (3.2,-0.10) {$x$};
  \node[below] at (4.6,-0.10) {$b$};
\end{tikzpicture}
\caption{}\label{fig:6.3.1}
\end{figure}

变限积分所定义的函数有着重要的性质. 由于
\[
\int_{x}^{b}f(t)\,\mathrm{d}t=-\int_{b}^{x}f(t)\,\mathrm{d}t,
\]
下面只讨论变上限积分.

\begin{theorem}\label{thm:6.3.1}
设 $f$ 在 $[a,b]$ 上 Riemann 可积, 则变上限积分
\[
\Phi(x)=\int_{a}^{x}f(t)\,\mathrm{d}t
\]
为满足 Lipschitz 条件的函数. 特别地, $\Phi(x)$ 在 $[a,b]$ 上一致连续.
\end{theorem}

\begin{proof}
因为 $f$ 在 $[a,b]$ 上 Riemann 可积, 根据定理~\ref{thm:6.1.2} 知, $f$ 在 $[a,b]$ 上必有界, 即 $|f(x)|\leqslant M$, $\forall x_{1},x_{2}\in[a,b]$. 由定理~\ref{thm:6.2.4} 及定理~\ref{thm:6.2.3} 得到
\begin{align*}
|\Phi(x_{1})-\Phi(x_{2})|&=\left|\int_{a}^{x_{1}}f(t)\,\mathrm{d}t-\int_{a}^{x_{2}}f(t)\,\mathrm{d}t\right|=\left|\int_{x_{2}}^{x_{1}}f(t)\,\mathrm{d}t\right|\\
&\leqslant\left|\int_{x_{2}}^{x_{1}}|f(t)|\,\mathrm{d}t\right|\leqslant M|x_{1}-x_{2}|.
\end{align*}
这就证明了 $\Phi$ 在 $[a,b]$ 上满足 Lipschitz 条件.

另一方面, $\forall\varepsilon>0$, 取 $\delta\in\left(0,\dfrac{\varepsilon}{M+1}\right)$, 则当 $x_{1},x_{2}\in[a,b]$, $|x_{1}-x_{2}|<\delta$ 时有
\[
|\Phi(x_{1})-\Phi(x_{2})|\leqslant M|x_{1}-x_{2}|\leqslant M\delta<M\cdot\frac{\varepsilon}{M+1}<\varepsilon,
\]
因此, $\Phi$ 在 $[a,b]$ 上是一致连续的.
\end{proof}

现在我们来研究积分学的最重要的微积分基本定理与微积分基本公式.

\begin{theorem}[微积分基本定理]\label{thm:6.3.2}
设 $f$ 在 $[a,b]$ 上 Riemann 可积, 且在 $x_{0}\in[a,b]$ 连续, 则
\[
\Phi(x)=\int_{a}^{x}f(t)\,\mathrm{d}t
\]
在 $x_{0}$ 处可导, 且
\[
\Phi'(x_{0})=f(x_{0}).
\]
\end{theorem}

\begin{proof}
$\forall\varepsilon>0$, 因 $f$ 在 $x_{0}$ 连续, 故 $\exists\delta>0$, 当 $|x-x_{0}|<\delta$, $x\in[a,b]$ 时, 有 $|f(x)-f(x_{0})|<\varepsilon$. 于是
\[
\begin{aligned}
\left|\frac{\Phi(x)-\Phi(x_{0})}{x-x_{0}}-f(x_{0})\right|&=\left|\frac{\int_{a}^{x}f(t)\,\mathrm{d}t-\int_{a}^{x_{0}}f(t)\,\mathrm{d}t}{x-x_{0}}-f(x_{0})\right|\\
&=\left|\frac{\int_{x_{0}}^{x}f(t)\,\mathrm{d}t-\int_{x_{0}}^{x}f(x_{0})\,\mathrm{d}t}{x-x_{0}}\right|=\left|\frac{\int_{x_{0}}^{x}[f(t)-f(x_{0})]\,\mathrm{d}t}{x-x_{0}}\right|\\
&\leqslant\frac{\int_{x_{0}}^{x}|f(t)-f(x_{0})|\,\mathrm{d}t}{|x-x_{0}|}<\frac{\varepsilon\left|\int_{x_{0}}^{x}\mathrm{d}t\right|}{|x-x_{0}|}=\varepsilon,
\end{aligned}
\]
即 $\Phi(x)$ 在 $x_{0}$ 处可导, 且
\[
\Phi'(x_{0})=\lim_{x\to x_{0}}\frac{\Phi(x)-\Phi(x_{0})}{x-x_{0}}=f(x_{0}).
\]
\end{proof}

\begin{corollary}\label{cor:6.3.1}
设 $f$ 在 $[a,b]$ 上连续, $\Phi(x)=\int_{a}^{x}f(t)\,\mathrm{d}t$, $x\in[a,b]$, 则 $\Phi'(x)=f(x)$, 即 $\Phi(x)$ 为 $f(x)$ 在 $[a,b]$ 上的一个原函数. 从而
\[
\int f(x)\,\mathrm{d}x=\Phi(x)+C=\int_{a}^{x}f(t)\,\mathrm{d}t+C.
\]
\end{corollary}

\begin{proof}[证法 1]
由定理~\ref{thm:6.3.2} 推得.
\end{proof}

\begin{proof}[证法 2]
对 $\forall x\in[a,b]$, 当 $\Delta x\neq 0$ 且 $x+\Delta x\in[a,b]$ 时, 有
\[
\begin{aligned}
&\frac{\Delta\Phi}{\Delta x}=\frac{\Phi(x+\Delta x)-\Phi(x)}{\Delta x}=\frac{1}{\Delta x}\int_{x}^{x+\Delta x}f(t)\,\mathrm{d}t\\
&\quad\xlongequal{\text{积分第一中值定理}}f(x+\theta\Delta x),\quad 0\leqslant\theta\leqslant 1.
\end{aligned}
\]
由于 $f$ 在点 $x$ 连续, 故有
\[
\Phi'(x)=\lim_{\Delta x\to 0}\frac{\Delta\Phi}{\Delta x}=\lim_{\Delta x\to 0}f(x+\theta\Delta x)=f(x).
\]
\end{proof}

\begin{corollary}\label{cor:6.3.2}
设 $f$ 在 $[a,b]$ 上连续, $u(x):(c,d)\to\mathbb{R}$ 与 $v(x):(c,d)\to\mathbb{R}$ 可导, $a\leqslant u(x)\leqslant b$, $a\leqslant v(x)\leqslant b$, 则
\[
\left(\int_{v(x)}^{u(x)}f(t)\,\mathrm{d}t\right)'=f(u(x))u'(x)-f(v(x))v'(x).
\]
\end{corollary}

\begin{proof}
应用复合函数求导的链规则, 有
\begin{align*}
\left(\int_{v(x)}^{u(x)}f(t)\,\mathrm{d}t\right)'&=\left(\int_{v(x)}^{a}f(t)\,\mathrm{d}t+\int_{a}^{u(x)}f(t)\,\mathrm{d}t\right)'\\
&=\left(\int_{a}^{u(x)}f(t)\,\mathrm{d}t\right)'-\left(\int_{a}^{v(x)}f(t)\,\mathrm{d}t\right)'\\
&=\left.\left(\int_{a}^{u}f(t)\,\mathrm{d}t\right)'\right|_{u=u(x)}\cdot u'(x)-\left.\left(\int_{a}^{v}f(t)\,\mathrm{d}t\right)'\right|_{v=v(x)}\cdot v'(x)\\
&=f(u(x))u'(x)-f(v(x))v'(x).
\end{align*}
\end{proof}

\begin{theorem}[微积分基本公式或 Newton-Leibniz 公式]\label{thm:6.3.3}
设 $f$ 在 $[a,b]$ 上连续, $F$ 为 $f$ 在 $[a,b]$ 上的任一原函数, 即 $F'(x)=f(x)$, $\forall x\in[a,b]$. 则
\[
\int_{a}^{b}f(x)\,\mathrm{d}x=\int_{a}^{b}F'(x)\,\mathrm{d}x=F(b)-F(a)=F(x)\Big|_{a}^{b}.
\]
\end{theorem}

微积分基本公式是体现导数 (微分) 与积分有着密切联系的重要公式.

\begin{proof}[证法 1]
设 $\Phi(x)=\int_{a}^{x}f(t)\,\mathrm{d}t$, 则由推论 6.1.1 知,
\[
\Phi'(x)=f(x)=F'(x),\quad \forall x\in[a,b].
\]
再根据推论 3.2.2 得到 $F(x)=\Phi(x)+C$, 所以,
\[
\begin{aligned}
\int_{a}^{b}f(x)\,\mathrm{d}x&=\int_{a}^{b}f(x)\,\mathrm{d}x-\int_{a}^{a}f(x)\,\mathrm{d}x=\Phi(b)-\Phi(a)\\
&=(\Phi(b)+C)-(\Phi(a)+C)=F(b)-F(a).
\end{aligned}
\]
\end{proof}

\begin{proof}[证法 2]
因为 $f$ 在 $[a,b]$ 上连续, 从而一致连续, 即 $\forall\varepsilon>0$, $\exists\delta>0$, 当 $x',x''\in[a,b]$ 且 $|x'-x''|<\delta$ 时, 有
\[
|f(x')-f(x'')|<\frac{\varepsilon}{b-a}.
\]
对上述 $\varepsilon>0$, 及 $[a,b]$ 的任一分割 $T:a=x_{0}<x_{1}<\cdots<x_{n}=b$, 当 $\|T\|<\delta$ 与 $\forall\xi_{i}\in[x_{i-1},x_{i}]$ 时, 有
\begin{align*}
&\left|\sum_{i=1}^{n}f(\xi_{i})\Delta x_{i}-[F(b)-F(a)]\right|\\
&=\left|\sum_{i=1}^{n}f(\xi_{i})\Delta x_{i}-\sum_{i=1}^{n}[F(x_{i})-F(x_{i-1})]\right|\\
&=\left|\sum_{i=1}^{n}f(\xi_{i})\Delta x_{i}-\sum_{i=1}^{n}f(\eta_{i})\Delta x_{i}\right|\\
&=\left|\sum_{i=1}^{n}[f(\xi_{i})-f(\eta_{i})]\Delta x_{i}\right|\\
&\leqslant\sum_{i=1}^{n}|f(\xi_{i})-f(\eta_{i})|\Delta x_{i}\\
&<\frac{\varepsilon}{b-a}\sum_{i=1}^{n}\Delta x_{i}=\frac{\varepsilon}{b-a}(b-a)=\varepsilon,
\end{align*}
其中第二个等式是应用了 Lagrange 中值定理, $\exists\eta_{i}\in(x_{i-1},x_{i})$, s. t.
\[
F(x_{i})-F(x_{i-1})=F'(\eta_{i})(x_{i}-x_{i-1})=f(\eta_{i})\Delta x_{i}.
\]
由上推得 $f$ 在 $[a,b]$ 上 Riemann 可积, 且有
\[
\int_{a}^{b}f(x)\,\mathrm{d}x=F(b)-F(a).
\]
\end{proof}

\begin{theorem}[推广的微积分基本公式或推广的 Newton-Leibniz 公式]\label{thm:6.3.4}
设 $f$ 在 $[a,b]$ 上 Riemann 可积, $F$ 在 $[a,b]$ 上连续, 且在 $[a,b]$ 上除有限个点外,
\[
F'(x)=f(x)
\]
成立. 则
\[
\int_{a}^{b}f(x)\,\mathrm{d}x=F(b)-F(a)=F(x)\Big|_{a}^{b}.
\]
\end{theorem}

\begin{proof}
作 $[a,b]$ 的分割 $T:a=x_{0}<x_{1}<\cdots<x_{n}=b$, 使 $F'(x)=f(x)$ 不成立的点全在上述分点中. 所以, $F'(x)=f(x)$, $\forall x\in(x_{i-1},x_{i})$, 且 $F$ 在 $[x_{i-1},x_{i}]$ 上连续. 由 Lagrange 中值定理知, $\exists\xi_{i}\in(x_{i-1},x_{i})$, s. t.
\[
f(\xi_{i})=F'(\xi_{i})=\frac{F(x_{i})-F(x_{i-1})}{x_{i}-x_{i-1}}.
\]
由此推得
\[
F(x_{i})-F(x_{i-1})=f(\xi_{i})(x_{i}-x_{i-1})=f(\xi_{i})\Delta x_{i}.
\]
于是,
\[
\begin{aligned}
\int_{a}^{b}f(x)\,\mathrm{d}x&=\lim_{\|T\|\to 0}\sum_{i=1}^{n}f(\xi_{i})\Delta x_{i}=\lim_{\|T\|\to 0}\sum_{i=1}^{n}[F(x_{i})-F(x_{i-1})]\\
&=\lim_{\|T\|\to 0}[F(b)-F(a)]=F(b)-F(a).
\end{aligned}
\]
\end{proof}

\begin{remark}\label{rem:6.3.1}
现在我们用定积分与 Newton-Leibniz 公式来描述匀加速运动与例~\ref{ex:5.2.11} 中的半生期问题. 关于匀加速运动有
\[
v(t)=\int_{t_{0}}^{t}a\,\mathrm{d}t+v(t_{0})=a(t-t_{0})+v(t_{0}),
\]
\[
\begin{aligned}
x(t)&=\int_{t_{0}}^{t}[a(t-t_{0})+v(t_{0})]\,\mathrm{d}t+x(t_{0})\\
&=\frac{a}{2}(t-t_{0})^{2}+v(t_{0})(t-t_{0})+x(t_{0}).
\end{aligned}
\]
关于半生期有
\[
\frac{\mathrm{d}m}{m}=-k\,\mathrm{d}t,
\]
\[
\int_{m(0)}^{m(t)}\frac{\mathrm{d}m}{m}=\int_{0}^{t}(-k)\,\mathrm{d}t,
\]
\[
\ln\frac{m(t)}{m(0)}=-kt,
\]
\[
m(t)=m(0)\mathrm{e}^{-kt}.
\]
\end{remark}

\begin{example}\label{ex:6.3.1}
已经知道, 如果 $f(x)$ 连续, 则变上限积分 $\int_{a}^{x}f(t)\,\mathrm{d}t$ 为 $f(x)$ 的一个原函数, 且
\[
\int f(x)\,\mathrm{d}x=\int_{a}^{x}f(t)\,\mathrm{d}t+C,
\]
还有 Newton-Leibniz 公式
\[
\int_{a}^{b}f(x)\,\mathrm{d}x=F(x)\Big|_{a}^{b}=F(b)-F(a)
\]
成立, 其中 $F(x)=\int_{a}^{x}f(t)\,\mathrm{d}t+C$.

值得注意的是,
\[
f(x)=\begin{cases}\frac{3}{2}x^{\frac{1}{2}}\sin\frac{1}{x}-x^{-\frac{1}{2}}\cos\frac{1}{x},&x\neq 0,\\0,&x=0\end{cases}
\]
有原函数
\[
F(x)=\begin{cases}x^{\frac{3}{2}}\sin\frac{1}{x},&x\neq 0,\\0,&x=0.\end{cases}
\]
但 $f(x)$ 在 $[0,1]$ 上无界, 根据定理~\ref{thm:6.1.2} 推得它在 $[0,1]$ 上不是 Riemann 可积的, 从而, Newton-Leibniz 公式不成立, 即
\[
\int_{0}^{1}f(x)\,\mathrm{d}x\neq F(x)\Big|_{0}^{1}.
\]
即使加上 $f(x)$ 在 $[a,b]$ 上 Riemann 可积, 推广的 Newton-Leibniz 公式也未必成立. 如
\[
f(x)=0,
\]
\[
F(x)=\begin{cases}1,&x\geqslant 0,\\0,&x<0,\end{cases}
\]
则 $F'(x)=0=f(x)$, $\forall x\neq 0$, 且 $f(x)$ 在 $[-1,1]$ 上 Riemann 可积, 但
\[
\int_{-1}^{1}f(x)\,\mathrm{d}x=\int_{-1}^{1}0\,\mathrm{d}x=0\neq 1=1-0=F(1)-F(-1).
\]
$F(x)$ 在 $0$ 不连续是造成推广的 Newton-Leibniz 公式不成立的主要原因.
\end{example}

\begin{example}\label{ex:6.3.2}
应用 Newton-Leibniz 公式计算下列 Riemann 积分:

(1) $\int_{a}^{b}x\,\mathrm{d}x$;\qquad (2) $\int_{a}^{b}x^{2}\,\mathrm{d}x$;\qquad (3) $\int_{a}^{b}\sin x\,\mathrm{d}x$;

(4) $\int_{a}^{b}\mathrm{e}^{x}\,\mathrm{d}x$;\qquad (5) $\int_{0}^{2}x\sqrt{4-x^{2}}\,\mathrm{d}x$.
\end{example}

\begin{solution}
(1) $\int_{a}^{b}x\,\mathrm{d}x=\left.\frac{x^{2}}{2}\right|_{a}^{b}=\frac{b^{2}-a^{2}}{2}$.

(2) $\int_{a}^{b}x^{2}\,\mathrm{d}x=\left.\frac{x^{3}}{3}\right|_{a}^{b}=\frac{b^{3}-a^{3}}{3}$.

(3) $\int_{a}^{b}\sin x\,\mathrm{d}x=-\cos x\Big|_{a}^{b}=\cos a-\cos b$.

(4) $\int_{a}^{b}\mathrm{e}^{x}\,\mathrm{d}x=\mathrm{e}^{x}\Big|_{a}^{b}=\mathrm{e}^{b}-\mathrm{e}^{a}$.

(5) 由
\[
\begin{aligned}
\int x\sqrt{4-x^{2}}\,\mathrm{d}x&=-\frac{1}{2}\int\sqrt{4-x^{2}}\,\mathrm{d}(4-x^{2})\\
&=-\frac{1}{3}(4-x^{2})^{\frac{3}{2}}+C
\end{aligned}
\]
得到
\[
\int_{0}^{2}x\sqrt{4-x^{2}}\,\mathrm{d}x=-\left.\frac{1}{3}(4-x^{2})^{\frac{3}{2}}\right|_{0}^{2}=\frac{8}{3}.
\]
比较例~\ref{ex:6.3.2}(2) 与例~\ref{ex:6.1.13}, 再比较例~\ref{ex:6.3.2}(3) 与例~\ref{ex:6.1.14} 就可知道, 用 Newton-Leibniz 公式计算 Riemann 积分比用定义及定理~\ref{thm:6.1.4} 要简捷得多, 这充分显示了 Newton-Leibniz 公式的强大威力, 也显示了求原函数或求不定积分的重要性.
\end{solution}

\begin{example}\label{ex:6.3.3}
设 $m,n=0,1,2,\cdots$, 证明:

(1)
\[
\int_{-\pi}^{\pi}\cos mx\cos nx\,\mathrm{d}x=\int_{-\pi}^{\pi}\sin mx\sin nx\,\mathrm{d}x=\begin{cases}0,&m\neq n,\\\pi,&m=n;\end{cases}
\]

(2) $\int_{-\pi}^{\pi}\sin mx\cos nx\,\mathrm{d}x=0$.
\end{example}

\begin{proof}
应用 Newton-Leibniz 公式, 有

(1)
\[
\begin{aligned}
\int_{-\pi}^{\pi}\cos mx\cos nx\,\mathrm{d}x&=\frac{1}{2}\int_{-\pi}^{\pi}[\cos(m+n)x+\cos(m-n)x]\,\mathrm{d}x\\
&=\begin{cases}\frac{1}{2}\left.\left[\frac{\sin(m+n)x}{m+n}+\frac{\sin(m-n)x}{m-n}\right]\right|_{-\pi}^{\pi}=0,&m\neq n,\\[2mm]\frac{1}{2}\left.\left[\frac{\sin(m+n)x}{m+n}+x\right]\right|_{-\pi}^{\pi}=\frac{1}{2}(0+2\pi)=\pi,&m=n.\end{cases}
\end{aligned}
\]
\[
\begin{aligned}
\int_{-\pi}^{\pi}\sin mx\sin nx\,\mathrm{d}x&=\frac{1}{2}\int_{-\pi}^{\pi}[\cos(m-n)x-\cos(m+n)x]\,\mathrm{d}x\\
&=\begin{cases}\frac{1}{2}\left.\left[\frac{\sin(m-n)x}{m-n}-\frac{\sin(m+n)x}{m+n}\right]\right|_{-\pi}^{\pi}=0,&m\neq n,\\[2mm]\frac{1}{2}\left.\left[x-\frac{\sin(m+n)x}{m+n}\right]\right|_{-\pi}^{\pi}=\frac{1}{2}(2\pi-0)=\pi,&m=n.\end{cases}
\end{aligned}
\]

(2)
\[
\begin{aligned}
\int_{-\pi}^{\pi}\sin mx\cos nx\,\mathrm{d}x&=\frac{1}{2}\int_{-\pi}^{\pi}[\sin(m+n)x+\sin(m-n)x]\,\mathrm{d}x\\
&=\begin{cases}\frac{1}{2}\left.\left[-\frac{\cos(m+n)x}{m+n}-\frac{\cos(m-n)x}{m-n}\right]\right|_{-\pi}^{\pi}=0,&m\neq n,\\[2mm]\frac{1}{2}\left.\left[-\frac{\cos(m+n)x}{m+n}\right]\right|_{-\pi}^{\pi}+0=\frac{1}{2}(0+0)=0,&m=n,\,m+n\neq 0,\\[2mm]0,&m=n=0.\end{cases}
\end{aligned}
\]
\end{proof}

\begin{example}\label{ex:6.3.4}
应用 Riemann 积分定义及 Newton-Leibniz 公式计算下面的数列极限:

(1) $\lim\limits_{n\to+\infty}\dfrac{1^{p}+2^{p}+\cdots+n^{p}}{n^{p+1}}$ ($p>-1$);

(2)
\[
\lim_{n\to+\infty}\frac{1}{n}\left(1+\sec^{2}\frac{\pi}{4n}+\sec^{2}\frac{2\pi}{4n}+\cdots+\sec^{2}\frac{n\pi}{4n}\right);
\]

(3)
\[
\lim_{n\to+\infty}\left(\frac{1}{n+1}+\frac{1}{n+2}+\cdots+\frac{1}{2n}\right).
\]
\end{example}

\begin{solution}
(1)
\[
\begin{aligned}
&\lim_{n\to+\infty}\frac{1^{p}+2^{p}+\cdots+n^{p}}{n^{p+1}}\\
&=\lim_{n\to+\infty}\frac{1}{n}\left[\left(\frac{1}{n}\right)^{p}+\left(\frac{2}{n}\right)^{p}+\cdots+\left(\frac{n}{n}\right)^{p}\right]\\
&=\int_{0}^{1}x^{p}\,\mathrm{d}x=\left.\frac{x^{p+1}}{p+1}\right|_{0}^{1}=\frac{1}{p+1},\quad p>-1.
\end{aligned}
\]
参阅例~\ref{ex:1.6.3}(1), 它应用 Stolz 公式对 $p$ 为自然数时求得极限值 $\dfrac{1}{p+1}$.

(2)
\[
\begin{aligned}
&\lim_{n\to+\infty}\frac{1}{n}\left(1+\sec^{2}\frac{\pi}{4n}+\sec^{2}\frac{2\pi}{4n}+\cdots+\sec^{2}\frac{n\pi}{4n}\right)\\
&=\frac{4}{\pi}\lim_{n\to+\infty}\frac{\pi}{4n}\left(1+\sec^{2}\frac{\pi}{4n}+\sec^{2}\frac{2\pi}{4n}+\cdots+\sec^{2}\frac{n\pi}{4n}\right)\\
&=\frac{4}{\pi}\int_{0}^{\frac{\pi}{4}}\sec^{2}x\,\mathrm{d}x=\left.\frac{4}{\pi}\tan x\right|_{0}^{\frac{\pi}{4}}=\frac{4}{\pi}(1-0)=\frac{4}{\pi}.
\end{aligned}
\]

(3)
\[
\begin{aligned}
&\lim_{n\to+\infty}\left(\frac{1}{n+1}+\frac{1}{n+2}+\cdots+\frac{1}{2n}\right)\\
&=\lim_{n\to+\infty}\frac{1}{n}\left(\frac{1}{1+\frac{1}{n}}+\cdots+\frac{1}{1+\frac{n}{n}}\right)\\
&=\int_{0}^{1}\frac{\mathrm{d}x}{1+x}=\left.\ln(1+x)\right|_{0}^{1}=\ln 2.
\end{aligned}
\]
另一方法参阅例~\ref{ex:1.4.9}(1).
\end{solution}

\begin{example}\label{ex:6.3.5}
设 $f$ 在 $[a,b]$ 上连续可导, $f(a)=0$. 证明:

(1) $\max\limits_{x\in[a,b]}f^{2}(x)\leqslant(b-a)\int_{a}^{b}[f'(x)]^{2}\,\mathrm{d}x$;

(2) $\int_{a}^{b}f^{2}(x)\,\mathrm{d}x\leqslant\dfrac{(b-a)^{2}}{2}\int_{a}^{b}[f'(x)]^{2}\,\mathrm{d}x$.
\end{example}

\begin{proof}
(1) 由 Cauchy-Schwarz 不等式、$f(a)=0$ 及 Newton-Leibniz 公式得到
\[
\begin{aligned}
f^{2}(x)&=[f(x)-f(a)]^{2}=\left[\int_{a}^{x}f'(t)\,\mathrm{d}t\right]^{2}=\left[\int_{a}^{x}f'(t)\cdot 1\,\mathrm{d}t\right]^{2}\\
&\leqslant\int_{a}^{x}[f'(x)]^{2}\,\mathrm{d}x\cdot\int_{a}^{x}1^{2}\,\mathrm{d}x=(x-a)\int_{a}^{x}[f'(x)]^{2}\,\mathrm{d}x\\
&\leqslant(b-a)\int_{a}^{b}[f'(x)]^{2}\,\mathrm{d}x,\quad\forall x\in[a,b],
\end{aligned}
\]
所以,
\[
\max_{x\in[a,b]}f^{2}(x)\leqslant(b-a)\int_{a}^{b}[f'(x)]^{2}\,\mathrm{d}x.
\]

(2) 由 (1) 推得了
\[
f^{2}(x)\leqslant(x-a)\int_{a}^{b}[f'(x)]^{2}\,\mathrm{d}x.
\]
两边积分得
\[
\begin{aligned}
\int_{a}^{b}f^{2}(x)\,\mathrm{d}x&\leqslant\int_{a}^{b}\left\{\int_{a}^{b}[f'(x)]^{2}\,\mathrm{d}x\cdot(x-a)\right\}\mathrm{d}x\\
&=\int_{a}^{b}[f'(x)]^{2}\,\mathrm{d}x\cdot\int_{a}^{b}(x-a)\,\mathrm{d}x\\
&=\int_{a}^{b}[f'(x)]^{2}\,\mathrm{d}x\cdot\left.\frac{(x-a)^{2}}{2}\right|_{a}^{b}\\
&=\frac{(b-a)^{2}}{2}\int_{a}^{b}[f'(x)]^{2}\,\mathrm{d}x.
\end{aligned}
\]
\end{proof}

\begin{example}\label{ex:6.3.6}
设 $f$ 在 $[0,+\infty)$ 上连续并恒取正值. 证明:
\[
\varphi(x)=\frac{\int_{0}^{x}tf(t)\,\mathrm{d}t}{\int_{0}^{x}f(t)\,\mathrm{d}t}
\]
是 $(0,+\infty)$ 上的严格增函数.
\end{example}

\begin{proof}
由微积分基本定理得到
\[
\begin{aligned}
\varphi'(x)&=\frac{xf(x)\int_{0}^{x}f(t)\,\mathrm{d}t-f(x)\int_{0}^{x}tf(t)\,\mathrm{d}t}{\left(\int_{0}^{x}f(t)\,\mathrm{d}t\right)^{2}}\\
&=\frac{f(x)\int_{0}^{x}(x-t)f(t)\,\mathrm{d}t}{\left(\int_{0}^{x}f(t)\,\mathrm{d}t\right)^{2}}>0
\end{aligned}
\]
($f(x)>0$, $x-t>0$, $\forall t\in(0,x)$), 根据定理~\ref{thm:3.5.1}, $\varphi(x)$ 为 $(0,+\infty)$ 上的严格增函数.
\end{proof}

\begin{example}\label{ex:6.3.7}
应用微积分基本定理求下列极限:

(1)
\[
\lim_{x\to+\infty}\frac{\int_{0}^{x}\sqrt{1+t^{4}}\,\mathrm{d}t}{x^{3}};
\]

(2)
\[
\lim_{x\to+\infty}\frac{\int_{0}^{x}t^{2}\mathrm{e}^{t^{2}}\,\mathrm{d}t}{\left(\int_{0}^{x}\mathrm{e}^{t^{2}}\,\mathrm{d}t\right)^{2}};
\]

(3)
\[
\lim_{x\to 0}\frac{\int_{0}^{x}t^{2}\mathrm{e}^{t^{2}}\,\mathrm{d}t}{\left(\int_{0}^{x}\mathrm{e}^{t^{2}}\,\mathrm{d}t\right)^{2}}.
\]
\end{example}

\begin{solution}
(1)
\[
\begin{aligned}
\lim_{x\to+\infty}\frac{\int_{0}^{x}\sqrt{1+t^{4}}\,\mathrm{d}t}{x^{3}}&\xlongequal[\text{L'Hospital 法则}]{\frac{\infty}{\infty}}\lim_{x\to+\infty}\frac{\sqrt{1+x^{4}}}{3x^{2}}\\
&=\lim_{x\to+\infty}\frac{1}{3}\sqrt{\frac{1}{x^{4}}+1}=\frac{1}{3}\sqrt{0+1}=\frac{1}{3}.
\end{aligned}
\]

(2)
\[
\begin{aligned}
\lim_{x\to+\infty}\frac{\int_{0}^{x}t^{2}\mathrm{e}^{t^{2}}\,\mathrm{d}t}{\left(\int_{0}^{x}\mathrm{e}^{t^{2}}\,\mathrm{d}t\right)^{2}}&\xlongequal{\frac{\infty}{\infty}}\lim_{x\to+\infty}\frac{x^{2}\mathrm{e}^{x^{2}}}{\left(\int_{0}^{x}\mathrm{e}^{t^{2}}\,\mathrm{d}t\right)\cdot\mathrm{e}^{x^{2}}}\\
&=\lim_{x\to+\infty}\frac{x^{2}}{2\int_{0}^{x}\mathrm{e}^{t^{2}}\,\mathrm{d}t}=\lim_{x\to+\infty}\frac{2x}{2\mathrm{e}^{x^{2}}}\\
&=\lim_{x\to+\infty}\frac{x}{\mathrm{e}^{x^{2}}}=\lim_{x\to+\infty}\frac{1}{2x\mathrm{e}^{x^{2}}}=0.
\end{aligned}
\]

(3)
\[
\begin{aligned}
\lim_{x\to 0}\frac{\int_{0}^{x}t^{2}\mathrm{e}^{t^{2}}\,\mathrm{d}t}{\left(\int_{0}^{x}\mathrm{e}^{t^{2}}\,\mathrm{d}t\right)^{2}}&\xlongequal{\frac{0}{0}}\lim_{x\to 0}\frac{x^{2}\mathrm{e}^{x^{2}}}{2\left(\int_{0}^{x}\mathrm{e}^{t^{2}}\,\mathrm{d}t\right)\cdot\mathrm{e}^{x^{2}}}\\
&=\lim_{x\to 0}\frac{x^{2}}{2\int_{0}^{x}\mathrm{e}^{t^{2}}\,\mathrm{d}t}=\lim_{x\to 0}\frac{x}{\mathrm{e}^{x^{2}}}=\frac{0}{\mathrm{e}^{x^{2}}}=0.
\end{aligned}
\]
\end{solution}

\begin{example}\label{ex:6.3.8}
设函数 $f$ 连续可导, $f(1)=1$, 且当 $x\geqslant 1$ 时, 有
\[
f'(x)=\frac{1}{x^{2}+f^{2}(x)}.
\]
证明: $\lim\limits_{x\to+\infty}f(x)$ 存在, 且 $\lim\limits_{x\to+\infty}f(x)\leqslant 1+\dfrac{\pi}{4}$.
\end{example}

\begin{proof}
由于 $f'(x)=\dfrac{1}{x^{2}+f^{2}(x)}>0$, 所以 $f$ 在 $[1,+\infty)$ 上严格增; 且当 $x\geqslant 1$ 时, 有 $f(x)>f(1)=1$, 于是, 由 Newton-Leibniz 公式得到
\[
\begin{aligned}
f(x)&=f(1)+\int_{1}^{x}f'(t)\,\mathrm{d}t=1+\int_{1}^{x}\frac{\mathrm{d}t}{t^{2}+f^{2}(t)}\\
&<1+\int_{1}^{x}\frac{\mathrm{d}t}{1+t^{2}}=1+\arctan x-\arctan 1\\
&<1+\frac{\pi}{2}-\frac{\pi}{4}=1+\frac{\pi}{4}.
\end{aligned}
\]
当 $x\to+\infty$ 时, $f$ 递增且有上界 $1+\dfrac{\pi}{4}$. 因此, $\lim\limits_{x\to+\infty}f(x)$ 存在有限, 并且
\[
\lim_{x\to+\infty}f(x)\leqslant 1+\frac{\pi}{4}.
\]
\end{proof}

\begin{example}\label{ex:6.3.9}
设 $f$ 连续可导, 且 $f(0)=0$, $f(1)=1$. 求证:
\[
\int_{0}^{1}|f(x)-f'(x)|\,\mathrm{d}x\geqslant\frac{1}{\mathrm{e}}.
\]
\end{example}

\begin{proof}
从 $\mathrm{e}^{-x}\leqslant 1$, $\forall x\in[0,1]$ 得到
\[
\begin{aligned}
\int_{0}^{1}|f(x)-f'(x)|\,\mathrm{d}x&\geqslant\int_{0}^{1}\mathrm{e}^{-x}|f(x)-f'(x)|\,\mathrm{d}x\\
&\geqslant\int_{0}^{1}\mathrm{e}^{-x}[f'(x)-f(x)]\,\mathrm{d}x=\int_{0}^{1}[\mathrm{e}^{-x}f(x)]'\,\mathrm{d}x\\
&\xlongequal{\text{Newton-Leibniz 公式}}\left.\mathrm{e}^{-x}f(x)\right|_{0}^{1}=\mathrm{e}^{-1}\cdot 1-\mathrm{e}^{-0}\cdot 0=\frac{1}{\mathrm{e}}.
\end{aligned}
\]
\end{proof}

\begin{example}\label{ex:6.3.10}
设 $f$ 在 $[0,1]$ 上有连续的 1 阶导数, $f(0)=0$, 并且 $0\leqslant f'(x)\leqslant 1$. 求证:
\[
\int_{0}^{1}f^{3}(x)\,\mathrm{d}x\leqslant\left(\int_{0}^{1}f(x)\,\mathrm{d}x\right)^{2}.
\]
\end{example}

\begin{proof}
令 $F(t)=\left(\int_{0}^{t}f(x)\,\mathrm{d}x\right)^{2}-\int_{0}^{t}f^{3}(x)\,\mathrm{d}x$, 则 $F(0)=0$,
\[
\begin{aligned}
F'(t)&=2f(t)\int_{0}^{t}f(x)\,\mathrm{d}x-f^{3}(t)\\
&=f(t)\left[2\int_{0}^{t}f(x)\,\mathrm{d}x-f^{2}(t)\right]=f(t)G(t),
\end{aligned}
\]
其中
\[
G(t)=2\int_{0}^{t}f(x)\,\mathrm{d}x-f^{2}(t),\quad G(0)=0-f^{2}(0)=0,
\]
\[
G'(t)=2f(t)-2f(t)f'(t)=2f(t)[1-f'(t)]\geqslant 0
\]
(因为 $f'(t)\geqslant 0$, $f(t)$ 单调增, $f(t)\geqslant f(0)=0$), 因此, $G$ 单调增, $G(t)\geqslant G(0)=0$. 由此可知, $F'(t)=f(t)G(t)\geqslant 0$. 这又得出 $F$ 单调增. 特别有
\[
\left(\int_{0}^{1}f(x)\,\mathrm{d}x\right)^{2}-\int_{0}^{1}f^{3}(x)\,\mathrm{d}x=F(1)\geqslant F(0)=0,
\]
\[
\int_{0}^{1}f^{3}(x)\,\mathrm{d}x\leqslant\left(\int_{0}^{1}f(x)\,\mathrm{d}x\right)^{2}.
\]
\end{proof}

\begin{example}\label{ex:6.3.11}
证明: $\lim\limits_{n\to+\infty}\int_{0}^{1}\dfrac{x^{n}}{1+x}\,\mathrm{d}x=0$.
\end{example}

\begin{solution}[证法 1]
因
\[
\begin{aligned}
0&\leqslant\int_{0}^{1}\frac{x^{n}}{1+x}\,\mathrm{d}x\leqslant\int_{0}^{1}x^{n}\,\mathrm{d}x\xlongequal{\text{Newton-Leibniz 公式}}\left.\frac{x^{n+1}}{n+1}\right|_{0}^{1}\\
&=\frac{1}{n+1}\to 0\ (n\to+\infty),
\end{aligned}
\]
故由夹逼定理知,
\[
\lim_{n\to+\infty}\int_{0}^{1}\frac{x^{n}}{1+x}\,\mathrm{d}x=0.
\]
\end{solution}

\begin{solution}[证法 2]
$\forall\varepsilon>0$, 取 $\delta=\min\left\{\dfrac{\varepsilon}{2},1\right\}$. 固定 $\delta$, 因 $\lim\limits_{n\to+\infty}(1-\delta)^{n}=0$, 故 $\exists N\in\mathbb{N}$, 当 $n>N$ 时, 有 $(1-\delta)^{n+1}<\dfrac{\varepsilon}{2}$. 于是,
\[
\begin{aligned}
0&\leqslant\int_{0}^{1}\frac{x^{n}}{1+x}\,\mathrm{d}x=\int_{0}^{1-\delta}\frac{x^{n}}{1+x}\,\mathrm{d}x+\int_{1-\delta}^{1}\frac{x^{n}}{1+x}\,\mathrm{d}x\\
&\leqslant\int_{0}^{1-\delta}x^{n}\,\mathrm{d}x+\int_{1-\delta}^{1}\,\mathrm{d}x\\
&<(1-\delta)^{n}(1-\delta)+[1-(1-\delta)]=(1-\delta)^{n+1}+\delta\\
&<\frac{\varepsilon}{2}+\frac{\varepsilon}{2}=\varepsilon,
\end{aligned}
\]
所以,
\[
\lim_{n\to+\infty}\int_{0}^{1}\frac{x^{n}}{1+x}\,\mathrm{d}x=0.
\]
\end{solution}

\begin{solution}[证法 3]
\[
\begin{aligned}
\int_{0}^{1}\frac{x^{n}}{1+x}\,\mathrm{d}x&=\begin{cases}\int_{0}^{1}\left(x^{n-1}-x^{n-2}+x^{n-3}+\cdots+x^{2}-x+1-\dfrac{1}{1+x}\right)\mathrm{d}x,&n\text{ 为奇数},\\[3mm]\int_{0}^{1}\left(x^{n-1}-x^{n-2}+x^{n-3}+\cdots+x-1+\dfrac{1}{1+x}\right)\mathrm{d}x,&n\text{ 为偶数},\end{cases}\\
&=\begin{cases}\dfrac{1}{n}-\dfrac{1}{n-1}+\dfrac{1}{n-2}+\cdots+\dfrac{1}{3}-\dfrac{1}{2}+1-\ln 2,&n\text{ 为奇数},\\[3mm]\dfrac{1}{n}-\dfrac{1}{n-1}+\dfrac{1}{n-2}+\cdots+\dfrac{1}{2}-1+\ln 2,&n\text{ 为偶数}.\end{cases}
\end{aligned}
\]
再根据例~\ref{ex:4.1.5} 立即推得
\[
\lim_{n\to+\infty}\int_{0}^{1}\frac{x^{n}}{1+x}\,\mathrm{d}x=0.
\]
\end{solution}

\begin{remark}\label{rem:6.3.2}
读者仔细检查下面的证明是否正确.
\[
\lim_{n\to+\infty}\int_{0}^{1}\frac{x^{n}}{1+x}\,\mathrm{d}x=\int_{0}^{1}\lim_{n\to+\infty}\frac{x^{n}}{1+x}\,\mathrm{d}x=0,
\]
其中
\[
\lim_{n\to+\infty}\frac{x^{n}}{1+x}=\begin{cases}0,&0\leqslant x<1,\\\dfrac{1}{2},&x=1.\end{cases}
\]
疑问之处是极限号与积分号为什么能交换. 似乎没有足够的理由. 但是, 这样做得到的答案也为 $0$, 也许是一种巧合. 为此, 我们举出极限号与积分号不可变换的例子.
\end{remark}

\begin{example}\label{ex:6.3.12}
设
\[
f_{n}(x)=\begin{cases}4n^{2}x,&0\leqslant x\leqslant\dfrac{1}{2n},\\-4n^{2}x+4n,&\dfrac{1}{2n}<x\leqslant\dfrac{1}{n},\\0,&x>\dfrac{1}{n},\end{cases}
\]
则 $f(x)=\lim\limits_{n\to+\infty}f_{n}(x)=0$. 于是,
\[
\begin{aligned}
\lim_{n\to+\infty}\int_{0}^{1}f_{n}(x)\,\mathrm{d}x&=\lim_{n\to+\infty}\frac{1}{2}\cdot 2n\cdot\frac{1}{n}=1\neq 0=\int_{0}^{1}0\,\mathrm{d}x\\
&=\int_{0}^{1}\lim_{n\to+\infty}f_{n}(x)\,\mathrm{d}x\quad(\text{图~\ref{fig:6.3.2}}).
\end{aligned}
\]
\end{example}

\begin{figure}[htbp]
\centering
\begin{tikzpicture}[x=1.00cm,y=1.00cm,>=Stealth]
  \draw[->] (-0.35,0) -- (6.7,0) node[below right] {$x$};
  \draw[->] (0,-0.35) -- (0,3.75) node[left] {$y$};
  \draw[thick] (0,0) -- (1.2,3.0) -- (2.4,0);
  \draw[thick] (2.4,0) -- (6.0,0);
  \draw[dotted,thick] (1.2,0) -- (1.2,3.0);
  \draw[dotted,thick] (0,3.0) -- (1.2,3.0);
  \draw (2.4,0) -- (2.4,-0.10);
  \draw (6.0,0) -- (6.0,-0.10);
  \node[below left] at (0,0) {$O$};
  \node[left] at (-0.06,3.0) {$2n$};
  \node[below] at (1.2,-0.14) {$\dfrac{1}{2n}$};
  \node[below] at (2.5,-0.14) {$\dfrac{1}{n}$};
  \node[below] at (6.0,-0.14) {$1$};
\end{tikzpicture}
\caption{}\label{fig:6.3.2}
\end{figure}

为了确保极限号与积分号可交换, 我们必须引进函数列一致收敛的概念.

\begin{definition}\label{def:6.3.1}
设 $X\subsetneqq\mathbb{R}$, $f_{n}:X\to\mathbb{R}$, $n=1,2,\cdots$ 为一个函数列, $f:X\to\mathbb{R}$, 如果 $\forall\varepsilon>0$, $\exists N=N(\varepsilon)\in\mathbb{N}$ ($N(\varepsilon)$ 只与 $\varepsilon$ 有关, 而与 $x\in X$ 无关), 当 $n>N=N(\varepsilon)$ 时, 有
\[
|f_{n}(x)-f(x)|<\varepsilon,\quad\forall x\in X,
\]
则称 $f_{n}(x)$ 在 $X$ 上一致收敛于 $f(x)$, 记作 $f_{n}(x)\Rightarrow f(x)$ ($n\to+\infty$), $x\in X$.
\end{definition}

\begin{theorem}\label{thm:6.3.5}
设 $f_{n}(x)\Rightarrow f(x)$ ($n\to+\infty$), $x\in X$, 如果 $f_{n}(x)$ 在 $x_{0}\in X$ 连续, $n=1,2,\cdots$, 则 $f(x)$ 在 $x_{0}$ 也连续. 特别当 $f_{n}(x)$ 在 $X$ 上为连续函数时, $f(x)$ 在 $X$ 上也为连续函数.
\end{theorem}

\begin{proof}
因 $f_{n}(x)\Rightarrow f(x)$ ($n\to+\infty$), $x\in X$, 故 $\exists N=N(\varepsilon)\in\mathbb{N}$, 当 $n>N$ 时, 有 $|f_{n}(x)-f(x)|<\dfrac{\varepsilon}{3}$, $\forall x\in X$. 于是, 对 $N+1$, 由 $f_{N+1}(x)$ 在 $x_{0}\in X$ 连续知, $\exists\delta>0$, 当 $x\in X$, $|x-x_{0}|<\delta$ 时, 有
\[
|f_{N+1}(x)-f_{N+1}(x_{0})|<\frac{\varepsilon}{3}.
\]
因此,
\[
\begin{aligned}
|f(x)-f(x_{0})|&\leqslant|f(x)-f_{N+1}(x)|+|f_{N+1}(x)-f_{N+1}(x_{0})|\\
&\quad+|f_{N+1}(x_{0})-f(x_{0})|\\
&<\frac{\varepsilon}{3}+\frac{\varepsilon}{3}+\frac{\varepsilon}{3}=\varepsilon,
\end{aligned}
\]
即 $\lim\limits_{x\to x_{0}}f(x)=f(x_{0})$, $f$ 在 $x_{0}$ 连续.
\end{proof}

\begin{theorem}\label{thm:6.3.6}
设 $f_{n}(x)$ 在 $[a,b]$ 上 Riemann 可积, $f_{n}(x)\Rightarrow f(x)$ ($n\to+\infty$), $x\in[a,b]$, 则 $f(x)$ 在 $[a,b]$ 上也 Riemann 可积, 且
\[
\lim_{n\to+\infty}\int_{a}^{b}f_{n}(x)\,\mathrm{d}x=\int_{a}^{b}\lim_{n\to+\infty}f_{n}(x)\,\mathrm{d}x=\int_{a}^{b}f(x)\,\mathrm{d}x.
\]
\end{theorem}

\begin{proof}
因为 $f_{n}(x)$ 在 $[a,b]$ 上 Riemann 可积, 根据定理~\ref{thm:6.1.2} 与定理~\ref{thm:6.1.4}(8), $f_{n}(x)$ 在 $[a,b]$ 上有界且几乎处处连续, 即 $D_{\text{不}}^{f_{n}}$ 为零测集. 又因为 $f_{n}(x)\Rightarrow f(x)$ ($n\to+\infty$), $x\in[a,b]$, 根据定理~\ref{thm:6.3.5} 知, $D_{\text{不}}^{f}\subset\bigcup\limits_{n=1}^{\infty}D_{\text{不}}^{f_{n}}$, 再由引理~\ref{lem:6.1.6}(4) 及 (2) 立即推得 $D_{\text{不}}^{f}$ 为零测集.

对 $\forall\varepsilon>0$, 因 $f_{n}(x)\Rightarrow f(x)$ ($n\to+\infty$), $x\in[a,b]$, 故 $\exists N=N(\varepsilon)\in\mathbb{N}$, 当 $n>N=N(\varepsilon)$ 时, 有
\[
|f_{n}(x)-f(x)|<\frac{\varepsilon}{b-a},\quad\forall x\in[a,b].
\]
于是,
\[
f_{N+1}(x)-\frac{\varepsilon}{b-a}<f(x)<f_{N+1}(x)+\frac{\varepsilon}{b-a},
\]
$f(x)$ 在 $[a,b]$ 上有界. 从定理~\ref{thm:6.1.4}(8) 得到 $f$ 在 $[a,b]$ 上 Riemann 可积, 进而, 有
\[
\begin{aligned}
&\left|\int_{a}^{b}f_{n}(x)\,\mathrm{d}x-\int_{a}^{b}f(x)\,\mathrm{d}x\right|=\left|\int_{a}^{b}[f_{n}(x)-f(x)]\,\mathrm{d}x\right|\\
&\leqslant\int_{a}^{b}|f_{n}(x)-f(x)|\,\mathrm{d}x<\frac{\varepsilon}{b-a}\cdot(b-a)=\varepsilon,
\end{aligned}
\]
\[
\lim_{n\to+\infty}\int_{a}^{b}f_{n}(x)\,\mathrm{d}x=\int_{a}^{b}f(x)\,\mathrm{d}x=\int_{a}^{b}\lim_{n\to+\infty}f_{n}(x)\,\mathrm{d}x.
\]
\end{proof}

回顾注~\ref{rem:6.3.2}, 由于
\[
f(x)=\lim_{n\to+\infty}\frac{x^{n}}{1+x}=\begin{cases}0,&0\leqslant x<1,\\\dfrac{1}{2},&x=1\end{cases}
\]
在 $x=1$ 不连续, 因此, 从定理~\ref{thm:6.3.5} 及反证法可以推断 $f_{n}(x)=\dfrac{x^{n}}{1+x}$ 在 $[0,1]$ 上不一致收敛于 $f(x)$. 由此知, 对上述的 $f_{n}$, $f$ 不能应用定理~\ref{thm:6.3.6}. 而 $\lim\limits_{n\to+\infty}\int_{0}^{1}\dfrac{x^{n}}{1+x}\,\mathrm{d}x=\int_{0}^{1}\lim\limits_{n\to+\infty}\dfrac{x^{n}}{1+x}\,\mathrm{d}x$ 只是一种巧合.

例~\ref{ex:6.3.12} 中, $f_{n}(x)$ 在 $[0,1]$ 上连续, 而 $f(x)\equiv 0$ 在 $[0,1]$ 上也连续, 这不能对 $f_{n}(x)$ 在 $[0,1]$ 上是否一致收敛于 $f(x)\equiv 0$ 给出任何信息. 但是, 由\linebreak[4]$\lim\limits_{n\to+\infty}\int_{0}^{1}f_{n}(x)\,\mathrm{d}x\neq\int_{0}^{1}\lim\limits_{n\to+\infty}f_{n}(x)\,\mathrm{d}x$ 与定理~\ref{thm:6.3.6} 表明 $f_{n}(x)$ 在 $[0,1]$ 上不一致收敛于 $f(x)=0$.

\begin{exblock}{练习题 6.3}

  \item 应用 Newton-Leibniz 公式计算下列 Riemann 积分:

  (1) $\int_{-1}^{1}\frac{x^{2}}{1+x^{2}}\,\mathrm{d}x$;\qquad (2) $\int_{0}^{1}\frac{x^{n}}{1+x}\,\mathrm{d}x$;\qquad (3) $\int_{1}^{4}\frac{x+1}{\sqrt{x}}\,\mathrm{d}x$;

  (4) $\int_{0}^{\frac{\pi}{2}}\cos mx\sin nx\,\mathrm{d}x$, $m,n\in\mathbb{Z}$;\qquad (5) $\int_{\mathrm{e}}^{\mathrm{e}^{2}}\frac{\mathrm{d}x}{x\ln x}$;

  (6) $\int_{0}^{1}\frac{\mathrm{e}^{x}-\mathrm{e}^{-x}}{2}\,\mathrm{d}x$;\qquad (7) $\int_{0}^{\frac{\pi}{3}}\tan^{2}x\,\mathrm{d}x$;

  (8) $\int_{4}^{9}\left(\sqrt{x}+\frac{1}{\sqrt{x}}\right)\mathrm{d}x$;\qquad (9) $\int_{\frac{1}{\mathrm{e}}}^{\mathrm{e}}\frac{1}{x}(\ln x)^{2}\,\mathrm{d}x$.

  \item 利用 Riemann 积分求下列极限:

  (1) $\lim\limits_{n\to+\infty}\frac{1}{n}\sum\limits_{k=1}^{n}\sin\frac{k\pi}{n}$;\qquad (2) $\lim\limits_{n\to+\infty}\frac{1}{n^{4}}(1+2^{3}+\cdots+n^{3})$;

  (3) $\lim\limits_{n\to+\infty}n\left[\frac{1}{(n+1)^{2}}+\frac{1}{(n+2)^{2}}+\cdots+\frac{1}{(n+n)^{2}}\right]$;

  (4) $\lim\limits_{n\to+\infty}\left(\frac{n}{n^{2}+1^{2}}+\frac{n}{n^{2}+2^{2}}+\cdots+\frac{n}{n^{2}+n^{2}}\right)$;

  (5) $\lim\limits_{n\to+\infty}\left(\frac{1}{\sqrt{n^{2}}}+\frac{1}{\sqrt{n(n+1)}}+\cdots+\frac{1}{\sqrt{n(2n-1)}}\right)$.

  \item 证明:

  (1) $\int_{0}^{1}\frac{(1+x)^{4}}{1+x^{2}}\,\mathrm{d}x=\frac{22}{3}-\pi$;

  (2) $\int_{0}^{1}\frac{x^{4}(1-x)^{4}}{1+x^{2}}\,\mathrm{d}x=\frac{22}{7}-\pi$.

  \item 设 $f$ 为 $[0,+\infty)$ 上的连续函数, 且它满足 $\int_{0}^{x}f(t)\,\mathrm{d}t=\frac{1}{2}xf(x)$, $x>0$. 证明: $f=cx$, 其中 $c$ 为常数.

  \item 设 $f$ 为 $(0,+\infty)$ 上的连续函数, 且对 $\forall a>0$ 有
  \[
  g(x)=\int_{x}^{ax}f(t)\,\mathrm{d}t\equiv\text{常数},\quad x\in(0,+\infty).
\]
  证明: $f(x)=\frac{c}{x}$, $x\in(0,+\infty)$, $c$ 为常数.

  \item 应用微积分基本定理求下列极限:

  (1) $\lim\limits_{x\to 0}\frac{1}{x}\int_{0}^{x}\cos t^{2}\,\mathrm{d}t$;

  (2)
  \[
  \lim_{x\to+\infty}\frac{\left(\int_{0}^{x}\mathrm{e}^{t^{2}}\,\mathrm{d}t\right)^{2}}{\int_{0}^{x}\mathrm{e}^{2t^{2}}\,\mathrm{d}t}.
\]

  \item 设 $f$ 为连续可导的函数, 试求
  \[
  \frac{\mathrm{d}}{\mathrm{d}x}\int_{a}^{x}(x-t)f'(t)\,\mathrm{d}t.
\]

  \item 设 $y=f(x)$ 为 $[a,b]$ 上严格增的连续曲线 (图~\ref{fig:6.3.3}). 试分别用 (1) 积分第二中值定理, (2) 连续函数 $F(x)=f(a)(x-a)+(b-x)f(b)$ 的介值定理, 证明: 存在 $\xi\in(a,b)$, 使得图~\ref{fig:6.3.3} 中两阴影部分面积相等.

  \item 设 $f$ 在 $[a,b]$ 上连续且单调增,
  \[
  F(x)=\begin{cases}\frac{1}{x-a}\int_{a}^{x}f(t)\,\mathrm{d}t,&x\in(a,b],\\f(a),&x=a.\end{cases}
\]
  证明: $F$ 为 $[a,b]$ 上的增函数.

  \item 应用微积分基本定理证明: 连续的奇函数的一切原函数皆为偶函数; 连续的偶函数的原函数中只有一个奇函数 (比较练习题 5.1 中题 5 与题 6).

  \item 设 $f$ 在 $[a,b]$ 上连续, 且 $f(x)>0$. 证明:
  \[
  F(x)=\int_{a}^{x}f(t)\,\mathrm{d}t-\int_{x}^{b}\frac{\mathrm{d}t}{f(t)}
\]
  在 $(a,b)$ 内有一根且仅有一根.
\end{exblock}

\begin{figure}[htbp]
\centering
\begin{tikzpicture}[x=1.10cm,y=1.00cm,>=Stealth]
  \draw[->] (-0.45,0) -- (5.85,0) node[below right] {$x$};
  \draw[->] (0,-0.45) -- (0,3.95) node[left] {$y$};
  \fill[pattern=north east lines] (1.2,1.4) .. controls (2.45,1.52) and (3.25,1.78) .. (3.6,2.2)
    -- (3.6,1.4) -- cycle;
  \fill[pattern=dots] (3.6,3.2) -- (5.0,3.2) .. controls (4.62,3.06) and (4.10,2.62) .. (3.6,2.2) -- cycle;
  \draw[thick] (1.2,1.4) .. controls (2.45,1.52) and (3.25,1.78) .. (3.6,2.2)
    .. controls (4.05,2.55) and (4.60,3.06) .. (5.0,3.2);
  \draw (1.2,1.4) -- (3.6,1.4);
  \draw (1.2,0) -- (1.2,1.4);
  \draw (3.6,0) -- (3.6,3.2);
  \draw (3.6,3.2) -- (5.0,3.2);
  \draw (5.0,0) -- (5.0,3.2);
  \node[below left] at (0,0) {$O$};
  \node[below] at (1.2,-0.10) {$a$};
  \node[below] at (3.6,-0.10) {$\xi$};
  \node[below] at (5.0,-0.10) {$b$};
  \node at (2.20,2.10) {$y=f(x)$};
\end{tikzpicture}
\caption{}\label{fig:6.3.3}
\end{figure}

\begin{exblock}{思考题 6.3}

  \item 设函数 $f$ 在 $(-\infty,+\infty)$ 上连续, 并且 $g(x)=f(x)\int_{0}^{x}f(t)\,\mathrm{d}t$ 为 $(-\infty,+\infty)$ 上的减函数. 证明: $f(x)\equiv 0$.

  \item 求出所有 $[0,+\infty)$ 上的正值连续函数 $g(x)$, 使得对 $\forall x>0$ 有
  \[
  \frac{1}{2}\int_{0}^{x}[g(t)]^{2}\,\mathrm{d}t=\frac{1}{x}\left(\int_{0}^{x}g(t)\,\mathrm{d}t\right)^{2}.
\]
\end{exblock}

% ====================================================================
%  徐森林、薛春华《数学分析》（第一册）
%  \S6.4  Riemann 积分的换元与分部积分
%  源：原书印刷页 386--398（PDF 物理页 392--404）
%  底本：PaddleOCR-VL-1.6，见 chapters/_ch06_4_raw.md
%  约定：全西文标点；行间公式用一对方括号；多行公式用 align 或 aligned；
%        跨册引用一律用 extref（写作「定理 x.y.z」原文）；数学模式内的中文必须包text
% ====================================================================
\section{Riemann 积分的换元与分部积分}\label{sec:6.4}

Newton-Leibniz 公式给了我们计算 Riemann 积分的简捷办法, 它的关键在于求出原函数. 换元 (变量代换) 与分部积分是求原函数或不定积分的两个重要方法. 用它们先求原函数再应用 Newton-Leibniz 公式计算 Riemann 积分, 但有时应用起来不免有些繁琐. 这里将要介绍的是不先求原函数, 而是套用换元与分部积分的方法与步骤直接计算 Riemann 积分.

\begin{theorem}[Riemann 积分的换元 (变量代换)]\label{thm:6.4.1}
设 $f(x)$ 在 $[a,b]$ 上连续, $x=\varphi(t)$ 在 $[\alpha,\beta]$ 上连续可导, 且 $\varphi([\alpha,\beta])\subset[a,b]$, $\varphi(a)=a$, $\varphi(\beta)=b$, 则
\[
\int_{a}^{b}f(x)\,\mathrm{d}x=\int_{a}^{\beta}f(\varphi(t))\varphi'(t)\,\mathrm{d}t=\int_{a}^{\beta}f\circ\varphi(t)\varphi'(t)\,\mathrm{d}t.
\]
\end{theorem}

\begin{proof}
因为 $f$ 在 $[a,b]$ 上连续, 根据微积分基本定理知, $f$ 在 $[a,b]$ 上有原函数 $F$, 即 $F'(x)=f(x)$, 所以,
\[
[F(\varphi(t))]'=F'(\varphi(t))\cdot\varphi'(t)=f(\varphi(t))\cdot\varphi'(t).
\]
再根据 Newton-Leibniz 公式得到
\begin{align*}
\int_{a}^{\beta}f(\varphi(t))\varphi'(t)\,\mathrm{d}t&=F(\varphi(t))\Big|_{\alpha}^{\beta}=F(\varphi(\beta))-F(\varphi(\alpha))\\
&=F(b)-F(a)=\int_{a}^{b}f(x)\,\mathrm{d}x.
\end{align*}
\end{proof}

如果在定理~\ref{thm:6.4.1} 中, 当 $t$ 从 $\alpha$ 变到 $\beta$ 时, $\varphi(t)$ 严格增地从 $a$ 变到 $b$, 则我们有不同特色的另一证明.

\begin{theoremp}[Riemann 积分的换元 (变量代换)]\label{thm:6.4.1p}
设 $f(x)$ 在 $[a,b]$ 上连续, $x=\varphi(t)$ 在 $[\alpha,\beta]$ 上连续可导, 且当 $t$ 从 $\alpha$ 到 $\beta$ 时, $\varphi(t)$ 严格递增地从 $a$ 变到 $b$, 则有
\[
\int_{a}^{b}f(x)\,\mathrm{d}x=\int_{a}^{\beta}f(\varphi(t))\varphi'(t)\,\mathrm{d}t=\int_{a}^{\beta}f\circ\varphi(t)\varphi'(t)\,\mathrm{d}t.
\]
\end{theoremp}

\begin{proof}
由于上式两边的被积函数都是连续的, 因此, Riemann 积分都存在.

作 $[\alpha,\beta]$ 上的分割 $W:\alpha=t_{0}<t_{1}<\cdots<t_{n}=\beta$. 令 $x_{i}=\varphi(t_{i})$, $i=0,1,\cdots,n$. 这就产生了 $[a,b]$ 的一个分割 $T:a=x_{0}<x_{1}<\cdots<x_{n}=b$. 由 Lagrange 中值定理, 我们有
\[
\Delta x_{i}=x_{i}-x_{i-1}=\varphi(t_{i})-\varphi(t_{i-1})=\varphi'(\tau_{i})(t_{i}-t_{i-1})=\varphi'(\tau_{i})\Delta t_{i},
\]
其中 $\tau_{i}\in(t_{i-1},t_{i})$, $i=1,2,\cdots,n$. 设 $K$ 为 $|\varphi'|$ 在 $[\alpha,\beta]$ 上的一个上界, 便有
\[
\Delta x_{i}\leqslant K\Delta t_{i},\quad i=1,2,\cdots,n.
\]
由此可知 $\|T\|\leqslant K\|W\|$. 所以, $\|W\|\to0$ 蕴涵着 $\|T\|\to0$.

令 $\xi_{i}=\varphi(\tau_{i})$, $i=1,2,\cdots,n$, 则易见 $x_{i-1}=\varphi(t_{i-1})<\xi_{i}=\varphi(\tau_{i})<\varphi(t_{i})=x_{i}$, 即 $\xi_{i}\in(x_{i-1},x_{i})$, $i=1,2,\cdots,n$. 于是,
\begin{align*}
\int_{a}^{b}f(x)\,\mathrm{d}x&=\lim_{\|T\|\to0}\sum_{i=1}^{n}f(\xi_{i})\Delta x_{i}=\lim_{\|W\|\to0}\sum_{i=1}^{n}f(\varphi(\tau_{i}))\varphi'(\tau_{i})\Delta t_{i}\\
&=\int_{a}^{\beta}f\circ\varphi(t)\varphi'(t)\,\mathrm{d}t.
\end{align*}
\end{proof}

\begin{theorempp}\label{thm:6.4.1pp}
在定理~\ref{thm:6.4.1p} 中, 如果将 $f(x)$ 在 $[a,b]$ 上「连续」改为「Riemann 可积」, 则结论仍正确.
\end{theorempp}

\begin{proof}
$\forall\varepsilon>0$, 因为 $\varphi'(t)$ 连续, 故 $\exists\delta_{1}>0$, 当 $t',t''\in[\alpha,\beta]$, $|t'-t''|<\delta_{1}$ 时, $|\varphi'(t')-\varphi'(t'')|<\dfrac{\varepsilon}{3M+1}$, 其中 $M=\sup\limits_{x\in[a,b]}|f(x)|$. 又因为 $f$ 在 $[a,b]$ 上 Riemann 可积, 故 $\exists\delta_{2}>0$, 当 $[a,b]$ 的分割 $T$, $\|T\|<\delta_{2}$ 时, $\sum\limits_{i=1}^{n}\omega_{i}^{f}\Delta x_{i}<\dfrac{\varepsilon}{3}$, 且
\[
\left|\sum_{i=1}^{n}f(\xi_{i})\Delta x_{i}-\int_{a}^{b}f(x)\,\mathrm{d}x\right|<\frac{\varepsilon}{3},\quad\xi_{i}\in[x_{i-1},x_{i}].
\]
于是, 对 $[\alpha,\beta]$ 的任一分割 $W:\alpha=t_{0}<t_{1}<\cdots<t_{n}=\beta$, $\forall\tilde{\tau}_{i}\in[t_{i-1},t_{i}]$, 令 $x_{i}=\varphi(t_{i})$, $x_{i}-x_{i-1}=\varphi(t_{i})-\varphi(t_{i-1})=\varphi'(\tau_{i})(t_{i}-t_{i-1})$, 当 $\|W\|<\delta=\min\left\{\delta_{1},\dfrac{\delta_{2}}{K+1}\right\}$ 时, 有 $\|T\|\leqslant K\|W\|<K\cdot\dfrac{\delta_{2}}{K+1}<\delta_{2}$, 其中 $K$ 为 $|\varphi'|$ 在 $[\alpha,\beta]$ 上的一个上界.
\begin{align*}
&\left|\sum_{i=1}^{n}f(\varphi(\tilde{\tau}_{i}))\varphi'(\tilde{\tau}_{i})\Delta t_{i}-\int_{a}^{b}f(x)\,\mathrm{d}x\right|\\
\leqslant{}&\sum_{i=1}^{n}\left|f(\varphi(\tilde{\tau}_{i}))\right|\bullet\left|\varphi'(\tilde{\tau}_{i})-\varphi'(\tau_{i})\right|\Delta t_{i}\\
&+\sum_{i=1}^{n}\left|f(\varphi(\tilde{\tau}_{i}))-f(\varphi(\tau_{i}))\right|\varphi'(\tau_{i})\Delta t_{i}\\
&+\left|\sum_{i=1}^{n}f(\varphi(\tau_{i}))\varphi'(\tau_{i})\Delta t_{i}-\int_{a}^{b}f(x)\,\mathrm{d}x\right|\\
\leqslant{}&M\frac{\varepsilon}{3M+1}+\sum_{i=1}^{n}\omega_{i}^{f}\Delta x_{i}+\left|\sum_{i=1}^{n}f(\xi_{i})\Delta x_{i}-\int_{a}^{b}f(x)\,\mathrm{d}x\right|\\
<{}&\frac{\varepsilon}{3}+\frac{\varepsilon}{3}+\frac{\varepsilon}{3}=\varepsilon.
\end{align*}
由此可知
\[
\int_{a}^{\beta}f(\varphi(t))\varphi'(t)\,\mathrm{d}t=\lim_{\|W\|\to0}\sum_{i=1}^{n}f(\varphi(\tilde{\tau}_{i}))\varphi'(\tilde{\tau}_{i})\Delta t_{i}=\int_{a}^{b}f(x)\,\mathrm{d}x.
\]
\end{proof}

\begin{theorem}[Riemann 积分的分部积分]\label{thm:6.4.2}
设 $u(x)$, $v(x)$ 为 $[a,b]$ 上的连续可导函数, 则有 Riemann 积分的分部积分公式
\[
\int_{a}^{b}u(x)v'(x)\,\mathrm{d}x=u(x)v(x)\Big|_{a}^{b}-\int_{a}^{b}u'(x)v(x)\,\mathrm{d}x,
\]
即
\[
\int_{a}^{b}u(x)\,\mathrm{d}v(x)=u(x)v(x)\Big|_{a}^{b}-\int_{a}^{b}v(x)\,\mathrm{d}u(x).
\]
\end{theorem}

\begin{proof}[证法 1]
\begin{align*}
&\int_{a}^{b}u(x)v'(x)\,\mathrm{d}x\\
&\xlongequal{\text{Newton-Leibniz 公式}}\left[u(x)v(x)-\int u'(x)v(x)\,\mathrm{d}x\right]\Big|_{a}^{b}\\
&=u(x)v(x)\Big|_{a}^{b}-\left[\int u'(x)v(x)\,\mathrm{d}x\right]\Big|_{a}^{b}\\
&=u(x)v(x)\Big|_{a}^{b}-\int_{a}^{b}u'(x)v(x)\,\mathrm{d}x.
\end{align*}
\end{proof}

\begin{proof}[证法 2]
因为 $uv$ 为 $uv'+u'v$ 在 $[a,b]$ 上的一个原函数, 故
\begin{align*}
&\int_{a}^{b}u(x)v'(x)\,\mathrm{d}x+\int_{a}^{b}u'(x)v(x)\,\mathrm{d}x\\
&=\int_{a}^{b}[u(x)v'(x)+u'(x)v(x)]\,\mathrm{d}x\\
&\xlongequal{\text{Newton-Leibniz}}u(x)v(x)\Big|_{a}^{b},
\end{align*}
从而
\[
\int_{a}^{b}u(x)v'(x)\,\mathrm{d}x=u(x)v(x)\Big|_{a}^{b}-\int_{a}^{b}u'(x)v(x)\,\mathrm{d}x.
\]
\end{proof}

进而, 有下面的定理.

\begin{theorem}\label{thm:6.4.3}
设 $u(x)$, $v(x)$ 在 $[a,b]$ 上具有 $n+1$ 阶连续导数, 则有
\begin{align*}
\int_{a}^{b}u(x)v^{(n+1)}(x)\,\mathrm{d}x&=\left[u(x)v^{(n)}(x)-u'(x)v^{(n-1)}(x)+\cdots+(-1)^{n}u^{(n)}(x)v(x)\right]\Big|_{a}^{b}\\
&\quad+(-1)^{n+1}\int_{a}^{b}u^{(n+1)}(x)v(x)\,\mathrm{d}x,\quad n=0,1,2,\cdots.
\end{align*}
\end{theorem}

\begin{proof}
\begin{align*}
\int_{a}^{b}u(x)v^{(n+1)}(x)\,\mathrm{d}x&=u(x)v^{(n)}(x)\Big|_{a}^{b}-\int_{a}^{b}u'(x)v^{(n)}(x)\,\mathrm{d}x\\
&=u(x)v^{(n)}(x)\Big|_{a}^{b}-u'(x)v^{(n-1)}(x)\Big|_{a}^{b}+\int_{a}^{b}u''(x)v^{(n-1)}(x)\,\mathrm{d}x\\
&=\cdots\\
&=\left[u(x)v^{(n)}(x)-u'(x)v^{(n-1)}(x)+\cdots+(-1)^{n}u^{(n)}(x)v(x)\right]\Bigg|_{a}^{b}\\
&\quad+(-1)^{(n+1)}\int_{a}^{b}u^{(n+1)}(x)v(x)\,\mathrm{d}x.
\end{align*}
\end{proof}

\begin{example}\label{ex:6.4.1}
Taylor 公式的积分型余项.

设函数 $f$ 在点 $x_{0}$ 的某开邻域 $U(x_{0})$ 内有 $n+1$ 阶连续导数. 令 $x\in U(x_{0})$, $u(t)=(x-t)^{n}$, $v(t)=f(t)$, $t\in[x_{0},x]\ (\text{或}\ t\in[x,x_{0}])$. 根据定理~\ref{thm:6.4.3}, 有
\begin{align*}
&\int_{x_{0}}^{x}(x-t)^{n}f^{(n+1)}(t)\,\mathrm{d}t\\
&=\left[(x-t)^{n}f^{(n)}(t)+n(x-t)^{n-1}f^{(n-1)}(t)+\cdots+n!f(t)\right]\Bigg|_{x_{0}}^{x}+\int_{x_{0}}^{x}0\cdot f(t)\,\mathrm{d}t\\
&=n!f(x)-n!\left[f(x_{0})+f'(x_{0})(x-x_{0})+\cdots+\frac{f^{(n)}(x_{0})}{n!}(x-x_{0})^{n}\right]\\
&=n!R_{n}(x),
\end{align*}
由此求得
\[
R_{n}(x)=\frac{1}{n!}\int_{x_{0}}^{x}f^{(n+1)}(t)(x-t)^{n}\,\mathrm{d}t,
\]
并称它为 Taylor 公式的积分型余项.

由于 $f^{(n+1)}(t)$ 连续, $(x-t)^{n}$ 在 $[x_{0},x]$ (或 $[x,x_{0}]$) 上保持同号, 因此, 由推广的积分第一中值定理, 有
\begin{align*}
R_{n}(x)&=\frac{1}{n!}\int_{x_{0}}^{x}f^{(n+1)}(t)(x-t)^{n}\,\mathrm{d}t\\
&=\frac{1}{n!}f^{(n+1)}(\xi)\int_{x_{0}}^{x}(x-t)^{n}\,\mathrm{d}t\\
&=\frac{1}{n!}f^{(n+1)}(\xi)\frac{-(x-t)^{n+1}}{n+1}\Bigg|_{x_{0}}^{x}\\
&=\frac{f^{(n+1)}(\xi)}{(n+1)!}(x-x_{0})^{n+1},\quad\xi\in(x_{0},x)\ (\text{或}\ (x,x_{0})),
\end{align*}
其中 $\xi=x_{0}+\theta(x-x_{0})$, $0<\theta<1$. 这就是熟悉的 Taylor 公式的 Lagrange 型余项.

如果直接应用积分第一中值定理, 就得到 Taylor 公式的 Cauchy 余项
\begin{align*}
R_{n}(x)&=\frac{1}{n!}\int_{x_{0}}^{x}f^{(n+1)}(t)(x-t)^{n}\,\mathrm{d}t\\
&=\frac{1}{n!}f^{(n+1)}(\xi)(x-\xi)^{n}\int_{x_{0}}^{x}\mathrm{d}t\\
&=\frac{1}{n!}f^{(n+1)}(\xi)(x-\xi)^{n}(x-x_{0})\\
&=\frac{1}{n!}f^{(n+1)}(x+\theta(x-x_{0}))(1-\theta)^{n}(x-x_{0})^{n+1},
\end{align*}
其中 $\xi=x_{0}+\theta(x-x_{0})$, $x-\xi=(x-x_{0})-\theta(x-x_{0})=(1-\theta)(x-x_{0})$.
\end{example}

\begin{example}\label{ex:6.4.2}
求下列 Riemann 积分:
\[
(1)\ \int_{0}^{1}\sqrt{1-x^{2}}\,\mathrm{d}x;\qquad(2)\ \int_{0}^{\frac{\pi}{2}}\sin t\cos^{2}t\,\mathrm{d}t;\qquad(3)\ \int_{1}^{\mathrm{e}}x^{2}\ln x\,\mathrm{d}x.
\]
\end{example}

\begin{solution}
(1)
\begin{align*}
\int_{0}^{1}\sqrt{1-x^{2}}\,\mathrm{d}x&\xlongequal{x=\sin t}\int_{0}^{\frac{\pi}{2}}\sqrt{1-\sin^{2}t}\cos t\,\mathrm{d}t\\
&=\int_{0}^{\frac{\pi}{2}}\cos^{2}t\,\mathrm{d}t=\frac{1}{2}\int_{0}^{\frac{\pi}{2}}(1+\cos2t)\,\mathrm{d}t\\
&=\frac{1}{2}\left(t+\frac{1}{2}\sin2t\right)\Bigg|_{0}^{\frac{\pi}{2}}=\frac{\pi}{4}.
\end{align*}
或者由下面的例~\ref{ex:6.4.5} 的结论, 得
\[
\int_{0}^{1}\sqrt{1-x^{2}}\,\mathrm{d}x=\int_{0}^{\frac{\pi}{2}}\cos^{2}t\,\mathrm{d}t=\frac{1}{2}\cdot\frac{\pi}{2}=\frac{\pi}{4}.
\]
(2)
\[
\int_{0}^{\frac{\pi}{2}}\sin t\cos^{2}t\,\mathrm{d}t\xlongequal{x=\cos t}\int_{1}^{0}x^{2}(-\mathrm{d}x)=-\frac{x^{3}}{3}\Bigg|_{1}^{0}=\frac{1}{3}.
\]
或者
\begin{align*}
\int_{0}^{\frac{\pi}{2}}\sin t\cos^{2}t\,\mathrm{d}t&=\int_{0}^{\frac{\pi}{2}}\sin t(1-\sin^{2}t)\,\mathrm{d}t\\
&=\int_{0}^{\frac{\pi}{2}}\sin t\,\mathrm{d}t-\int_{0}^{\frac{\pi}{2}}\sin^{3}t\,\mathrm{d}t=-\cos t\Bigg|_{0}^{\frac{\pi}{2}}-\frac{2!!}{3!!}=1-\frac{2}{3}=\frac{1}{3}.
\end{align*}
(3)
\begin{align*}
\int_{1}^{\mathrm{e}}x^{2}\ln x\,\mathrm{d}x&\xlongequal{\text{升幂}}\int_{1}^{\mathrm{e}}\ln x\,\mathrm{d}\left(\frac{x^{3}}{3}\right)\\
&\xlongequal{\text{分部积分}}\frac{1}{3}\left(x^{3}\ln x\Bigg|_{1}^{\mathrm{e}}-\int_{1}^{\mathrm{e}}x^{3}\cdot\frac{1}{x}\,\mathrm{d}x\right)\\
&=\frac{1}{3}\left(\mathrm{e}^{3}-\frac{x^{3}}{3}\Bigg|_{1}^{\mathrm{e}}\right)=\frac{1}{3}\left(\mathrm{e}^{3}-\frac{\mathrm{e}^{3}}{3}+\frac{1}{3}\right)=\frac{1}{9}(2\mathrm{e}^{3}+1).
\end{align*}
\end{solution}

\begin{example}\label{ex:6.4.3}
求 $\displaystyle\int_{0}^{1}\frac{\ln(1+x)}{1+x^{2}}\,\mathrm{d}x$.
\end{example}

\begin{solution}[解法 1]
\begin{align*}
\int_{0}^{1}\frac{\ln(1+x)}{1+x^{2}}\,\mathrm{d}x&\xlongequal{x=\tan t}\int_{0}^{\frac{\pi}{4}}\frac{\ln(1+\tan t)}{1+\tan^{2}t}\sec^{2}t\,\mathrm{d}t\\
&=\int_{0}^{\frac{\pi}{4}}\ln\frac{\cos t+\sin t}{\cos t}\,\mathrm{d}t=\int_{0}^{\frac{\pi}{4}}\ln\frac{\sqrt{2}\cos\left(\frac{\pi}{4}-t\right)}{\cos t}\,\mathrm{d}t\\
&=\int_{0}^{\frac{\pi}{4}}\ln\sqrt{2}\,\mathrm{d}t+\int_{0}^{\frac{\pi}{4}}\ln\cos\left(\frac{\pi}{4}-t\right)\mathrm{d}t-\int_{0}^{\frac{\pi}{4}}\ln\cos t\,\mathrm{d}t\\
&\xlongequal{u=\frac{\pi}{4}-t}\frac{\pi}{8}\ln2+\int_{\frac{\pi}{4}}^{0}\ln\cos u(-\mathrm{d}u)-\int_{0}^{\frac{\pi}{4}}\ln\cos t\,\mathrm{d}t\\
&=\frac{\pi}{8}\ln2+\int_{0}^{\frac{\pi}{4}}\ln\cos t\,\mathrm{d}t-\int_{0}^{\frac{\pi}{4}}\ln\cos t\,\mathrm{d}t=\frac{\pi}{8}\ln2.
\end{align*}
上述计算的妙处在于套用了换元的方法与步骤并利用定理~\ref{thm:6.4.1} 导致两个 Riemann 积分 $\int_{0}^{\frac{\pi}{4}}\ln\cos u\,\mathrm{d}u$ 与 $\int_{0}^{\frac{\pi}{4}}\ln\cos t\,\mathrm{d}t$ 相消, 从而得到 $\int_{0}^{1}\frac{\ln(1+x)}{1+x^{2}}\,\mathrm{d}x=\frac{\pi}{8}\ln2$. 但是, 无法用同样的换元来求出 $\int\frac{\ln(1+x)}{1+x^{2}}\,\mathrm{d}x$, 当然也无法用 Newton-Leibniz 公式.
\end{solution}

\begin{solution}[解法 2]
令
\[
I(\alpha)=\int_{0}^{1}\frac{\ln(1+\alpha x)}{1+x^{2}}\,\mathrm{d}x,
\]
其中 $x$ 为积分变量, $\alpha$ 为参变量. 则
\begin{align*}
I'(\alpha)&=\int_{0}^{1}\left(\frac{\ln(1+\alpha x)}{1+x^{2}}\right)'_{\alpha}\mathrm{d}x=\int_{0}^{1}\frac{x}{(1+\alpha x)(1+x^{2})}\,\mathrm{d}x\\
&=\frac{1}{1+\alpha^{2}}\int_{0}^{1}\left(\frac{-\alpha}{1+\alpha x}+\frac{x+\alpha}{1+x^{2}}\right)\mathrm{d}x\\
&=\frac{1}{1+\alpha^{2}}\left[-\ln(1+\alpha x)+\frac{1}{2}\ln(1+x^{2})+\alpha\arctan x\right]\Bigg|_{0}^{1}\\
&=-\frac{\ln(1+\alpha)}{1+\alpha^{2}}+\frac{\ln2}{2(1+\alpha^{2})}+\frac{\pi\alpha}{4(1+\alpha^{2})}.
\end{align*}
因此,
\begin{align*}
\int_{0}^{1}\frac{\ln(1+x)}{1+x^{2}}\,\mathrm{d}x&=I(1)=I(1)-I(0)\\
&\xlongequal{\text{Newton-Leibniz 公式}}\int_{0}^{1}I'(\alpha)\,\mathrm{d}\alpha\\
&=-\int_{0}^{1}\frac{\ln(1+\alpha)}{1+\alpha^{2}}\,\mathrm{d}\alpha+\frac{\ln2}{2}\int_{0}^{1}\frac{\mathrm{d}\alpha}{1+\alpha^{2}}+\frac{\pi}{4}\int_{0}^{1}\frac{\alpha}{1+\alpha^{2}}\,\mathrm{d}\alpha\\
&=-\int_{0}^{1}\frac{\ln(1+x)}{1+x^{2}}\,\mathrm{d}x+\frac{\ln2}{2}\arctan\alpha\Bigg|_{0}^{1}+\frac{\pi}{8}\ln(1+\alpha^{2})\Bigg|_{0}^{1}\\
&=-\int_{0}^{1}\frac{\ln(1+x)}{1+x^{2}}\,\mathrm{d}x+\frac{\pi\ln2}{8}+\frac{\pi}{8}\ln2,
\end{align*}
移项化简后得到
\[
\int_{0}^{1}\frac{\ln(1+x)}{1+x^{2}}\,\mathrm{d}x=\frac{1}{2}\cdot\frac{\pi}{4}\ln2=\frac{\pi}{8}\ln2.
\]
这个方法的妙处在于等号右边出现了 $-\int_{0}^{1}\frac{\ln(1+x)}{1+x^{2}}\,\mathrm{d}x$, 从而可解出 $\int_{0}^{1}\frac{\ln(1+x)}{1+x^{2}}\,\mathrm{d}x$.

但是, 遗憾的是
\[
I'(\alpha)=\int_{0}^{1}\left(\frac{\ln(1+\alpha x)}{1+x^{2}}\right)'_{\alpha}\mathrm{d}x,
\]
即求导与积分号可变换, 这里没有足够的理由, 读者学会了参变量积分求导的定理后再来检验它, 这等式确实是成立的.

我们给出解法 2, 是为了告诉读者除了解法 1 外还有一个很巧妙的方法, 还想告诉读者要思路开阔、大胆设想、不要受到束缚, 当然也不要忘了严格论证.
\end{solution}

\begin{example}\label{ex:6.4.4}
求 $\displaystyle\int_{0}^{\pi}\frac{x\sin x}{1+\cos^{2}x}\,\mathrm{d}x$.
\end{example}

\begin{solution}
\begin{align*}
\int_{0}^{\pi}\frac{x\sin x}{1+\cos^{2}x}\,\mathrm{d}x&\xlongequal{x=\pi-t}\int_{\pi}^{0}\frac{(\pi-t)\sin(\pi-t)}{1+\cos^{2}(\pi-t)}\,\mathrm{d}(\pi-t)\\
&=\int_{0}^{\pi}\frac{\pi\sin t}{1+\cos^{2}t}\,\mathrm{d}t+\int_{0}^{\pi}\frac{-t\sin t}{1+\cos^{2}t}\,\mathrm{d}t,
\end{align*}
移项化简得到
\begin{align*}
\int_{0}^{\pi}\frac{x\sin x}{1+\cos^{2}x}\,\mathrm{d}x&=\frac{\pi}{2}\int_{0}^{\pi}\frac{\sin x}{1+\cos^{2}x}\,\mathrm{d}x=-\frac{\pi}{2}\int_{0}^{\pi}\frac{\mathrm{d}\cos x}{1+\cos^{2}x}\\
&=-\frac{\pi}{2}\arctan(\cos x)\Bigg|_{0}^{\pi}=-\frac{\pi}{2}\left(\arctan(\cos\pi)-\arctan(\cos0)\right)\\
&=-\frac{\pi}{2}\left(-\frac{\pi}{4}-\frac{\pi}{4}\right)=\frac{\pi^{2}}{4}.
\end{align*}
\end{solution}

\begin{example}\label{ex:6.4.5}
证明:

(1) Wallis 公式
\[
\int_{0}^{\frac{\pi}{2}}\sin^{n}x\,\mathrm{d}x=\int_{0}^{\frac{\pi}{2}}\cos^{n}x\,\mathrm{d}x=\begin{cases}\dfrac{(n-1)!!}{n!!}\dfrac{\pi}{2}, & n\ \text{为偶数},\\[3mm]\dfrac{(n-1)!!}{n!!}, & n\ \text{为奇数};\end{cases}
\]
(2)
\[
\frac{\pi}{2}=\lim_{m\to+\infty}\left[\frac{(2m)!!}{(2m-1)!!}\right]^{2}\cdot\frac{1}{2m+1},
\]
其中
\[
(2m)!!=(2m)\cdot(2m-2)\cdots2,\quad(2m-1)!!=(2m-1)(2m-3)\cdots3\cdot1.
\]
\end{example}

\begin{proof}
(1)
\begin{align*}
I_{n}&=\int_{0}^{\frac{\pi}{2}}\sin^{n}x\,\mathrm{d}x=-\int_{0}^{\frac{\pi}{2}}\sin^{n-1}x\,\mathrm{d}\cos x\\
&\xlongequal{\text{分部积分}}-\sin^{n-1}x\cos x\Bigg|_{0}^{\frac{\pi}{2}}+\int_{0}^{\frac{\pi}{2}}(n-1)\sin^{n-2}x\cos^{2}x\,\mathrm{d}x\\
&=(n-1)\int_{0}^{\frac{\pi}{2}}\sin^{n-2}x(1-\sin^{2}x)\,\mathrm{d}x\\
&=(n-1)I_{n-2}-(n-1)I_{n},
\end{align*}
移项整理后得到递推公式
\[
I_{n}=\frac{n-1}{n}I_{n-2}\quad(n\geqslant2).
\]
由于
\[
I_{0}=\int_{0}^{\frac{\pi}{2}}\sin^{0}x\,\mathrm{d}x=\int_{0}^{\frac{\pi}{2}}\mathrm{d}x=\frac{\pi}{2},
\]
\[
I_{1}=\int_{0}^{\frac{\pi}{2}}\sin x\,\mathrm{d}x=-\cos x\Bigg|_{0}^{\frac{\pi}{2}}=\cos0-\cos\frac{\pi}{2}=1-0=1.
\]
重复应用递推公式便得
\begin{align*}
I_{2m}&=\frac{2m-1}{2m}I_{2m-2}=\frac{2m-1}{2m}\cdot\frac{2m-3}{2m-2}I_{2m-4}\\
&=\cdots\\
&=\frac{2m-1}{2m}\cdot\frac{2m-3}{2m-2}\cdots\frac{1}{2}I_{0}\\
&=\frac{(2m-1)!!}{(2m)!!}\frac{\pi}{2},
\end{align*}
\begin{align*}
I_{2m-1}&=\frac{2m-2}{2m-1}I_{2m-3}=\frac{2m-2}{2m-1}\cdot\frac{2m-4}{2m-3}I_{2m-5}\\
&=\cdots\\
&=\frac{2m-2}{2m-1}\cdot\frac{2m-4}{2m-3}\cdots\frac{2}{3}I_{1}\\
&=\frac{(2m-2)!!}{(2m-1)!!}.
\end{align*}
而
\begin{align*}
\int_{0}^{\frac{\pi}{2}}\cos^{n}x\,\mathrm{d}x&\xlongequal{x=\frac{\pi}{2}-t}-\int_{\frac{\pi}{2}}^{0}\cos^{n}\left(\frac{\pi}{2}-t\right)\mathrm{d}t\\
&=\int_{0}^{\frac{\pi}{2}}\sin^{n}t\,\mathrm{d}t=\int_{0}^{\frac{\pi}{2}}\sin^{n}x\,\mathrm{d}x\\
&=\begin{cases}\dfrac{(n-1)!!}{n!!}\cdot\dfrac{\pi}{2}, & n\ \text{为偶数},\\[
3mm]\dfrac{(n-1)!!}{n!!}, & n\ \text{为奇数}.\end{cases}
\end{align*}
或类似 $\int_{0}^{\frac{\pi}{2}}\sin^{n}x\,\mathrm{d}x$, 应用分部积分法得到递推公式求得上述公式.
(2) 因为
\[
\int_{0}^{\frac{\pi}{2}}\sin^{2m+1}x\,\mathrm{d}x<\int_{0}^{\frac{\pi}{2}}\sin^{2m}x\,\mathrm{d}x<\int_{0}^{\frac{\pi}{2}}\sin^{2m-1}x\,\mathrm{d}x.
\]
所以由 (1), 上式就是
\[
\frac{(2m)!!}{(2m+1)!!}<\frac{(2m-1)!!}{(2m)!!}\cdot\frac{\pi}{2}<\frac{(2m-2)!!}{(2m-1)!!}.
\]
由此得到
\[
A_{m}=\frac{1}{2m+1}\left[\frac{(2m)!!}{(2m-1)!!}\right]^{2}<\frac{\pi}{2}<\left[\frac{(2m)!!}{(2m-1)!!}\right]^{2}\cdot\frac{1}{2m}=B_{m}.
\]
因为
\begin{align*}
0<B_{m}-A_{m}&=\left[\frac{(2m)!!}{(2m-1)!!}\right]^{2}\left(\frac{1}{2m}-\frac{1}{2m+1}\right)\\
&=\left[\frac{(2m)!!}{(2m-1)!!}\right]^{2}\frac{1}{2m(2m+1)}<\frac{1}{2m}\cdot\frac{\pi}{2}\to0\quad(m\to+\infty),
\end{align*}
故由夹逼定理立即有
\[
\lim_{m\to+\infty}(B_{m}-A_{m})=0.
\]
再由 $0<\frac{\pi}{2}-A_{m}<B_{m}-A_{m}\to0$ $(m\to+\infty)$ 得到
\[
\lim_{m\to+\infty}\left(\frac{\pi}{2}-A_{m}\right)=0,
\]
\begin{align*}
\lim_{m\to+\infty}\left[\frac{(2m)!!}{(2m-1)!!}\right]^{2}\frac{1}{2m+1}&=\lim_{m\to+\infty}A_{m}=\lim_{m\to+\infty}\left[\frac{\pi}{2}-\left(\frac{\pi}{2}-A_{m}\right)\right]\\
&=\frac{\pi}{2}-0=\frac{\pi}{2}.
\end{align*}
\end{proof}

\begin{example}\label{ex:6.4.6}
计算 $\displaystyle I(m,n)=\int_{0}^{\frac{\pi}{2}}\cos^{m}x\sin^{n}x\,\mathrm{d}x$.
\end{example}

\begin{solution}
\begin{align*}
I(m,n)&=\frac{1}{n+1}\int_{0}^{\frac{\pi}{2}}\cos^{m-1}x\,\mathrm{d}\sin^{n+1}x\\
&=\frac{1}{n+1}\left[\cos^{m-1}x\sin^{n+1}x\Bigg|_{0}^{\frac{\pi}{2}}-\int_{0}^{\frac{\pi}{2}}(m-1)\cos^{m-2}x(-\sin^{n+2}x)\,\mathrm{d}x\right]\\
&=\frac{m-1}{n+1}\int_{0}^{\frac{\pi}{2}}\cos^{m-2}x\sin^{n}x(1-\cos^{2}x)\,\mathrm{d}x\\
&=\frac{m-1}{n+1}\left[I(m-2,n)-I(m,n)\right],
\end{align*}
移项化简得
\begin{align*}
I(m,n)&=\frac{m-1}{m+n}I(m-2,n)\\
&=\begin{cases}\dfrac{(m-1)(m-3)\cdots1}{(m+n)(m+n-2)\cdots(n+2)}I(0,n), & m\ \text{为偶数},\\[
3mm]\dfrac{(m-1)(m-3)\cdots2}{(m+n)(m+n-2)\cdots(n+3)}I(1,n), & m\ \text{为奇数},\end{cases}\\
&=\begin{cases}\dfrac{(m-1)!!(n-1)!!}{(m+n)!!}\dfrac{\pi}{2}, & m\ \text{与}\ n\ \text{均为偶数},\\[3mm]\dfrac{(m-1)!!(n-1)!!}{(m+n)!!}, & m\ \text{与}\ n\ \text{中有一个奇数}.\end{cases}
\end{align*}
\end{solution}
\begin{example}\label{ex:6.4.7}
求 $\displaystyle\lim_{n\to+\infty}\left[\left(1+\frac{1}{n}\right)\left(1+\frac{2}{n}\right)\cdots\left(1+\frac{n}{n}\right)\right]^{\frac{1}{n}}$.
\end{example}
\begin{solution}
因为
\begin{align*}
\lim_{n\to+\infty}\frac{1}{n}\sum_{k=1}^{n}\left[\ln\left(1+\frac{k}{n}\right)\right]&=\int_{0}^{1}\ln(1+x)\,\mathrm{d}x=\int_{0}^{1}\ln(1+x)\,\mathrm{d}(1+x)\\
&\xlongequal{\text{分部积分}}(1+x)\ln(1+x)\Bigg|_{0}^{1}-\int_{0}^{1}(1+x)\cdot\frac{1}{1+x}\,\mathrm{d}x\\
&=2\ln2-\int_{0}^{1}\mathrm{d}x=2\ln2-1,
\end{align*}
所以
\begin{align*}
&\lim_{n\to+\infty}\left[\left(1+\frac{1}{n}\right)\left(1+\frac{2}{n}\right)+\cdots+\left(1+\frac{n}{n}\right)\right]^{\frac{1}{n}}\\
&=\exp\left[\lim_{n\to+\infty}\frac{1}{n}\sum_{k=1}^{n}\ln\left(1+\frac{k}{n}\right)\right]=\mathrm{e}^{2\ln2-1}=\frac{\mathrm{e}^{\ln4}}{\mathrm{e}}=\frac{4}{\mathrm{e}}.
\end{align*}
\end{solution}
\begin{example}\label{ex:6.4.8}
计算 $\displaystyle\int_{0}^{\frac{\pi}{4}}\frac{\mathrm{d}x}{\cos x}$.
\end{example}
\begin{solution}
\begin{align*}
\int_{0}^{\frac{\pi}{4}}\frac{\mathrm{d}x}{\cos x}&=\int_{0}^{\frac{\pi}{4}}\frac{\mathrm{d}\sin x}{1-\sin^{2}x}\xlongequal{t=\sin x}\int_{0}^{\frac{\sqrt{2}}{2}}\frac{\mathrm{d}t}{1-t^{2}}\\
&=\frac{1}{2}\int_{0}^{\frac{\sqrt{2}}{2}}\left(\frac{1}{1+t}+\frac{1}{1-t}\right)\mathrm{d}t\\
&=\frac{1}{2}\ln\frac{1+t}{1-t}\Bigg|_{0}^{\frac{\sqrt{2}}{2}}=\frac{1}{2}\ln\frac{1+\frac{\sqrt{2}}{2}}{1-\frac{\sqrt{2}}{2}}=\ln(\sqrt{2}+1).
\end{align*}
\end{solution}
\begin{example}\label{ex:6.4.9}
(1) 设 $f$ 为连续的奇函数, 则对 $\forall a\in\mathbb{R}$, 有
\[
\int_{-a}^{a}f(x)\,\mathrm{d}x=0;
\]
(2) 设 $f$ 为连续的偶函数, 则对 $\forall a\in\mathbb{R}$, 有
\[
\int_{-a}^{a}f(x)\,\mathrm{d}x=2\int_{0}^{a}f(x)\,\mathrm{d}x;
\]
(3) 设 $f$ 是周期为 $T$ 的连续函数, 则对 $\forall a\in\mathbb{R}$, 有
\[
\int_{a}^{a+T}f(x)\,\mathrm{d}x=\int_{0}^{T}f(x)\,\mathrm{d}x.
\]
\end{example}

\begin{proof}
(1)
\begin{align*}
\int_{-a}^{a}f(x)\,\mathrm{d}x&=\int_{-a}^{0}f(x)\,\mathrm{d}x+\int_{0}^{a}f(x)\,\mathrm{d}x\xlongequal{x=-t}\int_{a}^{0}f(-t)\,\mathrm{d}(-t)+\int_{0}^{a}f(x)\,\mathrm{d}x\\
&=\int_{0}^{a}-f(t)\,\mathrm{d}t+\int_{0}^{a}f(x)\,\mathrm{d}x=0.
\end{align*}
(2)
\begin{align*}
\int_{-a}^{a}f(x)\,\mathrm{d}x&=\int_{-a}^{0}f(x)\,\mathrm{d}x+\int_{0}^{a}f(x)\,\mathrm{d}x\\
&\xlongequal{x=-t}\int_{a}^{0}f(-t)\,\mathrm{d}(-t)+\int_{0}^{a}f(x)\,\mathrm{d}x\\
&=\int_{0}^{a}f(t)\,\mathrm{d}t+\int_{0}^{a}f(x)\,\mathrm{d}x=2\int_{0}^{a}f(x)\,\mathrm{d}x.
\end{align*}
(3) 因为
\begin{align*}
\int_{T}^{a+T}f(x)\,\mathrm{d}x&\xlongequal{x=t+T}\int_{0}^{a}f(t+T)\,\mathrm{d}(t+T)=\int_{0}^{a}f(t)\,\mathrm{d}t\\
&=\int_{0}^{a}f(x)\,\mathrm{d}x,
\end{align*}
所以
\begin{align*}
\int_{a}^{a+T}f(x)\,\mathrm{d}x&=\int_{a}^{0}f(x)\,\mathrm{d}x+\int_{0}^{T}f(x)\,\mathrm{d}x+\int_{T}^{a+T}f(x)\,\mathrm{d}x\\
&=-\int_{0}^{a}f(x)\,\mathrm{d}x+\int_{0}^{T}f(x)\,\mathrm{d}x+\int_{0}^{a}f(x)\,\mathrm{d}x\\
&=\int_{0}^{T}f(x)\,\mathrm{d}x.
\end{align*}
\end{proof}

\begin{remark}\label{rem:6.4.1}
在例~\ref{ex:6.4.9} 中, 如果将 $f$「连续」改为「Riemann 可积」, 则其结论仍正确. 只需应用 Riemann 积分定义证明:
\[
(1)\ \int_{0}^{a}f(x)\,\mathrm{d}x=-\int_{-a}^{0}f(x)\,\mathrm{d}x;\qquad(2)\ \int_{0}^{a}f(x)\,\mathrm{d}x=\int_{-a}^{0}f(x)\,\mathrm{d}x;
\]
\[
(3)\ \int_{T}^{a+T}f(x)\,\mathrm{d}x=\int_{0}^{a}f(x)\,\mathrm{d}x.
\]
或者应用定理~\ref{thm:6.4.1pp} 及例~\ref{ex:6.4.9} 中的方法.
\end{remark}

\begin{exblock}{练习题 6.4}

  \item 证明: $\lim\limits_{n\to+\infty}\prod\limits_{k=1}^{n}\left(1+\dfrac{k}{n}\right)^{\frac{k}{n^{2}}}=\mathrm{e}^{\frac{1}{4}}$.

  (提示: $\lim\limits_{n\to+\infty}\dfrac{1}{n}\sum\limits_{k=1}^{n}\dfrac{k}{n}\ln\left(1+\dfrac{k}{n}\right)=\int_{0}^{1}x\ln(1+x)\,\mathrm{d}x$.)

  \item 求下列极限:
  \begin{enumerate}[label=(\arabic*)]
    \item $\lim\limits_{n\to+\infty}\sin\dfrac{\pi}{n}\sum\limits_{k=1}^{n}\left(2+\cos\dfrac{k\pi}{n}\right)^{-1}$;
    \item $\lim\limits_{n\to+\infty}\sum\limits_{i=1}^{n}\dfrac{i}{n^{2}}\sin\left(\dfrac{i\pi}{2n}\right)^{2}$.
  \end{enumerate}

  \item 求下列 Riemann 积分:
  \begin{enumerate}[label=(\arabic*)]
    \item $\int_{0}^{1}x(2-x^{2})^{\frac{1}{2}}\,\mathrm{d}x$;
    \item $\int_{0}^{\frac{\pi}{2}}\dfrac{\mathrm{d}x}{1+\varepsilon\cos x}$.
  \end{enumerate}

  \item 证明: $\int_{\mathrm{e}^{-2n\pi}}^{1}\left|\left[\cos\left(\ln\dfrac{1}{x}\right)\right]'\right|\mathrm{d}x=4n$, $n\in\mathbb{N}$.

  \item 设 $I$ 为一个开区间, 函数 $f$ 在 $I$ 内连续, $a<b$ 且 $a,b\in I$. 证明:
  \[
  \lim_{h\to0}\frac{1}{h}\int_{a}^{b}[f(x+h)-f(x)]\,\mathrm{d}x=f(b)-f(a).
\]
  (提示: 先作换元 $t=x+h$, 再应用积分中值定理.)

  \item 计算下列 Riemann 积分:
  \begin{enumerate}[label=(\arabic*)]
    \item $\int_{0}^{\pi}\sin^{3}x\,\mathrm{d}x$;
    \item $\int_{-\pi}^{\pi}x^{2}\cos x\,\mathrm{d}x$;
    \item $\int_{0}^{\sqrt{3}}x\arctan x\,\mathrm{d}x$;
    \item $\int_{0}^{3}\dfrac{x}{1+\sqrt{1+x}}\,\mathrm{d}x$;
    \item $\int_{-1}^{0}(2x+1)\sqrt{1-x-x^{2}}\,\mathrm{d}x$;
    \item $\int_{\mathrm{e}^{-1}}^{\mathrm{e}}|\ln x|\,\mathrm{d}x$;
    \item $\int_{0}^{5}[x]\sin\dfrac{\pi x}{5}\,\mathrm{d}x$;
    \item $\int_{0}^{a}x^{2}\sqrt{a^{2}-x^{2}}\,\mathrm{d}x$;
    \item $\int_{0}^{\ln2}\sqrt{\mathrm{e}^{x}-1}\,\mathrm{d}x$;
    \item $\int_{0}^{1}x^{n}\ln x\,\mathrm{d}x$, $n\in\mathbb{N}$;
    \item $\int_{0}^{\frac{\pi}{2}}\dfrac{\cos x\sin x}{a^{2}\sin^{2}x+b^{2}\cos^{2}x}\,\mathrm{d}x$, $ab\neq0$;
    \item $\int_{0}^{a}\ln\left(x+\sqrt{a^{2}+x^{2}}\right)\mathrm{d}x$, $a>0$.
  \end{enumerate}

  \item 证明: $\lim\limits_{n\to+\infty}\int_{0}^{\frac{\pi}{2}}\sin^{n}x\,\mathrm{d}x=0$. 由此推证
  \begin{enumerate}[label=(\arabic*)]
    \item $\lim\limits_{n\to+\infty}\dfrac{(2n)!!}{(2n+1)!!}=0$;
    \item $\lim\limits_{n\to+\infty}\dfrac{(2n-1)!!}{(2n)!!}=0$.
  \end{enumerate}

  \item 设 $f$ 为连续函数, 证明:
  \begin{enumerate}[label=(\arabic*)]
    \item $\int_{0}^{\frac{\pi}{2}}f(\cos x)\,\mathrm{d}x=\int_{0}^{\frac{\pi}{2}}f(\sin x)\,\mathrm{d}x$;
    \item $\int_{0}^{\pi}xf(\sin x)\,\mathrm{d}x=\dfrac{\pi}{2}\int_{0}^{\pi}f(\sin x)\,\mathrm{d}x$.
  \end{enumerate}
  并由 (2) 证明 $\int_{0}^{\pi}\dfrac{x\sin x}{1+\cos^{2}x}\,\mathrm{d}x=\dfrac{\pi^{2}}{4}$.

  \item 计算下列 Riemann 积分:
  \begin{enumerate}[label=(\arabic*)]
    \item $\int_{0}^{\frac{\pi}{2}}\dfrac{\cos^{2}x}{\cos x+\sin x}\,\mathrm{d}x$;
    \item $\int_{0}^{\frac{\pi}{2}}\dfrac{\sin^{2}x}{\cos x+\sin x}\,\mathrm{d}x$.
  \end{enumerate}

  \item 设 $f$ 在 $[a,b]$ 上连续可导, 应用分部积分证明:
  \begin{enumerate}[label=(\arabic*)]
    \item $\lim\limits_{\lambda\to+\infty}\int_{a}^{b}f(x)\cos\lambda x\,\mathrm{d}x=0$;
    \item $\lim\limits_{\lambda\to+\infty}\int_{a}^{b}f(x)\sin\lambda x\,\mathrm{d}x=0$.
  \end{enumerate}
  若将 $f$ 放宽成「在 $[a,b]$ 上 Riemann 可积」, 则结论仍然正确. 那就是著名的 Riemann 引理. 本书第 3 册中将讨论它.

  \item 设 $f$ 在 $[0,1]$ 上 2 阶连续可导, 证明:
  \[
  \int_{0}^{1}f(x)\,\mathrm{d}x=\frac{1}{2}\left[f(0)+f(1)-\int_{0}^{1}x(1-x)f''(x)\,\mathrm{d}x\right].
\]

  \item 应用分部积分证明: $\int_{0}^{2\pi}\left(\int_{x}^{2\pi}\frac{\sin t}{t}\,\mathrm{d}t\right)\mathrm{d}x=0$.

  \item 设 $f$ 在 $[a,+\infty)$ 中连续, 对 $\forall a>0$, 证明:
  \[
  \int_{0}^{a}\left(\int_{0}^{x}f(t)\,\mathrm{d}t\right)\mathrm{d}x=\int_{0}^{a}f(x)(a-x)\,\mathrm{d}x.
\]

  \item 计算下列 Riemann 积分:
  \begin{enumerate}[label=(\arabic*)]
    \item $\int_{0}^{\frac{\pi}{2}}\cos^{5}x\sin2x\,\mathrm{d}x$;
    \item $\int_{0}^{1}\sqrt{4-x^{2}}\,\mathrm{d}x$;
    \item $\int_{0}^{a}x^{2}\sqrt{a^{2}-x^{2}}\,\mathrm{d}x$ $(a>0)$;
    \item $\int_{0}^{1}\dfrac{\mathrm{d}x}{(x^{2}-x+1)^{\frac{3}{2}}}$;
    \item $\int_{0}^{1}\dfrac{\mathrm{d}x}{\mathrm{e}^{x}+\mathrm{e}^{-x}}$;
    \item $\int_{0}^{\frac{\pi}{2}}\dfrac{\cos x}{1+\sin^{2}x}\,\mathrm{d}x$;
    \item $\int_{0}^{1}\arcsin x\,\mathrm{d}x$;
    \item $\int_{0}^{\frac{\pi}{2}}\mathrm{e}^{x}\sin x\,\mathrm{d}x$;
    \item $\int_{0}^{1}\mathrm{e}^{\sqrt{x}}\,\mathrm{d}x$;
    \item $\int_{0}^{a}x^{2}\sqrt{\dfrac{a-x}{a+x}}\,\mathrm{d}x$;
    \item $\int_{0}^{\frac{\pi}{2}}\dfrac{\cos\theta}{\sin\theta+\cos\theta}\,\mathrm{d}\theta$.
  \end{enumerate}

\end{exblock}

\begin{exblock}{思考题 6.4}

  \item 设 $m,n\in\mathbb{N}$, $B(m,n)=\sum\limits_{k=0}^{n}C_{n}^{k}\dfrac{(-1)^{k}}{m+k+1}$. 证明:
  \begin{enumerate}[label=(\arabic*)]
    \item $B(m,n)=\int_{0}^{1}x^{m}(1-x)^{n}\,\mathrm{d}x$;
    \item $B(m,n)=B(n,m)$;
    \item $B(m,n)=\dfrac{m!n!}{(m+n+1)!}$.
  \end{enumerate}

  \item 对于 $\alpha\in(0,1]$, 定义 $f_{\alpha}(x)=\left[\dfrac{\alpha}{x}\right]-\alpha\left[\dfrac{1}{x}\right]$. 证明:
  \[
  \int_{0}^{1}f_{\alpha}(x)\,\mathrm{d}x=\alpha\ln\alpha.
\]
  (提示: $f_{\alpha}(x)=-\left(\dfrac{\alpha}{x}-\left[\dfrac{\alpha}{x}\right]\right)+\alpha\left(\dfrac{1}{x}-\left[\dfrac{1}{x}\right]\right)$, 对积分 $\int_{0}^{\alpha}\left(\dfrac{\alpha}{x}-\left[\dfrac{\alpha}{x}\right]\right)\mathrm{d}x$ 作换元 $x=\alpha t$.)

  \item 设 $I(m,n)=\int_{0}^{\frac{\pi}{2}}\cos^{m}x\sin^{n}x\,\mathrm{d}x$ ($m,n\in\mathbb{N}$). 证明:
  \[
  I(m,n)=\frac{m-1}{m+n}I(m-2,n)=\frac{n-1}{m+n}I(m,n-2).
\]
  并求 $I(2m,2n)$.

  \item 设 $x>0$, 证明不等式:
  \[
  \left|\int_{x}^{x+c}\sin t^{2}\,\mathrm{d}t\right|<\frac{1}{x},\quad\text{其中}\ c>0.
\]

\end{exblock}

\section{广义积分}\label{sec:6.5}

在 Riemann 积分中, 积分区间是一个有限的闭区间, 而且 Riemann 可积函数必有界函数. 但是, 对于许多来自数学本身以及其他学科的问题, 这两条限制显得过于苛刻, 因此, 有必要推广已有的积分概念. 如果说将已经详细讨论过的 Riemann 积分称作正常积分, 那么, 下面要推广的积分统称为广义积分或反常积分.

广义积分大致可以分为两大类, 一类是积分区间无界, 简称为无穷积分, 另一类是无界函数的积分, 称为瑕积分.

\begin{definition}\label{def:6.5.1}
设函数 $f$ 在无穷区间 $[a,+\infty)$ 上有定义, 且在任何有限区间 $[a,u]$ 上 Riemann 可积. 如果存在极限
\[
\lim_{u\to+\infty}\int_{a}^{u}f(x)\,\mathrm{d}x=J,
\]
则称此极限 $J$ 为函数 $f$ 在 $[a,+\infty)$ 上的无穷积分, 记作
\[
J=\lim_{u\to+\infty}\int_{a}^{u}f(x)\,\mathrm{d}x=\int_{a}^{+\infty}f(x)\,\mathrm{d}x.
\]
几何上, 它为 $x$ 轴, 直线 $x=a$ 与曲线 $y=f(x)$ 所围图形的开口面积 (图~\ref{fig:6.5.1}).

如果 $J$ 为实数, 我们就称无穷积分 $\int_{0}^{+\infty}f(x)\,\mathrm{d}x$ 收敛; 否则, 称 $\int_{0}^{+\infty}f(x)\,\mathrm{d}x$ 发散.

类似地, 可定义
\[
\int_{-\infty}^{b}f(x)\,\mathrm{d}x=\lim_{v\to-\infty}\int_{v}^{b}f(x)\,\mathrm{d}x.
\]
\[
\int_{-\infty}^{+\infty}f(x)\,\mathrm{d}x=\int_{-\infty}^{a}f(x)\,\mathrm{d}x+\int_{a}^{+\infty}f(x)\,\mathrm{d}x=\lim_{v\to-\infty}\int_{v}^{a}f(x)\,\mathrm{d}x+\lim_{u\to+\infty}\int_{a}^{u}f(x)\,\mathrm{d}x,
\]
其中 $a$ 为任一实数, $f$ 在任何有限区间 $[v,u]\subset(-\infty,+\infty)$ 上 Riemann 可积. 则当且仅当右边两个无穷积分都收敛时它才是收敛的.

易见, 无穷积分 $\int_{-\infty}^{+\infty}f(x)\,\mathrm{d}x$ 的收敛性与收敛的值都与实数 $a$ 的选取无关.
\end{definition}

\begin{figure}[htbp]
\centering
\begin{tikzpicture}[x=0.95cm,y=0.80cm,>=Stealth]
  \draw[->] (-0.3,0) -- (7.3,0) node[below] {$x$};
  \draw[->] (0,-0.3) -- (0,4.2) node[left] {$y$};
  \node[below left] at (0,0) {$O$};
  \fill[pattern=dots,pattern color=black]
    (1.3,0) -- plot[smooth] coordinates {(1.3,3.45)(1.6,2.95)(1.9,2.55)(2.3,2.15)(2.7,1.85)(3.1,1.60)(3.5,1.40)(3.9,1.24)(4.3,1.11)(4.7,1.00)(5.3,0.86)} -- (5.3,0) -- cycle;
  \draw[thick] plot[smooth] coordinates {(1.3,3.45)(1.6,2.95)(1.9,2.55)(2.3,2.15)(2.7,1.85)(3.1,1.60)(3.5,1.40)(3.9,1.24)(4.3,1.11)(4.7,1.00)(5.3,0.86)(5.9,0.72)(6.5,0.62)(7.0,0.55)};
  \node[right] at (4.6,1.62) {$y=f(x)$};
  \node[below] at (1.3,-0.06) {$a$};
  \node[below] at (5.3,-0.06) {$u$};
\end{tikzpicture}
\caption{}\label{fig:6.5.1}
\end{figure}

\begin{definition}\label{def:6.5.2}
设函数 $f$ 在区间 $(a,b]$ 上有定义, 在点 $a$ 的任一右开邻域内无界, 但在任何内闭区间 $[v,b]\subset(a,b]$ 上 Riemann 可积. 如果存在极限
\[
\lim_{v\to a^{+}}\int_{v}^{b}f(x)\,\mathrm{d}x=J,
\]
则称此极限 $J$ 为无界函数 $f$ 在 $(a,b]$ 上的瑕积分, 记作
\[
J=\lim_{v\to a^{+}}\int_{v}^{b}f(x)\,\mathrm{d}x=\int_{a}^{b}f(x)\,\mathrm{d}x.
\]
几何上, 它为 $x$ 轴, 直线 $x=b$, $x=a$ 与曲线 $y=f(x)$ 所围成图形的开口面积 (图~\ref{fig:6.5.2}).

如果 $J$ 为实数, 则称 $\int_{a}^{b}f(x)\,\mathrm{d}x$ 收敛; 否则, 称 $\int_{a}^{b}f(x)\,\mathrm{d}x$ 发散. 上述点 $a$ 称为 $f$ 的瑕点.
\end{definition}

\begin{figure}[htbp]
\centering
\begin{tikzpicture}[x=0.85cm,y=0.80cm,>=Stealth]
  \draw[->] (-0.3,0) -- (7.4,0) node[below] {$x$};
  \draw[->] (0,-0.3) -- (0,4.0) node[left] {$y$};
  \node[below left] at (0,0) {$O$};
  \fill[pattern=dots,pattern color=black]
    (2.2,0) -- (2.2,4.0) -- plot[smooth] coordinates {(2.35,3.5)(2.55,2.95)(2.85,2.45)(3.2,2.05)(3.6,1.72)(4.1,1.42)(4.6,1.20)(5.1,1.02)(5.6,0.90)(6.1,0.80)(6.5,0.74)} -- (6.5,0) -- cycle;
  \draw[thick] plot[smooth] coordinates {(2.25,3.9)(2.45,3.1)(2.75,2.55)(3.1,2.13)(3.5,1.78)(4.0,1.47)(4.5,1.24)(5.0,1.05)(5.5,0.93)(6.0,0.82)(6.5,0.74)};
  \draw (2.2,0) -- (2.2,4.0);
  \draw (3.5,0) -- (3.5,1.78);
  \draw (6.5,0) -- (6.5,0.74);
  \node[below] at (2.2,-0.06) {$a$};
  \node[below] at (3.5,-0.06) {$v$};
  \node[below] at (6.5,-0.06) {$b$};
  \node[right] at (4.5,1.95) {$y=f(x)$};
\end{tikzpicture}
\caption{}\label{fig:6.5.2}
\end{figure}

类似地, 可定义瑕点为 $b$ 的瑕积分:
\[
\int_{a}^{b}f(x)\,\mathrm{d}x=\lim_{u\to b^{-}}\int_{a}^{u}f(x)\,\mathrm{d}x,
\]
其中 $f$ 在 $[a,b)$ 上有定义, 在点 $b$ 的任何左开邻域内无界, 但在任何内闭区间 $[a,u]\subset[a,b)$ 上 Riemann 可积 (图~\ref{fig:6.5.3}).

\begin{figure}[htbp]
\centering
\begin{tikzpicture}[x=0.85cm,y=0.78cm,>=Stealth]
  \draw[->] (-0.3,0) -- (6.6,0) node[below] {$x$};
  \draw[->] (0,-0.3) -- (0,4.2) node[left] {$y$};
  \node[below left] at (0,0) {$O$};
  \fill[pattern=dots,pattern color=black]
    (1.6,0) -- (1.6,1.55) -- plot[smooth] coordinates {(1.9,1.72)(2.3,1.98)(2.7,2.30)(3.1,2.68)(3.5,3.15)(3.9,3.75)(4.2,4.25)} -- (4.2,0) -- cycle;
  \draw[thick] plot[smooth] coordinates {(1.6,1.55)(1.9,1.72)(2.3,1.98)(2.7,2.30)(3.1,2.68)(3.5,3.15)(3.9,3.75)(4.2,4.3)};
  \draw (4.2,0) -- (4.2,4.3);
  \draw (3.3,0) -- (3.3,2.83);
  \node[below] at (1.6,-0.06) {$a$};
  \node[below] at (3.3,-0.06) {$u$};
  \node[rotate=40] at (2.62,2.60) {$y=f(x)$};
\end{tikzpicture}
\caption{}\label{fig:6.5.3}
\end{figure}

如果 $f$ 的瑕点 $c\in(a,b)$, 则定义瑕积分为
\begin{align*}
\int_{a}^{b}f(x)\,\mathrm{d}x&=\int_{a}^{c}f(x)\,\mathrm{d}x+\int_{c}^{b}f(x)\,\mathrm{d}x\\
&=\lim_{u\to c^{-}}\int_{a}^{u}f(x)\,\mathrm{d}x+\lim_{v\to c^{+}}\int_{v}^{b}f(x)\,\mathrm{d}x,
\end{align*}
其中 $f$ 在 $[a,c)\cup(c,b]$ 上有定义, 在点 $c$ 的任一开邻域内无界, 但在任何 $[a,u]\subset[a,c)$ 与 $[v,b]\subset(c,b]$ 上都 Riemann 可积. 当且仅当上面右边两个瑕积分都收敛时, 左边的瑕积分才是收敛的.

又若 $a$, $b$ 两点都是 $f$ 的瑕点, 而 $f$ 在任何内闭区间 $[u,v]\subset(a,b)$ 上 Riemann 可积, 则瑕积分定义为
\begin{align*}
\int_{a}^{b}f(x)\,\mathrm{d}x&=\int_{a}^{c}f(x)\,\mathrm{d}x+\int_{c}^{b}f(x)\,\mathrm{d}x\\
&=\lim_{v\to a^{+}}\int_{v}^{c}f(x)\,\mathrm{d}x+\lim_{u\to b^{-}}\int_{c}^{u}f(x)\,\mathrm{d}x,
\end{align*}
其中 $c$ 为 $(a,b)$ 内任一实数. 同样, 当且仅当上面右边两个瑕积分都收敛时, 左边的瑕积分才是收敛的.

显然, 该瑕积分的收敛性和积分值与 $c$ 的选取无关.

如果 $a$, $b$ 中有一个 $-\infty$ 或 $+\infty$, 且 $(a,b)$ 中还有有限个瑕点, 则可将 $(a,b)$ 分成有限部分, 使在每一部分中只含瑕点、$-\infty$, $+\infty$ 之一, 则当且仅当每一部分的积分收敛时, $\int_{a}^{b}f(x)\,\mathrm{d}x$ 才是收敛的.

关于广义积分, 同样有 Newton-Leibniz 公式, 变量代换公式及分部积分公式.

\begin{theorem}\label{thm:6.5.1}
设函数 $f$ 在 $[a,+\infty)$ 上有定义, 在任何内闭区间 $[a,u]$ 上 Riemann 可积, $f$ 有原函数 $F$, 且 $F(+\infty)=\lim\limits_{u\to+\infty}F(u)$ 存在, 则
\[
\int_{a}^{+\infty}f(x)\,\mathrm{d}x=F(+\infty)-F(a).
\]
\end{theorem}

\begin{proof}
\[
\int_{a}^{+\infty}f(x)\,\mathrm{d}x=\lim_{u\to+\infty}\int_{a}^{u}f(x)\,\mathrm{d}x=\lim_{u\to+\infty}[F(u)-F(a)]=F(+\infty)-F(a).
\]
\end{proof}

其他公式读者可自行给出.

\begin{example}\label{ex:6.5.1}
讨论无穷积分
\[
\int_{1}^{+\infty}\frac{\mathrm{d}x}{x^{p}}
\]
的敛散性.
\end{example}

\begin{solution}
\[
\int_{1}^{+\infty}\frac{\mathrm{d}x}{x^{p}}=\lim_{u\to+\infty}\int_{1}^{u}\frac{\mathrm{d}x}{x^{p}}=\begin{cases}
\lim\limits_{u\to+\infty}\dfrac{1}{1-p}\left(u^{1-p}-1\right), & p\neq1,\\[3mm]
\lim\limits_{u\to+\infty}\ln u, & p=1
\end{cases}=\begin{cases}
\dfrac{1}{p-1}, & p>1,\ \text{收敛},\\[3mm]
+\infty, & p\leqslant1,\ \text{发散}.
\end{cases}
\]
上述极限过程也可简单表达为
\[
\int_{1}^{+\infty}\frac{\mathrm{d}x}{x^{p}}=\begin{cases}
\dfrac{1}{1-p}\left(u^{1-p}-1\right)\bigg|_{1}^{+\infty}, & p\neq1,\\[3mm]
\ln u\bigg|_{1}^{+\infty}, & p=1
\end{cases}=\begin{cases}
\dfrac{1}{p-1}, & p>1,\\[3mm]
+\infty, & p\leqslant1.
\end{cases}
\]
\end{solution}

\begin{example}\label{ex:6.5.2}
讨论下列无穷积分的敛散性.
\begin{enumerate}[label=(\arabic*)]
\item $\int_{2}^{+\infty}\dfrac{\mathrm{d}x}{x\left(\ln x\right)^{p}}$;
\item $\int_{-\infty}^{+\infty}\dfrac{\mathrm{d}x}{1+x^{2}}$.
\end{enumerate}
\end{example}

\begin{solution}
由于无穷积分是通过变限积分的极限来定义的, 因此有关 Riemann 积分的换元与分部积分法一般都可用到广义 (反常) 积分 (无穷积分与瑕积分) 中来 (请思考理由).

(1)
\[
\int_{2}^{+\infty}\frac{\mathrm{d}x}{x\left(\ln x\right)^{p}}=\int_{2}^{+\infty}\frac{\mathrm{d}\left(\ln x\right)}{\left(\ln x\right)^{p}}\xlongequal{t=\ln x}\int_{\ln2}^{+\infty}\frac{\mathrm{d}t}{t^{p}},
\]
从例~\ref{ex:6.5.1} 知, 当 $p>1$ 时收敛; 当 $p\leqslant1$ 时发散.

(2)
\begin{align*}
\int_{-\infty}^{+\infty}\frac{\mathrm{d}x}{1+x^{2}}&=\int_{-\infty}^{0}\frac{\mathrm{d}x}{1+x^{2}}+\int_{0}^{+\infty}\frac{\mathrm{d}x}{1+x^{2}}\\
&=\lim_{v\to-\infty}\int_{v}^{0}\frac{\mathrm{d}x}{1+x^{2}}+\lim_{u\to+\infty}\int_{0}^{u}\frac{\mathrm{d}x}{1+x^{2}}\\
&=\lim_{v\to-\infty}(\arctan0-\arctan v)+\lim_{u\to+\infty}(\arctan u-\arctan0)\\
&=-\left(-\frac{\pi}{2}\right)+\frac{\pi}{2}=\pi.
\end{align*}
上述描述很烦琐, 简化如下:
\[
\int_{-\infty}^{+\infty}\frac{\mathrm{d}x}{1+x^{2}}=\arctan x\bigg|_{-\infty}^{+\infty}=\frac{\pi}{2}-\left(-\frac{\pi}{2}\right)=\pi.
\]
\end{solution}

\begin{example}\label{ex:6.5.3}
求极限 $\lim\limits_{n\to+\infty}\sin\dfrac{\pi}{n}\sum\limits_{k=1}^{n}\dfrac{1}{3+\sin\dfrac{k\pi}{n}}$.
\end{example}

\begin{solution}
\begin{align*}
\lim_{n\to+\infty}\sin\frac{\pi}{n}\sum_{k=1}^{n}\frac{1}{3+\sin\frac{k\pi}{n}}&=\lim_{n\to+\infty}\frac{\sin\frac{\pi}{n}}{\frac{\pi}{n}}\frac{\pi}{n}\sum_{k=1}^{n}\frac{1}{3+\sin\frac{k\pi}{n}}\\
&=\int_{0}^{\pi}\frac{\mathrm{d}x}{3+\sin x}\xlongequal{u=\tan\frac{x}{2}}\int_{0}^{+\infty}\frac{\frac{2\mathrm{d}u}{1+u^{2}}}{3+\frac{2u}{1+u^{2}}}=2\int_{0}^{+\infty}\frac{\mathrm{d}u}{3+2u+3u^{2}}\\
&=\frac{2}{3}\int_{0}^{+\infty}\frac{\mathrm{d}\left(u+\frac{1}{3}\right)}{\left(u+\frac{1}{3}\right)^{2}+\frac{8}{9}}=\frac{2}{3}\cdot\frac{3}{2\sqrt{2}}\arctan\frac{u+\frac{1}{3}}{\frac{2\sqrt{2}}{3}}\bigg|_{0}^{+\infty}\\
&=\frac{1}{\sqrt{2}}\left(\frac{\pi}{2}-\arctan\frac{1}{2\sqrt{2}}\right).
\end{align*}
\end{solution}

\begin{example}\label{ex:6.5.4}
证明: $\int_{0}^{+\infty}\dfrac{\mathrm{d}x}{1+x^{4}}=\dfrac{\pi}{2\sqrt{2}}$.
\end{example}

\begin{proof}[证法 1]
由例~\ref{ex:3.3.7}, 有
\begin{align*}
\int_{0}^{+\infty}\frac{\mathrm{d}x}{1+x^{4}}&=\frac{1}{4\sqrt{2}}\ln\frac{x^{2}+\sqrt{2}x+1}{x^{2}-\sqrt{2}x+1}\bigg|_{0}^{+\infty}+\frac{1}{2\sqrt{2}}\big[\arctan(\sqrt{2}x+1)\\
&\quad+\arctan(\sqrt{2}x-1)\big]\bigg|_{0}^{+\infty}\\
&=\frac{1}{2\sqrt{2}}\left(\frac{\pi}{2}+\frac{\pi}{2}-\frac{\pi}{4}+\frac{\pi}{4}\right)=\frac{\pi}{2\sqrt{2}}.
\end{align*}
\end{proof}

\begin{proof}[证法 2]
\begin{align*}
\int_{0}^{+\infty}\frac{\mathrm{d}x}{1+x^{4}}&\xlongequal[x=\left(\frac{1}{t}-1\right)^{\frac{1}{4}}]{t=\frac{1}{1+x^{4}}}\int_{1}^{0}t\cdot\frac{1}{4}\left(\frac{1}{t}-1\right)^{-\frac{3}{4}}\cdot\frac{-1}{t^{2}}\,\mathrm{d}t\\
&=\frac{1}{4}\int_{0}^{1}t^{-\frac{1}{4}}(1-t)^{-\frac{3}{4}}\,\mathrm{d}t=\frac{1}{4}\int_{0}^{1}t^{\frac{3}{4}-1}(1-t)^{\frac{1}{4}-1}\,\mathrm{d}t=\frac{1}{4}\mathrm{B}\left(\frac{3}{4},\frac{1}{4}\right)\\
&=\frac{1}{4}\frac{\Gamma\left(\frac{3}{4}\right)\Gamma\left(\frac{1}{4}\right)}{\Gamma\left(\frac{3}{4}+\frac{1}{4}\right)}=\frac{1}{4}\Gamma\left(\frac{3}{4}\right)\Gamma\left(\frac{1}{4}\right)\xlongequal{\text{余元公式}}\frac{\pi}{4\sin\frac{\pi}{4}}=\frac{\pi}{2\sqrt{2}}.
\end{align*}
\end{proof}

注: 这里为了给出一种新的计算方法, 我们不得不借用第 3 册中的 B 函数与 $\Gamma$ 函数的公式:
\[
\mathrm{B}(p,q)=\int_{0}^{1}t^{p-1}(1-t)^{q-1}\,\mathrm{d}t,
\]
\[
\Gamma(s)=\int_{0}^{+\infty}t^{s-1}\mathrm{e}^{-t}\,\mathrm{d}t,\quad\Gamma(s+1)=s!,\quad\Gamma(1)=1,
\]
\[
\mathrm{B}(p,q)=\frac{\Gamma(p)\Gamma(q)}{\Gamma(p+q)},\quad\text{余元公式}\ \Gamma(s)\Gamma(1-s)=\frac{\pi}{\sin\pi s}.
\]

\begin{example}\label{ex:6.5.5}
计算 $I=\int_{0}^{+\infty}\sin ax\,\mathrm{e}^{-bx}\,\mathrm{d}x$, $J=\int_{0}^{+\infty}\cos ax\,\mathrm{e}^{-bx}\,\mathrm{d}x$, $a>0$, $b>0$.
\end{example}

\begin{solution}
\[
I=\int_{0}^{+\infty}\sin ax\,\mathrm{e}^{-bx}\,\mathrm{d}x=\int_{0}^{+\infty}\sin ax\,\mathrm{d}\left(-\frac{\mathrm{e}^{-bx}}{b}\right)
\]
\begin{align*}
&=\frac{-1}{b}\sin ax\,\mathrm{e}^{-bx}\bigg|_{0}^{+\infty}+\frac{a}{b}\int_{0}^{+\infty}\cos ax\,\mathrm{e}^{-bx}\,\mathrm{d}x\\
&=\frac{a}{b}\int_{0}^{+\infty}\cos ax\,\mathrm{d}\left(-\frac{\mathrm{e}^{-bx}}{b}\right)\\
&=\frac{a}{b^{2}}\left(-\cos ax\,\mathrm{e}^{-bx}\bigg|_{0}^{+\infty}-a\int_{0}^{+\infty}\sin ax\,\mathrm{e}^{-bx}\,\mathrm{d}x\right),\\
&=\frac{a}{b^{2}}\left(1-a\int_{0}^{+\infty}\sin ax\,\mathrm{e}^{-bx}\,\mathrm{d}x\right)=\frac{a}{b^{2}}-\frac{a^{2}}{b^{2}}I,
\end{align*}
\[
I=\frac{\frac{a}{b^{2}}}{1+\frac{a^{2}}{b^{2}}}=\frac{a}{a^{2}+b^{2}},
\]
\[
J=\frac{b}{a}I=\frac{b}{a^{2}+b^{2}}.
\]
\end{solution}

\begin{example}\label{ex:6.5.6}
求 Poisson 积分 $(0<r<1)$
\[
\int_{-\pi}^{\pi}\frac{1-r^{2}}{1-2r\cos x+r^{2}}\,\mathrm{d}x.
\]
\end{example}

\begin{solution}
\begin{align*}
&\int_{-\pi}^{\pi}\frac{1-r^{2}}{1-2r\cos x+r^{2}}\,\mathrm{d}x\\
&\xlongequal{t=\tan\frac{x}{2}}\int_{-\infty}^{+\infty}\frac{1-r^{2}}{1-2r\frac{1-t^{2}}{1+t^{2}}+r^{2}}\frac{2\,\mathrm{d}t}{1+t^{2}}\\
&=\int_{-\infty}^{+\infty}\frac{2(1-r^{2})}{(1-r)^{2}+(1+r)^{2}t^{2}}\,\mathrm{d}t=2\arctan\frac{1+r}{1-r}t\bigg|_{-\infty}^{+\infty}\\
&=2\left[\frac{\pi}{2}-\left(-\frac{\pi}{2}\right)\right]=2\pi.
\end{align*}
\end{solution}

\begin{example}\label{ex:6.5.7}
证明: Euler-Poisson 积分
\[
\int_{0}^{+\infty}\mathrm{e}^{-x^{2}}\,\mathrm{d}x=\frac{\sqrt{\pi}}{2}
\]
(此积分在概率统计中非常重要).
\end{example}

\begin{proof}
令 $\varphi(x)=\mathrm{e}^{-x^{2}}-(1-x^{2})$, 则 $\varphi'(x)=2x(1-\mathrm{e}^{-x^{2}})>0$ $(x>0)$, $\varphi$ 严格增, 故 $\mathrm{e}^{-x^{2}}-(1-x^{2})=\varphi(x)>\varphi(0)=0$, $\mathrm{e}^{-x^{2}}>1-x^{2}$ $(x>0)$. 于是, $\mathrm{e}^{-nx^{2}}>(1-x^{2})^{n}$, $x\in(0,1)$. 此外, 由
\[
\mathrm{e}^{x^{2}}>1+x^{2},\quad\mathrm{e}^{-x^{2}}<\frac{1}{1+x^{2}},\quad\mathrm{e}^{-nx^{2}}<\frac{1}{\left(1+x^{2}\right)^{n}},\quad x\in(0,+\infty),
\]
得到
\[
\frac{(2n)!!}{(2n+1)!!}=\int_{0}^{\frac{\pi}{2}}\cos^{2n+1}t\,\mathrm{d}t\xlongequal{x=\sin t}\int_{0}^{1}(1-x^{2})^{n}\,\mathrm{d}x
\]
\begin{align*}
&<\int_{0}^{1}\mathrm{e}^{-nx^{2}}\,\mathrm{d}x<\int_{0}^{+\infty}\mathrm{e}^{-nx^{2}}\,\mathrm{d}x<\int_{0}^{+\infty}\frac{\mathrm{d}x}{(1+x^{2})^{n}}\\
&\xlongequal{x=\tan t}\int_{0}^{\frac{\pi}{2}}\cos^{2n-2}t\,\mathrm{d}t=\frac{(2n-3)!!}{(2n-2)!!}\frac{\pi}{2},
\end{align*}
再由 $\int_{0}^{+\infty}\mathrm{e}^{-nx^{2}}\,\mathrm{d}x\xlongequal{x=\frac{t}{\sqrt{n}}}\int_{0}^{+\infty}\mathrm{e}^{-t^{2}}\frac{\mathrm{d}t}{\sqrt{n}}$ 得到
\[
n\left[\frac{(2n)!!}{(2n+1)!!}\right]^{2}<\left(\int_{0}^{+\infty}\mathrm{e}^{-x^{2}}\,\mathrm{d}x\right)^{2}<n\left[\frac{(2n-3)!!}{(2n-2)!!}\right]^{2}\cdot\frac{\pi^{2}}{4}.
\]
根据 Wallis 公式及例~\ref{ex:6.4.5}(2) 推得
\begin{align*}
\lim_{n\to+\infty}n\left[\frac{(2n)!!}{(2n+1)!!}\right]^{2}&=\lim_{n\to+\infty}\frac{n}{2n+1}\left[\frac{(2n)!!}{(2n-1)!!}\right]^{2}\cdot\frac{1}{2n+1}\\
&=\frac{1}{2}\cdot\frac{\pi}{2}=\frac{\pi}{4},
\end{align*}
\begin{align*}
\lim_{n\to+\infty}n\left[\frac{(2n-3)!!}{(2n-2)!!}\right]^{2}\cdot\frac{\pi^{2}}{4}&=\lim_{n\to+\infty}\frac{n\cdot(2n)^{2}}{(2n+1)(2n-1)^{2}}\left[\frac{(2n-1)!!}{(2n)!!}\right]^{2}(2n+1)\cdot\frac{\pi^{2}}{4}\\
&=\frac{1}{2}\cdot\frac{2}{\pi}\cdot\frac{\pi^{2}}{4}=\frac{\pi}{4}.
\end{align*}
由夹逼定理知
\[
\left(\int_{0}^{+\infty}\mathrm{e}^{-x^{2}}\,\mathrm{d}x\right)^{2}=\frac{\pi}{4},\quad\int_{0}^{+\infty}\mathrm{e}^{-x^{2}}\,\mathrm{d}x=\frac{\sqrt{\pi}}{2}.
\]
以后应用二重积分或 $\Gamma$ 函数的余元公式也可推得 $\int_{0}^{+\infty}\mathrm{e}^{-x^{2}}\,\mathrm{d}x=\frac{\sqrt{\pi}}{2}$.
\end{proof}

\begin{example}\label{ex:6.5.8}
计算瑕积分:
\begin{enumerate}[label=(\arabic*)]
\item $\int_{0}^{1}\dfrac{\mathrm{d}x}{\sqrt{1-x^{2}}}$;
\item $\int_{0}^{1}\dfrac{x^{n}\,\mathrm{d}x}{\sqrt{1-x^{2}}}$, $n\in\mathbb{N}$.
\end{enumerate}
\end{example}

\begin{solution}
(1) 被积函数 $f(x)=\dfrac{1}{\sqrt{1-x^{2}}}$ 在 $[0,1)$ 上连续, 从而在任何 $[0,u]\subset[0,1)$ 上 Riemann 可积, $x=1$ 为其瑕点. 因此
\begin{align*}
\int_{0}^{1}\frac{\mathrm{d}x}{\sqrt{1-x^{2}}}&=\lim_{u\to 1^{-}}\int_{0}^{u}\frac{\mathrm{d}x}{\sqrt{1-x^{2}}}\\
&=\lim_{u\to 1^{-}}(\arcsin u-\arcsin0)=\arcsin1=\frac{\pi}{2}.
\end{align*}
简单表达为
\[
\int_{0}^{1}\frac{\mathrm{d}x}{\sqrt{1-x^{2}}}=\arcsin x\bigg|_{0}^{1}=\frac{\pi}{2}-0=\frac{\pi}{2}.
\]
(2) $x=1$ 为瑕点. 作变换 $x=\sin t$, 则有
\begin{align*}
\int_{0}^{1}\frac{x^{n}\,\mathrm{d}x}{\sqrt{1-x^{2}}}&=\int_{0}^{\frac{\pi}{2}}\frac{\sin^{n}t\cos t}{\sqrt{1-\sin^{2}t}}\,\mathrm{d}t=\int_{0}^{\frac{\pi}{2}}\sin^{n}t\,\mathrm{d}t\\
&=\begin{cases}
\dfrac{(n-1)!!}{n!!}\dfrac{\pi}{2}, & n\ \text{为偶数},\\[
3mm]
\dfrac{(n-1)!!}{n!!}, & n\ \text{为奇数}.
\end{cases}
\end{align*}
\end{solution}
\begin{example}\label{ex:6.5.9}
讨论瑕积分 $\int_{0}^{1}\dfrac{\mathrm{d}x}{x^{p}}$ 的敛散性.
\end{example}
\begin{solution}
当 $p\leqslant0$ 时, 这是正常积分, 视作是收敛的.
当 $p>0$ 时, 被积函数 $\dfrac{1}{x^{p}}$ 在 $(0,1]$ 上连续, $x=0$ 为其瑕点, 于是,
\begin{align*}
\int_{0}^{1}\frac{\mathrm{d}x}{x^{p}}&=\begin{cases}
\dfrac{1}{1-p}x^{-p+1}\bigg|_{0}^{1}=\dfrac{1}{1-p}, & 0<p<1,\\[3mm]
\dfrac{1}{1-p}x^{-p+1}\bigg|_{0}^{1}=+\infty, & p>1,\\[3mm]
\ln x\bigg|_{0}^{1}=+\infty, & p=1
\end{cases}\\
&=\begin{cases}
\dfrac{1}{1-p}, & 0<p<1,\ \text{收敛},\\[3mm]
+\infty, & p\geqslant1,\ \text{发散}.
\end{cases}
\end{align*}
类似地, $\int_{a}^{b}\dfrac{\mathrm{d}x}{(x-a)^{p}}$ 与 $\int_{a}^{b}\dfrac{\mathrm{d}x}{(b-x)^{p}}$ 当 $p<1$ 时收敛; 当 $p\geqslant1$ 时发散.
考察广义积分
\[
\int_{0}^{+\infty}\frac{\mathrm{d}x}{x^{p}}=\int_{0}^{1}\frac{\mathrm{d}x}{x^{p}}+\int_{1}^{+\infty}\frac{\mathrm{d}x}{x^{p}},
\]
当且仅当 $p<1$ 时, $\int_{0}^{1}\dfrac{\mathrm{d}x}{x^{p}}$ 收敛; 当且仅当 $p>1$ 时, $\int_{1}^{+\infty}\dfrac{\mathrm{d}x}{x^{p}}$ 收敛. 因此, 右边两个积分对任何 $p$ 不能同时收敛, 故 $\int_{0}^{+\infty}\dfrac{\mathrm{d}x}{x^{p}}$ 对任何实数 $p$ 都是发散的.
\end{solution}

如果只停留在数学概念上, 那么肯定学不深, 钻不透, 也走不远. 因此, 我们有必要来讨论广义积分的一些性质与各种收敛判别法则.

\begin{theorem}[无穷积分收敛的 Cauchy 准则]\label{thm:6.5.2}
无穷积分 $\int_{a}^{+\infty}f(x)\,\mathrm{d}x$ 收敛 $\Longleftrightarrow\forall\varepsilon>0$, $\exists\Delta>\max\{a,0\}$, 当 $u_{1},u_{2}>\Delta$ 时, 便有
\[
\left|\int_{a}^{u_{1}}f(x)\,\mathrm{d}x-\int_{a}^{u_{2}}f(x)\,\mathrm{d}x\right|=\left|\int_{u_{1}}^{u_{2}}f(x)\,\mathrm{d}x\right|<\varepsilon.
\]
\end{theorem}

\begin{proof}
令 $F(u)=\int_{a}^{u}f(x)\,\mathrm{d}x$, 则
\begin{align*}
\int_{a}^{+\infty}f(x)\,\mathrm{d}x\ \text{收敛} &\Longleftrightarrow\lim_{u\to+\infty}\int_{a}^{u}f(x)\,\mathrm{d}x=\lim_{u\to+\infty}F(u)\ \text{存在且有限}\\
&\xLeftrightarrow[\text{定理~\ref{thm:2.1.4}}]{\text{函数收敛的 Cauchy 准则}}\forall\varepsilon>0,\exists\Delta>\max\{a,0\},\ \text{当}\ u_{1},u_{2}>\Delta\ \text{时, 便有}\\
&\left|\int_{u_{1}}^{u_{2}}f(x)\,\mathrm{d}x\right|=\left|F(u_{1})-F(u_{2})\right|<\varepsilon.
\end{align*}
\end{proof}

\begin{lemma}\label{lem:6.5.1}
设 $\int_{a}^{+\infty}f(x)\,\mathrm{d}x$ 与 $\int_{a}^{+\infty}g(x)\,\mathrm{d}x$ 都收敛, $\lambda$, $\mu$ 为任意实常数, 则 $\int_{a}^{+\infty}[\lambda f(x)+\mu g(x)]\,\mathrm{d}x$ 也收敛, 且
\[
\int_{a}^{+\infty}\left[\lambda f(x)+\mu g(x)\right]\mathrm{d}x=\lambda\int_{a}^{+\infty}f(x)\,\mathrm{d}x+\mu\int_{a}^{+\infty}g(x)\,\mathrm{d}x.
\]
\end{lemma}

\begin{proof}
\begin{align*}
\int_{a}^{+\infty}\left[\lambda f(x)+\mu g(x)\right]\mathrm{d}x&=\lim_{u\to+\infty}\int_{a}^{u}\left[\lambda f(x)+\mu g(x)\right]\mathrm{d}x\\
&=\lim_{u\to+\infty}\left[\lambda\int_{a}^{u}f(x)\,\mathrm{d}x+\mu\int_{a}^{u}g(x)\,\mathrm{d}x\right]\\
&=\lambda\int_{a}^{+\infty}f(x)\,\mathrm{d}x+\mu\int_{a}^{+\infty}g(x)\,\mathrm{d}x.
\end{align*}
\end{proof}

\begin{lemma}\label{lem:6.5.2}
设 $f$ 在任何有限区间 $[v,u]$ 上 Riemann 可积, $a<b$. 则 $\int_{a}^{+\infty}f(x)\,\mathrm{d}x$ 与 $\int_{b}^{+\infty}f(x)\,\mathrm{d}x$ 同敛散. 如果收敛, 则
\[
\int_{a}^{+\infty}f(x)\,\mathrm{d}x=\int_{a}^{b}f(x)\,\mathrm{d}x+\int_{b}^{+\infty}f(x)\,\mathrm{d}x,
\]
其中 $\int_{a}^{b}f(x)\,\mathrm{d}x$ 为 Riemann 积分.
\end{lemma}

\begin{proof}
设 $a<b<u$, 则由定理~\ref{thm:6.2.4}, 有
\[
\int_{a}^{u}f(x)\,\mathrm{d}x=\int_{a}^{b}f(x)\,\mathrm{d}x+\int_{b}^{u}f(x)\,\mathrm{d}x.
\]
因此, $\lim\limits_{u\to+\infty}\int_{a}^{u}f(x)\,\mathrm{d}x$ 与 $\lim\limits_{u\to+\infty}\int_{b}^{u}f(x)\,\mathrm{d}x$ 同时存在或同时不存在, 故 $\int_{a}^{+\infty}f(x)\,\mathrm{d}x$ 与 $\int_{b}^{+\infty}f(x)\,\mathrm{d}x$ 同敛散. 当它们收敛时, 有
\begin{align*}
\int_{a}^{+\infty}f(x)\,\mathrm{d}x&=\lim_{u\to+\infty}\int_{a}^{u}f(x)\,\mathrm{d}x=\lim_{u\to+\infty}\left[\int_{a}^{b}f(x)\,\mathrm{d}x+\int_{b}^{u}f(x)\,\mathrm{d}x\right]\\
&=\int_{a}^{b}f(x)\,\mathrm{d}x+\int_{b}^{+\infty}f(x)\,\mathrm{d}x.
\end{align*}
\end{proof}

\begin{lemma}\label{lem:6.5.3}
设 $f$ 在任何有限区间 $[a,u]$ 上 Riemann 可积, 且 $\int_{a}^{+\infty}|f(x)|\,\mathrm{d}x$ 收敛, 则 $\int_{a}^{+\infty}f(x)\,\mathrm{d}x$ 也收敛, 且有
\[
\left|\int_{a}^{+\infty}f(x)\,\mathrm{d}x\right|\leqslant\int_{a}^{+\infty}\left|f(x)\right|\mathrm{d}x.
\]
\end{lemma}

\begin{proof}
由 $\int_{a}^{+\infty}|f(x)|\,\mathrm{d}x$ 收敛, 根据无穷积分的收敛准则 (定理~\ref{thm:6.5.2}) 的必要性, $\forall\varepsilon>0$, $\exists\Delta>\max\{a,0\}$, 当 $u_{2}\geqslant u_{1}>\Delta$ 时, 总有
\[
\left|\int_{u_{1}}^{u_{2}}|f(x)|\,\mathrm{d}x\right|=\int_{u_{1}}^{u_{2}}|f(x)|\,\mathrm{d}x<\varepsilon.
\]
于是,
\[
\left|\int_{u_{1}}^{u_{2}}f(x)\,\mathrm{d}x\right|\leqslant\int_{u_{1}}^{u_{2}}|f(x)|\,\mathrm{d}x<\varepsilon.
\]
再由无穷积分的收敛准则 (定理~\ref{thm:6.5.2}) 的充分性推得 $\int_{a}^{+\infty}f(x)\,\mathrm{d}x$ 收敛.

又因为
\[
\left|\int_{a}^{u}f(x)\,\mathrm{d}x\right|\leqslant\int_{a}^{u}\left|f(x)\right|\mathrm{d}x\quad(u>a),
\]
令 $u\to+\infty$, 得到
\[
\left|\int_{a}^{+\infty}f(x)\,\mathrm{d}x\right|\leqslant\int_{a}^{+\infty}\left|f(x)\right|\mathrm{d}x.
\]
\end{proof}

\begin{definition}\label{def:6.5.3}
设 $f$ 在任何有限区间 $[a,u]$ 上 Riemann 可积, 如果 $\int_{a}^{+\infty}|f(x)|\,\mathrm{d}x$ 收敛, 则称 $\int_{a}^{+\infty}f(x)\,\mathrm{d}x$ 绝对收敛.
\end{definition}

引理~\ref{lem:6.5.3} 指出, 若 $\int_{a}^{+\infty}f(x)\,\mathrm{d}x$ 绝对收敛, 则 $\int_{a}^{+\infty}f(x)\,\mathrm{d}x$ 也收敛, 但反之不成立.

反例: 设
\[
f(x)=\begin{cases}
\dfrac{x}{2^{n-2}}-\dfrac{n}{2^{n-3}}, & 2n\leqslant x<2n+\dfrac{1}{2},\\[3mm]
-\dfrac{x}{2^{n-2}}+\dfrac{2n+1}{2^{n-2}}, & 2n+\dfrac{1}{2}\leqslant x<2n+\dfrac{3}{2},\\[3mm]
\dfrac{x}{2^{n-2}}-\dfrac{n+1}{2^{n-3}}, & 2n+\dfrac{3}{2}\leqslant x<2n+2,
\end{cases}
\]

\begin{figure}[htbp]
\centering
\begin{tikzpicture}[x=0.78cm,y=0.72cm,>=Stealth]
  \draw[->] (-0.4,0) -- (8.6,0) node[below] {$x$};
  \draw[->] (0,-2.7) -- (0,2.8) node[left] {$y$};
  \node[below left] at (0,0) {$O$};
  \node[left] at (0,2) {2};
  \node[left] at (0,1) {1};
  \node[left] at (0,-1) {-1};
  \node[left] at (0,-2) {-2};
  \draw[thick] (0,0) -- (1,2) -- (2,-2) -- (3,1) -- (3.4,0);
  \draw[thick] (3.75,0.18) -- (3.95,-0.18);
  \draw[thick] (4.15,0.18) -- (4.35,-0.18);
  \draw[thick] (4.6,0) -- (5,0.5) -- (6,-0.5) -- (7,0) -- (7.4,0.55);
  \node[below] at (1,-0.15) {1};
  \node[below] at (2,-0.15) {2};
  \node[below] at (3,-0.15) {3};
  \node[below] at (5,-0.15) {$2n$};
  \node[below=10pt] at (6,-0.15) {$2n+1$};
  \node[below] at (7,-0.15) {$2n+2$};
\end{tikzpicture}
\caption{}\label{fig:6.5.4}
\end{figure}

由图~\ref{fig:6.5.4} 立刻知道, $\int_{0}^{+\infty}f(x)\,\mathrm{d}x=0$, $\int_{0}^{+\infty}|f(x)|\,\mathrm{d}x=2\sum\limits_{n=1}^{\infty}\dfrac{1}{n}=+\infty$, 即 $\int_{0}^{+\infty}f(x)\,\mathrm{d}x$ 收敛, 但 $\int_{0}^{+\infty}|f(x)|\,\mathrm{d}x$ 发散. 我们称这种收敛但非绝对收敛的无穷积分为条件收敛.

\begin{theorem}[无穷积分的比较判别法]\label{thm:6.5.3}
设定义在 $[a,+\infty)$ 上的两个函数 $f$ 与 $g$ 都在任何有限区间 $[a,u]$ 上 Riemann 可积, 且满足
\[
|f(x)|\leqslant g(x),\quad x\in[a,+\infty).
\]
则
\begin{enumerate}[label=(\arabic*)]
\item 当 $\int_{a}^{+\infty}g(x)\,\mathrm{d}x$ 收敛时, $\int_{a}^{+\infty}|f(x)|\,\mathrm{d}x$ 必收敛, 当然 $\int_{a}^{+\infty}f(x)\,\mathrm{d}x$ 也收敛;
\item 当 $\int_{a}^{+\infty}|f(x)|\,\mathrm{d}x$ 发散时, $\int_{a}^{+\infty}g(x)\,\mathrm{d}x$ 必发散.
\end{enumerate}
(简单说即是, 大的收敛, 小的也收敛; 小的发散, 大的也发散.)
\end{theorem}

\begin{proof}
因为 $0\leqslant|f(x)|\leqslant g(x)$, 所以 $\int_{a}^{u}|f(x)|\,\mathrm{d}x$ 与 $\int_{a}^{u}g(x)\,\mathrm{d}x$ 都为 $u$ 的增函数.

(1) 如果 $\int_{a}^{+\infty}g(x)\,\mathrm{d}x$ 收敛, 则 $\forall u\in[a,+\infty)$, 有 $\int_{a}^{u}g(x)\,\mathrm{d}x\leqslant M$(常数), 从而
\[
\int_{a}^{u}|f(x)|\,\mathrm{d}x\leqslant\int_{a}^{u}g(x)\,\mathrm{d}x\leqslant M,\quad\forall u\in[a,+\infty).
\]
因此, $\int_{a}^{+\infty}|f(x)|\,\mathrm{d}x=\lim\limits_{u\to+\infty}\int_{a}^{u}|f(x)|\,\mathrm{d}x$ 存在且有限, 这就证明了 $\int_{a}^{+\infty}|f(x)|\,\mathrm{d}x$ 收敛.

根据引理~\ref{lem:6.5.3}, $\int_{a}^{+\infty}f(x)\,\mathrm{d}x$ 也收敛.

(2) 如果 $\int_{a}^{+\infty}|f(x)|\,\mathrm{d}x$ 发散, 则 $\lim\limits_{u\to+\infty}\int_{a}^{u}|f(x)|\,\mathrm{d}x=+\infty$. 再由
\[
\int_{a}^{u}|f(x)|\,\mathrm{d}x\leqslant\int_{a}^{u}g(x)\,\mathrm{d}x
\]
得到
\[
+\infty=\lim_{u\to+\infty}\int_{a}^{u}|f(x)|\,\mathrm{d}x\leqslant\lim_{u\to+\infty}\int_{a}^{u}g(x)\,\mathrm{d}x=\int_{a}^{+\infty}g(x)\,\mathrm{d}x,
\]
即
\[
\int_{a}^{+\infty}g(x)\,\mathrm{d}x=+\infty,
\]
是发散的.
\end{proof}

由无穷积分的比较判别法可派生出用起来更方便的比较判别法的极限形式.

\begin{theoremp}[无穷积分的比较判别法的极限形式]\label{thm:6.5.3p}
设 $f$, $g$ 在任何有穷区间 $[a,u]$ 上 Riemann 可积, $g(x)>0$, 且
\[
\lim_{x\to+\infty}\frac{|f(x)|}{g(x)}=c,
\]
则有
\begin{enumerate}[label=(\arabic*)]
\item 当 $0<c<+\infty$ 时, $\int_{a}^{+\infty}|f(x)|\,\mathrm{d}x$ 与 $\int_{a}^{+\infty}g(x)\,\mathrm{d}x$ 同敛散;
\item 当 $c=0$ 时, $\int_{a}^{+\infty}g(x)\,\mathrm{d}x$ 收敛蕴涵着 $\int_{a}^{+\infty}|f(x)|\,\mathrm{d}x$ 也收敛;
\item 当 $c=+\infty$ 时, $\int_{a}^{+\infty}g(x)\,\mathrm{d}x$ 发散蕴涵着 $\int_{a}^{+\infty}|f(x)|\,\mathrm{d}x$ 也发散.
\end{enumerate}
\end{theoremp}

\begin{proof}
(1) 因为 $\lim\limits_{x\to+\infty}\dfrac{|f(x)|}{g(x)}=c\in(0,+\infty)$, 所以 $\exists\Delta>\max\{a,0\}$, 当 $x>\Delta$ 时, 有
\[
\frac{3c}{2}=c+\frac{c}{2}>\frac{|f(x)|}{g(x)}>c-\frac{c}{2}=\frac{c}{2}>0,
\]
\[
\frac{3c}{2}g(x)>|f(x)|>\frac{c}{2}g(x).
\]
再由定理~\ref{thm:6.5.3} 知, $\int_{a}^{+\infty}|f(x)|\,\mathrm{d}x$ 与 $\int_{a}^{+\infty}g(x)\,\mathrm{d}x$ 同敛散.

(2) 当 $c=0$ 时, 由于 $\lim\limits_{x\to+\infty}\dfrac{|f(x)|}{g(x)}=c=0$, 故对 $\varepsilon_{0}=1$, $\exists\Delta>\max\{a,0\}$, 当 $x>\Delta$ 时, 有
\[
\frac{|f(x)|}{g(x)}<\varepsilon_{0}=1,
\]
故
\[
|f(x)|<g(x).
\]
根据定理~\ref{thm:6.5.3}(1) 知, $\int_{a}^{+\infty}g(x)\,\mathrm{d}x$ 收敛蕴涵着 $\int_{a}^{+\infty}|f(x)|\,\mathrm{d}x$ 收敛.

(3) 当 $c=+\infty$ 时, 由于 $\lim\limits_{x\to+\infty}\dfrac{|f(x)|}{g(x)}=c=+\infty$, 故对 $A_{0}=1$, $\exists\Delta>\max\{a,0\}$, 当 $x>\Delta$ 时, 有
\[
\frac{|f(x)|}{g(x)}>A_{0}=1,
\]
故
\[
|f(x)|>g(x).
\]
根据定理~\ref{thm:6.5.3}(2) 知, $\int_{a}^{+\infty}g(x)\,\mathrm{d}x$ 发散蕴涵着 $\int_{a}^{+\infty}|f(x)|\,\mathrm{d}x$ 发散.
\end{proof}

\begin{corollary}\label{cor:6.5.1}
$\left(\text{与}\int_{a}^{+\infty}\dfrac{\mathrm{d}x}{x^{p}}\text{比较}\right)$ 设函数 $f$ 定义于 $[a,+\infty)$ ($a>0$) 上, 且在任何有限区间 $[a,u]$ 上 Riemann 可积, 则有
\begin{enumerate}[label=(\arabic*)]
\item 当 $|f(x)|\leqslant\dfrac{1}{x^{p}}$, $x\in[a,+\infty)$, 且 $p>1$ 时, $\int_{a}^{+\infty}|f(x)|\,\mathrm{d}x$ 收敛;
\item 当 $|f(x)|\geqslant\dfrac{1}{x^{p}}$, $x\in[a,+\infty)$, 且 $p\leqslant1$ 时, $\int_{a}^{+\infty}|f(x)|\,\mathrm{d}x$ 发散.
\end{enumerate}
\end{corollary}

\begin{proof}
由定理~\ref{thm:6.5.3} 与例~\ref{ex:6.5.1} 可立即推得.
\end{proof}

\begin{corollaryp}\label{cor:6.5.1p}
$\left(\text{与}\int_{a}^{+\infty}\dfrac{\mathrm{d}x}{x^{p}}\text{比较的极限形式}\right)$ 设 $f$ 定义于 $[a,+\infty)$ ($a>0$) 上, 且在任何有限区间 $[a,u]$ 上 Riemann 可积, 且
\[
\lim_{x\to+\infty}x^{p}|f(x)|=\lim_{x\to+\infty}\frac{|f(x)|}{\frac{1}{x^{p}}}=\lambda.
\]
则有
\begin{enumerate}[label=(\arabic*)]
\item 当 $p>1$, $0\leqslant\lambda<+\infty$ 时, $\int_{a}^{+\infty}|f(x)|\,\mathrm{d}x$ 收敛;
\item 当 $p\leqslant1$, $0<\lambda\leqslant+\infty$ 时, $\int_{a}^{+\infty}|f(x)|\,\mathrm{d}x$ 发散.
\end{enumerate}
\end{corollaryp}

\begin{proof}
由定理~\ref{thm:6.5.3p} 与例~\ref{ex:6.5.1} 可立即推得.
\end{proof}

\begin{example}\label{ex:6.5.10}
讨论 $\int_{a}^{+\infty}\dfrac{\sin x}{1+x^{2}}\,\mathrm{d}x$ 的敛散性.
\end{example}

\begin{solution}
由 $\left|\dfrac{\sin x}{1+x^{2}}\right|\leqslant\dfrac{1}{1+x^{2}}$, $x\in[0,+\infty)$ 以及 $\int_{0}^{+\infty}\dfrac{\mathrm{d}x}{1+x^{2}}=\dfrac{\pi}{2}$ 收敛, 根据比较判别法知, $\int_{a}^{+\infty}\dfrac{\sin x}{1+x^{2}}\,\mathrm{d}x$ 绝对收敛.
\end{solution}

\begin{example}\label{ex:6.5.11}
讨论下列无穷积分的敛散性:
\begin{enumerate}[label=(\arabic*)]
\item $\int_{1}^{+\infty}x^{\alpha}\mathrm{e}^{-x}\,\mathrm{d}x$;
\item $\int_{0}^{+\infty}\dfrac{x^{2}}{\sqrt{x^{5}+1}}\,\mathrm{d}x$.
\end{enumerate}
\end{example}

\begin{solution}
上面两个无穷积分中的被积函数都是非负的, 故收敛就是绝对收敛.

(1) 因为
\[
\lim_{x\to+\infty}\frac{x^{\alpha}\mathrm{e}^{-x}}{\frac{1}{x^{2}}}=\lim_{x\to+\infty}\frac{x^{\alpha+2}}{\mathrm{e}^{x}}\xlongequal{\text{L'Hospital 法则}}0,
\]
所以, 由推论~\ref{cor:6.5.1p}(1) 及 $\int_{1}^{+\infty}\dfrac{\mathrm{d}x}{x^{2}}$ 收敛推得对任何实数 $\alpha$, $\int_{1}^{+\infty}x^{\alpha}\mathrm{e}^{-x}\,\mathrm{d}x$ 收敛.

(2) 因为
\[
\lim_{x\to+\infty}\frac{\frac{x^{2}}{\sqrt{x^{5}+1}}}{\frac{1}{x^{\frac{1}{2}}}}=\lim_{x\to+\infty}\frac{x^{\frac{5}{2}}}{\sqrt{x^{5}+1}}=\lim_{x\to+\infty}\frac{1}{\sqrt{1+\frac{1}{x^{5}}}}=\frac{1}{\sqrt{1+0}}=1,
\]
所以, 由推论~\ref{cor:6.5.1p}(2) 及 $\int_{1}^{+\infty}\dfrac{\mathrm{d}x}{x^{\frac{1}{2}}}$ 发散推得 $\int_{0}^{+\infty}\dfrac{x^{2}}{\sqrt{x^{5}+1}}\,\mathrm{d}x$ 是发散的.
\end{solution}

关于 $\int_{-\infty}^{b}|f(x)|\,\mathrm{d}x$ 的比较判别法也可类似讨论.

下面介绍关于 $\int_{a}^{+\infty}f(x)g(x)\,\mathrm{d}x$ 敛散性的 Dirichlet 判别法与 Abel 判别法.

\begin{theorem}[无穷积分的 Dirichlet 判别法]\label{thm:6.5.4}
如果函数 $f$, $g$ 满足:
\begin{enumerate}[label=(\arabic*)]
\item $g$ 在 $[a,+\infty)$ 上单调, 且 $\lim\limits_{x\to+\infty}g(x)=0$;
\item $F(u)=\int_{a}^{u}f(x)\,\mathrm{d}x$ 在 $[a,+\infty)$ 上有界.
\end{enumerate}
则 $\int_{a}^{+\infty}f(x)g(x)\,\mathrm{d}x$ 收敛.
\end{theorem}

\begin{proof}
设 $\left|\int_{a}^{u}f(x)\,\mathrm{d}x\right|<M$, $\forall u>a$.

$\forall\varepsilon>0$, 由于 $\lim\limits_{x\to+\infty}g(x)=0$, 则 $\exists\Delta>\max\{a,0\}$, 当 $u_{2}>u_{1}>\Delta$ 时, 有
\[
|g(u_{2})|<\frac{\varepsilon}{4M},\quad|g(u_{1})|<\frac{\varepsilon}{4M}.
\]
因此, 由一般积分第二中值定理可知, $\exists\xi\in[u_{1},u_{2}]$, s.t.
\begin{align*}
\left|\int_{u_{1}}^{u_{2}}f(x)g(x)\,\mathrm{d}x\right|&=\left|g(u_{1})\int_{u_{1}}^{\xi}f(x)\,\mathrm{d}x+g(u_{2})\int_{\xi}^{u_{2}}f(x)\,\mathrm{d}x\right|\\
&\leqslant|g(u_{1})|\cdot\left|\int_{u_{1}}^{\xi}f(x)\,\mathrm{d}x\right|+|g(u_{2})|\cdot\left|\int_{\xi}^{u_{2}}f(x)\,\mathrm{d}x\right|\\
&=|g(u_{1})|\cdot\left|\int_{a}^{\xi}f(x)\,\mathrm{d}x-\int_{a}^{u_{1}}f(x)\,\mathrm{d}x\right|\\
&\quad+|g(u_{2})|\cdot\left|\int_{a}^{u_{2}}f(x)\,\mathrm{d}x-\int_{a}^{\xi}f(x)\,\mathrm{d}x\right|\\
&<2M\cdot\frac{\varepsilon}{4M}+2M\cdot\frac{\varepsilon}{4M}=\varepsilon.
\end{align*}
根据无穷积分的 Cauchy 收敛准则, $\int_{a}^{+\infty}f(x)g(x)\,\mathrm{d}x$ 收敛.
\end{proof}

\begin{theorem}[无穷积分的 Abel 判别法]\label{thm:6.5.5}
如果函数 $f$, $g$ 满足:
\begin{enumerate}[label=(\arabic*)]
\item $g$ 在 $[a,+\infty)$ 上单调有界;
\item $\int_{a}^{+\infty}f(x)\,\mathrm{d}x$ 收敛.
\end{enumerate}
则 $\int_{a}^{+\infty}f(x)g(x)\,\mathrm{d}x$ 收敛.
\end{theorem}

\begin{proof}[证法 1]
设 $|g(x)|<M$, $\forall x\in[a,+\infty)$.

$\forall\varepsilon>0$, 因为 $\int_{a}^{+\infty}f(x)\,\mathrm{d}x$ 收敛, 故 $\exists\Delta>\max\{a,0\}$, 当 $u_{2}>u_{1}>\Delta$ 时, 有
\[
\left|\int_{u_{1}}^{u_{2}}f(x)\,\mathrm{d}x\right|<\frac{\varepsilon}{2M}.
\]
因此, 由一般积分第二中值定理可知, $\exists\xi\in[u_{1},u_{2}]$, s.t.
\begin{align*}
&\left|\int_{u_{1}}^{u_{2}}f(x)g(x)\,\mathrm{d}x\right|=\left|g(u_{1})\int_{u_{1}}^{\xi}f(x)\,\mathrm{d}x+g(u_{2})\int_{\xi}^{u_{2}}f(x)\,\mathrm{d}x\right|\\
&\leqslant\left|g(u_{1})\right|\cdot\left|\int_{u_{1}}^{\xi}f(x)\,\mathrm{d}x\right|+\left|g(u_{2})\right|\cdot\left|\int_{\xi}^{u_{2}}f(x)\,\mathrm{d}x\right|\\
&<M\cdot\frac{\varepsilon}{2M}+M\cdot\frac{\varepsilon}{2M}=\varepsilon.
\end{align*}
根据无穷积分的 Cauchy 收敛准则, $\int_{a}^{+\infty}f(x)g(x)\,\mathrm{d}x$ 收敛.
\end{proof}

\begin{proof}[证法 2]
因 $\int_{a}^{+\infty}f(x)\,\mathrm{d}x$ 收敛, 故 $\int_{a}^{u}f(x)\,\mathrm{d}x$ 在 $[a,+\infty)$ 上有界. 又因为 $g(x)-g(+\infty)$ 在 $[a,+\infty)$ 上单调趋于 $0$ $(x\to+\infty)$, 根据 Dirichlet 判别法知,
\[
\int_{a}^{+\infty}f(x)[g(x)-g(+\infty)]\,\mathrm{d}x
\]
收敛. 因此,
\begin{align*}
\int_{a}^{+\infty}f(x)g(x)\,\mathrm{d}x=&\int_{a}^{+\infty}f(x)\big[g(x)-g(+\infty)\big]\mathrm{d}x\\
&+g(+\infty)\int_{a}^{+\infty}f(x)\,\mathrm{d}x
\end{align*}
也收敛.
\end{proof}

\begin{example}\label{ex:6.5.12}
讨论 $\int_{1}^{+\infty}\dfrac{\sin x}{x^{p}}\,\mathrm{d}x$ 与 $\int_{1}^{+\infty}\dfrac{\cos x}{x^{p}}\,\mathrm{d}x$ 的敛散性.
\end{example}

\begin{solution}
(1) 当 $p>1$ 时, 因为
\[
\left|\frac{\sin x}{x^{p}}\right|\leqslant\frac{1}{x^{p}},\quad x\in[1,+\infty),
\]
而 $\int_{1}^{+\infty}\dfrac{\mathrm{d}x}{x^{p}}$ 当 $p>1$ 时收敛, 根据比较判别法知, $\int_{1}^{+\infty}\left|\dfrac{\sin x}{x^{p}}\right|\mathrm{d}x$ 收敛, 从而 $\int_{1}^{+\infty}\dfrac{\sin x}{x^{p}}\,\mathrm{d}x$ 绝对收敛.

(2) 当 $0<p\leqslant1$ 时, 因为 $\forall u\geqslant1$, 有
\[
\left|\int_{1}^{u}\sin x\,\mathrm{d}x\right|=\left|\cos1-\cos u\right|=2,
\]
而 $\dfrac{1}{x^{p}}$ 当 $p>0$ 时单调趋于 $0$ $(x\to+\infty)$, 故由 Dirichlet 判别法推得无穷积分 $\int_{1}^{+\infty}\dfrac{\sin x}{x^{p}}\,\mathrm{d}x$ 当 $p>0$ 时总是收敛的.

另一方面, 由于
\[
\left|\frac{\sin x}{x^{p}}\right|\geqslant\frac{\sin^{2}x}{x}=\frac{1-\cos2x}{2x}=\frac{1}{2x}-\frac{\cos2x}{2x},\quad x\in[1,+\infty),
\]
其中 $\int_{1}^{+\infty}\dfrac{\cos2x}{2x}\,\mathrm{d}x\xlongequal{t=2x}\dfrac{1}{2}\int_{2}^{+\infty}\dfrac{\cos t}{t}\,\mathrm{d}t$ 满足 Dirichlet 判别法的条件, 它是收敛的. 而 $\int_{1}^{+\infty}\dfrac{\mathrm{d}x}{2x}=\dfrac{1}{2}\ln x\bigg|_{1}^{+\infty}=+\infty$ 发散. 因此, 当 $0<p\leqslant1$ 时, 由于 $\int_{1}^{+\infty}\left[\dfrac{1}{2x}-\dfrac{\cos2x}{2x}\right]\mathrm{d}x$ 发散, 因而 $\int_{1}^{+\infty}\left|\dfrac{\sin x}{x^{p}}\right|\mathrm{d}x$ 发散. 故 $\int_{1}^{+\infty}\dfrac{\sin x}{x^{p}}\,\mathrm{d}x$ 当 $0<p\leqslant1$ 时是条件收敛的.

(3) 当 $p\leqslant0$ 时, $\forall n\in\mathbb{N}$, 有
\[
\left|\int_{2n\pi}^{2n\pi+\pi}\frac{\sin x}{x^{p}}\,\mathrm{d}x\right|=\int_{2n\pi}^{2n\pi+\pi}\frac{\sin x}{x^{p}}\,\mathrm{d}x\geqslant\int_{0}^{\pi}\sin x\,\mathrm{d}x=2>1=\varepsilon_{0}.
\]
综合以上讨论可知, $\int_{1}^{+\infty}\dfrac{\sin x}{x^{p}}\,\mathrm{d}x$ 当 $p>1$ 时绝对收敛; 当 $0<p\leqslant1$ 时条件收敛; 当 $p\leqslant0$ 时发散.

类似地可证 $\int_{1}^{+\infty}\dfrac{\cos x}{x^{p}}\,\mathrm{d}x$ 当 $p>1$ 时绝对收敛; 当 $0<p\leqslant1$ 时条件收敛; 当 $p\leqslant0$ 时发散.
\end{solution}

\begin{example}\label{ex:6.5.13}
证明下列无穷积分都是条件收敛的.
\begin{enumerate}[label=(\arabic*)]
\item $\int_{1}^{+\infty}\sin x^{2}\,\mathrm{d}x$;
\item $\int_{1}^{+\infty}\cos x^{2}\,\mathrm{d}x$;
\item $\int_{1}^{+\infty}x\sin x^{4}\,\mathrm{d}x$.
\end{enumerate}
\end{example}

\begin{proof}
由
\[
\int_{1}^{+\infty}\sin x^{2}\,\mathrm{d}x\xlongequal{t=x^{2}}\int_{1}^{+\infty}\frac{\sin t}{2\sqrt{t}}\,\mathrm{d}t,
\]
\[
\int_{1}^{+\infty}\cos x^{2}\,\mathrm{d}x\xlongequal{t=x^{2}}\int_{1}^{+\infty}\frac{\cos t}{2\sqrt{t}}\,\mathrm{d}t,
\]
\[
\int_{1}^{+\infty}x\sin x^{4}\,\mathrm{d}x\xlongequal{t=x^{4}}\int_{1}^{+\infty}t^{\frac{1}{4}}\sin t\cdot\frac{\mathrm{d}t}{4t^{\frac{3}{4}}}=\int_{1}^{+\infty}\frac{\sin t}{4\sqrt{t}}\,\mathrm{d}t
\]
及例~\ref{ex:6.5.12} 立即推得它们都是条件收敛的.
\end{proof}

以上例~\ref{ex:6.5.13} 中三个无穷积分, 当 $x\to+\infty$ 时, 被积函数并不趋于 $0$, 甚至是无界的 (如 $x\sin x^{4}$), 但无穷积分收敛. 因此, 我们有必要研究 $\int_{a}^{+\infty}f(x)\,\mathrm{d}x$ 收敛与 $\lim\limits_{x\to+\infty}f(x)=0$ 的关系问题.

\begin{example}\label{ex:6.5.14}
设 $\int_{a}^{+\infty}f(x)\,\mathrm{d}x$ 收敛, 且 $\lim\limits_{x\to+\infty}f(x)=A\in\mathbb{R}$, 则 $A=0$.
\end{example}

\begin{proof}
(反证) 假设 $A\neq0$, 不妨设 $\lim\limits_{x\to+\infty}f(x)=A>0$, 则 $\exists\Delta>\max\{a,0\}$, 当 $x>\Delta$ 时, $f(x)>A-\dfrac{A}{2}=\dfrac{A}{2}$, 从而对 $u>\Delta$(图~\ref{fig:6.5.5}) 有
\begin{align*}
\int_{a}^{u}f(x)\,\mathrm{d}x&=\int_{a}^{\Delta}f(x)\,\mathrm{d}x+\int_{\Delta}^{u}f(x)\,\mathrm{d}x>\int_{a}^{\Delta}f(x)\,\mathrm{d}x+\int_{\Delta}^{u}\frac{A}{2}\,\mathrm{d}x\\
&=\int_{a}^{\Delta}f(x)\,\mathrm{d}x+\frac{A}{2}(u-\Delta)\to+\infty\quad(u\to+\infty),
\end{align*}
所以,
\[
\int_{a}^{+\infty}f(x)\,\mathrm{d}x=\lim_{u\to+\infty}\int_{a}^{u}f(x)\,\mathrm{d}x=+\infty.
\]
这与 $\int_{a}^{+\infty}f(x)\,\mathrm{d}x$ 收敛相矛盾.
\end{proof}

\begin{figure}[htbp]
\centering
\begin{tikzpicture}[x=0.95cm,y=0.90cm,>=Stealth]
  \draw[->] (-0.3,0) -- (7.6,0) node[below] {$x$};
  \draw[->] (0,-0.3) -- (0,3.7) node[left] {$y$};
  \node[below left] at (0,0) {$O$};
  \draw (0.2,1) -- (7.2,1);
  \draw (0.2,2) -- (7.2,2);
  \draw (0.2,3) -- (7.2,3);
  \draw (1.2,0) -- (1.2,3.2);
  \draw (3.0,0) -- (3.0,3.2);
  \draw (6.4,0) -- (6.4,3.2);
  \node[left] at (0,1) {$\dfrac{A}{2}$};
  \node[left] at (0,2) {$A$};
  \node[left] at (0,3) {$\dfrac{3A}{2}$};
  \fill[pattern=dots,pattern color=black] (3.0,0) rectangle (6.4,1);
  \draw[thick] (1.2,1.0) .. controls (1.9,1.06) and (2.3,1.20) .. (3.0,1.35)
    .. controls (3.7,1.55) and (4.2,2.16) .. (4.65,2.12)
    .. controls (4.95,2.02) and (5.05,1.92) .. (5.45,2.00)
    .. controls (5.9,2.06) and (6.1,2.00) .. (6.4,2.00);
  \node[right] at (3.4,2.45) {$y=f(x)$};
  \node[below] at (1.2,-0.06) {$a$};
  \node[below] at (3.0,-0.06) {$\Delta$};
  \node[below] at (6.4,-0.06) {$u$};
\end{tikzpicture}
\caption{}\label{fig:6.5.5}
\end{figure}

\begin{example}\label{ex:6.5.15}
设 $\int_{a}^{+\infty}f(x)\,\mathrm{d}x$ 收敛, $f$ 在 $[a,+\infty)$ 中一致连续, 则 $\lim\limits_{x\to+\infty}f(x)=0$.
\end{example}

\begin{proof}
(反证) 假设 $\lim\limits_{x\to+\infty}f(x)\neq0$, 则 $\exists\varepsilon_{0}>0$, s.t. $\forall\Delta>0$, $\exists x_{1}>\Delta$, 有 $|f(x_{1})|\geqslant\varepsilon_{0}$. 又因为 $f(x)$ 在 $[a,+\infty)$ 上一致连续, 故 $\exists\delta>0$, 当 $|x'-x''|<\delta$ 时, 有
\[
|f(x')-f(x'')|<\frac{\varepsilon_{0}}{2}.
\]
于是, 当 $x\in[x_{1},x_{1}+\delta]$ 时,
\[
|f(x)|=|f(x_{1})-[f(x_{1})-f(x)]|\geqslant|f(x_{1})|-|f(x_{1})-f(x)|>\varepsilon_{0}-\frac{\varepsilon_{0}}{2}=\frac{\varepsilon_{0}}{2}.
\]
并且 $f(x)$ 与 $f(x_{1})$ 同号 (否则, $|f(x)-f(x_{1})|>|f(x_{1})|\geqslant\varepsilon_{0}$ 与 $|f(x)-f(x_{1})|<\dfrac{\varepsilon_{0}}{2}$ 矛盾). 如果 $f(x_{1})>0$, 则 $f(x)>0$, 从而
\[
f(x)>\frac{\varepsilon_{0}}{2}.
\]
故
\[
\left|\int_{x_{1}}^{x_{1}+\delta}f(x)\,\mathrm{d}x\right|\geqslant\frac{\varepsilon_{0}}{2}\int_{x_{1}}^{x_{1}+\delta}\mathrm{d}x=\frac{\varepsilon_{0}}{2}\delta.
\]
同理, 如果 $f(x_{1})<0$, 则 $f(x)<0$, 也有
\[
\left|\int_{x_{1}}^{x_{1}+\delta}f(x)\,\mathrm{d}x\right|=\int_{x_{1}}^{x_{1}+\delta}|f(x)|\,\mathrm{d}x\geqslant\frac{\varepsilon_{0}}{2}\int_{x_{1}}^{x_{1}+\delta}\mathrm{d}x=\frac{\varepsilon_{0}}{2}\delta.
\]
这就证明了, 对于 $\dfrac{\varepsilon_{0}}{2}\delta>0$, $\forall\Delta>0$, $\exists x_{1}+\delta>x_{1}>\Delta$, s.t.
\[
\left|\int_{x_{1}}^{x_{1}+\delta}f(x)\,\mathrm{d}x\right|\geqslant\frac{\varepsilon_{0}}{2}\delta.
\]
根据无穷积分的 Cauchy 准则, $\int_{a}^{+\infty}f(x)\,\mathrm{d}x$ 发散, 这与题设 $\int_{a}^{+\infty}f(x)\,\mathrm{d}x$ 收敛矛盾.
\end{proof}

\begin{example}\label{ex:6.5.16}
设 $\int_{a}^{+\infty}f(x)\,\mathrm{d}x$ 收敛, $f$ 在 $[a,+\infty)$ 上连续, 且 $f(x)\geqslant0$, $x\in[a,+\infty)$. 问: 是否有 $\lim\limits_{x\to+\infty}f(x)=0$?
\end{example}

\begin{solution}
不一定有 $\lim\limits_{x\to+\infty}f(x)=0$.

反例: 令
\[
f(x)=\begin{cases}
n^{2}\cdot2^{n}\left[x-\left(n-\dfrac{1}{2^{n}n}\right)\right], & x\in\left[n-\dfrac{1}{2^{n}n},n\right],\\[3mm]
-n^{2}\cdot2^{n}\left[x-\left(n+\dfrac{1}{2^{n}n}\right)\right], & x\in\left[n,n+\dfrac{1}{2^{n}n}\right],\\[3mm]
0, & x\ \text{为}\ [0,+\infty)\ \text{中的其他点}.
\end{cases}
\]
则 $f(x)\geqslant0$, 且 $\lim\limits_{n\to+\infty}f(n)=\lim\limits_{n\to+\infty}n=+\infty$, 故 $\lim\limits_{x\to+\infty}f(x)\neq0$, 但
\[
\int_{a}^{+\infty}f(x)\,\mathrm{d}x=\sum_{n=1}^{\infty}\frac{1}{2^{n}}=1,
\]
收敛 (图~\ref{fig:6.5.6}).
\end{solution}

\begin{figure}[htbp]
\centering
\begin{tikzpicture}[x=0.82cm,y=0.60cm,>=Stealth]
  \draw[->] (-0.4,0) -- (10.6,0) node[below] {$x$};
  \draw[->] (0,-0.3) -- (0,3.9) node[left] {$y$};
  \node[below left] at (0,0) {$O$};
  \node[left] at (0,1) {1};
  \node[left] at (0,2) {2};
  \node[left] at (0,2.65) {$\vdots$};
  \node[left] at (0,3.3) {$n$};
  \draw[thick] (0,0) -- (1.05,0) -- (1.2,1) -- (1.35,0) -- (2.6,0) -- (2.8,2) -- (3.0,0) -- (7.3,0);
  \draw[thick] (7.3,0) -- (7.65,3.3) -- (8.0,0) -- (10.4,0);
  \node[below] at (1.2,-0.14) {1};
  \node[below] at (2.8,-0.14) {2};
  \node[below] at (7.65,-0.14) {$n$};
  \node[anchor=north east] at (7.15,-0.32) {$n-\dfrac{1}{2^{n}n}$};
  \node[anchor=north west] at (8.15,-0.32) {$n+\dfrac{1}{2^{n}n}$};
  \draw[thin] (7.15,-0.32) -- (7.29,-0.06);
  \draw[thin] (8.15,-0.32) -- (8.01,-0.06);
  \draw (7.65,3.3) -- (8.6,3.3);
  \node[right] at (8.6,3.3) {\small 面积 $S_{n}=\dfrac{1}{2}\cdot n\cdot\dfrac{1}{2^{n-1}n}=\dfrac{1}{2^{n}}$};
\end{tikzpicture}
\caption{}\label{fig:6.5.6}
\end{figure}

类似无穷积分, 关于瑕积分有相应的性质与判别法.

\begin{theorem}[瑕积分收敛的 Cauchy 准则]\label{thm:6.5.6}
瑕积分 $\int_{a}^{b}f(x)\,\mathrm{d}x$ ($a$ 为瑕点) 收敛 $\Longleftrightarrow\forall\varepsilon>0$, $\exists\delta>0$, 当 $v_{1},v_{2}\in(a,a+\delta)$ 时, 便有
\[
\left|\int_{v_{1}}^{b}f(x)\,\mathrm{d}x-\int_{v_{2}}^{b}f(x)\,\mathrm{d}x\right|=\left|\int_{v_{1}}^{v_{2}}f(x)\,\mathrm{d}x\right|<\varepsilon.
\]
\end{theorem}

\begin{proof}
设 $F(v)=\int_{v}^{b}f(x)\,\mathrm{d}x$, 则

$\int_{a}^{b}f(x)\,\mathrm{d}x$ 收敛 $\Longleftrightarrow$ $\lim\limits_{v\to a^{+}}\int_{v}^{b}f(x)\,\mathrm{d}x=\lim\limits_{v\to a^{+}}F(v)$ 存在且有限
\begin{align*}
&\xLeftrightarrow{\text{函数收敛的 Cauchy 准则}}\forall\varepsilon>0,\exists\delta>0,\ \text{当}\ v_{1},v_{2}\in(a,a+\delta)\ \text{时, 便有}\\
&\left|\int_{v_{1}}^{v_{2}}f(x)\,\mathrm{d}x\right|=\left|\int_{a}^{v_{2}}f(x)\,\mathrm{d}x-\int_{a}^{v_{1}}f(x)\,\mathrm{d}x\right|\\
&=\left|F(v_{2})-F(v_{1})\right|<\varepsilon.
\end{align*}
\end{proof}

\begin{lemma}\label{lem:6.5.4}
设函数 $f$, $g$ 的瑕点同为 $a$, $\lambda$ 与 $\mu$ 为常数, 如果瑕积分 $\int_{a}^{b}f(x)\,\mathrm{d}x$ 与 $\int_{a}^{b}g(x)\,\mathrm{d}x$ 都收敛, 则瑕积分 $\int_{a}^{b}[\lambda f(x)+\mu g(x)]\,\mathrm{d}x$ 也收敛 (此时它也许为正常积分, $a$ 未必为瑕点, 如 $k_{1}=k_{2}=0$, 而正常积分视作收敛), 并有
\[
\int_{a}^{b}\left[\lambda f(x)+\mu g(x)\right]\mathrm{d}x=\lambda\int_{a}^{b}f(x)\,\mathrm{d}x+\mu\int_{a}^{b}g(x)\,\mathrm{d}x.
\]
\end{lemma}

\begin{proof}
\begin{align*}
\int_{a}^{b}\left[\lambda f(x)+\mu g(x)\right]\mathrm{d}x&=\lim_{v\to a^{+}}\int_{v}^{b}\left[\lambda f(x)+\mu g(x)\right]\mathrm{d}x\\
&=\lim_{v\to a^{+}}\left[\lambda\int_{v}^{b}f(x)\,\mathrm{d}x+\mu\int_{v}^{b}g(x)\,\mathrm{d}x\right]\\
&=\lambda\int_{a}^{b}f(x)\,\mathrm{d}x+\mu\int_{a}^{b}g(x)\,\mathrm{d}x.
\end{align*}
\end{proof}

\begin{lemma}\label{lem:6.5.5}
设函数 $f$ 的瑕点为 $a$, 则瑕积分 $\int_{a}^{b}f(x)\,\mathrm{d}x$ 与 $\int_{a}^{c}f(x)\,\mathrm{d}x$ ($c\in(a,b)$) 同敛散. 如果收敛, 则
\[
\int_{a}^{b}f(x)\,\mathrm{d}x=\int_{a}^{c}f(x)\,\mathrm{d}x+\int_{c}^{b}f(x)\,\mathrm{d}x,
\]
其中 $\int_{c}^{b}f(x)\,\mathrm{d}x$ 为 Riemann 积分.
\end{lemma}

\begin{proof}
设 $a<v<c$, 则
\[
\int_{v}^{b}f(x)\,\mathrm{d}x=\int_{v}^{c}f(x)\,\mathrm{d}x+\int_{c}^{b}f(x)\,\mathrm{d}x.
\]
因此, $\lim\limits_{v\to a^{+}}\int_{v}^{b}f(x)\,\mathrm{d}x$ 与 $\lim\limits_{v\to a^{+}}\int_{v}^{c}f(x)\,\mathrm{d}x$ 同时存在或同时不存在. 故瑕积分 $\int_{a}^{b}f(x)\,\mathrm{d}x$ 与 $\int_{a}^{c}f(x)\,\mathrm{d}x$ 同敛散. 当它们收敛时, 有
\begin{align*}
\int_{a}^{b}f(x)\,\mathrm{d}x&=\lim_{v\to a^{+}}\int_{v}^{b}f(x)\,\mathrm{d}x=\lim_{v\to a^{+}}\left[\int_{v}^{c}f(x)\,\mathrm{d}x+\int_{c}^{b}f(x)\,\mathrm{d}x\right]\\
&=\int_{a}^{c}f(x)\,\mathrm{d}x+\int_{c}^{b}f(x)\,\mathrm{d}x.
\end{align*}
\end{proof}

\begin{lemma}\label{lem:6.5.6}
设函数 $f$ 的瑕点为 $a$, $f$ 在 $(a,b]$ 的任一内闭区间 $[u,b]$ 上 Riemann 可积, 则当 $\int_{a}^{b}|f(x)|\,\mathrm{d}x$ 收敛时, $\int_{a}^{b}f(x)\,\mathrm{d}x$ 也必收敛, 并有
\[
\left|\int_{a}^{b}f(x)\,\mathrm{d}x\right|\leqslant\int_{a}^{b}\left|f(x)\right|\mathrm{d}x.
\]
\end{lemma}

\begin{proof}
由 $\int_{a}^{b}|f(x)|\,\mathrm{d}x$ 收敛, 根据瑕积分收敛的 Cauchy 准则的必要性知, $\forall\varepsilon>0$, $\exists\delta>0$, 当 $a<v_{1}<v_{2}<a+\delta$ 时, 总有
\[
\left|\int_{v_{1}}^{v_{2}}|f(x)|\,\mathrm{d}x\right|=\int_{v_{1}}^{v_{2}}|f(x)|\,\mathrm{d}x<\varepsilon.
\]
于是,
\[
\left|\int_{v_{1}}^{v_{2}}f(x)\,\mathrm{d}x\right|\leqslant\int_{v_{1}}^{v_{2}}|f(x)|\,\mathrm{d}x<\varepsilon.
\]
再由瑕积分收敛的 Cauchy 准则的充分性推得 $\int_{a}^{b}f(x)\,\mathrm{d}x$ 收敛.

又因为
\[
\left|\int_{v}^{b}f(x)\,\mathrm{d}x\right|\leqslant\int_{v}^{b}|f(x)|\,\mathrm{d}x\quad(a<v\leqslant b),
\]
令 $v\to a^{+}$, 得到
\[
\left|\int_{a}^{b}f(x)\,\mathrm{d}x\right|\leqslant\int_{a}^{b}|f(x)|\,\mathrm{d}x.
\]
\end{proof}

\begin{definition}\label{def:6.5.4}
设 $f$ 在任何内闭区间 $[v,b]\subset(a,b]$ 上 Riemann 可积, $a$ 为其瑕点, 如果 $\int_{a}^{b}|f(x)|\,\mathrm{d}x$ 收敛, 则称 $\int_{a}^{b}f(x)\,\mathrm{d}x$ 绝对收敛. 引理~\ref{lem:6.5.6} 指出 $\int_{a}^{b}f(x)\,\mathrm{d}x$ 也收敛, 但反之不成立. 读者可仿照定义~\ref{def:6.5.3} 中的反例构造瑕积分的反例.
\end{definition}

我们称收敛而非绝对收敛的瑕积分为条件收敛的.

\begin{theorem}[瑕积分的比较判别法]\label{thm:6.5.7}
设定义在 $(a,b]$ 上的两个函数 $f$ 与 $g$, 其瑕点同为 $a$, 在任何内闭区间 $[v,b]\subset(a,b]$ 上都 Riemann 可积, 且满足
\[
|f(x)|\leqslant g(x),\quad x\in(a,b].
\]
则
\begin{enumerate}[label=(\arabic*)]
\item 当 $\int_{a}^{b}g(x)\,\mathrm{d}x$ 收敛时, $\int_{a}^{b}|f(x)|\,\mathrm{d}x$ 必收敛, 当然 $\int_{a}^{b}f(x)\,\mathrm{d}x$ 也收敛;
\item 当 $\int_{a}^{b}|f(x)|\,\mathrm{d}x$ 发散时, $\int_{a}^{b}g(x)\,\mathrm{d}x$ 必发散.
\end{enumerate}
(简单说即是: 大的收敛, 小的也收敛; 小的发散, 大的也发散.)
\end{theorem}

\begin{proof}
因为 $0\leqslant|f(x)|\leqslant g(x)$, 所以 $\int_{v}^{b}|f(x)|\,\mathrm{d}x$ 与 $\int_{v}^{b}g(x)\,\mathrm{d}x$ 都为 $v$ 当 $v\to a^{+}$ 时的增函数.

(1) 如果 $\int_{a}^{b}g(x)\,\mathrm{d}x$ 收敛, 则 $\forall v\in(a,b]$, 有 $\int_{v}^{b}g(x)\,\mathrm{d}x\leqslant M$(常数), 从而
\[
\int_{v}^{b}|f(x)|\,\mathrm{d}x\leqslant\int_{v}^{b}g(x)\,\mathrm{d}x\leqslant M,\quad\forall v\in(a,b],
\]
$\int_{a}^{b}|f(x)|\,\mathrm{d}x=\lim\limits_{v\to a^{+}}\int_{v}^{b}|f(x)|\,\mathrm{d}x$ 存在且有限, 即 $\int_{a}^{b}|f(x)|\,\mathrm{d}x$ 收敛.

(2) 如果 $\int_{a}^{b}|f(x)|\,\mathrm{d}x$ 发散, 则 $\lim\limits_{v\to a^{+}}\int_{v}^{b}|f(x)|\,\mathrm{d}x=\int_{a}^{b}|f(x)|\,\mathrm{d}x=+\infty$, 则由
\[
\int_{v}^{b}|f(x)|\,\mathrm{d}x\leqslant\int_{v}^{b}g(x)\,\mathrm{d}x
\]
得到
\[
+\infty=\lim_{v\to a^{+}}\int_{v}^{b}|f(x)|\,\mathrm{d}x\leqslant\lim_{v\to a^{+}}\int_{v}^{b}g(x)\,\mathrm{d}x=\int_{a}^{b}g(x)\,\mathrm{d}x,
\]
即
\[
\int_{a}^{b}g(x)\,\mathrm{d}x=+\infty,
\]
发散.
\end{proof}

\begin{theoremp}[瑕积分的比较判别法的极限形式]\label{thm:6.5.7p}
设 $f$, $g$ 在任何有限区间 $[v,b]\subset(a,b]$ 上都 Riemann 可积, $g(x)>0$, 且
\[
\lim_{x\to a^{+}}\frac{|f(x)|}{g(x)}=c.
\]
则有
\begin{enumerate}[label=(\arabic*)]
\item 当 $0<c<+\infty$ 时, $\int_{a}^{b}|f(x)|\,\mathrm{d}x$ 与 $\int_{a}^{b}g(x)\,\mathrm{d}x$ 同敛散;
\item 当 $c=0$ 时, $\int_{a}^{b}g(x)\,\mathrm{d}x$ 收敛蕴涵着 $\int_{a}^{b}|f(x)|\,\mathrm{d}x$ 也收敛;
\item 当 $c=+\infty$ 时, $\int_{a}^{b}g(x)\,\mathrm{d}x$ 发散蕴涵着 $\int_{a}^{b}|f(x)|\,\mathrm{d}x$ 也发散.
\end{enumerate}
\end{theoremp}

\begin{proof}
(1) 因为 $\lim\limits_{x\to a^{+}}\dfrac{|f(x)|}{g(x)}=c\in(0,+\infty)$, 故 $\exists\delta>0$, 当 $x\in(a,a+\delta)$ 时, 有
\[
\frac{3c}{2}=c+\frac{c}{2}>\frac{|f(x)|}{g(x)}>c-\frac{c}{2}=\frac{c}{2}>0,
\]
\[
\frac{3c}{2}g(x)>|f(x)|>\frac{c}{2}g(x).
\]
再由定理~\ref{thm:6.5.7} 知, $\int_{a}^{b}|f(x)|\,\mathrm{d}x$ 与 $\int_{a}^{b}g(x)\,\mathrm{d}x$ 同敛散.

(2) 当 $c=0$ 时, 由 $\lim\limits_{x\to a^{+}}\dfrac{|f(x)|}{g(x)}=c=0$, 故对 $\varepsilon_{0}=1$, $\exists\delta>0$, 当 $x\in(a,a+\delta)$ 时, 有
\[
\frac{|f(x)|}{g(x)}<1,
\]
故
\[
|f(x)|<g(x).
\]
于是, 根据定理~\ref{thm:6.5.7}(1) 知, $\int_{a}^{b}g(x)\,\mathrm{d}x$ 收敛蕴涵着 $\int_{a}^{b}|f(x)|\,\mathrm{d}x$ 也收敛.

(3) 当 $c=+\infty$ 时, 由于 $\lim\limits_{x\to a^{+}}\dfrac{|f(x)|}{g(x)}=c=+\infty$, 故对 $A_{0}=1$, $\exists\delta>0$, 当 $x\in(a,a+\delta)$ 时, 有
\[
\frac{|f(x)|}{g(x)}>A_{0}=1,
\]
故
\[
|f(x)|>g(x).
\]
于是, 根据定理~\ref{thm:6.5.7}(2) 知, $\int_{a}^{b}g(x)\,\mathrm{d}x$ 发散蕴涵着 $\int_{a}^{b}|f(x)|\,\mathrm{d}x$ 也发散.
\end{proof}

\begin{corollary}\label{cor:6.5.2}
$\left(\text{与}\int_{a}^{b}\dfrac{\mathrm{d}x}{(x-a)^{p}}\text{比较}\right)$ 设 $f$ 定义于 $(a,b]$, $a$ 为其瑕点, 且在任何内闭区间 $[u,b]\subset(a,b]$ 上 Riemann 可积. 则有
\begin{enumerate}[label=(\arabic*)]
\item 当 $|f(x)|\leqslant\dfrac{1}{(x-a)^{p}}$ 且 $p<1$ 时, $\int_{a}^{b}|f(x)|\,\mathrm{d}x$ 收敛;
\item 当 $|f(x)|\geqslant\dfrac{1}{(x-a)^{p}}$, 且 $p\geqslant1$ 时, $\int_{a}^{b}|f(x)|\,\mathrm{d}x$ 发散.
\end{enumerate}
\end{corollary}

\begin{proof}
由定理~\ref{thm:6.5.7} 与例~\ref{ex:6.5.9} 可立即推得.
\end{proof}

\begin{corollaryp}\label{cor:6.5.2p}
$\left(\text{与}\int_{a}^{b}\dfrac{\mathrm{d}x}{(x-a)^{p}}\text{比较的极限形式}\right)$ 设 $f$ 定义于 $(a,b]$, $a$ 为其瑕点, 且在任何内闭区间 $[u,b]\subset(a,b]$ 上 Riemann 可积. 如果
\[
\lim_{x\to a^{+}}(x-a)^{p}|f(x)|=\lim_{x\to a^{+}}\frac{|f(x)|}{\frac{1}{(x-a)^{p}}}=\lambda,
\]
则有
\begin{enumerate}[label=(\arabic*)]
\item 当 $p<1$, $0\leqslant\lambda<+\infty$ 时, $\int_{a}^{b}|f(x)|\,\mathrm{d}x$ 收敛;
\item 当 $p\geqslant1$, $0<\lambda\leqslant+\infty$ 时, $\int_{a}^{b}|f(x)|\,\mathrm{d}x$ 发散.
\end{enumerate}
\end{corollaryp}

\begin{proof}
由定理~\ref{thm:6.5.7p} 与例~\ref{ex:6.5.9} 可立即推得.
\end{proof}

\begin{example}\label{ex:6.5.17}
讨论广义积分
\[
\int_{0}^{+\infty}\frac{\ln(1+x)}{x^{p}}\,\mathrm{d}x
\]
的敛散性.
\end{example}

\begin{solution}
\begin{align*}
\int_{0}^{+\infty}\frac{\ln(1+x)}{x^{p}}\,\mathrm{d}x&=\int_{0}^{1}\frac{\ln(1+x)}{x^{p}}\,\mathrm{d}x+\int_{1}^{+\infty}\frac{\ln(1+x)}{x^{p}}\,\mathrm{d}x\\
&=I_{1}+I_{2}.
\end{align*}
首先讨论 $I_{1}=\int_{0}^{1}\dfrac{\ln(1+x)}{x^{p}}\,\mathrm{d}x$ 的敛散性, $x=0$ 疑似瑕点.

因为 $\dfrac{\ln(1+x)}{x^{p}}\sim\dfrac{x}{x^{p}}=\dfrac{1}{x^{p-1}}$ $(x\to0^{+})$, 所以, 根据推论~\ref{cor:6.5.2p}, 当 $p-1<1$, 即 $p<2$ 时, $I_{1}$ 收敛; 当 $p-1\geqslant1$, 即 $p\geqslant2$ 时发散.

再来讨论 $I_{2}=\int_{1}^{+\infty}\dfrac{\ln(1+x)}{x^{p}}\,\mathrm{d}x$ 的敛散性.

当 $p\leqslant1$ 时, 由 $\lim\limits_{x\to+\infty}\dfrac{\ln(1+x)}{x^{p}}\Big/\dfrac{1}{x^{p}}=\lim\limits_{x\to+\infty}\ln(1+x)=+\infty$ 及推论~\ref{cor:6.5.1p}(2) 知, $I_{2}$ 发散.

当 $p>1$ 时, 令 $p=1+2\delta$, 则 $\delta>0$, 于是,
\begin{align*}
\lim_{x\to+\infty}\frac{\ln(1+x)}{x^{p}}\Big/\frac{1}{x^{1+\delta}}&=\lim_{x\to+\infty}\frac{\ln(1+x)}{x^{\delta}}\\
&=\lim_{x\to+\infty}\frac{1}{\delta x^{\delta-1}(1+x)}=0.
\end{align*}
再由 $\int_{1}^{+\infty}\dfrac{\mathrm{d}x}{x^{1+\delta}}$ 收敛及推论~\ref{cor:6.5.1p}(1) 立刻知道, $I_{2}$ 收敛.

综合上述讨论, 当且仅当 $1<p<2$ 时, $\int_{0}^{+\infty}\dfrac{\ln(1+x)}{x^{p}}\,\mathrm{d}x$ 收敛.
\end{solution}

\begin{example}\label{ex:6.5.18}
证明: 瑕积分 $\int_{0}^{\frac{\pi}{2}}\ln\sin x\,\mathrm{d}x$ 与 $\int_{0}^{\frac{\pi}{2}}\ln\cos x\,\mathrm{d}x$ 均收敛, 且
\[
\int_{0}^{\frac{\pi}{2}}\ln\sin x\,\mathrm{d}x=\int_{0}^{\frac{\pi}{2}}\ln\cos x\,\mathrm{d}x=-\frac{\pi}{2}\ln2.
\]
\end{example}

\begin{proof}[证法 1]
\begin{align*}
\int_{0}^{\frac{\pi}{2}}\ln\cos x\,\mathrm{d}x&\xlongequal{t=\frac{\pi}{2}-x}\int_{\frac{\pi}{2}}^{0}\ln\sin t\,(-\mathrm{d}t)=\int_{0}^{\frac{\pi}{2}}\ln\sin t\,\mathrm{d}t\\
&=\int_{0}^{\frac{\pi}{2}}\ln\left(2\sin\frac{t}{2}\cos\frac{t}{2}\right)\mathrm{d}t=\int_{0}^{\frac{\pi}{2}}\left(\ln2+\ln\sin\frac{t}{2}+\ln\cos\frac{t}{2}\right)\mathrm{d}t\\
&\xlongequal{u=\frac{t}{2}}\frac{\pi}{2}\ln2+2\left(\int_{0}^{\frac{\pi}{4}}\ln\sin u\,\mathrm{d}u+\int_{0}^{\frac{\pi}{4}}\ln\cos u\,\mathrm{d}u\right)\\
&\xlongequal{v=\frac{\pi}{2}-u}\frac{\pi}{2}\ln2+2\left[\int_{0}^{\frac{\pi}{4}}\ln\sin u\,\mathrm{d}u+\int_{\frac{\pi}{2}}^{\frac{\pi}{4}}\ln\sin v\,(-\mathrm{d}v)\right]\\
&=\frac{\pi}{2}\ln2+2\int_{0}^{\frac{\pi}{2}}\ln\sin t\,\mathrm{d}t,
\end{align*}
移项化简得
\[
\int_{0}^{\frac{\pi}{2}}\ln\cos x\,\mathrm{d}x=\int_{0}^{\frac{\pi}{2}}\ln\sin x\,\mathrm{d}x=-\frac{\pi}{2}\ln2.
\]
\end{proof}

\begin{proof}[证法 2]
由证法 1, 有
\begin{align*}
\int_{0}^{\frac{\pi}{2}}\ln\sin x\,\mathrm{d}x&=\frac{1}{2}\int_{0}^{\frac{\pi}{2}}(\ln\sin x+\ln\cos x)\,\mathrm{d}x\\
&=\frac{1}{2}\int_{0}^{\frac{\pi}{2}}\ln\left(\frac{1}{2}\sin2x\right)\mathrm{d}x\\
&=-\frac{1}{2}\ln2\cdot\frac{\pi}{2}+\frac{1}{2}\int_{0}^{\frac{\pi}{2}}\ln\sin2x\,\mathrm{d}x\\
&\xlongequal{u=2x}-\frac{\pi}{4}\ln2+\frac{1}{2}\int_{0}^{\pi}\ln\sin u\,\frac{\mathrm{d}u}{2}\\
&=-\frac{\pi}{4}\ln2+\frac{1}{4}\int_{0}^{\frac{\pi}{2}}\ln\sin u\,\mathrm{d}u+\frac{1}{4}\int_{\frac{\pi}{2}}^{\pi}\ln\sin u\,\mathrm{d}u\\
&\xlongequal{u=\pi-x}-\frac{\pi}{4}\ln2+\frac{1}{4}\int_{0}^{\frac{\pi}{2}}\ln\sin x\,\mathrm{d}x+\frac{1}{4}\int_{\frac{\pi}{2}}^{0}\ln\sin x\,(-\mathrm{d}x)\\
&=-\frac{\pi}{4}\ln2+\frac{1}{2}\int_{0}^{\frac{\pi}{2}}\ln\sin x\,\mathrm{d}x,
\end{align*}
移项化简得
\[
\int_{0}^{\frac{\pi}{2}}\ln\sin x\,\mathrm{d}x=-\frac{\pi}{2}\ln2.
\]
\end{proof}

\begin{example}\label{ex:6.5.19}
设 $f$ 在 $[0,+\infty)$ 上连续, $0<a<b$. 证明:
\begin{enumerate}[label=(\arabic*)]
\item 如果 $\lim\limits_{x\to+\infty}f(x)=k$, 则
\[
\int_{0}^{+\infty}\frac{f(ax)-f(bx)}{x}\,\mathrm{d}x=\left[f(0)-k\right]\ln\frac{b}{a};
\]
\item 如果 $\int_{0}^{+\infty}\frac{f(x)}{x}\,\mathrm{d}x$ 收敛, 则
\[
\int_{0}^{+\infty}\frac{f(ax)-f(bx)}{x}\,\mathrm{d}x=f(0)\ln\frac{b}{a}.
\]
\end{enumerate}
\end{example}

\begin{proof}
(1)
\begin{align*}
\int_{\varepsilon}^{A}\frac{f(ax)-f(bx)}{x}\,\mathrm{d}x&=\int_{\varepsilon}^{A}\frac{f(ax)}{x}\,\mathrm{d}x-\int_{\varepsilon}^{A}\frac{f(bx)}{x}\,\mathrm{d}x\\
&\xlongequal[{t=bx}]{t=ax}\int_{a\varepsilon}^{aA}\frac{f(t)}{\frac{t}{a}}\frac{\mathrm{d}t}{a}-\int_{b\varepsilon}^{bA}\frac{f(t)}{\frac{t}{b}}\frac{\mathrm{d}t}{b}=\int_{a\varepsilon}^{aA}\frac{f(t)}{t}\,\mathrm{d}t-\int_{b\varepsilon}^{bA}\frac{f(t)}{t}\,\mathrm{d}t\\
&=\int_{a\varepsilon}^{b\varepsilon}\frac{f(t)}{t}\,\mathrm{d}t-\int_{aA}^{bA}\frac{f(t)}{t}\,\mathrm{d}t\\
&\xlongequal[{t=Au}]{t=\varepsilon u}\int_{a}^{b}\frac{f(\varepsilon u)}{\varepsilon u}\varepsilon\,\mathrm{d}u-\int_{a}^{b}\frac{f(Au)}{Au}A\,\mathrm{d}u\\
&=\int_{a}^{b}\frac{\left[f(\varepsilon u)-f(Au)\right]}{u}\,\mathrm{d}u\\
&\xlongequal[{\exists\xi\in[a,b]}]{\text{积分第一中值定理}}\left[f(\varepsilon\xi)-f(A\xi)\right]\int_{a}^{b}\frac{\mathrm{d}u}{u}=\left[f(\varepsilon\xi)-f(A\xi)\right]\ln\frac{b}{a}.
\end{align*}
于是,
\begin{align*}
\int_{0}^{+\infty}\frac{f(ax)-f(bx)}{x}\,\mathrm{d}x&=\lim_{\substack{\varepsilon\to0^{+}\\ A\to+\infty}}\int_{\varepsilon}^{A}\frac{f(ax)-f(bx)}{x}\,\mathrm{d}x\\
&=\lim_{\substack{\varepsilon\to0^{+}\\ A\to+\infty}}\left[f(\varepsilon\xi)-f(A\xi)\right]\ln\frac{b}{a}=\left[f(0)-f(+\infty)\right]\ln\frac{b}{a}\\
&=\left(f(0)-k\right)\ln\frac{b}{a}.
\end{align*}
(2) 由于 $\int_{0}^{+\infty}\dfrac{f(x)}{x}\,\mathrm{d}x$ 收敛, 故对 $\forall\varepsilon>0$, 有
\[
\int_{\varepsilon}^{+\infty}\frac{f(ax)}{x}\,\mathrm{d}x\xlongequal{t=ax}\int_{\varepsilon a}^{+\infty}\frac{f(t)}{\frac{t}{a}}\frac{\mathrm{d}t}{a}=\int_{\varepsilon a}^{+\infty}\frac{f(t)}{t}\,\mathrm{d}t.
\]
于是,
\begin{align*}
\int_{\varepsilon}^{+\infty}\frac{f(ax)-f(bx)}{x}\,\mathrm{d}x&=\int_{\varepsilon a}^{+\infty}\frac{f(t)}{t}\,\mathrm{d}t-\int_{\varepsilon b}^{+\infty}\frac{f(t)}{t}\,\mathrm{d}t\\
&=\int_{\varepsilon a}^{\varepsilon b}\frac{f(t)}{t}\,\mathrm{d}t\xlongequal{t=\varepsilon x}\int_{a}^{b}\frac{f(\varepsilon x)}{\varepsilon x}\varepsilon\,\mathrm{d}x=\int_{a}^{b}\frac{f(\varepsilon x)}{x}\,\mathrm{d}x\\
&\xlongequal[{\exists\xi\in[a,b]}]{\text{积分第一中值定理}}f(\varepsilon\xi)\int_{a}^{b}\frac{\mathrm{d}x}{x}=f(\varepsilon\xi)\ln\frac{b}{a}.
\end{align*}
再令 $\varepsilon\to0^{+}$, 得到
\begin{align*}
\int_{0}^{+\infty}\frac{f(ax)-f(bx)}{x}\,\mathrm{d}x&=\lim_{\varepsilon\to0^{+}}\int_{\varepsilon}^{+\infty}\frac{f(ax)-f(bx)}{x}\,\mathrm{d}x\\
&=\lim_{\varepsilon\to0^{+}}f(\varepsilon\xi)\ln\frac{b}{a}=f(0)\ln\frac{b}{a}.
\end{align*}
\end{proof}

\begin{exblock}{练习题 6.5}

  \item 证明: $\lim\limits_{n\to+\infty}\dfrac{\sqrt[n]{n!}}{n}=\dfrac{1}{\mathrm{e}}$.

  (提示: (1) 参阅例~\ref{ex:1.4.11}; (2) 对 $\dfrac{\sum\limits_{k=1}^{n}(\ln k-\ln n)}{n}$ 应用 Stolz 公式; (3) $\lim\limits_{n\to+\infty}\dfrac{1}{n}\sum\limits_{k=1}^{n}\ln\dfrac{k}{n}=\int_{0}^{1}\ln x\,\mathrm{d}x$.)

  \item 在 $[-1,+\infty)$ 上定义函数
  \[
  f(x)=\int_{-1}^{x}\frac{\mathrm{e}^{\frac{1}{t}}}{t^{2}\left(1+\mathrm{e}^{\frac{1}{t}}\right)^{2}}\,\mathrm{d}t,
\]
  试写出函数 $f$ 的简单表达式.

  \item 计算下列无穷积分:
  \begin{enumerate}[label=(\arabic*)]
    \item $\int_{2}^{+\infty}\dfrac{\mathrm{d}x}{x(\ln x)^{p}}$, $p>1$;
    \item $\int_{0}^{+\infty}\mathrm{e}^{-\sqrt{x}}\,\mathrm{d}x$;
    \item $\int_{-\infty}^{0}x\mathrm{e}^{x}\,\mathrm{d}x$;
    \item $\int_{0}^{+\infty}x^{5}\mathrm{e}^{-x^{2}}\,\mathrm{d}x$;
    \item $\int_{1}^{+\infty}\dfrac{\mathrm{d}x}{x(1+x)}$;
    \item $\int_{0}^{+\infty}\dfrac{\mathrm{d}x}{1+x^{3}}$;
    \item $\int_{-\infty}^{+\infty}\dfrac{\mathrm{d}x}{x^{2}+2x+2}$;
    \item $\int_{-\infty}^{+\infty}\dfrac{\mathrm{d}x}{(x^{2}+x+1)^{2}}$;
    \item $\int_{0}^{+\infty}x^{n-1}\mathrm{e}^{-x}\,\mathrm{d}x$, $n\in\mathbb{N}$;
    \item $\int_{0}^{+\infty}\dfrac{\mathrm{d}x}{(x^{2}+a^{2})^{n}}$, $n\in\mathbb{N}$;
    \item $\int_{0}^{+\infty}x^{2n-1}\mathrm{e}^{-x^{2}}\,\mathrm{d}x$, $n\in\mathbb{N}$;
    \item $\int_{0}^{+\infty}\mathrm{e}^{-ax}\cos bx\,\mathrm{d}x$, 其中 $a$, $b$ 为常数, $a>0$;
    \item $\int_{0}^{+\infty}\dfrac{1+x^{2}}{1+x^{4}}\,\mathrm{d}x$;
    \item $\int_{2}^{+\infty}\dfrac{x\ln x}{1-x^{2}}\,\mathrm{d}x$;
    \item $\int_{0}^{+\infty}x\mathrm{e}^{-x^{2}}\,\mathrm{d}x$;
    \item $\int_{-\infty}^{+\infty}x\mathrm{e}^{-x^{2}}\,\mathrm{d}x$;
    \item $\int_{0}^{+\infty}\dfrac{\mathrm{d}x}{\sqrt{\mathrm{e}^{x}}}$;
    \item $\int_{0}^{+\infty}\mathrm{e}^{-x}\sin x\,\mathrm{d}x$;
    \item $\int_{1}^{+\infty}\dfrac{\mathrm{d}x}{x^{2}(1+x)}$;
    \item $\int_{-\infty}^{+\infty}\dfrac{\mathrm{d}x}{4x^{2}+4x+5}$.
  \end{enumerate}

  \item 判断下列无穷积分的敛散性:
  \begin{enumerate}[label=(\arabic*)]
    \item $\int_{1}^{+\infty}\dfrac{\mathrm{d}x}{\sqrt[3]{x^{4}+1}}$;
    \item $\int_{1}^{+\infty}\dfrac{x}{1-\mathrm{e}^{x}}\,\mathrm{d}x$;
    \item $\int_{0}^{+\infty}\dfrac{\mathrm{d}x}{1+\sqrt{x}}$;
    \item $\int_{1}^{+\infty}\dfrac{x\arctan x}{1+x^{3}}\,\mathrm{d}x$;
    \item $\int_{1}^{+\infty}\dfrac{\ln(1+x)}{x^{n}}\,\mathrm{d}x$;
    \item $\int_{0}^{+\infty}\dfrac{x^{m}}{1+x^{n}}\,\mathrm{d}x$;
    \item $\int_{1}^{+\infty}\dfrac{\mathrm{d}x}{x(1+x)^{\frac{1}{3}}}$;
    \item $\int_{0}^{+\infty}\dfrac{\cos x}{1+x^{2}}\,\mathrm{d}x$;
    \item $\int_{0}^{+\infty}\dfrac{x^{3}}{\sqrt{1+x^{7}}}\,\mathrm{d}x$;
    \item $\int_{1}^{+\infty}\dfrac{x^{2}}{\sqrt{1+x^{6}}}\,\mathrm{d}x$.
  \end{enumerate}

  \item 讨论下列无穷积分是绝对收敛还是条件收敛:
  \begin{enumerate}[label=(\arabic*)]
    \item $\int_{1}^{+\infty}\dfrac{\sin\sqrt{x}}{x}\,\mathrm{d}x$;
    \item $\int_{0}^{+\infty}\dfrac{\operatorname{sgn}(\sin x)}{1+x^{2}}\,\mathrm{d}x$;
    \item $\int_{0}^{+\infty}\dfrac{\sqrt{x}\cos x}{100+x}\,\mathrm{d}x$;
    \item $\int_{e}^{+\infty}\dfrac{\ln(\ln x)}{\ln x}\sin x\,\mathrm{d}x$.
  \end{enumerate}

  \item 计算下列瑕积分:
  \begin{enumerate}[label=(\arabic*)]
    \item $\int_{-1}^{1}\dfrac{\mathrm{d}x}{\sqrt{1-x^{2}}}$;
    \item $\int_{-1}^{1}\dfrac{\arcsin x}{\sqrt{1-x^{2}}}\,\mathrm{d}x$;
    \item $\int_{0}^{1}\dfrac{\mathrm{d}x}{(2-x)\sqrt{1-x}}$;
    \item $\int_{-1}^{1}\dfrac{\mathrm{d}x}{(2-x^{2})\sqrt{1-x^{2}}}$;
    \item $\int_{0}^{1}\dfrac{\arcsin\sqrt{x}}{\sqrt{x(1-x)}}\,\mathrm{d}x$;
    \item $\int_{0}^{1}(\ln x)^{n}\,\mathrm{d}x$, $n\in\mathbb{N}$;
    \item $\int_{0}^{1}\dfrac{(1-x)^{n}}{\sqrt{x}}\,\mathrm{d}x$, $n\in\mathbb{N}$;
    \item $\int_{0}^{1}\dfrac{x^{n}}{\sqrt{1-x}}\,\mathrm{d}x$, $n\in\mathbb{N}$.
  \end{enumerate}

  \item 讨论下列瑕积分的敛散性:
  \begin{enumerate}[label=(\arabic*)]
    \item $\int_{0}^{2}\dfrac{\mathrm{d}x}{(x-1)^{2}}$;
    \item $\int_{0}^{\pi}\dfrac{\sin x}{x^{\frac{3}{2}}}\,\mathrm{d}x$;
    \item $\int_{0}^{1}\dfrac{\mathrm{d}x}{\sqrt{x}\ln x}$;
    \item $\int_{0}^{1}\dfrac{\ln x}{1-x}\,\mathrm{d}x$;
    \item $\int_{0}^{1}\dfrac{\arctan x}{1-x^{3}}\,\mathrm{d}x$;
    \item $\int_{0}^{\frac{\pi}{2}}\dfrac{1-\cos x}{x^{m}}\,\mathrm{d}x$;
    \item $\int_{0}^{1}\dfrac{1}{x^{\alpha}}\sin\dfrac{1}{x}\,\mathrm{d}x$;
    \item $\int_{0}^{\frac{\pi}{2}}\mathrm{e}^{-x}\ln x\,\mathrm{d}x$;
    \item $\int_{0}^{1}\dfrac{\mathrm{d}x}{\sqrt{1-x^{4}}}$;
    \item $\int_{0}^{1}\dfrac{\mathrm{d}x}{\ln x}$;
    \item $\int_{0}^{1}\ln x\,\mathrm{d}x$;
    \item $\int_{0}^{1}\dfrac{\mathrm{d}x}{x(\ln x)^{p}}$.
  \end{enumerate}

  \item 判断下列广义积分的敛散性:
  \begin{enumerate}[label=(\arabic*)]
    \item $\int_{0}^{+\infty}\dfrac{|\sin x|}{x^{\frac{4}{3}}}\,\mathrm{d}x$;
    \item $\int_{0}^{+\infty}\dfrac{\mathrm{d}x}{\mathrm{e}^{x}\sqrt{x}}$.
  \end{enumerate}

  \item (1) 设函数 $f$ 在 $[0,+\infty)$ 上连续且 $f\geqslant0$, $\int_{0}^{+\infty}f(x)\,\mathrm{d}x=0$. 证明: $f=0$.

  (2) 设函数 $f$ 在任何 $[0,A]\subset[0,+\infty)$ 上 Riemann 可积且 $f\geqslant0$, $\int_{0}^{+\infty}f(x)\,\mathrm{d}x=0$. 证明: $f\xrightarrow{\text{a.e.}}0$, $x\in[0,+\infty)$.

  \item 设函数 $f$ 在 $[0,+\infty)$ 上一致连续, 且 $f\geqslant0$, $\int_{0}^{+\infty}f(x)\,\mathrm{d}x$ 收敛. 证明: $\lim\limits_{x\to+\infty}f(x)=0$.

  (此题比例~\ref{ex:6.5.15} 多了条件 $f\geqslant0$, 因此要求给出比例~\ref{ex:6.5.15} 更简单的论述.)

  \item 设 $\int_{0}^{+\infty}f(x)\,\mathrm{d}x$ 收敛, 且 $\left|f'(x)\right|<M$(常数), $x>0$. 证明: $\lim\limits_{x\to+\infty}f(x)=0$.

  \item 设 $f$ 与 $g$ 为定义在 $[a,+\infty)$ 上的函数, 对 $\forall u>a$, 它们在 $[a,u]$ 上都 Riemann 可积. 证明: 如果 $\int_{a}^{+\infty}f^{2}(x)\,\mathrm{d}x$ 与 $\int_{a}^{+\infty}g^{2}(x)\,\mathrm{d}x$ 收敛, 则 $\int_{a}^{+\infty}f(x)g(x)\,\mathrm{d}x$ 与 $\int_{a}^{+\infty}[f(x)+g(x)]^{2}\,\mathrm{d}x$ 也都收敛.

  \item 设 $f$, $g$, $h$ 为定义在 $[a,+\infty)$ 上的三个连续函数, 且 $h(x)\leqslant f(x)\leqslant g(x)$. 证明:
  \begin{enumerate}[label=(\arabic*)]
    \item 若 $\int_{a}^{+\infty}h(x)\,\mathrm{d}x$ 与 $\int_{a}^{+\infty}g(x)\,\mathrm{d}x$ 都收敛, 则 $\int_{a}^{+\infty}f(x)\,\mathrm{d}x$ 也收敛;
    \item 若 $\int_{a}^{+\infty}h(x)\,\mathrm{d}x=\int_{a}^{+\infty}g(x)\,\mathrm{d}x=A$, 则 $\int_{a}^{+\infty}f(x)\,\mathrm{d}x=A$.
  \end{enumerate}

  \item 证明:
  \begin{enumerate}[label=(\arabic*)]
    \item $\int_{0}^{+\infty}\dfrac{1-\mathrm{e}^{-x^{2}}}{x^{2}}\,\mathrm{d}x=\sqrt{\pi}$;
    \item $\int_{0}^{1}\left[\int_{x}^{\sqrt{x}}\dfrac{\sin t}{t}\,\mathrm{d}t\right]\mathrm{d}x=1-\sin1$;
    \item $\int_{\frac{\pi}{4}}^{\frac{\pi}{2}}\dfrac{1+\sin x}{1+\cos x}\mathrm{e}^{x}\,\mathrm{d}x=\mathrm{e}^{\frac{\pi}{2}}-(\sqrt{2}-1)\mathrm{e}^{\frac{\pi}{4}}$;
    \item $\int_{0}^{1}\left(\ln\dfrac{1}{x}\right)^{\frac{1}{2}}\mathrm{d}x=\dfrac{\sqrt{\pi}}{2}$. (提示: 作变换 $u=\left(\ln\dfrac{1}{x}\right)^{\frac{1}{2}}$.)
  \end{enumerate}

\end{exblock}

\begin{exblock}{思考题 6.5}

  \item 计算下列两个积分的比值:
  \[
  \int_{0}^{1}\frac{\mathrm{d}t}{\sqrt{1-t^{4}}},\quad\int_{0}^{1}\frac{\mathrm{d}t}{\sqrt{1+t^{4}}}.
\]

  \item 令 $\varphi(x)=-\int_{0}^{x}\ln\cos t\,\mathrm{d}t$, $|x|\leqslant\dfrac{\pi}{2}$. 证明:
  \[
  \varphi(x)=-x\ln2+2\varphi\left(\frac{\pi}{4}+\frac{x}{2}\right)-2\varphi\left(\frac{\pi}{4}-\frac{x}{2}\right).
\]
  并计算 $\varphi\left(\dfrac{\pi}{2}\right)$. 再计算下列积分:
  \begin{enumerate}[label=(\arabic*)]
    \item $\int_{0}^{\frac{\pi}{2}}\ln\sin x\,\mathrm{d}x$;
    \item $\int_{0}^{\pi}x\ln\sin x\,\mathrm{d}x$;
    \item $\int_{0}^{\frac{\pi}{2}}x\cot x\,\mathrm{d}x$;
    \item $\int_{0}^{1}\dfrac{\ln x}{\sqrt{1-x^{2}}}\,\mathrm{d}x$;
    \item $\int_{0}^{1}\dfrac{\arcsin x}{x}\,\mathrm{d}x$.
  \end{enumerate}

  \item 分别作变换 $x=\tan t$ 与 $x=\dfrac{1}{u}$, 计算无穷积分
  \[
  \int_{0}^{+\infty}\frac{\mathrm{d}x}{(1+x^{2})(1+x^{\alpha})},\quad\alpha\in\mathbb{R}.
\]

  \item 应用 $\int_{0}^{\frac{\pi}{2}}\ln\sin x\,\mathrm{d}x=\int_{0}^{\frac{\pi}{2}}\ln\cos x\,\mathrm{d}x=-\dfrac{\pi}{2}\ln2$, 证明:
  \begin{enumerate}[label=(\arabic*)]
    \item $\int_{0}^{\pi}\theta\ln\sin\theta\,\mathrm{d}\theta=-\dfrac{\pi^{2}}{2}\ln2$;
    \item $\int_{0}^{\pi}\dfrac{\theta\sin\theta}{1-\cos\theta}\,\mathrm{d}\theta=2\pi\ln2$.
  \end{enumerate}

  \item 举例说明: $\int_{a}^{+\infty}f(x)\,\mathrm{d}x$ 收敛时, $\int_{a}^{+\infty}f^{2}(x)\,\mathrm{d}x$ 不一定收敛; $\int_{a}^{+\infty}f(x)\,\mathrm{d}x$ 绝对收敛时, $\int_{a}^{+\infty}f^{2}(x)\,\mathrm{d}x$ 也不一定收敛.

  \item 证明: 若 $\int_{a}^{+\infty}f(x)\,\mathrm{d}x$ 绝对收敛, 且 $\lim\limits_{x\to+\infty}f(x)=0$, 则 $\int_{a}^{+\infty}f^{2}(x)\,\mathrm{d}x$ 必定收敛.

  \item 证明: $f$ 为 $[a,+\infty)$ 上的单调函数, 且 $\int_{a}^{+\infty}f(x)\,\mathrm{d}x$ 收敛, 则 $\lim\limits_{x\to+\infty}f(x)=0$, 且 $f(x)=o\left(\dfrac{1}{x}\right)$, $x\to+\infty$.

  \item 举例说明: 瑕积分 $\int_{a}^{b}f(x)\,\mathrm{d}x$ 收敛时, $\int_{a}^{b}f^{2}(x)\,\mathrm{d}x$ 不一定收敛; $\int_{a}^{b}f(x)\,\mathrm{d}x$ 绝对收敛时, $\int_{a}^{b}f^{2}(x)\,\mathrm{d}x$ 也不一定收敛.

  \item 证明: 若 $f$ 在 $[a,+\infty)$ 上可导, 且 $\int_{a}^{+\infty}f(x)\,\mathrm{d}x$ 与 $\int_{a}^{+\infty}f'(x)\,\mathrm{d}x$ 都收敛, 则 $\lim\limits_{x\to+\infty}f(x)=0$.

\end{exblock}

% ====================================================================
%  徐森林、薛春华《数学分析》（第一册）
%  \S6.6  Riemann 积分与广义积分的应用
%  源：原书印刷页 427--448（PDF 物理页 433--454）
%  底本：PaddleOCR-VL-1.6，见 chapters/_ch06_6_raw.md
%  约定：全西文标点；行间公式用一对方括号；多行公式用 align 或 aligned；
%        跨册引用一律用 extref（写作「定理 x.y.z」原文）；数学模式内的中文必须包text
% ====================================================================

\section{Riemann 积分与广义积分的应用}\label{sec:6.6}

Riemann 积分与广义积分来源于几何问题与物理问题. 由此我们建立了一整套的积分理论与计算积分的方法. 这一节主要研究一些具体的实例.

\subsection*{1. 平面图形的面积}

(1) 由连续曲线 $y=f(x)\geqslant0$, 直线 $x=a$, $x=b$ ($a<b$) 与 $x$ 轴所围的曲边梯形面积为

\[
S=\int_{a}^{b}f(x)\,\mathrm{d}x.
\]

如果 $y=f(x)$ 为 $[a,b]$ 上的一般连续曲线, 则所围的面积为

\[
S=\int_{a}^{b}|f(x)|\,\mathrm{d}x.
\]

而

\[
\int_{a}^{b}f(x)\,\mathrm{d}x
\]

为所围面积的代数和, 其中 $x$ 轴上方图形的面积为正, 下方图形的面积为负.

\begin{figure}[htbp]
\centering
\begin{tikzpicture}[x=1.0cm,y=0.82cm,>=Stealth]
  \draw[->] (-0.4,0) -- (6.5,0) node[below right] {$x$};
  \draw[->] (0,-1.7) -- (0,2.6) node[left] {$y$};
  \node[below left] at (0,0) {$O$};
  \fill[pattern=north east lines]
    (0.9,1.55) .. controls (1.8,2.15) and (2.7,2.10) .. (3.5,1.80)
    .. controls (4.4,1.40) and (5.3,0.80) .. (5.9,0.45)
    -- (5.9,-0.90) .. controls (5.2,-1.25) and (4.3,-1.05) .. (3.4,-0.30)
    .. controls (2.5,0.40) and (1.7,1.00) .. (0.9,1.55) -- cycle;
  \draw[very thick] (0.9,1.55) .. controls (1.8,2.15) and (2.7,2.10) .. (3.5,1.80)
    .. controls (4.4,1.40) and (5.3,0.80) .. (5.9,0.45);
  \draw[very thick] (0.9,1.55) .. controls (1.7,1.00) and (2.5,0.40) .. (3.4,-0.30)
    .. controls (4.3,-1.05) and (5.2,-1.25) .. (5.9,-0.90);
  \node[above] at (4.5,1.32) {$y=f_{2}(x)$};
  \node[below] at (2.05,-0.72) {$y=f_{1}(x)$};
  \node[below] at (0.9,-0.06) {$a$};
  \node[below] at (5.9,-0.06) {$b$};
\end{tikzpicture}
\caption{}\label{fig:6.6.1}
\end{figure}

更一般地, 由上、下两条连续曲线 $y=f_{2}(x)$ 与 $y=f_{1}(x)$ 以及两条直线 $x=a$ 与 $x=b$ ($a<b$) 所围图形的面积 (图~\ref{fig:6.6.1}) 计算公式为

\[
S=\int_{a}^{b}|f_{2}(x)-f_{1}(x)|\,\mathrm{d}x.
\]

(2) 如果曲线 $C$ 由极坐标方程

\[
r=r(\theta),\quad\theta\in[\alpha,\beta]
\]

给出, 其中 $r(\theta)$ 在 $[\alpha,\beta]$ 上连续, $\beta-\alpha\leqslant2\pi$. 由曲线 $C$ 与两条射线 $\theta=\alpha$, $\theta=\beta$ 所围图形 (扇形) 的面积 (图~\ref{fig:6.6.2}) 计算公式为

\[
S=\lim_{\|T\|\to0}\sum_{i=1}^{n}\frac{1}{2}r^{2}(\xi_{i})\Delta\theta_{i}=\frac{1}{2}\int_{\alpha}^{\beta}r^{2}(\theta)\,\mathrm{d}\theta.
\]

其中 $T:\alpha=\theta_{0}<\theta_{1}<\cdots<\theta_{n}=\beta$ 为 $[\alpha,\beta]$ 的一个分割, 而射线 $\theta=\theta_{i}$ ($i=1,\cdots,n-1$) 将扇形分成 $n$ 个小扇形. 第 $i$ 个小扇形面积近似为 $\dfrac{1}{2}r^{2}(\xi_{i})\Delta\theta_{i}$, $\xi_{i}\in[\theta_{i-1},\theta_{i}]$.

\begin{figure}[htbp]
\centering
\begin{tikzpicture}[x=1.0cm,y=1.0cm,>=Stealth]
  \draw[->] (-0.3,0) -- (6.4,0) node[below right] {$x$};
  \node[below left] at (0,0) {$O$};
  \coordinate (O) at (0,0);
  \coordinate (Lo) at (10:5.6);
  \coordinate (Hi) at (63:5.0);
  \draw[thick] (O) -- (Lo);
  \draw[thick] (O) -- (Hi);
  \draw[very thick] (Lo) .. controls (2.6,1.35) and (2.0,3.2) .. (Hi);
  \fill[pattern=north east lines]
    (O) -- (24:5.6) .. controls (2.9,1.85) and (2.6,2.4) .. (30:4.65)
    -- (O) -- cycle;
  \draw[thin] (O) -- (24:4.2);
  \draw (1.45,0) arc (0:20:1.45);
  \node[right] at (1.5,0.18) {$\alpha$};
  \node[above] at (Hi) {$\alpha=\beta$};
  \node[right] at (Lo) {$\theta=\alpha$};
  \node[right] at (30:4.72) {$\theta=\theta_{i-1}$};
  \node[right] at (27:4.35) {$\theta=\xi_{i}$};
  \node[right] at (24:4.0) {$\theta=\theta_{i}$};
\end{tikzpicture}
\caption{}\label{fig:6.6.2}
\end{figure}

(3) 如果曲线 $C$ 由参数方程

\[
\begin{cases}
x=x(t),\\
y=y(t),
\end{cases}\quad t\in[\alpha,\beta]
\]

给出, 其中 $x=x(t)$, $y=y(t)$ 在 $[\alpha,\beta]$ 上连续可导. 记 $a=x(\alpha)$, $b=x(\beta)$ ($a<b$ 或 $b<a$), 则由曲线 $C$ 与直线 $x=a$, $x=b$ 及 $x$ 轴所围图形的面积公式为 (请证之)

\[
S=\int_{\alpha}^{\beta}|y(t)x'(t)|\,\mathrm{d}t
\]

或

\[
S=\int_{\alpha}^{\beta}|x(t)y'(t)|\,\mathrm{d}t.
\]

如果由参数方程

\[
\begin{cases}
x=x(t),\\
y=y(t),
\end{cases}\quad\alpha\leqslant t\leqslant\beta
\]

所表示的曲线是封闭的, 且 $x(t)$, $y(t)$ 在 $[\alpha,\beta]$ 上连续可导, 即有

\[
(x(\alpha),y(\alpha))=(x(\beta),y(\beta)),
\]

且在 $(\alpha,\beta)$ 内曲线自身不再相交, 则由曲线自身所围图形的面积计算公式为

\[
S=\left|\int_{\alpha}^{\beta}y(t)x'(t)\,\mathrm{d}t\right|
\]

或

\[
S=\left|\int_{\alpha}^{\beta}x(t)y'(t)\,\mathrm{d}t\right|.
\]

\begin{example}\label{ex:6.6.1}
求由抛物线 $y^{2}=x$ 与直线 $x-2y-3=0$ 所围图形的面积.
\end{example}

\begin{solution}[解法 1]
易见抛物线 $y^{2}=x$ 与直线 $x-2y-3=0$ 的交点为 $(1,-1)$ 与 $(9,3)$. 于是, 所围图形的面积为 (图~\ref{fig:6.6.3})
\[
\begin{aligned}
S&=\int_{-1}^{3}\left[(2y+3)-y^{2}\right]\mathrm{d}y=\left.\left(y^{2}+3y-\frac{y^{3}}{3}\right)\right|_{-1}^{3}\\
&=(9+9-9)-\left(1-3+\frac{1}{3}\right)=\frac{32}{3}.
\end{aligned}
\]
\end{solution}

\begin{solution}[解法 2]
$S=S_{1}+S_{2}$
\[
\begin{aligned}
&=\int_{0}^{1}\left[\sqrt{x}-(-\sqrt{x})\right]\mathrm{d}x+\int_{1}^{9}\left(\sqrt{x}-\frac{x-3}{2}\right)\mathrm{d}x\\
&=2\int_{0}^{1}\sqrt{x}\,\mathrm{d}x+\int_{1}^{9}\left(\sqrt{x}-\frac{x-3}{2}\right)\mathrm{d}x\\
&=\left.\frac{4}{3}x^{\frac{3}{2}}\right|_{0}^{1}+\left.\left(\frac{2}{3}x^{\frac{3}{2}}-\frac{x^{2}}{4}+\frac{3}{2}x\right)\right|_{1}^{9}=\frac{32}{3}.
\end{aligned}
\]
\end{solution}

\begin{figure}[htbp]
\centering
\begin{tikzpicture}[x=0.62cm,y=0.72cm,>=Stealth]
  \draw[->] (-1.6,0) -- (10.6,0) node[below right] {$x$};
  \draw[->] (0,-2.5) -- (0,3.9) node[left] {$y$};
  \node[below left] at (0,0) {$O$};
  \fill[pattern=north east lines]
    (1,-1) -- (9,3) .. controls (7.2,2.72) and (3.0,1.62) .. (1,1) -- cycle;
  \fill[pattern=north east lines]
    (0,0) .. controls (0.5,0.7) and (0.8,0.9) .. (1,1) -- (1,-1)
    .. controls (0.7,-0.85) and (0.4,-0.6) .. (0,0) -- cycle;
  \draw[very thick] (0,0) .. controls (0.5,0.7) and (0.8,0.9) .. (1,1)
    .. controls (2.0,1.42) and (4.0,2.0) .. (9,3);
  \draw[very thick] (0,0) .. controls (0.4,-0.6) and (0.7,-0.85) .. (1,-1)
    .. controls (2.0,-1.42) and (4.0,-2.0) .. (9,-3);
  \draw[very thick] (-0.5,-1.75) -- (10.0,3.25);
  \draw[thin] (1,-1) -- (1,1);
  \draw[dotted] (9,3) -- (0,3);
  \draw[dotted] (9,3) -- (9,0);
  \draw[dotted] (1,-1) -- (0,-1);
  \node[above right] at (9,3) {$Q$};
  \node[below] at (1,-1.25) {$P$};
  \node at (0.48,0.62) {$S_{1}$};
  \node at (3.4,1.15) {$S_{2}$};
  \node[below] at (1,-0.06) {1};
  \node[below] at (3,-0.06) {3};
  \node[below] at (6,-0.06) {6};
  \node[below] at (9,-0.06) {9};
  \node[left] at (0,3) {3};
  \node[left] at (0,-1) {-1};
\end{tikzpicture}
\caption{}\label{fig:6.6.3}
\end{figure}

\begin{example}\label{ex:6.6.2}
求由摆线
\[
\begin{cases}
x=a(t-\sin t),\\
y=a(1-\cos t),
\end{cases}\quad(a>0)
\]
的一拱与 $x$ 轴所围图形的面积 $S$ (图~\ref{fig:6.6.4}).
\end{example}

\begin{solution}
\[
\begin{aligned}
S&=\int_{0}^{2\pi a}y(x)\,\mathrm{d}x
\xrightarrow[y=a(1-\cos t)]{x=a(t-\sin t)}\int_{0}^{2\pi}a(1-\cos t)\cdot a(1-\cos t)\,\mathrm{d}t\\
&=a^{2}\int_{0}^{2\pi}\left(1-2\cos t+\frac{1+\cos2t}{2}\right)\mathrm{d}t\\
&=a^{2}\left.\left(\frac{3}{2}t-2\sin t+\frac{\sin2t}{4}\right)\right|_{0}^{2\pi}\\
&=a^{2}\left(\frac{3}{2}\cdot2\pi-0-0\right)=3\pi a^{2}.
\end{aligned}
\]
\end{solution}

\begin{figure}[htbp]
\centering
\begin{tikzpicture}[x=1.0cm,y=1.0cm,>=Stealth]
  \draw[->] (-0.4,0) -- (7.6,0) node[below right] {$x$};
  \draw[->] (0,-0.4) -- (0,2.9) node[left] {$y$};
  \node[below left] at (0,0) {$O$};
  \fill[pattern=north east lines]
    plot[smooth] coordinates {(0.0000,0.0000) (0.0004,0.0086) (0.0030,0.0341) (0.0100,0.0761) (0.0236,0.1340) (0.0457,0.2066) (0.0783,0.2929) (0.1229,0.3912) (0.1812,0.5000) (0.2542,0.6173) (0.3431,0.7412) (0.4485,0.8695) (0.5708,1.0000) (0.7103,1.1305) (0.8667,1.2588) (1.0396,1.3827) (1.2284,1.5000) (1.4319,1.6088) (1.6491,1.7071) (1.8783,1.7934) (2.1180,1.8660) (2.3662,1.9239) (2.6210,1.9659) (2.8802,1.9914) (3.1416,2.0000) (3.4030,1.9914) (3.6622,1.9659) (3.9170,1.9239) (4.1652,1.8660) (4.4049,1.7934) (4.6341,1.7071) (4.8512,1.6088) (5.0548,1.5000) (5.2436,1.3827) (5.4165,1.2588) (5.5729,1.1305) (5.7124,1.0000) (5.8347,0.8695) (5.9401,0.7412) (6.0290,0.6173) (6.1020,0.5000) (6.1602,0.3912) (6.2049,0.2929) (6.2374,0.2066) (6.2596,0.1340) (6.2732,0.0761) (6.2802,0.0341) (6.2828,0.0086) (6.2832,0.0000)}
    -- (6.2832,0) -- (0,0) -- cycle;
  \draw[very thick] plot[smooth] coordinates {(0.0000,0.0000) (0.0004,0.0086) (0.0030,0.0341) (0.0100,0.0761) (0.0236,0.1340) (0.0457,0.2066) (0.0783,0.2929) (0.1229,0.3912) (0.1812,0.5000) (0.2542,0.6173) (0.3431,0.7412) (0.4485,0.8695) (0.5708,1.0000) (0.7103,1.1305) (0.8667,1.2588) (1.0396,1.3827) (1.2284,1.5000) (1.4319,1.6088) (1.6491,1.7071) (1.8783,1.7934) (2.1180,1.8660) (2.3662,1.9239) (2.6210,1.9659) (2.8802,1.9914) (3.1416,2.0000) (3.4030,1.9914) (3.6622,1.9659) (3.9170,1.9239) (4.1652,1.8660) (4.4049,1.7934) (4.6341,1.7071) (4.8512,1.6088) (5.0548,1.5000) (5.2436,1.3827) (5.4165,1.2588) (5.5729,1.1305) (5.7124,1.0000) (5.8347,0.8695) (5.9401,0.7412) (6.0290,0.6173) (6.1020,0.5000) (6.1602,0.3912) (6.2049,0.2929) (6.2374,0.2066) (6.2596,0.1340) (6.2732,0.0761) (6.2802,0.0341) (6.2828,0.0086) (6.2832,0.0000)};
  \draw[dashed] (0,1) circle (1);
  \node at (3.6,0.72) {$S$};
  \node[left] at (0,1) {$a$};
  \node[left] at (0,2) {$2a$};
  \node[below] at (3.1416,-0.06) {$\pi a$};
  \node[below] at (6.2832,-0.06) {$2\pi a$};
\end{tikzpicture}
\caption{}\label{fig:6.6.4}
\end{figure}

\begin{example}\label{ex:6.6.3}
求椭圆 $\dfrac{x^{2}}{a^{2}}+\dfrac{y^{2}}{b^{2}}=1$ 所围图形的面积 $S$ ($a>0$, $b>0$).
\end{example}

\begin{solution}[解法 1]
椭圆的参数方程为
\[
\begin{cases}
x=a\cos t,\\
y=b\sin t,
\end{cases}\quad t\in[0,2\pi].
\]
于是,
\[
\begin{aligned}
S&=\left|\int_{0}^{2\pi}b\sin t\cdot(a\cos t)'\mathrm{d}t\right|=ab\int_{0}^{2\pi}\sin^{2}t\,\mathrm{d}t\\
&=ab\int_{0}^{2\pi}\frac{1-\cos2t}{2}\,\mathrm{d}t=\left.\frac{ab}{2}\left(t-\frac{\sin2t}{2}\right)\right|_{0}^{2\pi}=\frac{ab}{2}(2\pi-0)=\pi ab.
\end{aligned}
\]
\end{solution}

\begin{solution}[解法 2]
\[
\begin{aligned}
S&\xlongequal{\text{对称性}}4\int_{0}^{a}b\sqrt{1-\frac{x^{2}}{a^{2}}}\,\mathrm{d}x
\xlongequal{x=a\sin t}4b\int_{0}^{\frac{\pi}{2}}\sqrt{1-\sin^{2}t}\,(a\cos t)\,\mathrm{d}t\\
&=4ab\int_{0}^{\frac{\pi}{2}}\cos^{2}t\,\mathrm{d}t=4ab\cdot\frac{1}{2}\cdot\frac{\pi}{2}=\pi ab.
\end{aligned}
\]
\end{solution}

\begin{solution}[解法 3]
椭圆 $\dfrac{x^{2}}{a^{2}}+\dfrac{y^{2}}{b^{2}}=1$ 的极坐标方程为 $\dfrac{r^{2}\cos^{2}\theta}{a^{2}}+\dfrac{r^{2}\sin^{2}\theta}{b^{2}}=1$, 即
\[
r^{2}=r^{2}(\theta)=\frac{1}{\dfrac{\cos^{2}\theta}{a^{2}}+\dfrac{\sin^{2}\theta}{b^{2}}}.
\]
于是,
\[
\begin{aligned}
S&=\frac{1}{2}\int_{0}^{2\pi}r^{2}(\theta)\,\mathrm{d}\theta=\frac{1}{2}\int_{0}^{2\pi}\frac{\mathrm{d}\theta}{\dfrac{\cos^{2}\theta}{a^{2}}+\dfrac{\sin^{2}\theta}{b^{2}}}\\
&=2\int_{0}^{\frac{\pi}{2}}\frac{\mathrm{d}\tan\theta}{\dfrac{1}{a^{2}}+\dfrac{1}{b^{2}}\tan^{2}\theta}
\xlongequal{u=\tan\theta}2a^{2}b^{2}\int_{0}^{+\infty}\frac{\mathrm{d}u}{b^{2}+a^{2}u^{2}}\\
&=\left.2ab\arctan\frac{au}{b}\right|_{0}^{+\infty}=2ab\cdot\frac{\pi}{2}=\pi ab.
\end{aligned}
\]
\end{solution}

\begin{example}\label{ex:6.6.4}
求双纽线 $(x^{2}+y^{2})^{2}=a^{2}(x^{2}-y^{2})$ 所围图形的面积 $S$.
\end{example}

\begin{solution}
双纽线的极坐标方程为
\[
r^{4}=a^{2}r^{2}\left(\cos^{2}\theta-\sin^{2}\theta\right)=a^{2}r^{2}\cos2\theta,\quad r^{2}=a^{2}\cos2\theta.
\]
因为 $r^{2}\geqslant0$, 故 $\cos2\theta\geqslant0$, 且 $\theta$ 的取值范围为 $\left[-\dfrac{\pi}{4},\dfrac{\pi}{4}\right]$ 与 $\left[\dfrac{3\pi}{4},\dfrac{5\pi}{4}\right]$. 再由图形的对称性 (图~\ref{fig:6.6.5}) 立即有
\[
S=4\cdot\frac{1}{2}\int_{0}^{\frac{\pi}{4}}r^{2}(\theta)\,\mathrm{d}\theta=2\int_{0}^{\frac{\pi}{4}}a^{2}\cos2\theta\,\mathrm{d}\theta=\left.a^{2}\sin2\theta\right|_{0}^{\frac{\pi}{4}}=a^{2}.
\]
\end{solution}

\begin{figure}[htbp]
\centering
\begin{tikzpicture}[x=1.55cm,y=1.35cm,>=Stealth]
  \draw[->] (-2.1,0) -- (2.3,0) node[below right] {$x$};
  \draw[->] (0,-1.25) -- (0,1.35) node[left] {$y$};
  \node[below left] at (0,0) {$O$};
  \fill[pattern=north east lines]
    plot[smooth] coordinates {(1.0000,0.0000) (0.9936,0.0651) (0.9744,0.1283) (0.9427,0.1875) (0.8989,0.2409) (0.8434,0.2863) (0.7769,0.3218) (0.6998,0.3451) (0.6124,0.3536) (0.5144,0.3437) (0.4036,0.3097) (0.2716,0.2382) (0.0000,0.0000)}
    -- (0,0) -- cycle;
  \draw[very thick] plot[smooth] coordinates {(0.0000,-0.0000) (0.2716,-0.2382) (0.4036,-0.3097) (0.5144,-0.3437) (0.6124,-0.3536) (0.6998,-0.3451) (0.7769,-0.3218) (0.8434,-0.2863) (0.8989,-0.2409) (0.9427,-0.1875) (0.9744,-0.1283) (0.9936,-0.0651) (1.0000,0.0000) (0.9936,0.0651) (0.9744,0.1283) (0.9427,0.1875) (0.8989,0.2409) (0.8434,0.2863) (0.7769,0.3218) (0.6998,0.3451) (0.6124,0.3536) (0.5144,0.3437) (0.4036,0.3097) (0.2716,0.2382) (0.0000,0.0000)};
  \draw[very thick] plot[smooth] coordinates {(-0.0000,0.0000) (-0.2716,0.2382) (-0.4036,0.3097) (-0.5144,0.3437) (-0.6124,0.3536) (-0.6998,0.3451) (-0.7769,0.3218) (-0.8434,0.2863) (-0.8989,0.2409) (-0.9427,0.1875) (-0.9744,0.1283) (-0.9936,0.0651) (-1.0000,-0.0000) (-0.9936,-0.0651) (-0.9744,-0.1283) (-0.9427,-0.1875) (-0.8989,-0.2409) (-0.8434,-0.2863) (-0.7769,-0.3218) (-0.6998,-0.3451) (-0.6124,-0.3536) (-0.5144,-0.3437) (-0.4036,-0.3097) (-0.2716,-0.2382) (-0.0000,-0.0000)};
  \draw[thin] (0,0) -- (0.7854:1.4);
  \draw[thin] (0,0) -- (-0.7854:1.4);
  \node[right] at (0.7854:1.42) {$\theta=\dfrac{\pi}{4}$};
  \node[right] at (-0.7854:1.42) {$\theta=-\dfrac{\pi}{4}$};
  \node[below] at (1,-0.05) {$a$};
  \node[anchor=west] at (0.75,1.15) {$(x^{2}+y^{2})^{2}=a^{2}(x^{2}-y^{2})$};
\end{tikzpicture}
\caption{}\label{fig:6.6.5}
\end{figure}

\subsection*{2. 平行截面间立体的体积}

设 $\Omega$ 为 $\mathbb{R}^{3}$ 中一个立体, 它夹在垂直于 $x$ 轴的两平面 $x=a$ 与 $x=b$ 之间 ($a<b$), 并称 $\Omega$ 为位于 $[a,b]$ 上的立体. 记 $S(x)$ 为 $x\in[a,b]$ 点处垂直于 $x$ 轴的平面截 $\Omega$ 的截面面积函数, 如果它连续, 则立体 $\Omega$ 的体积公式为

\[
V=\int_{a}^{b}S(x)\,\mathrm{d}x.
\]

事实上, 设 $T:a=x_{0}<x_{1}<\cdots<x_{n}=b$ 为 $[a,b]$ 的一个分割, 过各个分点作垂直于 $x$ 轴的平面 $x=x_{i}$, $i=1,2,\cdots,n$, 它们将 $\Omega$ 切成 $n$ 个薄片 (就如切萝卜片). 设 $S(x)$ 在每个小区间 $\Delta_{i}=[x_{i-1},x_{i}]$ 上的最大值与最小值分别为 $M_{i}$ 与 $m_{i}$, 则每一薄片的体积 $\Delta V_{i}$ 满足

\[
m_{i}\Delta x_{i}\leqslant\Delta V_{i}\leqslant M_{i}\Delta x_{i},
\]

因此,

\[
\sum_{i=1}^{n}m_{i}\Delta x_{i}\leqslant\sum_{i=1}^{n}\Delta V_{i}\leqslant\sum_{i=1}^{n}M_{i}\Delta x_{i}.
\]

从而

\[
V=\lim_{\|T\|\to0}\sum_{i=1}^{n}M_{i}\Delta x_{i}=\lim_{\|T\|\to0}\sum_{i=1}^{n}m_{i}\Delta x_{i}=\int_{a}^{b}S(x)\,\mathrm{d}x.
\]

或者

\[
V=\lim_{\|T\|\to0}\sum_{i=1}^{n}S(\xi_{i})\Delta x_{i}=\int_{a}^{b}S(x)\,\mathrm{d}x,
\]

其中 $\xi_{i}\in[x_{i-1},x_{i}]$, $i=1,2,\cdots,n$ (图~\ref{fig:6.6.6}).

\begin{figure}[htbp]
\centering
\begin{tikzpicture}[x=1.0cm,y=1.0cm,>=Stealth]
  \draw[->] (-0.4,0) -- (8.6,0) node[below right] {$x$};
  \node[below left] at (0,0) {$O$};
  \draw[thick]
    (1.4,0.15) .. controls (2.2,1.35) and (4.4,1.45) .. (5.6,0.95)
    .. controls (6.5,0.55) and (7.0,0.75) .. (7.6,1.05);
  \draw[thick]
    (1.4,-0.15) .. controls (2.2,-1.35) and (4.4,-1.45) .. (5.6,-0.95)
    .. controls (6.5,-0.55) and (7.0,-0.75) .. (7.6,-1.05);
  \draw[thick] (1.4,0.15) ellipse (0.34 and 0.16);
  \draw[thick] (7.6,1.05) .. controls (7.9,0.6) and (7.9,-0.6) .. (7.6,-1.05);
  \draw[thick] (7.6,1.05) .. controls (8.05,0.7) and (8.05,-0.7) .. (7.6,-1.05);
  \draw[thick,dashed] (5.6,0.95) -- (5.6,-0.95);
  \draw[thick] (5.6,0.95) ellipse (0.30 and 0.72);
  \fill[pattern=north east lines] (5.15,0.95) .. controls (5.15,0.0) and (5.15,0.0) .. (5.15,-0.95)
    -- (5.45,-0.95) .. controls (5.45,0.0) and (5.45,0.0) .. (5.45,0.95) -- cycle;
  \draw[thin] (3.2,1.62) -- (3.2,0.15);
  \draw[thin] (0,1.62) -- (3.2,1.62);
  \draw[thin] (1.4,1.62) -- (1.4,0.15);
  \node[above] at (2.3,1.66) {$S(x)$};
  \node at (4.6,1.15) {$\Omega$};
  \node[below] at (1.4,-0.06) {$a$};
  \node[below] at (3.2,-0.06) {$x$};
  \node[below] at (7.6,-0.06) {$b$};
\end{tikzpicture}
\caption{}\label{fig:6.6.6}
\end{figure}

\begin{figure}[htbp]
\centering
\begin{tikzpicture}[x=1.0cm,y=1.0cm,>=Stealth]
  \draw[->] (0,0) -- (0,4.1) node[left] {$z$};
  \draw[->] (0,0) -- (4.9,0) node[below right] {$y$};
  \draw[->] (0,0) -- (-2.1,-1.55) node[below left] {$x$};
  \node[below right] at (0,0) {$O$};
  \draw[thick] (-1.0,-0.72) -- (-1.0,2.3) -- (2.5,2.3) -- (2.5,-0.72) -- cycle;
  \draw[thick,dashed] (-1.0,2.3) -- (0.35,3.0) -- (3.85,3.0) -- (2.5,2.3);
  \draw[thick] (-1.0,-0.72) -- (0.35,0.02) -- (3.85,0.02) -- (2.5,-0.72);
  \draw[thick] (0.35,0.02) -- (0.35,3.0);
  \draw[thick] (3.85,0.02) -- (3.85,3.0);
  \draw[thick] (0.0,0.02) .. controls (0.55,1.0) and (1.45,1.95) .. (1.35,3.0);
  \draw[thick] (0.0,0.02) .. controls (-0.55,1.0) and (-1.45,1.95) .. (-0.55,3.0);
  \draw[thick,dotted] (1.35,3.0) .. controls (1.7,2.2) and (1.7,0.9) .. (1.35,0.02);
  \fill[pattern=north east lines] (1.15,1.62) rectangle (1.55,2.55);
  \draw (1.15,1.62) rectangle (1.55,2.55);
  \node[below] at (0.35,-0.06) {$a$};
  \node[below right] at (1.35,0.0) {$a$};
\end{tikzpicture}
\caption{}\label{fig:6.6.7}
\end{figure}

\begin{example}\label{ex:6.6.5}
求由两个圆柱面 $x^{2}+y^{2}=a^{2}$, $z^{2}+x^{2}=a^{2}$ 所围立体的体积 $V$.
\end{example}

\begin{solution}
由对称性 (图~\ref{fig:6.6.7}), 有
\[
\begin{aligned}
V&=8\int_{0}^{a}S(x)\,\mathrm{d}x\\
&=8\int_{0}^{a}\sqrt{a^{2}-x^{2}}\cdot\sqrt{a^{2}-x^{2}}\,\mathrm{d}x=8\int_{0}^{a}(a^{2}-x^{2})\,\mathrm{d}x\\
&=\left.8\left(a^{2}x-\frac{x^{3}}{3}\right)\right|_{0}^{a}=\frac{16}{3}a^{3}.
\end{aligned}
\]
\end{solution}

\begin{example}\label{ex:6.6.6}
求由椭球面 $\dfrac{x^{2}}{a^{2}}+\dfrac{y^{2}}{b^{2}}+\dfrac{z^{2}}{c^{2}}=1$ 所围椭球体的体积.
\end{example}

\begin{solution}
以平面 $x=x_{0}$ ($|x_{0}|\leqslant a$) 截椭球面得椭圆 (它在 $yOz$ 平面上的正投影) 为
\[
\frac{y^{2}}{b^{2}\left(1-\dfrac{x_{0}^{2}}{a^{2}}\right)}+\frac{z^{2}}{c^{2}\left(1-\dfrac{x_{0}^{2}}{a^{2}}\right)}=1.
\]
根据例~\ref{ex:6.6.3} 知, 该截面 (椭圆片) 的面积为
\[
S(x)=\pi\left(b\sqrt{1-\frac{x^{2}}{a^{2}}}\right)\left(c\sqrt{1-\frac{x^{2}}{a^{2}}}\right)=\pi bc\left(1-\frac{x^{2}}{a^{2}}\right),\quad x\in[-a,a].
\]
于是, 椭球体体积为
\[
\begin{aligned}
V&=\int_{-a}^{a}S(x)\,\mathrm{d}x=\int_{-a}^{a}\pi bc\left(1-\frac{x^{2}}{a^{2}}\right)\mathrm{d}x\\
&=\left.\pi bc\left(x-\frac{x^{3}}{3a^{2}}\right)\right|_{-a}^{a}=\frac{4}{3}\pi abc.
\end{aligned}
\]
\end{solution}

当 $a=b=c$ 时, 椭球体就为球体, 其体积为 $V=\dfrac{4}{3}\pi a^{3}$.

\subsection*{3. 旋转体的体积}

设 $f$ 为 $[a,b]$ 上的连续函数, $\Omega$ 是由平面图形

\[
\{(x,y)\mid a\leqslant x\leqslant b,\ 0\leqslant|y|\leqslant|f(x)|\}
\]

绕 $x$ 轴旋转一周所得的旋转体. 易知截面 (半径为 $|f(x)|$ 的圆片) 面积

\[
S(x)=\pi[f(x)]^{2},\quad x\in[a,b].
\]

于是, 旋转体 $\Omega$ 的体积公式 (图~\ref{fig:6.6.8}) 为

\[
V=\int_{a}^{b}S(x)\,\mathrm{d}x=\pi\int_{a}^{b}[f(x)]^{2}\,\mathrm{d}x.
\]

\begin{figure}[htbp]
\centering
\begin{tikzpicture}[x=1.0cm,y=1.0cm,>=Stealth]
  \draw[->] (-0.4,0) -- (7.6,0) node[below right] {$x$};
  \draw[->] (0,-1.7) -- (0,2.1) node[left] {$y$};
  \node[below left] at (0,0) {$O$};
  \draw[thick,dashed] (1.3,0) ellipse (0.34 and 1.28);
  \draw[thick] (1.3,0) ellipse (0.34 and 1.28);
  \draw[thick] (5.6,0) ellipse (0.30 and 0.62);
  \draw[thick] (1.3,1.28) .. controls (3.0,1.72) and (4.3,0.95) .. (5.6,0.62);
  \draw[thick] (1.3,-1.28) .. controls (3.0,-1.72) and (4.3,-0.95) .. (5.6,-0.62);
  \draw[thick,dashed] (1.3,-1.28) arc (-90:90:0.34 and 1.28);
  \draw[thick] (1.3,1.28) arc (90:270:0.34 and 1.28);
  \draw[thick,dashed] (5.6,-0.62) arc (-90:90:0.30 and 0.62);
  \draw[thick] (5.6,0.62) arc (90:270:0.30 and 0.62);
  \draw[thick] (3.35,1.55) ellipse (0.26 and 1.05);
  \fill[pattern=north east lines] (3.09,1.55) .. controls (3.09,0.5) and (3.09,0.5) .. (3.09,-0.55)
    -- (3.61,-0.55) .. controls (3.61,0.5) and (3.61,0.5) .. (3.61,1.55) -- cycle;
  \draw[thick] (3.35,0) ellipse (0.26 and 1.05);
  \node[above] at (3.35,1.62) {$S(x)$};
  \node[above] at (4.7,1.20) {$y=f(x)$};
  \node[below] at (1.3,-0.06) {$a$};
  \node[below] at (3.35,-0.06) {$x$};
  \node[below] at (3.98,-0.06) {$x+\mathrm{d}x$};
  \node[below] at (5.6,-0.06) {$b$};
\end{tikzpicture}
\caption{}\label{fig:6.6.8}
\end{figure}

\begin{example}\label{ex:6.6.7}
证明: 高为 $h$, 底半径为 $r$ 的圆锥体的体积 $V=\dfrac{1}{3}\pi r^{2}h$.
\end{example}

\begin{proof}
设正圆锥的高为 $h$, 底圆半径为 $r$ (见图~\ref{fig:6.6.9}). 该正圆锥体可由平面图形 $\{(x,y)\mid0\leqslant x\leqslant h,\allowbreak\ 0\leqslant|y|\leqslant\frac{r}{h}x\}$ 绕 $x$ 轴旋转一周得到, 所以其体积为
\[
V=\pi\int_{0}^{h}\left(\frac{r}{h}x\right)^{2}\mathrm{d}x=\left.\pi\frac{r^{2}}{h^{2}}\frac{x^{3}}{3}\right|_{0}^{h}=\frac{1}{3}\pi r^{2}h.
\]
\end{proof}

进而, 任一高为 $h$, 底半径为 $r$ 的 (正或斜) 的圆锥, 它们在相同高处的截面为半径相同的圆片, 则其体积恒为 $\dfrac{1}{3}\pi r^{2}h$.

\begin{figure}[htbp]
\centering
\begin{tikzpicture}[x=1.0cm,y=1.0cm,>=Stealth]
  \coordinate (O) at (0,0);
  \coordinate (Apex) at (0,0);
  \coordinate (BaseC) at (4.6,0);
  \draw[->] (-0.4,0) -- (5.8,0) node[below left] {$x$};
  \draw[->] (0,-0.4) -- (0,2.05) node[left] {$y$};
  \draw[->] (0,0) -- (-1.5,-1.15) node[below left] {$z$};
  \node[above left] at (0,0) {$O$};
  \draw[thick] (Apex) -- (4.6,1.62);
  \draw[thick] (Apex) -- (4.6,-1.62);
  \draw[thick] (4.6,0) ellipse (0.40 and 1.62);
  \draw[thick] (4.6,0) -- (4.6,-1.62);
  \draw[thick,dashed] (2.6,0) ellipse (0.24 and 0.92);
  \draw[thick,dotted] (4.6,0) -- (1.9,1.55);
  \node[right] at (4.72,1.05) {$r$};
  \node[below] at (2.3,-0.06) {$h$};
\end{tikzpicture}
\caption{}\label{fig:6.6.9}
\end{figure}

\begingroup\xeCJKsetup{CJKglue={\hskip 0pt plus 1pt}}
上述只是一个特殊结果, 更一般的结果是截面面积相等则体积也相等的原理 (图~\ref{fig:6.6.10}): 设 $\Omega_{A}$, $\Omega_{B}$ 为位于同一区间 $[a,b]$ 上的两个立体, 其体积分别为 $V_{A}$, $V_{B}$. 如果在 $[a,b]$ 上它们的截面函数 $S_{A}(x)$ 与 $S_{B}(x)$ 皆连续, 且 $S_{A}(x)=S_{B}(x)$. 则\endgroup

\[
V_{A}=\int_{a}^{b}S_{A}(x)\,\mathrm{d}x=\int_{a}^{b}S_{B}(x)\,\mathrm{d}x=V_{B}.
\]

这个原理早已为我国齐梁时代的数学家祖{\zmbhan 暅} (祖冲之 (429—500) 之子, 生卒年代约在公元 5 世纪末至 6 世纪初) 在计算球的体积时所发现. 17 世纪意大利数学家 Cavalieri (1598—1647) 也提出了类似的原理, 但要比祖{\zmbhan 暅}晚 1100 多年.

\begin{figure}[htbp]
\centering
\begin{tikzpicture}[x=1.0cm,y=1.0cm,>=Stealth]
  \draw[->] (-0.4,0) -- (8.4,0) node[below right] {$x$};
  \draw[->] (0,-3.0) -- (0,3.2) node[left] {$y$};
  \node[below left] at (0,0) {$O$};
  \draw[thick] (1.2,1.35) .. controls (2.4,2.35) and (4.6,2.5) .. (5.7,1.7)
    .. controls (6.6,1.05) and (6.9,1.25) .. (7.5,1.55);
  \draw[thick] (1.2,1.35) .. controls (2.1,0.85) and (3.2,0.72) .. (4.3,0.85)
    .. controls (5.4,0.98) and (6.2,1.05) .. (7.5,1.55);
  \draw[thick] (1.2,1.35) ellipse (0.30 and 0.14);
  \draw[thick] (3.1,0.78) ellipse (0.28 and 0.46);
  \draw[thick,dashed] (3.1,1.24) -- (3.1,0.32);
  \draw[thick] (1.2,-1.35) .. controls (2.4,-2.35) and (4.6,-2.5) .. (5.7,-1.7)
    .. controls (6.6,-1.05) and (6.9,-1.25) .. (7.5,-1.55);
  \draw[thick] (1.2,-1.35) .. controls (2.1,-0.85) and (3.2,-0.72) .. (4.3,-0.85)
    .. controls (5.4,-0.98) and (6.2,-1.05) .. (7.5,-1.55);
  \draw[thick] (1.2,-1.35) ellipse (0.30 and 0.14);
  \draw[thick] (3.1,-0.78) ellipse (0.28 and 0.46);
  \draw[thick,dashed] (3.1,-0.32) -- (3.1,-1.24);
  \draw[dotted] (3.1,1.24) -- (3.1,-0.32);
  \node[above] at (4.4,2.12) {$S_{A}(x)$};
  \node[above] at (4.5,-0.44) {$S_{B}(x)$};
  \node[below] at (1.2,-0.06) {$a$};
  \node[below] at (3.1,-0.06) {$x$};
  \node[below] at (7.5,-0.06) {$b$};
\end{tikzpicture}
\caption{}\label{fig:6.6.10}
\end{figure}

\begin{example}\label{ex:6.6.8}
求由圆 $x^{2}+(y-R)^{2}\leqslant r^{2}$ ($0<r<R$) 绕 $x$ 轴旋转一周所得环状立体 (称为圆环体, 其边界为圆环面 (如救生圈)) 的体积 (图~\ref{fig:6.6.11}).
\end{example}

\begin{figure}[htbp]
\centering
\begin{tikzpicture}[x=1.0cm,y=1.0cm,>=Stealth]
  \draw[->] (-1.6,0) -- (1.6,0) node[below right] {$y$};
  \draw[->] (0,-1.9) -- (0,2.6) node[left] {$x$};
  \node[below left] at (0,0) {$O$};
  \draw[thick] (0,0.85) circle (0.82);
  \draw[thick] (0,0.85) -- (0.82,0.85);
  \node[right] at (0.42,0.95) {$R$};
  \node[above right] at (0.30,1.62) {$x^{2}+(y-R)^{2}\leqslant r^{2}$};
\end{tikzpicture}
\hspace{0.6cm}
\begin{tikzpicture}[x=1.0cm,y=1.0cm,>=Stealth]
  \draw[->] (0,0) -- (3.9,0) node[right] {$z$};
  \draw[->] (0,0) -- (-2.1,-0.85) node[below left] {$y$};
  \draw[->] (0,0) -- (0,2.5) node[above] {$x$};
  \node[below right] at (0,0) {$O$};
  \draw[thick] (0,0.45) ellipse (1.75 and 0.68);
  \draw[thick] (0,-0.45) ellipse (1.75 and 0.68);
  \draw[thick] (-1.75,0.45) arc (180:360:1.75 and 0.68);
  \draw[thick] (1.75,0.45) arc (0:180:1.75 and 0.68);
  \draw[thick,dashed] (-1.75,-0.45) arc (180:360:1.75 and 0.68);
  \draw[thick] (1.75,-0.45) arc (0:180:1.75 and 0.68);
  \draw[thick] (1.75,0.45) arc (90:-90:0.62 and 0.45);
  \draw[thick] (1.75,-0.45) arc (-90:90:0.62 and 0.45);
\end{tikzpicture}
\caption{}\label{fig:6.6.11}
\end{figure}

\begin{solution}
如图~\ref{fig:6.6.11} 所示, 圆 $x^{2}+(y-R)^{2}=r^{2}$ 的上、下半圆分别为
\[
y=f_{2}(x)=R+\sqrt{r^{2}-x^{2}},
\]
\[
y=f_{1}(x)=R-\sqrt{r^{2}-x^{2}},\quad|x|\leqslant r.
\]
故圆环体的截面面积函数为
\[
S(x)=\pi\left\{[f_{2}(x)]^{2}-[f_{1}(x)]^{2}\right\}=4\pi R\sqrt{r^{2}-x^{2}},\quad x\in[-r,r],
\]
于是, 圆环体的体积为
\[
\begin{aligned}
V&=V_{1}-V_{2}=\int_{-r}^{r}\pi[f_{2}(x)]^{2}\,\mathrm{d}x-\int_{-r}^{r}\pi[f_{1}(x)]^{2}\,\mathrm{d}x\\
&\xlongequal{\text{对称性}}8\pi R\int_{0}^{r}\sqrt{r^{2}-x^{2}}\,\mathrm{d}x
\xlongequal{x=r\sin\theta}8\pi R\int_{0}^{\frac{\pi}{2}}\sqrt{r^{2}-r^{2}\sin^{2}\theta}\,r\cos\theta\,\mathrm{d}\theta\\
&=8\pi Rr^{2}\int_{0}^{\frac{\pi}{2}}\cos^{2}\theta\,\mathrm{d}\theta=8\pi Rr^{2}\cdot\frac{1}{2}\cdot\frac{\pi}{2}\\
&=2\pi^{2}r^{2}R.
\end{aligned}
\]
\end{solution}

\subsection*{4. 平面曲线与空间曲线的弧长}

\begin{definition}\label{def:6.6.1}
设 $C=\overset{\frown}{AB}$ 为 $\mathbb{R}^{3}$ 中的空间曲线, 在 $C$ 上依次从 $A$ 到 $B$ 取分点
\[
A=P_{0},\ P_{1},\ \cdots,\ P_{n-1},\ P_{n}=B,
\]
它们成为对曲线 $C$ 的一个分割, 记为 $T$. 然后, 用线段连结 $T$ 中每相邻两点得到 $C$ 的 $n$ 条弦 $\overline{P_{i-1}P_{i}}$ ($i=1,\cdots,n$), 这 $n$ 条弦又组成 $C$ 的一条内接折线 (图~\ref{fig:6.6.12}). 记
\[
\|T\|=\max_{1\leqslant i\leqslant n}|P_{i-1}P_{i}|,
\]
\[
s_{T}=\sum_{i=1}^{n}|P_{i-1}P_{i}|,
\]
它们分别表示最长弦的长度与折线的总长度.
如果
\[
\lim_{\|T\|\to0}s_{T}=s,
\]
即 $\forall\varepsilon>0$, $\exists\delta>0$, 当 $\|T\|<\delta$ 时, 有
\[
|s_{T}-s|<\varepsilon,
\]
则称 $C$ 是可求长的, 并将极限 $s$ 称为曲线 $C$ 的弧长.
\end{definition}

\begin{figure}[htbp]
\centering
\begin{tikzpicture}[x=1.0cm,y=1.0cm,>=Stealth]
  \draw[->] (0,0) -- (0,3.5) node[above] {$z$};
  \draw[->] (0,0) -- (5.2,0) node[right] {$y$};
  \draw[->] (0,0) -- (-1.7,-1.2) node[below left] {$x$};
  \node[below right] at (0,0) {$O$};
  \draw[very thick] (0.9,0.55) .. controls (1.6,2.7) and (3.6,3.0) .. (4.5,1.35);
  \draw[thick] (0.9,0.55) -- (1.95,1.95) -- (3.1,2.2) -- (4.5,1.35);
  \fill (0.9,0.55) circle (0.05); \fill (1.95,1.95) circle (0.05);
  \fill (3.1,2.2) circle (0.05); \fill (4.5,1.35) circle (0.05);
  \node[left] at (0.82,0.5) {$P_{0}=A(t=\alpha)$};
  \node[above left] at (1.95,1.98) {$P_{1}(t=t_{1})$};
  \node[above] at (3.1,2.26) {$P_{i-1}(t=t_{i-1})$};
  \node[right] at (3.18,2.0) {$P_{i}(t=t_{i})$};
  \node[above right] at (4.52,1.4) {$B=P_{n}(t=\beta)$};
  \node[above] at (2.15,2.95) {$C$};
\end{tikzpicture}
\caption{}\label{fig:6.6.12}
\end{figure}

\begin{theorem}[可求长的充分条件]\label{thm:6.6.1}
设 $\mathbb{R}^{3}$ 中的空间曲线 $C$ 由参数方程
\[
(x,y,z)=(x(t),y(t),z(t)),\quad t\in[\alpha,\beta]
\]
给出, 并且它为一条 $C^{1}$ 光滑的正则曲线 (即 $x(t)$, $y(t)$, $z(t)$ 均连续可导, 且 $x'^{2}(t)+y'^{2}(t)+z'^{2}(t)\neq0$, 满足此条件的曲线 $C$ 的点称为正则点), 则 $C$ 是可求长的, 且弧长为
\[
s=\int_{\alpha}^{\beta}\sqrt{x'^{2}(t)+y'^{2}(t)+z'^{2}(t)}\,\mathrm{d}t.
\]
\end{theorem}

\begin{proof}
设
\[
T':\alpha=t_{0}<t_{1}<\cdots<t_{n-1}<t_{n}=\beta
\]
为区间 $[\alpha,\beta]$ 的一个分割, 相应于曲线 $C$ 的一个分割为
\[
T=\{P_{0},P_{1},\cdots,P_{n}\},
\]
其中 $P_{i}=(x(t_{i}),y(t_{i}),z(t_{i}))$, $i=0,1,\cdots,n$ (图~\ref{fig:6.6.12}). 在 $T'$ 所属的每个小区间 $[t_{i-1},t_{i}]$ 上, 由 Lagrange 中值定理得
\[
\Delta x_{i}=x(t_{i})-x(t_{i-1})=x'(\xi_{i})\Delta t_{i},\quad\xi_{i}\in[t_{i-1},t_{i}],
\]
\[
\Delta y_{i}=y(t_{i})-y(t_{i-1})=y'(\eta_{i})\Delta t_{i},\quad\eta_{i}\in[t_{i-1},t_{i}],
\]
\[
\Delta z_{i}=z(t_{i})-z(t_{i-1})=z'(\zeta_{i})\Delta t_{i},\quad\zeta_{i}\in[t_{i-1},t_{i}].
\]
又因为 $C$ 为光滑曲线, 当 $x'(t)\neq0$ 时, 在 $t$ 的某开邻域内 $x=x(t)$ 有连续的反函数, 故当 $\Delta x\to0$ 时 $\Delta t\to0$; 类似地, 当 $y'(t)\neq0$ 或 $z'(t)\neq0$ 时, 也能从 $\Delta y\to0$ 或 $\Delta z\to0$ 推知 $\Delta t\to0$. 所以,
\[
|P_{i-1}P_{i}|=\sqrt{\Delta x_{i}^{2}+\Delta y_{i}^{2}+\Delta z_{i}^{2}}\to0\Longleftrightarrow\Delta t\to0.
\]
由此知, 当曲线 $C$ 为光滑曲线时,
\[
\|T\|\to0\Longleftrightarrow\|T'\|\to0.
\]
由于 $\sqrt{x'^{2}(t)+y'^{2}(t)+z'^{2}(t)}$ 在 $[\alpha,\beta]$ 上连续, 从而 Riemann 可积, 所以 $\forall\varepsilon>0$, $\exists\delta_{1}>0$, 当 $\|T'\|<\delta_{1}$ 时,
\[
\left|\sum_{i=1}^{n}\sqrt{x'^{2}(\xi_{i})+y'^{2}(\xi_{i})+z'^{2}(\xi_{i})}\Delta t_{i}-\int_{\alpha}^{\beta}\sqrt{x'^{2}(t)+y'^{2}(t)+z'^{2}(t)}\,\mathrm{d}t\right|<\frac{\varepsilon}{2}.
\]
又因 $y(t)$, $z(t)$ 在 $[\alpha,\beta]$ 上连续, 故一致连续. 从而 $\exists\delta\in(0,\delta_{1})$, 当 $\|T'\|<\delta$ 时, 有
\[
\begin{aligned}
&\left|\sqrt{x'^{2}(\xi_{i})+y'^{2}(\eta_{i})+z'^{2}(\zeta_{i})}-\sqrt{x'^{2}(\xi_{i})+y'^{2}(\xi_{i})+z'^{2}(\xi_{i})}\right|\\
\leqslant&\frac{|y'(\eta_{i})+y'(\xi_{i})||y'(\eta_{i})-y'(\xi_{i})|+|z'(\zeta_{i})+z'(\xi_{i})||z'(\zeta_{i})-z'(\xi_{i})|}{\sqrt{x'^{2}(\xi_{i})+y'^{2}(\eta_{i})+z'^{2}(\zeta_{i})}+\sqrt{x'^{2}(\xi_{i})+y'^{2}(\xi_{i})+z'^{2}(\xi_{i})}}\\
\leqslant&|y'(\eta_{i})-y'(\xi_{i})|+|z'(\zeta_{i})-z'(\xi_{i})|<\frac{\varepsilon}{2(\beta-\alpha)}.
\end{aligned}
\]
于是
\[
\begin{aligned}
&\left|s_{T}-\int_{\alpha}^{\beta}\sqrt{x'^{2}(t)+y'^{2}(t)+z'^{2}(t)}\,\mathrm{d}t\right|\\
=&\left|\sum_{i=1}^{n}\sqrt{\Delta x_{i}^{2}+\Delta y_{i}^{2}+\Delta z_{i}^{2}}-\int_{\alpha}^{\beta}\sqrt{x'^{2}(t)+y'^{2}(t)+z'^{2}(t)}\,\mathrm{d}t\right|\\
=&\left|\sum_{i=1}^{n}\sqrt{x'^{2}(\xi_{i})+y'^{2}(\eta_{i})+z'^{2}(\zeta_{i})}\Delta t_{i}-\int_{\alpha}^{\beta}\sqrt{x'^{2}(t)+y'^{2}(t)+z'^{2}(t)}\,\mathrm{d}t\right|\\
\leqslant&\sum_{i=1}^{n}\left|\sqrt{x'^{2}(\xi_{i})+y'^{2}(\eta_{i})+z'^{2}(\zeta_{i})}-\sqrt{x'^{2}(\xi_{i})+y'^{2}(\xi_{i})+z'^{2}(\xi_{i})}\right|\Delta t_{i}\\
&+\left|\sum_{i=1}^{n}\sqrt{x'^{2}(\xi_{i})+y'^{2}(\xi_{i})+z'^{2}(\xi_{i})}\Delta t_{i}-\int_{\alpha}^{\beta}\sqrt{x'^{2}(t)+y'^{2}(t)+z'^{2}(t)}\,\mathrm{d}t\right|\\
<&\frac{\varepsilon}{2(\beta-\alpha)}(\beta-\alpha)+\frac{\varepsilon}{2}=\varepsilon,
\end{aligned}
\]
即
\[
\lim_{\|T\|\to0}s_{T}=\int_{\alpha}^{\beta}\sqrt{x'^{2}(t)+y'^{2}(t)+z'^{2}(t)}\,\mathrm{d}t.
\]
\end{proof}

如果空间曲线 $C$ 以直角坐标 $x$ 为参数, 即

\[
(x,y,z)=(x,y(x),z(x)),
\]

且 $y(x)$, $z(x)$ 在 $[a,b]$ 上连续可导, 则弧长公式为

\[
s=\int_{a}^{b}\sqrt{1+y'^{2}(x)+z'^{2}(x)}\,\mathrm{d}x.
\]

类似上面讨论, 或在上述公式中令 $z(t)=0$, 则有平面曲线的弧长公式为

\[
s=\int_{\alpha}^{\beta}\sqrt{x'^{2}(t)+y'^{2}(t)}\,\mathrm{d}t.
\]

如果直角坐标 $x$ 为参数, 即 $(x,y)=(x,y(x))$, 且 $y(x)$ 在 $[a,b]$ 上连续可导, 则弧长公式为

\[
s=\int_{a}^{b}\sqrt{1+y'^{2}(x)}\,\mathrm{d}x.
\]

如果平面曲线 $C$ 由极坐标方程

\[
r=r(\theta),\quad\theta\in[\alpha,\beta]
\]

给出, $r(\theta)$ 连续可导, 且

\[
(x,y)=(r(\theta)\cos\theta,r(\theta)\sin\theta),
\]

\[
\begin{aligned}
&x'^{2}(\theta)+y'^{2}(\theta)\\
=&[r'(\theta)\cos\theta-r(\theta)\sin\theta]^{2}+[r'(\theta)\sin\theta+r(\theta)\cos\theta]^{2}\\
=&r^{2}(\theta)+r'^{2}(\theta)\neq0,
\end{aligned}
\]

则弧长公式为

\[
s=\int_{\alpha}^{\beta}\sqrt{x'^{2}(\theta)+y'^{2}(\theta)}\,\mathrm{d}\theta=\int_{\alpha}^{\beta}\sqrt{r^{2}(\theta)+r'^{2}(\theta)}\,\mathrm{d}\theta.
\]

\begin{remark}\label{rem:6.6.1}
易见, 从曲线 $C$ 上定点 $P_{0}=(x(\alpha),y(\alpha),z(\alpha))$ 到动点 $P=(x(t),y(t),z(t))$ 的弧长为
\[
s(t)=\int_{\alpha}^{t}\sqrt{x'^{2}(\tau)+y'^{2}(\tau)+z'^{2}(\tau)}\,\mathrm{d}\tau.
\]
由于被积函数是连续的, 根据微积分基本定理, 有
\[
\frac{\mathrm{d}s}{\mathrm{d}t}=\sqrt{x'^{2}(t)+y'^{2}(t)+z'^{2}(t)},
\]
\[
\begin{aligned}
\mathrm{d}s&=\sqrt{x'^{2}(t)+y'^{2}(t)+z'^{2}(t)}\,\mathrm{d}t\\
&=\sqrt{\mathrm{d}x^{2}+\mathrm{d}y^{2}+\mathrm{d}z^{2}},
\end{aligned}
\]
并称 $s(t)$ 的微分 $\mathrm{d}s$ 为弧长元或弧微分 (图~\ref{fig:6.6.13}).
\end{remark}

\begin{figure}[htbp]
\centering
\begin{tikzpicture}[x=1.0cm,y=1.0cm,>=Stealth]
  \draw[->] (0,0) -- (0,3.2) node[above] {$z$};
  \draw[->] (0,0) -- (4.6,0) node[right] {$y$};
  \draw[->] (0,0) -- (-1.5,-1.05) node[below left] {$x$};
  \node[below right] at (0,0) {$O$};
  \draw[very thick] (0.8,0.5) .. controls (1.5,2.5) and (3.3,2.9) .. (4.2,1.3);
  \draw[thick] (0.8,0.5) -- (2.0,1.7) -- (3.0,1.95);
  \draw[thick,dashed] (3.0,1.95) -- (4.2,1.3);
  \fill (2.0,1.7) circle (0.05); \fill (3.0,1.95) circle (0.05);
  \draw[thick] (2.0,1.7) -- (3.0,1.70);
  \draw[thick] (3.0,1.95) -- (3.0,1.70);
  \draw[->,thick] (2.0,1.7) -- (3.0,1.95) node[midway,above left,inner sep=1pt] {$\mathrm{d}s$};
  \node[below] at (2.5,1.62) {$\mathrm{d}x$};
  \node[right] at (3.06,1.82) {$\mathrm{d}z$};
  \node[left] at (1.95,1.66) {$P$};
  \node[right] at (3.10,1.95) {$Q$};
  \node[above] at (1.42,2.58) {$T$};
\end{tikzpicture}
\caption{}\label{fig:6.6.13}
\end{figure}

特别当参数 $t$ 为弧长 $s$ 时, 有

\[
s(s)=\int_{\alpha}^{\beta}\sqrt{x'^{2}(\tau)+y'^{2}(\tau)+z'^{2}(\tau)}\,\mathrm{d}\tau,
\]

则

\[
\begin{aligned}
1&=\frac{\mathrm{d}s}{\mathrm{d}s}\\
&=\sqrt{x'^{2}(s)+y'^{2}(s)+z'^{2}(s)}\\
&=\|(x'(s),y'(s),z'(s))\|,
\end{aligned}
\]

即 $(x'(s),y'(s),z'(s))$ 为沿曲线 $C$ 的单位切向量场.

\subsection*{5. 旋转曲面的面积}

平面 $C^{1}$ 光滑正则曲线 $C$ (参数方程为 $x=x(t)$, $y=y(t)$, $t\in[\alpha,\beta]$, $y(t)\geqslant0$) 绕 $x$ 轴旋转所得旋转曲面的面积公式为

\[
\begin{aligned}
S&=\lim_{\|T\|\to0}\sum_{i=1}^{n}2\pi y(\xi_{i})\cdot\sqrt{x'^{2}(\xi_{i})+y'^{2}(\xi_{i})}\Delta t_{i}\\
&=\int_{\alpha}^{\beta}2\pi y(t)\sqrt{x'^{2}(t)+y'^{2}(t)}\,\mathrm{d}t=\int_{\alpha}^{\beta}2\pi y\,\mathrm{d}s.
\end{aligned}
\]

如果直角坐标 $x$ 为参数, $y(x)$ 为连续可导的函数, 则上述公式就成为

\[
S=\int_{a}^{b}2\pi y(x)\sqrt{1+y'^{2}(x)}\,\mathrm{d}x.
\]

\begin{example}\label{ex:6.6.9}
求摆线 $(x,y)=(a(t-\sin t),a(1-\cos t))$, $a>0$ 的一拱的弧长 $s$.
\end{example}

\begin{solution}
\[
\begin{aligned}
s&=\int_{0}^{2\pi}\sqrt{x'^{2}(t)+y'^{2}(t)}\,\mathrm{d}t\\
&=\int_{0}^{2\pi}\sqrt{a^{2}(1-\cos t)^{2}+a^{2}\sin^{2}t}\,\mathrm{d}t\\
&=\int_{0}^{2\pi}\sqrt{2a^{2}(1-\cos t)}\,\mathrm{d}t=2a\int_{0}^{2\pi}\sin\frac{t}{2}\,\mathrm{d}t\\
&=\left.-4a\cos\frac{t}{2}\right|_{0}^{2\pi}=8a.
\end{aligned}
\]
\end{solution}

\begin{example}\label{ex:6.6.10}
求悬链线 $y=\dfrac{\mathrm{e}^{x}+\mathrm{e}^{-x}}{2}$ 从 $x=0$ 到 $x=a>0$ 那一段的弧长 $s$.
\end{example}

\begin{solution}
\[
\begin{aligned}
s&=\int_{0}^{a}\sqrt{1+y'^{2}(x)}\,\mathrm{d}x=\int_{0}^{a}\sqrt{1+\left(\frac{\mathrm{e}^{x}-\mathrm{e}^{-x}}{2}\right)^{2}}\,\mathrm{d}x\\
&=\int_{0}^{a}\sqrt{\left(\frac{\mathrm{e}^{x}+\mathrm{e}^{-x}}{2}\right)^{2}}\,\mathrm{d}x=\int_{0}^{a}\frac{\mathrm{e}^{x}+\mathrm{e}^{-x}}{2}\,\mathrm{d}x\\
&=\left.\frac{\mathrm{e}^{x}-\mathrm{e}^{-x}}{2}\right|_{0}^{a}=\frac{\mathrm{e}^{a}-\mathrm{e}^{-a}}{2}.
\end{aligned}
\]
\end{solution}

\begin{example}\label{ex:6.6.11}
求心脏线 (或心形线) $r=a(1+\cos\theta)$, $a>0$, $0\leqslant\theta\leqslant2\pi$ 的周长 (图~\ref{fig:6.6.14}).
\end{example}

\begin{solution}
\[
\begin{aligned}
S&=\int_{0}^{2\pi}\sqrt{r^{2}(\theta)+r'^{2}(\theta)}\,\mathrm{d}\theta\\
&=\int_{0}^{2\pi}\sqrt{a^{2}(1+\cos\theta)^{2}+a^{2}(-\sin\theta)^{2}}\,\mathrm{d}\theta\\
&\xlongequal{\text{对称性}}2a\int_{0}^{\pi}\sqrt{2(1+\cos\theta)}\,\mathrm{d}\theta=4a\int_{0}^{\pi}\cos\frac{\theta}{2}\,\mathrm{d}\theta\\
&=\left.8a\sin\frac{\theta}{2}\right|_{0}^{\pi}=8a.
\end{aligned}
\]
\end{solution}

\begin{figure}[htbp]
\centering
\begin{tikzpicture}[x=1.15cm,y=1.15cm,>=Stealth]
  \draw[->] (-1.35,0) -- (2.5,0) node[below right] {$x$};
  \draw[->] (0,-1.35) -- (0,1.35) node[left] {$y$};
  \node[below left] at (0,0) {$O$};
  \draw[very thick] plot[smooth] coordinates {(2.0000,0.0000) (1.9886,0.1740) (1.9547,0.3447) (1.8989,0.5088) (1.8227,0.6634) (1.7277,0.8056) (1.6160,0.9330) (1.4902,1.0434) (1.3529,1.1352) (1.2071,1.2071) (1.0560,1.2584) (0.9026,1.2890) (0.7500,1.2990) (0.6012,1.2893) (0.4590,1.2611) (0.3258,1.2159) (0.2038,1.1558) (0.0948,1.0830) (0.0000,1.0000) (-0.0796,0.9094) (-0.1435,0.8138) (-0.1918,0.7159) (-0.2250,0.6183) (-0.2440,0.5233) (-0.2500,0.4330) (-0.2446,0.3493) (-0.2296,0.2736) (-0.2071,0.2071) (-0.1792,0.1504) (-0.1481,0.1037) (-0.1160,0.0670) (-0.0849,0.0396) (-0.0567,0.0206) (-0.0329,0.0088) (-0.0150,0.0026) (-0.0038,0.0003) (-0.0000,0.0000) (-0.0038,-0.0003) (-0.0150,-0.0026) (-0.0329,-0.0088) (-0.0567,-0.0206) (-0.0849,-0.0396) (-0.1160,-0.0670) (-0.1481,-0.1037) (-0.1792,-0.1504) (-0.2071,-0.2071) (-0.2296,-0.2736) (-0.2446,-0.3493) (-0.2500,-0.4330) (-0.2440,-0.5233) (-0.2250,-0.6183) (-0.1918,-0.7159) (-0.1435,-0.8138) (-0.0796,-0.9094) (-0.0000,-1.0000) (0.0948,-1.0830) (0.2038,-1.1558) (0.3258,-1.2159) (0.4590,-1.2611) (0.6012,-1.2893) (0.7500,-1.2990) (0.9026,-1.2890) (1.0560,-1.2584) (1.2071,-1.2071) (1.3529,-1.1352) (1.4902,-1.0434) (1.6160,-0.9330) (1.7277,-0.8056) (1.8227,-0.6634) (1.8989,-0.5088) (1.9547,-0.3447) (1.9886,-0.1740) (2.0000,-0.0000)};
  \draw[thick] (0,0) -- (0.9:1.0);
  \draw (0.55,0) arc (0:30:0.55);
  \node at (0.72,0.20) {$\theta$};
  \node[above right] at (1.6,0.62) {$r=a(1+\cos\theta)$};
  \node[below left] at (1.95,0.02) {$2a$};
\end{tikzpicture}
\caption{}\label{fig:6.6.14}
\end{figure}

\begin{example}\label{ex:6.6.12}
计算圆 $x^{2}+y^{2}=R^{2}$ 在 $[x_{1},x_{2}]\subset[-R,R]$ 上的弧段绕 $x$ 轴旋转所得球带的面积 $S$ (图~\ref{fig:6.6.15}).
\end{example}

\begin{figure}[htbp]
\centering
\begin{tikzpicture}[x=1.0cm,y=1.0cm,>=Stealth]
  \draw[->] (-0.3,0) -- (4.9,0) node[right] {$x$};
  \draw[->] (0,-2.2) -- (0,2.2) node[above] {$y$};
  \draw[->] (0,0) -- (-1.5,-0.95) node[below left] {$z$};
  \draw[thick] (0,0) circle (1.75);
  \draw[thick] (0,0) ellipse (0.55 and 1.75);
  \draw[thick] (0,1.75) arc (90:-90:0.55 and 1.75);
  \draw[thick,dashed] (0,-1.75) arc (-90:90:0.55 and 1.75);
  \fill[pattern=north east lines]
    (1.05,1.40) .. controls (1.55,1.05) and (1.60,0.45) .. (1.40,0.10)
    -- (0.72,0.06) .. controls (0.92,0.55) and (0.90,1.10) .. (0.60,1.55) -- cycle;
  \draw[thick] (0.60,1.55) .. controls (0.90,1.10) and (0.92,0.55) .. (0.72,0.06);
  \node[left] at (0.55,1.40) {$S$};
  \draw[dotted] (1.40,0) -- (1.40,0.10);
  \node[right] at (1.45,0.20) {$x_{2}$};
  \node[below left] at (0.70,0.0) {$x_{1}$};
  \node[below right] at (0,0) {$O$};
\end{tikzpicture}
\caption{}\label{fig:6.6.15}
\end{figure}

\begin{solution}
圆 $x^{2}+y^{2}=R^{2}$ 在 $[x_{1},x_{2}]\subset[-R,R]$ 上弧段的显函数表达式为
\[
y=f(x)=\sqrt{R^{2}-x^{2}},
\]
则
\[
y'=f'(x)=\frac{-x}{\sqrt{R^{2}-x^{2}}}.
\]
于是所求面积为
\[
\begin{aligned}
S&=\int_{x_{1}}^{x_{2}}2\pi f(x)\sqrt{1+f'^{2}(x)}\,\mathrm{d}x\\
&=2\pi\int_{x_{1}}^{x_{2}}\sqrt{R^{2}-x^{2}}\sqrt{1+\frac{x^{2}}{R^{2}-x^{2}}}\,\mathrm{d}x\\
&=2\pi R\int_{x_{1}}^{x_{2}}\mathrm{d}x=2\pi R(x_{2}-x_{1}).
\end{aligned}
\]
\end{solution}

特别地, 当 $x_{1}=-R$, $x_{2}=R$ 时, 所得球带的表面积为 $4\pi R^{2}$.

\begin{example}\label{ex:6.6.13}
计算由内摆线 $x^{\frac{2}{3}}+y^{\frac{2}{3}}=a^{\frac{2}{3}}$ (其参数表示为 $(x,y)=(a\cos^{3}t,a\sin^{3}t)$) 绕 $x$ 轴旋转所得旋转曲面的面积 $S$ (图~\ref{fig:6.6.16}).
\end{example}

\begin{solution}
\[
\begin{aligned}
S&\xlongequal{\text{对称性}}4\pi\int_{0}^{\frac{\pi}{2}}y(t)\sqrt{x'^{2}(t)+y'^{2}(t)}\,\mathrm{d}t\\
&=4\pi\int_{0}^{\frac{\pi}{2}}a\sin^{3}t\sqrt{(-3a\cos^{2}t\sin t)^{2}+(3a\sin^{2}t\cos t)^{2}}\,\mathrm{d}t\\
&=12\pi a^{2}\int_{0}^{\frac{\pi}{2}}\sin^{4}t\cos t\,\mathrm{d}t=\left.12\pi a^{2}\frac{\sin^{5}t}{5}\right|_{0}^{\frac{\pi}{2}}=\frac{12}{5}\pi a^{2}.
\end{aligned}
\]
\end{solution}

\begin{figure}[htbp]
\centering
\begin{tikzpicture}[x=1.5cm,y=1.5cm,>=Stealth]
  \draw[->] (-1.3,0) -- (1.5,0) node[below right] {$x$};
  \draw[->] (0,-1.3) -- (0,1.3) node[left] {$y$};
  \node[below left] at (0,0) {$O$};
  \fill[pattern=north east lines]
    plot[smooth] coordinates {(1.0000,0.0000) (0.9886,0.0007) (0.9551,0.0052) (0.9012,0.0173) (0.8298,0.0400) (0.7444,0.0755) (0.6495,0.1250) (0.5497,0.1887) (0.4495,0.2656) (0.3536,0.3536) (0.2656,0.4495) (0.1887,0.5497) (0.1250,0.6495) (0.0755,0.7444) (0.0400,0.8298) (0.0173,0.9012) (0.0052,0.9551) (0.0007,0.9886) (0.0000,1.0000)} -- (0,0) -- cycle;
  \draw[very thick] plot[smooth] coordinates {(1.0000,0.0000) (0.9886,0.0007) (0.9551,0.0052) (0.9012,0.0173) (0.8298,0.0400) (0.7444,0.0755) (0.6495,0.1250) (0.5497,0.1887) (0.4495,0.2656) (0.3536,0.3536) (0.2656,0.4495) (0.1887,0.5497) (0.1250,0.6495) (0.0755,0.7444) (0.0400,0.8298) (0.0173,0.9012) (0.0052,0.9551) (0.0007,0.9886) (0.0000,1.0000) (-0.0007,0.9886) (-0.0052,0.9551) (-0.0173,0.9012) (-0.0400,0.8298) (-0.0755,0.7444) (-0.1250,0.6495) (-0.1887,0.5497) (-0.2656,0.4495) (-0.3536,0.3536) (-0.4495,0.2656) (-0.5497,0.1887) (-0.6495,0.1250) (-0.7444,0.0755) (-0.8298,0.0400) (-0.9012,0.0173) (-0.9551,0.0052) (-0.9886,0.0007) (-1.0000,0.0000) (-0.9886,-0.0007) (-0.9551,-0.0052) (-0.9012,-0.0173) (-0.8298,-0.0400) (-0.7444,-0.0755) (-0.6495,-0.1250) (-0.5497,-0.1887) (-0.4495,-0.2656) (-0.3536,-0.3536) (-0.2656,-0.4495) (-0.1887,-0.5497) (-0.1250,-0.6495) (-0.0755,-0.7444) (-0.0400,-0.8298) (-0.0173,-0.9012) (-0.0052,-0.9551) (-0.0007,-0.9886) (-0.0000,-1.0000) (0.0007,-0.9886) (0.0052,-0.9551) (0.0173,-0.9012) (0.0400,-0.8298) (0.0755,-0.7444) (0.1250,-0.6495) (0.1887,-0.5497) (0.2656,-0.4495) (0.3536,-0.3536) (0.4495,-0.2656) (0.5497,-0.1887) (0.6495,-0.1250) (0.7444,-0.0755) (0.8298,-0.0400) (0.9012,-0.0173) (0.9551,-0.0052) (0.9886,-0.0007) (1.0000,-0.0000)};
  \node[below] at (1,-0.05) {$a$};
  \node[anchor=west] at (0.35,1.05) {$x^{\frac{2}{3}}+y^{\frac{2}{3}}=a^{\frac{2}{3}}$};
\end{tikzpicture}
\caption{}\label{fig:6.6.16}
\end{figure}

\begin{example}\label{ex:6.6.14}
证明: 圆柱螺线 $\boldsymbol{r}(t)=(x(t),y(t),z(t))=(a\cos t,a\sin t,bt)$, $t\in\mathbb{R}$, $(x^{2}(t)+y^{2}(t)=a^{2})$ 在每一点处的切向量 $\boldsymbol{r}'(t)=(x'(t),y'(t),z'(t))$ 与 $z$ 轴交成定角 $\theta$. 求出圆柱螺线的弧长公式, 讨论弧长 $s$ 与 $t$ 之间的关系 (图~\ref{fig:6.6.17}).
\end{example}

\begin{solution}
$\boldsymbol{r}'(t)=(x'(t),y'(t),z'(t))=(-a\sin t,a\cos t,b)$, 单位切向量为
\[
\begin{aligned}
\frac{\boldsymbol{r}'(t)}{\|\boldsymbol{r}'(t)\|}&=\frac{(-a\sin t,a\cos t,b)}{\sqrt{(-a\sin t)^{2}+(a\cos t)^{2}+b^{2}}}\\
&=\frac{1}{\sqrt{a^{2}+b^{2}}}(-a\sin t,a\cos t,b).
\end{aligned}
\]
于是,
\[
\cos\theta=\frac{\boldsymbol{r}'(t)}{\|\boldsymbol{r}'(t)\|}\cdot(0,0,1)=\frac{b}{\sqrt{a^{2}+b^{2}}}.
\]
\[
\theta=\arccos\frac{b}{\sqrt{a^{2}+b^{2}}}
\]
为一个定角, 它的弧长公式为
\[
\begin{aligned}
s&=\int_{0}^{t}\|\boldsymbol{r}'(t)\|\,\mathrm{d}t=\int_{0}^{t}\sqrt{x'^{2}(t)+y'^{2}(t)+z'^{2}(t)}\,\mathrm{d}t\\
&=\frac{1}{\sqrt{a^{2}+b^{2}}}\int_{0}^{t}\mathrm{d}t=\frac{t}{\sqrt{a^{2}+b^{2}}},\\
&\quad t=\sqrt{a^{2}+b^{2}}\,s.
\end{aligned}
\]
\end{solution}

\begin{figure}[htbp]
\centering
\begin{tikzpicture}[x=1.0cm,y=1.0cm,>=Stealth]
  \draw[->] (0,0) -- (0,3.7) node[above] {$z$};
  \draw[->] (0,0) -- (2.5,0) node[below right] {$y$};
  \draw[->] (0,0) -- (-1.6,-0.95) node[below left] {$x$};
  \node[below left] at (0,0) {$O$};
  \draw[thick] (0,-0.45) ellipse (0.95 and 0.42);
  \draw[thick] (0,0.9) ellipse (0.95 and 0.42);
  \draw[thick] (0,2.25) ellipse (0.95 and 0.42);
  \draw[thick] (0,3.15) ellipse (0.95 and 0.42);
  \draw[thick] (-0.95,-0.45) -- (-0.95,3.15);
  \draw[thick] (0.95,-0.45) -- (0.95,3.15);
  \draw[very thick] (0.95,3.15) .. controls (1.75,2.6) and (0.05,2.0) .. (0.95,1.45)
    .. controls (1.75,0.85) and (0.05,0.25) .. (0.95,-0.30);
  \draw[->,very thick] (0.95,0.32) -- (1.30,0.06);
  \node[right] at (1.06,0.32) {$\boldsymbol{r}(t)$};
\end{tikzpicture}
\caption{}\label{fig:6.6.17}
\end{figure}

\subsection*{6. 液体静压力}

\begin{example}\label{ex:6.6.15}
图~\ref{fig:6.6.18} 所示为一管道的圆形闸门 (半径为 3m). 问水平面齐及直径时, 闸门所受到的水的静压力为多大?
\end{example}

\begin{solution}
为方便起见, 取 $x$ 轴和 $y$ 轴如图, 此时圆的方程为
\[
x^{2}+y^{2}=9.
\]
由于在相同深度处水的静压强相同, 其值等于水的密度 $\nu$ 与深度 $x$ 的乘积, 故当 $\Delta x$ 很小时, 闸门上从深度 $x$ 到 $x+\Delta x$ 这一狭条 $\Delta A$ 上所受的静压力为
\[
\Delta P\approx\mathrm{d}P=2\nu x\sqrt{9-x^{2}}\,\mathrm{d}x.
\]
从而闸门上所受的总压力为
\[
P=\int_{0}^{3}2\nu x\sqrt{9-x^{2}}\,\mathrm{d}x=18\nu.
\]
\end{solution}

\begin{figure}[htbp]
\centering
\begin{tikzpicture}[x=1.0cm,y=1.0cm,>=Stealth]
  \draw[->] (0,0) -- (0,-3.5) node[below] {$x$};
  \draw[->] (0,0) -- (3.4,0) node[right] {$y$};
  \node[above] at (0,0) {$O$};
  \draw[thick] (0,0) circle (2.2);
  \draw[thin] (0,0.4) -- (2.0,0.4);
  \draw[thin] (0,0.4) -- (0,0.9);
  \draw[thin] (0,0.9) -- (1.95,0.9);
  \fill[pattern=north east lines] (-1.95,0.4) rectangle (1.95,0.9);
  \node[left] at (0,0.9) {$x$};
  \node[left] at (0,0.4) {$x+\Delta x$};
  \node[right] at (2.05,0.65) {$\Delta A$};
  \draw[<->,thin] (-2.55,-2.2) -- (-2.55,0);
  \node[left] at (-2.62,-1.1) {3};
\end{tikzpicture}
\caption{}\label{fig:6.6.18}
\end{figure}

\subsection*{7. 引力}

\begin{example}\label{ex:6.6.16}
一根长为 $l$ 的均匀细杆, 质量为 $M$, 在其中垂线上相距细杆 $a$ 处有一质量为 $m$ 的质点. 试求细杆对质点的万有引力.
\end{example}

\begin{solution}
如图~\ref{fig:6.6.19} 所示, 细杆位于 $x$ 轴上的 $\left[-\dfrac{l}{2},\dfrac{l}{2}\right]$, 质点位于 $y$ 轴上的点 $a$. 任取 $[x,x+\mathrm{d}x]\subset\left[-\dfrac{l}{2},\dfrac{l}{2}\right]$, 当 $\mathrm{d}x$ 很小时可把这一小段细杆看作一个质点, 其质量为 $\mathrm{d}M=\dfrac{M}{l}\mathrm{d}x$. 于是它对质点 $m$ 的引力为
\[
\mathrm{d}F=\frac{km\,\mathrm{d}M}{r^{2}}=\frac{km}{a^{2}+x^{2}}\cdot\frac{M}{l}\mathrm{d}x.
\]
由于细杆上各点对质点 $m$ 的引力方向各不相同, 因此不能直接对 $\mathrm{d}F$ 进行积分 (不符合代数可加的条件). 为此, 将 $\mathrm{d}F$ 分解到 $x$ 轴和 $y$ 轴两个方向上, 得到
\[
\mathrm{d}F_{x}=\mathrm{d}F\cdot\sin\theta,\quad\mathrm{d}F_{y}=-\mathrm{d}F\cdot\cos\theta.
\]
由于质点 $m$ 位于细杆的中垂线上, 必使水平合力为零, 即
\[
F_{x}=\int_{-\frac{l}{2}}^{\frac{l}{2}}\mathrm{d}F_{x}=0.
\]
又由 $\cos\theta=\dfrac{a}{\sqrt{a^{2}+x^{2}}}$, 得垂直方向合力为
\[
\begin{aligned}
F_{y}&=\int_{-\frac{l}{2}}^{\frac{l}{2}}\mathrm{d}F_{y}=-2\int_{0}^{\frac{l}{2}}\frac{kmMa}{l}(a^{2}+x^{2})^{-\frac{3}{2}}\,\mathrm{d}x\\
&=-\frac{2kmMa}{l}\cdot\frac{1}{a^{2}}\cdot\left.\frac{x}{\sqrt{a^{2}+x^{2}}}\right|_{0}^{\frac{l}{2}}\\
&=-\frac{2kmM}{a\sqrt{4a^{2}+l^{2}}},
\end{aligned}
\]
负号表示合力方向与 $y$ 轴方向相反.
\end{solution}

\begin{figure}[htbp]
\centering
\begin{tikzpicture}[x=1.0cm,y=1.0cm,>=Stealth]
  \draw[->] (-2.6,0) -- (3.0,0) node[below right] {$x$};
  \draw[->] (0,-0.5) -- (0,3.4) node[left] {$y$};
  \node[below left] at (0,0) {$O$};
  \draw[very thick] (-1.9,0) -- (1.9,0);
  \node[below] at (-1.9,-0.08) {$-\dfrac{l}{2}$};
  \node[below] at (1.9,-0.08) {$\dfrac{l}{2}$};
  \node[below] at (1.15,-0.08) {$x$};
  \node[below] at (1.72,-0.08) {$x+\mathrm{d}x$};
  \fill (2.0,0) circle (0.05);
  \draw[thin] (2.0,0) -- (2.0,-0.35);
  \fill (0,2.75) circle (0.06);
  \node[above right] at (0.04,2.78) {$a$};
  \draw[very thick,->] (0,2.75) -- (1.9,0);
  \draw[very thick,->] (0,2.75) -- (1.05,1.35);
  \draw[dotted] (1.05,1.35) -- (1.05,2.75);
  \draw[thin] (0,1.55) arc (90:56:1.55);
  \node at (0.35,1.68) {$\theta$};
  \node[right] at (1.02,2.30) {$\mathrm{d}F$};
  \node[right] at (1.10,1.30) {$\mathrm{d}F_{x}$};
  \node[below left] at (0.06,2.45) {$\mathrm{d}F_{y}$};
\end{tikzpicture}
\caption{}\label{fig:6.6.19}
\end{figure}

\subsection*{8. 功}

\begin{example}\label{ex:6.6.17}
一圆锥形水池, 池口直径 30m, 深 10m, 池中盛满了水. 试求将全部池水抽出池外需做的功.
\end{example}

\begin{solution}
为方便起见, 取坐标轴如图~\ref{fig:6.6.20} 所示. 由于抽出相同深度处单位体积的水需做相同的功 (等于水的密度 $\nu\times$ 深度), 因此首先考虑将池中深度为 $x$ 到 $x+\Delta x$ 的一薄层水 $\Delta\Omega$ 抽至池口需做的功 $\Delta W$. 当 $\Delta x$ 很小时, 把这一薄层水的深度都看作 $x$, 并作 $\Delta\Omega$ 的体积
\[
\Delta V\approx\pi\left[15\left(1-\frac{x}{10}\right)\right]^{2}\Delta x,
\]
这时有
\[
\Delta W\approx\mathrm{d}W=\pi\nu x\left[15\left(1-\frac{x}{10}\right)\right]^{2}\mathrm{d}x.
\]
从而将全部池水抽出池外需做的功为
\[
W=225\pi\nu\int_{0}^{10}x\left(1-\frac{x}{10}\right)^{2}\mathrm{d}x=1875\pi\nu.
\]
\end{solution}

\begin{figure}[htbp]
\centering
\begin{tikzpicture}[x=0.48cm,y=0.36cm,>=Stealth]
  \draw[->] (0,0) -- (8.4,0) node[right] {$y$};
  \draw[->] (0,0) -- (0,-7.4) node[below] {$x$};
  \node[above left] at (0,0) {$O$};
  \draw[thick] (0,0) -- (3.0,-6.0);
  \draw[thick] (0,0) -- (-3.0,-6.0);
  \draw[thick] (0,0) ellipse (3.0 and 0.63);
  \draw[thick] (-3.0,-6.0) arc (180:360:3.0 and 0.63);
  \draw[thick,dashed] (3.0,-6.0) arc (0:180:3.0 and 0.63);
  \draw[thick] (0,-6.0) ellipse (1.32 and 0.28);
  \draw[thick,dashed] (-1.32,-6.0) arc (180:360:1.32 and 0.28);
  \fill[pattern=north east lines] (-2.05,-2.75) rectangle (2.05,-3.25);
  \node[above] at (0,-0.45) {15};
  \node[below] at (2.55,-2.75) {$\Delta\Omega$};
  \node[left] at (0,-6.0) {$10$};
  \draw[<->,thin] (4.2,-2.75) -- (4.2,-3.25);
  \node[right] at (4.3,-3.0) {$\Delta x$};
  \draw[thin] (-2.05,-3.0) -- (4.2,-3.0);
\end{tikzpicture}
\caption{}\label{fig:6.6.20}
\end{figure}

\begin{example}\label{ex:6.6.18}
(第二宇宙速度) 在地球表面垂直向上发射火箭, 要使火箭克服地球引力无限远离地球. 试问: 初速度 $v_{0}$ 至少要多大?
\end{example}

\begin{solution}
根据万有引力定律, 在距地心 $x$ ($\geqslant R$) 处火箭所受的引力为
\[
F=\frac{GMm}{x^{2}}=\frac{mgR^{2}}{x^{2}},
\]
其中 $G$ 为引力常数, $M$ 为地球的质量, $m$ 为火箭质量, $R$ 为地球半径, $g$ 为重力加速度. 因此, 在地球表面有 $\dfrac{GMm}{R^{2}}=mg$, $GM=gR^{2}$.

于是, 火箭从地面上升到距离地心为 $r$ ($>R$) 处需做的功为 (图~\ref{fig:6.6.21})
\[
\int_{R}^{r}\frac{mgR^{2}}{x^{2}}\,\mathrm{d}x.
\]
从而火箭无限远离地球需做的功为
\[
\begin{aligned}
W&=\lim_{r\to+\infty}\int_{R}^{r}\frac{mgR^{2}}{x^{2}}\,\mathrm{d}x\\
&=\int_{R}^{+\infty}\frac{mgR^{2}}{x^{2}}\,\mathrm{d}x\\
&=\left.-\frac{mgR^{2}}{x}\right|_{R}^{+\infty}=\frac{mgR^{2}}{R}=mgR.
\end{aligned}
\]
再根据机械能守恒定律可求得初速度 $v_{0}$ 至少应使
\[
\frac{1}{2}mv_{0}^{2}=mgR,
\]
\[
\begin{aligned}
v_{0}&=\sqrt{2gR}=\sqrt{2\times9.81(\mathrm{m}/\mathrm{s}^{2})\times6.371\times10^{6}(\mathrm{m})}\\
&\approx11.2(\mathrm{km}/\mathrm{s}).
\end{aligned}
\]
\end{solution}

\begin{figure}[htbp]
\centering
\begin{tikzpicture}[x=1.0cm,y=1.0cm,>=Stealth]
  \draw[->] (0,0) -- (0,3.1) node[above] {$x$};
  \draw[->] (0,0) -- (3.1,0) node[below right] {$z$};
  \draw[->] (0,0) -- (-1.7,-1.15) node[below left] {$y$};
  \node[below left] at (0,0) {$O$};
  \draw[thick] (0,0) circle (1.75);
  \draw[thick] (0,0) ellipse (0.62 and 1.75);
  \draw[thick] (0,1.75) arc (90:-90:0.62 and 1.75);
  \draw[thick,dashed] (0,-1.75) arc (-90:90:0.62 and 1.75);
  \draw[thick,dashed] (1.75,0) .. controls (1.75,1.4) and (0.5,1.75) .. (0,1.75);
  \node[left] at (0,1.75) {$r$};
  \draw[thin] (0,0) -- (1.35,1.35);
  \node[right] at (0.62,0.78) {$R$};
  \draw[very thick,->] (0,1.75) -- (0,2.75);
\end{tikzpicture}
\caption{}\label{fig:6.6.21}
\end{figure}

我们称上述 $v_{0}$ 为第二宇宙速度.

\begin{exblock}{练习题 6.6}
\item 求由抛物线 $y=x^{2}$ 与 $y=2-x^{2}$ 所围图形的面积.
\item 求由曲线 $y=|\ln x|$ 与直线 $x=\dfrac{1}{10}$, $x=10$, $y=0$ 所围图形的面积.
\item 抛物线 $y^{2}=2x$ 将圆片 $x^{2}+y^{2}\leqslant8$ 分成两部分, 求这两部分的面积之比.
\item 求两椭圆 $\dfrac{x^{2}}{a^{2}}+\dfrac{y^{2}}{b^{2}}=1$ 与 $\dfrac{x^{2}}{b^{2}}+\dfrac{y^{2}}{a^{2}}=1$ ($a>0$, $b>0$) 所围公共部分的面积.
\item 求内摆线 $x=a\cos^{3}t$, $y=a\sin^{3}t$ ($a>0$, 即 $x^{\frac{2}{3}}+y^{\frac{2}{3}}=a^{\frac{2}{3}}$) 所围图形 (图~\ref{fig:6.6.16}) 的面积.
\item 求心形 (脏) 线 $r=a(1+\cos\theta)$ ($a>0$, 图~\ref{fig:6.6.14}) 所围图形的面积.
\item 求三叶形曲线 $r=a\sin3\theta$ ($a>0$) 所围图形的面积.
\item 求由曲线 $\sqrt{\dfrac{x}{a}}+\sqrt{\dfrac{y}{b}}=1$ ($a,b>0$) 与坐标轴所围图形的面积.
\item 求两曲线 $r=\sin\theta$ 与 $r=\sqrt{3}\cos\theta$ 所围公共部分的面积.
\item 求由曲线 $x=t-t^{3}$, $y=1-t^{4}$ 所围图形的面积.
\item 如图~\ref{fig:6.6.22} 所示, 直椭圆柱体被通过底面短轴的斜平面所截. 试求截得楔形体的体积.
\item 求下列平面曲线绕所给坐标轴旋转所围成立体的体积:
\begin{enumerate}[label=(\arabic*)]
\item $y=\sin x$, $0\leqslant x\leqslant\pi$, 绕 $x$ 轴;
\item $x=a(t-\sin t)$, $y=a(1-\cos t)$ ($a>0$), $0\leqslant t\leqslant2\pi$, 绕 $x$ 轴;
\item $r=a(1+\cos\theta)$ ($a>0$), 绕极轴;
\item $\dfrac{x^{2}}{a^{2}}+\dfrac{y^{2}}{b^{2}}=1$, 绕 $y$ 轴.
\end{enumerate}
\item 求曲线 $x=a\cos^{3}t$, $y=a\sin^{3}t$ 所围平面图形 (图~\ref{fig:6.6.16}) 绕 $x$ 轴旋转所得立体的体积.
\item 求 $0\leqslant y\leqslant\sin x$, $0\leqslant x\leqslant\pi$ 所示平面图形分别绕 $x$ 轴和 $y$ 轴旋转所得立体的体积.
\item 导出曲边梯形 $0\leqslant y\leqslant f(x)$, $a\leqslant x\leqslant b$ 绕 $y$ 轴旋转所得立体的体积公式
\[
V=2\pi\int_{a}^{b}xf(x)\,\mathrm{d}x.
\]
\item 设球的半径为 $r$, 证明: 高为 $h$ 的球缺 (图~\ref{fig:6.6.23}) 体积
\[
V=\pi h^{2}\left(r-\frac{h}{3}\right),\quad h\leqslant r.
\]
\item 求下列曲线的弧长:
\begin{enumerate}[label=(\arabic*)]
\item $y=x^{\frac{3}{2}}$, $0\leqslant x\leqslant4$;
\item $\sqrt{x}+\sqrt{y}=1$;
\item $x=a\cos^{3}t$, $y=a\sin^{3}t$, $0\leqslant t\leqslant2\pi$, 其中 $a>0$;
\item $x=\mathrm{e}^{t}\cos t$, $y=\mathrm{e}^{t}\sin t$, $0\leqslant t\leqslant2\pi$;
\item $x=a(\cos t+t\sin t)$, $y=a(\sin t-t\cos t)$, $0\leqslant t\leqslant2\pi$, 其中 $a>0$;
\item $x=\dfrac{c^{2}}{a^{2}}\cos^{3}t$, $y=\dfrac{c^{2}}{b^{2}}\sin^{3}t$, $0\leqslant t\leqslant2\pi$, 其中常数 $a>b>0$, 且 $c^{2}=a^{2}-b^{2}$;
\item $y^{2}=2px$, $0\leqslant x\leqslant a$, 其中 $p>0$;
\item $y=\ln\cos x$, $x\in\left[0,\dfrac{\pi}{3}\right]$;
\item $y=\sqrt{25-x^{2}}$, $x\in\left[0,\dfrac{5\sqrt{2}}{2}\right]$;
\item $r=a\sin^{3}\dfrac{\theta}{3}$, $0\leqslant\theta\leqslant3\pi$, 其中 $a>0$;
\item $r=a\theta$, $0\leqslant\theta\leqslant2\pi$, 其中 $a>0$;
\item $x=t$, $y=3t^{2}$, $z=6t^{3}$, $0\leqslant t\leqslant2$;
\item $x=\sin t$, $y=t$, $z=1-\cos t$, $0\leqslant t\leqslant2\pi$.
\end{enumerate}
\item 求 $a,b$ 的值, 使椭圆 $x=a\cos t$, $y=b\sin t$ 的周长等于正弦曲线 $y=\sin x$ 在 $0\leqslant x\leqslant2\pi$ 上一段的长.
\item 证明: 半径为 $a$ 的球的表面积为 $4\pi a^{2}$.
\item 求下列平面曲线绕指定轴旋转所得旋转曲面的面积:
\begin{enumerate}[label=(\arabic*)]
\item $y^{2}=x$, $0\leqslant x\leqslant1$, 绕 $x$ 轴;
\item $y=\sin x$, $0\leqslant x\leqslant\pi$, 绕 $x$ 轴;
\item $\dfrac{x^{2}}{a^{2}}+\dfrac{y^{2}}{b^{2}}=1$, 绕 $y$ 轴;
\item $x^{2}+(y-a)^{2}=r^{2}$ ($r<a$), 绕 $x$ 轴;
\item $x=a(t-\sin t)$, $y=a(1-\cos t)$, $0\leqslant t\leqslant2\pi$, 绕 $x$ 轴, 其中 $a>0$.
\end{enumerate}
\item 设平面 $C^{1}$ (连续可导) 光滑正则曲线由极坐标方程
\[
r=r(\theta),\quad\alpha\leqslant\theta\leqslant\beta\quad([\alpha,\beta]\subset[0,\pi],\ r(\theta)\geqslant0)
\]
给出. 试求它绕极轴旋转所得旋转曲面的面积计算公式.
\item 求下列极坐标曲线绕极轴旋转所得旋转曲面的面积:
\begin{enumerate}[label=(\arabic*)]
\item 心形 (脏) 线 $r=a(1+\cos\theta)$, 其中 $a>0$;
\item 双纽线 $r^{2}=2a^{2}\cos2\theta$, 其中 $a>0$.
\end{enumerate}
\item 有一等腰梯形闸门, 它的上、下两条底边各长为 10m 与 6m, 高为 20m. 计算当水面与上底边相齐时闸门一侧所受的静压力.
\item 边长为 $a$ 与 $b$ 的矩形薄板, 与液面成 $\alpha$ ($0<\alpha<90^{\circ}$) 角斜沉于液体中. 设 $a>b$, 长边平行于液面, 上沿位于深 $h$ 处, 液体的密度为 $\nu$. 试求薄板每侧所受的静压力.
\item 直径为 6m 的一个球浸入水中, 其球心在水平面下 10m 处, 求球面上所受的静压力.
\item 设在坐标系的原点有一质量为 $m$ 的质点, 在区间 $[a,a+l]$ ($a>0$) 上有一质量为 $M$ 的均匀细杆. 试求质点与细杆之间的万有引力.
\item 设有两条各长为 $l$ 的均匀细杆在同一直线上, 中间离开距离 $c$, 每根细杆的质量为 $M$. 试求它们之间的万有引力.
\item 一个半球形 (直径为 20m) 的容器内盛满了水. 试问: 将水抽尽需做多少功?
\item 长 10m 的铁索下垂于矿井中, 已知铁索每米的质量为 8kg. 问: 将此铁索提出地面需做多少功?
\item 一物体在某介质中按 $x=ct^{3}$ 作直线运动, 介质的阻力与速度 $\dfrac{\mathrm{d}x}{\mathrm{d}t}$ 的平方成正比. 计算物体由 $x=0$ 移至 $x=a$ 时克服介质阻力所做的功.
\item 半径为 $r$ 的球体沉入水中, 其比重与水相同. 试问: 将球体从水中捞出需做多少功?
\end{exblock}

\begin{figure}[htbp]
\centering
\begin{tikzpicture}[x=1.0cm,y=1.0cm,>=Stealth]
  \draw[thick] (-2.4,0.35) .. controls (-1.2,-1.05) and (1.2,-1.35) .. (2.4,-0.25);
  \draw[thick] (-2.4,-0.85) .. controls (-1.2,-2.25) and (1.2,-2.55) .. (2.4,-1.45);
  \draw[thick] (-2.4,0.35) .. controls (-2.85,-0.25) and (-2.85,-0.25) .. (-2.4,-0.85);
  \draw[thick] (2.4,-0.25) .. controls (2.85,-0.85) and (2.85,-0.85) .. (2.4,-1.45);
  \draw[thick] (-2.4,0.35) -- (0.0,-0.25);
  \draw[thick] (0.0,-0.25) -- (2.4,-0.25);
  \draw[thick] (-2.4,-0.85) -- (0.0,-1.45);
  \draw[thick] (0.0,-1.45) -- (2.4,-1.45);
  \draw[thick] (0.0,-0.25) -- (0.0,-1.45);
  \draw[thick,dashed] (0.0,-0.25) -- (0.0,1.05) -- (0.0,-0.25);
  \draw[thin] (2.55,-0.25) -- (2.55,0.95) -- (2.55,-0.25);
  \draw[dashed] (0.0,0.95) -- (2.55,0.95);
  \node[right] at (2.62,0.90) {6cm};
  \node[left] at (-2.55,-0.6) {8cm};
  \node[below] at (-1.2,-2.05) {10cm};
\end{tikzpicture}
\caption{}\label{fig:6.6.22}
\end{figure}

\begin{figure}[htbp]
\centering
\begin{tikzpicture}[x=1.0cm,y=1.0cm,>=Stealth]
  \draw[thick] (0,0) circle (1.7);
  \draw[thick,dashed] (0,0) ellipse (0.55 and 1.7);
  \draw[thick] (0,1.7) arc (90:-90:0.55 and 1.7);
  \draw[thick,dashed] (0,-1.7) arc (-90:90:0.55 and 1.7);
  \draw[thick] (-1.7,0.35) .. controls (-1.0,1.05) and (1.0,1.05) .. (1.7,0.35);
  \draw[thin,dashed] (0,0) -- (0,1.7);
  \node[below left] at (0,0) {$O$};
  \node[left] at (0,1.55) {$r$};
  \draw[thin] (0,0.95) -- (0.72,0.95);
  \node[above] at (1.05,0.78) {$b$};
\end{tikzpicture}
\caption{}\label{fig:6.6.23}
\end{figure}

% ==================================================================
%  徐森林、薛春华《数学分析》(第一册)
%  第 6 章复习题
% ==================================================================
%  源: 原书印刷页 444-449(PDF 物理页 450-455)
%  底本: PaddleOCR-VL-1.6 + 高清影印逐题核验校订
%  规范: exblock 容器; 全西文半角标点; 单一落盘出口 _texout.py
% ==================================================================
\begin{exblock}{复习题 6}

\item (1) 设 $f$ 在 $[a,b]$ 上 2 阶可导, 且 $f''(x)>0$. 证明:
① $f\left(\frac{a+b}{2}\right)\leqslant\frac{1}{b-a}\int_{a}^{b}f(x)\mathrm{d}x$;
② 若 $f(x)\leqslant 0$, $x\in[a,b]$, 则有
\[
f(x)\geqslant\frac{2}{b-a}\int_{a}^{b}f(x)\mathrm{d}x,\quad x\in[a,b];
\]
(2) 设凸函数 $f$ 在 $[a,b]$ 上 Riemann 可积, 证明
\[
f\left(\frac{a+b}{2}\right)\leqslant\frac{1}{b-a}\int_{a}^{b}f(x)\mathrm{d}x\leqslant\frac{f(a)+f(b)}{2}.
\]
(3) 设 $\varphi$ 在 $[0,a]$ 上连续, $f(x)$ 2 阶可导, 且 $f''(x)\geqslant 0$, 证明:
\[
\frac{1}{a}\int_{0}^{a}f(\varphi(t))\mathrm{d}t\geqslant f\left(\frac{1}{a}\int_{0}^{a}\varphi(t)\mathrm{d}t\right).
\]
(4) 设 $f$ 在 $[a,b]$ 上连续, 且 $f(x)>0$, 则
\[
\ln\left(\frac{1}{b-a}\int_{a}^{b}f(x)\mathrm{d}x\right)\geqslant\frac{1}{b-a}\int_{a}^{b}\ln f(x)\mathrm{d}x.
\]

\item (1) 设 $f(x)$ 在 $[a,b]$ 上具有 1 阶导函数, $f(a)=f(b)=0$. 证明: $\exists\,\xi\in(a,b)$, s.t.
\[
|f'(\xi)|\geqslant\frac{4}{(b-a)^{2}}\left|\int_{a}^{b}f(x)\mathrm{d}x\right|.
\]
(2) 设 $f(x)$ 在 $[a,b]$ 上具有 1 阶导函数, $f(a)=f(b)=0$, 且 $f(x)\not\equiv 0$. 证明: $\exists\,\xi\in(a,b)$, s.t.
\[
|f'(\xi)|>\frac{4}{(b-a)^{2}}\int_{a}^{b}|f(x)|\mathrm{d}x\geqslant\frac{4}{(b-a)^{2}}\left|\int_{a}^{b}f(x)\mathrm{d}x\right|.
\]
(3) 设 $f(x)$ 在 $[a,b]$ 上具有 1 阶导函数, $f(a)=f(b)=0$, 且 $f(x)\not\equiv 0$, 则 $\exists\,\xi\in(a,b)$, s.t.
\[
|f'(\xi)|>\frac{4}{(b-a)^{2}}\left|\int_{a}^{b}f(x)\mathrm{d}x\right|
\]
等价于设 $F(x)$ 在 $[a,b]$ 上 2 阶可导, $F'(a)=F'(b)=0$, 且 $F(x)\not\equiv$ 常数, 则 $\exists\,\xi\in(a,b)$, s.t.
\[
|F''(\xi)|>\frac{4}{(b-a)^{2}}|F(b)-F(a)|.
\]

\item (1) 设 $f$ 在 $[a,b]$ 上可导, $f\left(\frac{a+b}{2}\right)=0$. 证明: $\exists\,\xi\in(a,b)$, s.t.
\[
|f'(\xi)|\geqslant\frac{4}{(b-a)^{2}}\int_{a}^{b}|f(x)|\mathrm{d}x\geqslant\frac{4}{(b-a)^{2}}\left|\int_{a}^{b}f(x)\mathrm{d}x\right|.
\]
(2) 设 $f$ 在 $[a,b]$ 上可导, $f\left(\frac{a+b}{2}\right)=0$, $f(x)\not\equiv 0$. 证明: $\exists\,\xi\in(a,b)$, s.t.
\[
|f'(\xi)|>\frac{4}{(b-a)^{2}}\left|\int_{a}^{b}f(x)\mathrm{d}x\right|.
\]
举例说明: $|f'(\xi)|>\frac{4}{(b-a)^{2}}\int_{a}^{b}|f(x)|\mathrm{d}x$ 未必成立.
(3) 上述 (2) $\Leftrightarrow$ 设 $F(x)$ 在 $[a,b]$ 上 2 阶可导, $F'\left(\frac{a+b}{2}\right)=0$, $F(x)\not\equiv$ 常数, 则 $\exists\,\xi\in(a,b)$, s.t.
\[
|F''(\xi)|>\frac{4}{(b-a)^{2}}|F(b)-F(a)|.
\]

\item 设 $a,b>0$, $f$ 在 $[-a,b]$ 上连续. 又设 $f>0$ 且 $\int_{-a}^{b}xf(x)\mathrm{d}x=0$. 证明:
\[
\int_{-a}^{b}x^{2}f(x)\mathrm{d}x\leqslant ab\int_{-a}^{b}f(x)\mathrm{d}x.
\]

\item 证明: 不等式
\[
\int_{-1}^{1}(1-x^{2})^{n}\mathrm{d}x\geqslant\frac{4}{3\sqrt{n}},\quad n\in\mathbb{N},
\]
且当 $n\geqslant 2$ 时严格不等号成立.

\item 证明: $\int_{0}^{\pi}\cos nx\cos^{n}x\,\mathrm{d}x=\frac{\pi}{2^{n}}$, $n=0,1,2,\cdots$.

\item 设 $f$ 在 $[0,1]$ 上连续, 并且 $\int_{0}^{1}x^{k}f(x)\,\mathrm{d}x=0$, $k=0,1,2,\cdots,n-1$, $\int_{0}^{1}x^{n}f(x)\,\mathrm{d}x=1$. 证明: $\exists\,\xi\in(0,1)$, s.t. $|f(\xi)|\geqslant 2^{n}(n+1)$.

\item 设函数 $f$ 在 $[a,b]$ 上连续、非负且严格单调增, 由积分中值定理, $\forall n\in\mathbb{N}$, $\exists_{1}\,x_{n}\in[a,b]$, s.t.
\[
f^{n}(x_{n})=\frac{1}{b-a}\int_{a}^{b}f^{n}(t)\mathrm{d}t.
\]
证明: $\lim\limits_{n\to+\infty}x_{n}=b$.

\item 设 $f$ 是一个 $n$ 次多项式, 且满足 $\int_{0}^{1}x^{k}f(x)\,\mathrm{d}x=0$, $k=1,2,\cdots,n$. 证明:
\[
\int_{0}^{1}f^{2}(x)\mathrm{d}x=(n+1)^{2}\left(\int_{0}^{1}f(x)\mathrm{d}x\right)^{2}.
\]

\item 设函数 $f$ 在区间 $[0,+\infty)$ 上单调增, 证明: $\varphi(x)=\int_{0}^{x}f(t)\mathrm{d}t$ 在 $[0,+\infty)$ 上为凸函数. 如果 $f$ 在区间 $[0,+\infty]$ 上单调增且连续, 试用简单方法证明之.

\item 设 $f(x)$ 2 阶连续可导, $f\geqslant 0$, $f'\leqslant 0$. 证明: $\forall c\in[a,b]$, 有
\[
f(c)\leqslant\frac{2}{b-a}\int_{a}^{b}f(x)\mathrm{d}x.
\]

\item 对 $n\in\mathbb{N}$, 定义 $I_{n}=\int_{0}^{\frac{\pi}{2}}\frac{\sin^{2}nt}{\sin t}\,\mathrm{d}t$, 求极限 $\lim\limits_{n\to+\infty}\frac{I_{n}}{\ln n}$.

\item 设 $f$ 在 $[a,b]$ 上连续, 若 $\int_{a}^{b}f(x)g(x)\,\mathrm{d}x=0$ 对 $[a,b]$ 上一切满足条件 $g(a)=g(b)=0$ 的连续函数 $g$ 成立. 证明: 在 $[a,b]$ 上, $f(x)\equiv 0$.

\item 设函数 $f$ 在 $[-1,1]$ 上可导, $M=\sup|f'|$. 如果 $\exists\,a\in(0,1)$, s.t. $\int_{-a}^{a}f(x)\,\mathrm{d}x=0$. 证明:
\[
\left|\int_{-1}^{1}f(x)\mathrm{d}x\right|\leqslant M(1-a^{2}).
\]

\item 设 $f$ 在 $[a,b]$ 上连续可导. 证明:
\[
\max_{a\leqslant x\leqslant b}|f(x)|\leqslant\frac{1}{b-a}\left|\int_{a}^{b}f(x)\mathrm{d}x\right|+\int_{a}^{b}|f'(x)|\,\mathrm{d}x.
\]

\item 设 $f$ 在 $(0,+\infty)$ 上连续, 且 $\forall x>0$ 有
\[
\int_{x}^{x^{2}}f(t)\,\mathrm{d}t=\int_{1}^{x}f(t)\,\mathrm{d}t,
\]
试求出满足上述条件的一切函数 $f$.

\item 设 $f$ 在 $[0,1]$ 上 2 阶连续可导, $f(0)=f(1)=0$, 且当 $x\in(0,1)$ 时, $f(x)\neq 0$. 证明:
\[
\int_{0}^{1}\left|\frac{f''(x)}{f(x)}\right|\mathrm{d}x\geqslant 4.
\]

\item 设 $f$ 在 $[a,b]$ 上连续、单调增. 证明:
\[
\int_{a}^{b}xf(x)\,\mathrm{d}x\geqslant\frac{a+b}{2}\int_{a}^{b}f(x)\,\mathrm{d}x.
\]
(提示: 令 $F(t)=\int_{a}^{t}xf(x)\,\mathrm{d}x-\frac{a+t}{2}\int_{a}^{t}f(x)\,\mathrm{d}x$.)

\item 设 $b>a>0$, 证明不等式:
\[
\ln\frac{b}{a}>\frac{2(b-a)}{a+b}.
\]
(提示: 对 $f(x)=-\frac{1}{x}$ 应用题 18 的结果.)

\item 设函数 $f$ 在 $[0,1]$ 上 2 阶连续可导, 且 $f(0)=f(1)=f'(0)=0$, $f'(1)=1$. 证明:
\[
\int_{0}^{1}\left[f''(x)\right]^{2}\mathrm{d}x\geqslant 4.
\]
指出式中等号成立的条件.

\item 设函数 $f$ 在 $[0,1]$ 上 Riemann 可积, 且有正数 $m$ 与 $M$, 使得 $m\leqslant f(x)\leqslant M$, $\forall x\in[0,1]$. 证明:
\[
1\leqslant\int_{0}^{1}f(x)\,\mathrm{d}x\int_{0}^{1}\frac{\mathrm{d}x}{f(x)}\leqslant\frac{(m+M)^{2}}{4mM}.
\]

\item 设 $x(t)$ 在 $[0,a]$ 上连续且满足
\[
|x(t)|\leqslant M+k\int_{0}^{t}|x(t)|\,\mathrm{d}t,
\]
这里 $M$ 与 $k$ 为正常数. 证明: $|x(t)|\leqslant M\mathrm{e}^{kt}$, $t\in[0,a]$.

\item 设 $f$ 为 $[-1,1]$ 上的连续函数, 且对 $[-1,1]$ 上的任何偶函数 $g$, 积分 $\int_{-1}^{1}f(x)g(x)\,\mathrm{d}x=0$. 证明: $f$ 在 $[-1,1]$ 上为奇函数.

\item 设 $a,b,n\in\mathbb{N}$, 令 $f(x)=\frac{x^{n}(a-bx)^{n}}{n!}$. 证明:
\begin{enumerate}[label=(\arabic*)]
  \item $f\left(\frac{a}{b}-x\right)=f(x)$;
  \item $f^{(k)}(x)$ ($0\leqslant k\leqslant 2n$) 当 $x=0$, $x=\frac{a}{b}$ 时取得值为整数;
  \item 假设 $\pi$ 为有理数, 即 $\pi=\frac{a}{b}$, $a$ 与 $b$ 为既约正整数. 则可证 $\int_{0}^{\pi}f(x)\sin x\,\mathrm{d}x$ 为整数. 由此证明 $\pi$ 不可能为有理数, 即 $\pi$ 为无理数.
\end{enumerate}

\item 设函数 $f$ 在 $[0,2]$ 上连续可导, $f(0)=f(2)=1$, 且 $|f'|\leqslant 1$. 证明: $1\leqslant\int_{0}^{2}f(x)\,\mathrm{d}x\leqslant 3$.

\item 证明: $\int_{0}^{x}\mathrm{e}^{xt-t^{2}}\mathrm{d}t=\mathrm{e}^{\frac{x^{2}}{4}}\int_{0}^{x}\mathrm{e}^{-\frac{t^{2}}{4}}\mathrm{d}t$.

\item 证明: (1) $\forall n\in\mathbb{N}$, 积分
\[
I_{n}=\int_{0}^{1}\frac{1}{x}\left[\left(\ln x+\ln\frac{1+x}{1-x}\right)^{n}-\ln^{n}x\right]\mathrm{d}x
\]
是收敛的;
(2) 当 $n$ 为偶数时, $I_{n}=0$.

\item 证明: $\int_{0}^{+\infty}\left(\frac{\sin x}{x}\right)^{2}\mathrm{d}x=\int_{0}^{+\infty}\frac{\sin x}{x}\mathrm{d}x$.

\item 计算积分:
\begin{enumerate}[label=(\arabic*)]
  \item $\int_{0}^{+\infty}\frac{\ln x}{(1+x)\sqrt{x}}\,\mathrm{d}x$;
  \item $\int_{0}^{\frac{\pi}{2}}\frac{\mathrm{d}x}{1+\tan^{100}x}$;
  \item $\int_{1}^{+\infty}\left(\frac{1}{[x]}-\frac{1}{x}\right)\mathrm{d}x$;
  \item $\int_{0}^{+\infty}\frac{\mathrm{e}^{-x^{2}}}{\left(x^{2}+\frac{1}{2}\right)^{2}}\,\mathrm{d}x$.
\end{enumerate}

\item 设 $f$ 为 $[-1,1]$ 上的连续函数. 证明:
\[
\lim_{h\to 0^{+}}\int_{-1}^{1}\frac{h}{h^{2}+x^{2}}f(x)\,\mathrm{d}x=\pi f(0).
\]

\item 设 $f$ 为 $[-1,1]$ 上的连续函数. 证明:
\[
\lim_{\lambda\to 0^{+}}\frac{1}{2\lambda}\int_{-1}^{1}f(x)\mathrm{e}^{-\frac{|x|}{\lambda}}\,\mathrm{d}x=f(0).
\]

\item 设 $f$ 在 $[0,\pi]$ 上 Riemann 可积. 证明: 不能同时有
\[
\int_{0}^{\pi}|f(x)-\sin x|^{2}\,\mathrm{d}x\leqslant\frac{3}{4},\quad\int_{0}^{\pi}|f(x)-\cos x|^{2}\,\mathrm{d}x\leqslant\frac{3}{4}.
\]

\item 设 $f$ 在 $[a,b]$ 上 2 阶连续可导. 证明: $\exists\,\xi\in(a,b)$, s.t.
\[
\int_{a}^{b}f(x)\,\mathrm{d}x=(b-a)f\left(\frac{a+b}{2}\right)+\frac{(b-a)^{3}}{24}f''(\xi).
\]

\item 设
\[
F(x)=\begin{cases}\displaystyle\int_{0}^{x}\cos\frac{1}{t}\,\mathrm{d}t, & x\neq 0,\\[2mm] 0, & x=0.\end{cases}
\]
求 $F'(0)$.

\item 函数 $u$ 在 $[0,+\infty)$ 上满足积分--微分方程
\[
\begin{cases}
u'(t)=u(t)+\displaystyle\int_{0}^{1}u(x)\,\mathrm{d}x,\\
u(0)=1.
\end{cases}
\]
试确定函数 $u$.
\end{exblock}

